% =====================================================================
%  H. Brøchner
%  Den nyere Philosophies Historie i Grundrids.
%  (= Philosophiens Historie i Grundrids, 2. Deel.)
%  Kjøbenhavn: P. G. Philipsens Forlag, 1874.   291 pp.
%
%  "The History of Modern Philosophy in Outline" -- TRANSLATION (English).
%
%  Translated from transcription.tex (Danish) in this directory, the text
%  of record, as it stood at commit 75a01bf (2026-10-08), after its repair
%  of 2026-10-08 (ACCURACY-LEDGER.tsv).  If the transcription is repaired
%  again, re-check the English at every changed reading
%  (TRANSLATION-PLAYBOOK.md §6a).
%
%  STRUCTURE mirrors the transcription 1:1: the same hand-built heads
%  (\begin{center} blocks, \divrule, \dblrule, \markboth), no \section,
%  and every \opage{N} of the Danish once and in order, so a page of the
%  English can be cross-referenced to the Danish and the scan.  Footnotes
%  are marked *) and **) and the counter is reset at each printed page,
%  as in the transcription.
%
%  CONVENTIONS (see RESUME-NOTES.md, "Translation", for the glossary):
%   * Greek is kept as printed, in \gk{} (the print's cursive Greek
%     fount); a bracketed English gloss is added only where Brøchner
%     does not translate it himself.
%   * Latin, German and French are kept in the language and, as in the
%     print, set upright roman (the print reserves italic for Greek).
%   * The print's guillemets «...» become ``...''.
%   * ɔ: (the old Danish id est) becomes "i.e."
%   * Letterspacing (\emph{} in the Danish) stays \emph{} (italic).
%   * ERRATA: the print's list B ("2den Deel", 17 entries, back matter
%     of the transcription) is NOT applied in the Danish, which is
%     diplomatic.  The English follows Brøchner's correction at each of
%     its sites, with a % comment there giving the printed reading.
%   * Printer's defects logged in the transcription (e.g. p. 9 "Paton")
%     are translated in their evident sense, with a % comment at the site.
%
%  STATE: DRAFT COMPLETE (2026-10-08): title leaves, Contents, the whole body
%  (pp. 1--291) and the errata note.  Pp. 1--11 translated inline; pp. 11--291
%  by 25 translator agents in five waves from one brief and glossary (see
%  RESUME-NOTES.md, "TRANSLATION"), each part checked by script against the
%  Danish page by page (\opage, footnotes, counter resets, \emph, \gk,
%  paragraph counts).  Aand = mind for the knowing subject (Descartes,
%  Malebranche, Spinoza, Leibnitz, Kant), spirit for Hegel's and Schelling's
%  Geist and in the cultural sense.  The translators' odd spots in the Danish
%  are listed in RESUME-NOTES for the repair session; if the transcription is
%  repaired at any of them, re-check the English there (§6a).
% =====================================================================

\documentclass[12pt,a4paper,openany]{book}
\usepackage[utf8]{inputenc}
\usepackage[T1]{fontenc}
\usepackage{libertinus}
\usepackage{libertinust1math}
\usepackage{textalpha}
\usepackage[english]{babel}
\usepackage[protrusion=true,expansion=true]{microtype}
\usepackage{enumitem}
\usepackage{setspace}
\usepackage[margin=1.2in]{geometry}
\usepackage{fancyhdr}
\setlength{\headheight}{28pt}
\addtolength{\topmargin}{-13.5pt}
\usepackage{hyperref}
\hypersetup{pdftitle={The History of Modern Philosophy in Outline (1874)},
  pdfauthor={H. Brøchner}, hidelinks}

% Footnote marks as printed: *) and **), reset at each printed page by
% \setcounter{footnote}{0}, exactly as in the transcription.
\renewcommand{\thefootnote}{%
  \ifcase\value{footnote}\or *)\or **)\or ***)\or \dag)\or \dag\dag)\else
  \arabic{footnote})\fi}

% The short centred divisional rule and the double rule, as in the
% transcription (header §3 there).
\newcommand{\divrule}{\begin{center}\rule{2cm}{0.4pt}\end{center}}
\newcommand{\dblrule}{\begin{center}%
  \rule{2cm}{0.4pt}\\[1.2pt]\rule{2cm}{0.4pt}\end{center}}

% Greek, in the print's cursive fount (see the transcription's preamble).
\newcommand{\gk}[1]{\textit{#1}}

\pagestyle{fancy}
\fancyhf{}
\fancyhead[LE,RO]{\thepage}
\fancyhead[RE]{\textit{The History of Modern Philosophy in Outline}}
\fancyhead[LO]{\textit{\leftmark}}
\fancypagestyle{plain}{\fancyhf{}\fancyhead[LE,RO]{\thepage}}

\setstretch{1.3}

% ---- original-page markers: the Danish printed page begins here ----
\usepackage{marginnote}
\newcommand{\opage}[1]{\marginnote{\footnotesize\textit{[#1]}}}

\begin{document}

\chapter*{Title Pages}
\markboth{Title Pages}{Title Pages}
\thispagestyle{empty}

\begin{center}
  {\LARGE \emph{The History of Philosophy}}\\[8pt]
  {\LARGE \emph{in Outline.}}
\end{center}
\vspace{18pt}
\begin{center}
  \textbf{By}
\end{center}
\vspace{10pt}
\begin{center}
  \textbf{Dr.\ \emph{H. Brøchner,}}
\end{center}
\begin{center}
  \textbf{Professor of Philosophy.}
\end{center}
\vspace{14pt}
\dblrule
\vspace{14pt}
\begin{center}
  \textbf{Part 2.}
\end{center}
\begin{center}
  \textbf{\large The History of Philosophy in Modern Times.}
\end{center}
\vspace{14pt}
\dblrule
\vspace{18pt}
\begin{center}
  \textbf{Copenhagen.}
\end{center}
\begin{center}
  \textbf{P.\ G.\ \emph{Philipsen, Publisher.}}
\end{center}
\begin{center}
  \textbf{1874.}
\end{center}
\vfill
\begin{center}
  \small Translated from the Danish
\end{center}
\newpage
\thispagestyle{empty}

\begin{center}
  {\LARGE The History of Modern Philosophy}\\[8pt]
  {\LARGE in Outline.}
\end{center}
\vspace{18pt}
\begin{center}
  \textbf{By}
\end{center}
\vspace{10pt}
\begin{center}
  \textbf{Dr.\ \emph{H. Brøchner,}}
\end{center}
\begin{center}
  \textbf{Professor of Philosophy.}
\end{center}
\vspace{14pt}
\dblrule
\vspace{18pt}
\begin{center}
  \textbf{Copenhagen.}
\end{center}
\begin{center}
  \textbf{P.\ G.\ \emph{Philipsen, Publisher.}}
\end{center}
\begin{center}
  \textbf{1874.}
\end{center}
\newpage
\thispagestyle{empty}
\vspace*{0.62\textheight}
\begin{center}
  \textbf{Printed by Bianco Luno.}
\end{center}

\chapter*{Contents of the Second Part.}
\markboth{Contents}{Contents}

\divrule

{\small [The page numbers are those of the Danish original, given in the
margin of this translation.]\par}
\medskip

\begingroup
\parindent=0pt \parfillskip=0pt \tolerance=9999 \bfseries
\newcommand{\indholdleader}{%
  \nobreak\leaders\hbox{\hspace{0.22em}.\hspace{0.22em}}%
  \hskip 0.8em plus 1fill minus 0.8em\nobreak}
\newcommand{\indholdtop}[2]{\hangindent=2em\hangafter=1\noindent
  #1\indholdleader\mbox{#2}\par}
\newcommand{\indholdsub}[3]{\hangindent=2.4em\hangafter=1\noindent
  \hspace*{0.4em}\makebox[1.4em][r]{#1}\enspace #2\indholdleader\mbox{#3}\par}
\noindent\hfill{\small Page}\par
{\large\indholdtop{Second Main Division: The Philosophy of Modern Times}{1---275}}
\indholdtop{\emph{Division}}{1}
\indholdtop{First Section: \emph{The Philosophy of the Transitional Period}}{2---19}
\indholdtop{Second Section: \emph{The Philosophy of Modern Times in Its
Independent Formation}}{20---275}
\indholdtop{\emph{Chapter One:} The General Character of Modern Philosophy.
Its Forms of Development}{20---23}
\indholdtop{\emph{Chapter Two:} The \emph{Realistic} Line of Development}{23---76}
\indholdsub{1.}{Bacon}{23---37}
\indholdsub{2.}{Hobbes}{37---42}
\indholdsub{3.}{Locke}{42---51}
\indholdsub{4.}{Berkeley}{51---53}
\indholdsub{5.}{Hume}{54---58}
\indholdsub{6.}{Condillac}{58---66}
\indholdsub{7.}{Materialism (Système de la nature)}{66---76}
\indholdtop{\emph{Chapter Three:} The \emph{Idealistic} Line of Development}{77---275}
\indholdtop{I. \emph{Dogmatic} Idealism}{77---137}
\indholdsub{1.}{Descartes}{77---88}
\indholdsub{2.}{Occasionalism (de la Forge, Geulincx, Malebranche)}{88---94}
\indholdsub{3.}{Spinoza}{94---115}
\indholdsub{4.}{Leibnitz}{116---128}
\indholdsub{5.}{Idealism in the 18th Century}{128---137}
\indholdtop{II. \emph{Critical} Idealism}{137---275}
\indholdsub{1.}{Kant}{137---182}
\indholdsub{2.}{The Reaction against Criticism from the Philosophy of
Faith: Jacobi}{182---191}
\indholdsub{3.}{The Development of Criticism into Subjective Idealism:
Fichte}{192---204}
% The Indhold words this entry differently from the head on printed 204
% ("Reactionen mod Subjectivismen fra «Realismen»"); both are kept.
\indholdsub{4.}{The Reaction against Subjectivism from ``Realism'':
Herbart}{204---224}
\indholdsub{5.}{Speculative Idealism:  a) Schelling}{224---241}
\indholdsub{6.}{b) Hegel}{241---275}
\indholdtop{\emph{Conclusion}}{275---291}
\endgroup

\chapter*{Second Main Division.}
\markboth{The Philosophy of Modern Times}{The Philosophy of Modern Times}
\opage{1}
\divrule

\begin{center}
  \textbf{\large The Philosophy of Modern Times.}
\end{center}

\divrule

\begin{center}
  \textbf{Division.}
\end{center}

The philosophy of modern times, taken in the wider sense, comprises
what in the Introduction (Part 1, p. 8 ff.) was distinguished as
\emph{the philosophy of the transitional period} and as the
\emph{independently formed modern philosophy}. Common to these
forms is the tendency characterized at the close of the section on
Scholasticism, by which they differ in principle from the thought of
the Middle Ages. But whereas in the transitional period the basic
character of thought remained essentially the same as in Greek
philosophy, and something new was at most intimated and postulated,
the philosophy of modern times in the narrower sense, as it frees
itself from the authority of antiquity, develops a new and distinctive
fundamental form of thought, a new world-view and a new method, by
which it takes its place as the second main historical form of
philosophy, with equal right beside the philosophy of classical
antiquity.

\divrule

\opage{2}
\begin{center}
  \textbf{\large First Section.}
\end{center}

\markboth{First Section}{First Section}
\divrule

\begin{center}
  \textbf{The Philosophy of the Transitional Period.}
\end{center}

\divrule

By the transitional period we understand the span of roughly a
century and a half in the history of philosophy that lies between the
dissolution of Scholasticism and the independent development of modern
philosophy. The designation, transitional period, points partly to the
period's position as the middle term between the Middle Ages and modern
times, partly to its forming the \emph{philosophical} middle term
between the two great historical fundamental forms of philosophy. With
respect to this last relation it becomes the transitional period
\gk{κατ᾽ ἐξοχήν} [par excellence].

Such a transitional period was necessary at the close of the Middle
Ages before the new main form of philosophy could emerge in its
developed determinateness. In Scholasticism thought, since its
essential object was what thought cannot appropriate, had won only a
formal development, without being able to unfold an independent
content of thought. As it freed itself from the dominion of the Church
and, proclaiming a new world-view, turned in combat against
Scholasticism, it therefore had to seek the content necessary for this
struggle, which had to be such that in it the liberated human thought
could find itself again and recognize reality, where it lay ready as
given; and it found it in the philosophy of antiquity, with which it
felt itself akin. It then exchanged the enforced submission of
Scholasticism to a foreign authority for a self-chosen submission to a
kindred one, in which the thought of the new age, which could now have
reality and presence in its content, could gain fullness and
independence through the appropriation and working-out of this content
of reality, into which it had first \opage{3}to live its way. In all
the manifold and heterogeneous philosophical attempts that fall within
this span of time, the common distinguishing mark, beside the
opposition to the philosophy of the Middle Ages, is this dependence on
antiquity, despite the striving for independence and the demand for it
that gradually comes forward.

If we seek to grasp the essential peculiarity of the period through
the thoughts that are central in it, a great difficulty lies in the
multiplicity and the apparently total heterogeneity of the
philosophical attempts that fill the period. Consciously they all
agree only in the negative: in the opposition to the kind and
direction of the earlier philosophy; but the counter-position to it
appears in numerous shapes, --- all the more numerous as the taking-up
of the forms of thought of the past, conditioned by the lack of
independence, leads to various combinations on which the individual
striving of the thinkers gains a preponderant influence, --- and what
is positive in the spiritual tendency that runs through the various
attempts comes to clarity only successively. Only by fixing one's eye
on the peculiarity of the ages between which this period forms the
transition, and on the goal to which the development in it has
historically led, and by immersing oneself in the systems, as if
grasping their hidden thoughts, does it become possible to find the
guiding principle. In this age, as in every more strongly agitated
age, the task is to see, through the spiritual current that shows
itself on the surface and is moved by every breath of air, the current
that moves in the depths and easily escapes observation. Through the
former, which is the image of the conscious striving of individuals,
one must look out for the powers by which they are unconsciously led;
one must see the mysterious influence of the hidden powers of the past
and of the future, as restraining and as impelling; only in this way
does one understand such a time of ferment, rich in mighty forces, in
energetic endeavors, but \opage{4}poor in sober-mindedness, in
limitation and in clarity of consciousness.

If, then, we look to the peculiarity of thought in the ages that the
transitional period connects, Scholasticism has shown itself to be
characterized by the unfreedom of thought, by its want of reality and
by its formalism. The peculiarity of the independently formed modern
philosophy we shall see to lie in this: that, starting from doubt, in
which the transitional period ends because it cannot master the
oppositions it would unite, it sets itself the task, independently of
the authority of the Church and of antiquity, of developing out of the
interior of self-consciousness, or through experience directed at
nature, an independent, uniform content of thought, secured by the
exact method and the unshakable point of departure. The truth of
nature and the autonomy of self-consciousness become its watchword,
and with this watchword it comes into opposition in principle to the
tendency and the view of life of the Middle Ages.

By these two boundary points the character of the period that lies
between them is indirectly determined. It must point in the same
direction as what issues from it as the result of its development; in
its fundamental thoughts it must reveal its kinship with this, but in
the form of those fundamental thoughts, in their peculiar coloring, it
must at the same time show them as proceeding from the thoughts of the
past, and as not yet released from their connection with them.

This is confirmed, too, when, guided by that indication, we find, as
the fundamental thoughts which, more clearly or more obscurely, govern
the systems of the transitional period, and toward which the tendency
of those systems points even where the distance seems greatest, the
closely connected thoughts of \emph{the divine universe} and \emph{the
divine, microcosmic human being}. These thoughts, which are clearly
expressed in the most developed systems of the period, issue with
necessity from the tendency that calls the new thinking into life. As
nature again acquires significance and is no longer the merely
\opage{5}essenceless, it must be placed in an inner, essential relation
to the Godhead; the notion of an arbitrary creation in time must,
albeit with inconsistencies, pass over into the thought of an
\emph{eternal} creation, i.e.\ an eternal, necessary unfolding of the
life and essence of the Godhead, so that the Godhead and the universe
become the same thought, only seen from two sides. And in this divine
universe, in this essentially valid, inspirited nature, man, whose
spirit is the apprehending mirror of all the universe's forms of
existence, becomes the center, as a microcosmic being. The nimbus of
divinity that irradiates the universe as a whole irradiates human
nature, and with enthusiastic rapture, with the overweening boldness
of a young freedom, the age believes in the infinity of the human
being: in its infinity in power, which magically reaches out over all
the spheres of existence; in its infinity in being, which is bounded
by no time; in its infinity in knowledge, which is checked by no
barrier.

In these thoughts the central thoughts of modern philosophy come into
view, but on the one hand, as will appear later, as still influenced
by the notions of the Middle Ages; on the other, as determined through
a reflex of the thoughts of antiquity. Nature, whose unity is demanded
by thought anticipating experience, is intuited by imagination as an
infinitized cosmos, and throughout this whole period it becomes
essentially the object of abstract-metaphysical thought and of
imagination. And man, who forms the center in this cosmos, is
determined in the ancient manner, by way of metaphysical thinking, as
the harmonious \gk{τέλος} [end] of natural existence, as that which
concentrates the whole of existence in itself, yet in such a way that
the modern age's accentuation of the infinity of inwardness gives the
ancient thought a new nuance. For however much the side of nature is
accentuated, it is still man's \emph{spirit} that becomes the ideal
center in the infinite existence, which could no longer, like the
bounded cosmos of the Greeks, have the earth for its material center.
In general the ancient, as it is revived in the
mo\opage{6}dern age, is constantly modified, as if in passing through
refracting media.

While the two main thoughts of the transitional period indicate the
direction of modern philosophy, they form at the same time the
antithesis and the parallel to the central notions of the Middle Ages:
the God without nature and the historical God-man. The world-view and
view of life that proceed from them must differ in principle from
those of the preceding age; but only successively does the antithesis
become really clear, and in the form in which those fundamental
thoughts appear there shows itself from the beginning a theological
and theosophical element, beneath which the element of nature must, as
it were, work its way forward. It is precisely this process, this
development of what is striven for out of the swaddling in which it at
first lies wrapped, that provides the point of view for a distinction
between two main divisions of the transitional period, between which
the beginning of the ecclesiastical Reformation becomes the boundary.
In the first division the theological element had persisted from the
Middle Ages and veiled the element of nature, which, where it asserted
itself as such, did so under forms taken from the old philosophy. As
the Reformation gives the theological element an outlet and separates
it off by concentrating it in a domain of its own, the tendency toward
the natural emerges within philosophy with clearer consciousness and
with decisive preponderance. The rapid flowering of the natural
sciences in this age, which prepares independence from antiquity
through a view of nature widened and corrected by epoch-making
discoveries and inventions, and through the formation of a method that
philosophy can take up, at the same time frees it more completely from
its dependence on theology, and the remaining theological elements are
separated out more decisively, even in a hostile manner. The
opposition, which in the first division of the period had been
directed against the philosophy of the Middle Ages, comes in the
second to be directed also against the religion of the Middle Ages.
% "in the second": errata list B, p. 6 l. 3 from below, "dette, l. det
% andet"; printed "bliver i dette ogsaa" ("in this one").
Beside systems that are indifferent to Christianity or formally
\opage{7}subordinate themselves to its authority, there arise systems
in which the opposition to it appears with clear consciousness, even as
a passionate hatred, and a struggle arises in which philosophy gets its
martyrs (Bruno, Vanini).

\divrule

The most significant representative of the first span of time is also
the first in whom the new tendency is distinctly stamped: the German
Cardinal \emph{Nicolaus Cusanus} (born at Kues on the Moselle 1401,
died 1464). Through doubt about the capacity of our understanding for
a knowing of the divine he is led from ``positive'' theology to
``negative'' theology, to the avowal of our not-knowing; but in the
manner of the mystics and the Neoplatonists he gets beyond doubt and
the negative by assuming an immediate intuition of God in an intuitio
intellectualis, and his ``mystical'' theology determines God as the
infinite unity of all opposites. The world he conceives as having
issued from the higher matter that is in \emph{God}, as unbounded in
time and space, and he comes close, --- inasmuch as the determinations
concerning the world that state its contingency, finitude and
imperfection are taken as valid only from relative points of view,
while the world, as the expression of God's fullness of power, is
perfect and infinite, and has a beginning only in the sense that God is
its principle, --- to positing God and the universe as the same
concept, only seen from different sides. Using Platonic and
Pythagorean determinations, he designates the universe as God's
mirror, as the perfect, harmonious unfolding, determined by numerical
ratios, in multiplicity, of what in God is in unity, as the unfolded
greatest, while God is the absolutely greatest. It is a living
organism, in which the whole is in each part, all is in all, death
only a transformation. Man is in it the uniting center, to whose forms
of knowledge the forms of existence correspond. The unity of the
greatest within humanity and the absolutely greatest is Christ, the
God-\opage{8}man; in him the infinity of human nature is intuited.
Knowledge of the world is the condition of knowledge of God; but the
latter first becomes actual and complete in that beholding, through a
supernatural revelation and through faith.

In the Cusan, whom his mechanical and astronomical studies led to a
view of nature in its reality, we see the validity of nature asserted
by its ground being placed back in God's \emph{essence}, and he can
determine it by predicates that recall the Platonic conception of the
cosmos. Yet his essential interest is still the theological one; he
will not depart from the doctrine of the Church, and even if he lets
knowledge of nature be required as the condition of knowledge of God,
he places it in such a relation to supernatural revelation that its
significance is, as it were, revoked. Nevertheless the Cusan has
expressed in his conception of the universe thoughts that were later
taken up by Giordano Bruno and by Leibnitz.
% "taken up by Giordano Bruno": the print omits the first "af" ("optagne
% Giordano Bruno og af Leibnitz"), a defect logged in the transcription.

\emph{The Platonic Academy} in Florence represents the enthusiastic
adherence to Plato's doctrine that is characteristic of the beginning
of the transitional period, a doctrine which, however, was understood
essentially through the medium of Neoplatonism and was not kept free
of admixture with foreign elements, e.g.\ Aristotelian ones. Marsiglio
Ficino and Giovanni Pico of Mirandola are the most significant members
of this Academy. \emph{Marsiglio Ficino} (1433---1499) determines
passive matter, active quality, the rational soul, the eternal
intelligences and God as the grades of existence. God, the absolute,
free reason, independent of the eternal primal matter, is the eternal
primal form and primal force of things, who comprehends the world
\emph{from within} in intuiting himself, and has ordered everything so
that all individual things strive in concert after the good and partake
of it by virtue of the God who works in them. In the world, which has
life through the world-soul, man, on whom the main interest of the
system falls, is the uniting center through the union of soul and
body. He mirrors the universe microcosmically, his \opage{9}knowledge
is infinite, and his soul has immortality, like the world-soul and the
souls of the spheres. In \emph{Giovanni Pico of Mirandola} (1463---1495)
the same faith in the loftiness of human nature appears, in its
determination as universal, i.e.\ as including potentially all forms of
existence. Man is infinite in knowledge and existence; he embraces
everything in knowledge and love. Like Ficino, Pico takes up the
tradition of the past syncretistically. Plato he regards as the
preserver of an older tradition, which Plato expresses through symbols
and which is interpreted by the Neoplatonists.
% "Plato": printed "Paton" (defect logged in the transcription; the
% annotator inserted the l in ink).
From the Jewish Kabbalah too Pico takes up elements. His interest in
the physical is united with a theosophical tendency.

The other philosophers of this division are of lesser interest, since
either, like the \emph{Aristotelians}, with a still greater lack of
independence, they strive only to reproduce one particular form of
ancient thought, or, like the \emph{theosophists} proper, they give
only something unclear and fermenting, or, like the
\emph{philologists}, only something fragmentary, nothing really
coherent in a philosophical direction. Yet in all of them we can find
the fundamental thoughts indicated above at least intimated. In the
theosophists (Cornelius Agrippa, Reuchlin and others) they are
intimated through the essentiality ascribed to the material, the
natural, and by the doctrine of the threefold magic, by which man
stands in a relation of mastery to threefold existence. In the
Aristotelians they come forward inasmuch as the return to the genuine,
original Aristotelian doctrine leads again to the recognition of the
essentiality of nature and of the connection of the human spirit,
imperishable as universal, with nature as a whole. Precisely in them
the consciousness of the opposition to the Christian emerges most
clearly, even if this consciousness is accompanied by an, at least
formal, submission to theological truth, (as in Pomponazzo
1462---1524).
% "submission": printed "Underskastelse" (defect logged in the
% transcription).
And even in the last group in this division, in those philologists who
also had a, though subordinate, significance as philosophers
(Laurentius Valla, Rudolph Agricola and others), there is,
\opage{10}in the endeavor to return to naturalness in thinking, a
reference to the essentiality of nature and to the truth of thought in
its harmony with nature, which as it were distantly intimates the goal
of the basic tendency: the fundamental thoughts.

\divrule

The Reformation now forms, as already stated, the boundary between the
two divisions of the period. It has a likeness to the contemporary
philosophical phenomena which shows itself in the emphasis on
subjectivity, in the return to the original through a break with
tradition, and in a primitive relation to its essential object without
the mediation of a foreign authority; and by this likeness the
Reformation proves to be an expression of the whole age's tendency of
life. But it is only a limited, single expression of it, which, just
as it had been prepared not merely by the preceding religious
development but by the whole earlier spiritual development, by the
rebirth of science and art, by the awakening sense for reality and the
gradual organization of a secular ethical life, so also had to be
supplemented by the unconditional independence of philosophy and by
political and social reform, if the whole tendency of the spirit was
to be reformed in truth. And beside the likeness of the Reformation to
the contemporary spiritual phenomena, the unlikeness that separates it
from them must now be stressed no less. This unlikeness has its ground
in the peculiarity of the Reformation as a religious phenomenon. When
Protestantism pointed to subjectivity, it was not in order to assert
the independence of the human spirit, the autonomy of
self-consciousness, but precisely in order to let man, in the
innermost of his soul, give himself up in faith, so as to receive
everything from divine grace. And when it went back to the original,
to the historical source of the Christian faith, it was not in order to
let thought freely test the content of Scripture, but in order to bow
beneath what was revealed as an unshakable authority. It is therefore
unjustified when \opage{11}one designates the principle of
Protestantism without further ado as a principle of freedom, as the
principle of free thought. Protestantism has at most indirectly, by the
break with the authority of Catholicism and by the ferment it called
forth, given impulses toward a more independent thinking to many, also
within the domain of Catholicism; but it is itself just as unfree in
its thinking as Catholicism, having only exchanged the authority of the
Church for the authority of Scripture. It is so even to a higher
degree, inasmuch as, still more consistently than Catholicism, it
diminishes and abases the human for the glory of God, and inasmuch as,
having arrested tradition at a definite stage, it does not have in
itself the principle of development that lies in a constantly advancing
tradition, which at every time has its present, personal
representation. Luther speaks with the same scorn of philosophy and of
``Madam Reason'' as he speaks of the world-view of Copernicus. If, in
the course of the development of Protestantism, free thought has
penetrated its domain to a higher degree than that of Catholicism ---
where throughout the whole first century after the Reformation the
sciences, and philosophy in particular, flourished far more richly than
in the Protestant lands --- it has penetrated from other movements in
the life of the spirit; it has not developed out of Protestantism
itself. It has penetrated by an unconscious influence, which has been
supported by the lack of clarity about the relation of the humane forms
of life to the Christian, and by the fact that Protestantism did not,
like Catholicism, have the closed power of the Church to set against
it. Free thought, which unfolds in modern philosophy and in the
empirical sciences, goes its way independent of and unconcerned about
confessional differences.

\divrule

In the second division of the transitional period we could distinguish
between such endeavors as appear only as continuations of what was begun
earlier, and at any rate differ \opage{12}quantitatively from it by the
greater, if not consistently maintained, significance that is ascribed to
the material (the philosophizing philologists, the Platonic and
Aristotelian schools), and such as have a peculiarity characteristic of
this span of time, which comes forward through the attempt to find a new
principle of knowledge, or through a new determination of the grounds of
existence, or in both these ways (Paracelsus, Cardano, Telesio, Bruno,
Vanini, Campanella, Jacob Bøhm). In these last the relation to
Christianity is predominantly that of indifference or of opposition; only
exceptionally does a mediating or restorative tendency appear, as in
% printed "restáurerende" (defect logged in the transcription)
Campanella, whose philosophy's direction is influenced by the movement
that was called forth by the endeavors of the Catholic Church toward a
regeneration, or, as in Jacob Bøhm, a Christian-religious interest, into
which, however, an element of nature foreign to Christianity intrudes. At
the close of the transitional period the consciousness of the conditions
for the solution of the infinite task of knowledge, enlivened by the
contemplation of nature, in union with the strife between the different
elements of the systems and the inherited dualistic separation between
the spiritual and the corporeal, calls forth a skepticism (Montaigne,
Charron, Sanchez), which nevertheless preserves the age's direction
toward nature and points to a new point of departure in
self-consciousness.

In all the different forms the fundamental thoughts of the transitional
period come forward. Concerning the philosophical tendency that only
continues what was begun in the preceding division, what was remarked about this holds.
% errata list B, p. 12: printed "der", corrected "om dette"
The more independent attempts, which behave indifferently or with
hostility toward Christianity, point to the intuition of the infinite
% printed "Chrisendommen" (defect logged in the transcription)
universe as the totality of existence; they emphasize the significance
and essentiality of matter, they seek physical principles for the finite
forms of existence, they give the senses and experience a decisive
significance for knowledge and let the human mind have in itself the
source of truth and the capacity to comprehend existence. Campanella will
philosophize according to the testimony of God \opage{13}and of the
\emph{senses}, and sees in nature a divine revelation beside that of the
Bible, a living mirror of the Godhead; and in Jacob Bøhm there appears
distinctly the recognition of the significance of the natural, even if it
becomes a ``nature in God'', and the intuition of the infinity of
knowledge, even if it is as a knowledge of Christ in us. And the skeptics
show their spiritual connection with the age by pointing to
self-consciousness and to the law of nature as what is relatively certain
and as the only reliable norm and point of departure of thought.

In \emph{Giordano Bruno} (born c.~1548, burned in Rome 1600) the thought
of the age finds its clearest and most characteristic expression. In
appropriating the conception of God as the unity of opposites that had
been asserted at the beginning of the period by Nicolaus Cusanus, and in
determining it more precisely by conceiving God as the unity of the four
Aristotelian causes, he also places in God matter ``pregnant with all
forms'', and thereby brings the concept of God close to the concept of
self-active, all-comprehending nature, the ensouled world-all. From the
necessity of God's ``absolutely infinite'' essence the infinity of the
world eternally unfolds, in which the strife of opposites, amid the
uninterrupted coming-to-be and passing-away, motion and change of the
parts, produces the harmony of the infinite whole, eternally equal to
itself. The concept of the universe as this infinite, living,
all-comprehending, harmonious whole covers the concept of God so closely
that only the hidden influence of the notions of the past, which will be
shown later, prevents their complete identification, and that a
distinguishing mark can be produced only by force. In the world, ensouled
at every point, the ether, which corresponds to matter, is the universal
element, from which the physical fundamental oppositions: fire and water,
issue as conditions of the uninterrupted becoming of things. By the
unfolding of the divine unity the indivisible monads are posited, the
elementary parts of all that exists, ideal atoms, which have inner
infi\opage{14}nity and imperishability in the successive assumption of
all forms. By their combination the body is formed, which is held
together and governed by the soul as the active, omnipresent central
monad. Man is the bond of existence; his infinite and immortal being
comprehends all the spheres of existence and finds its highest form of
life in the knowledge of God, which is prepared by the gradual ascent
from the sensuous images of the ideas and is attained in contemplation,
through the ``heroic love'' which in ``divine passion'', in the ecstasy
of self-surrender, forgets the ever-enduring limitedness of finite
existence.

Even in Giordano Bruno, the greatest and freest thinker of the
transitional period, we nevertheless still see that influence of the
fundamental view of the Middle Ages which acts as a hindrance throughout
the whole period. The thoughts of the divine universe and of the divine,
microcosmic, spiritual-natural human being had, consistently maintained
and developed, to lead to seeing the ensouled universe as all existence
and as eternally being through itself, and to seeing the spiritual and
the corporeal in it and in man as indissolubly and inseparably bound
together. And these consequences we also see expressed in Giordano
Bruno. For him matter, which bore all forms in itself, was itself
something divine; it was determined as the very universal force or power
of God, which is one with God's producing activity and the expression of
the necessity of God's essence. The many alternating forms that matter,
filled with force, unfolds from within point to one universal form: the
world-soul, which is the source of the forms and as such in unity with
matter. This active unity of matter and form is designated as
\emph{nature}. As principle it is an eternal being; in its unfolding it
is divine and infinite, everywhere ensouled. And this natural, living
world-all, which is one with the world-soul (cf.\ the doctrine of the
Stoics), is posited as the goal and limit of all natural philosophy; as
the living, all-comprehending unity it is called: ``God in things.''
\opage{15}

In these determinations the standpoint of a consistent pantheism is
reached: what Greek philosophy strove for in its culminating period and
in Stoicism is expressed, and the opposition in principle to the
Christian is established; and with these determinations in Bruno the
statements of the most significant representatives of the period agree.
But at the moment when the goal seems reached, an after-effect of the
medieval fundamental view, against which the struggle was waged, asserts
itself: there again presses forward a thought of a God separated from
the universe, we find again the notion of creation, we see a
supernatural separated from a natural, a purely spiritual from a
corporeal, we see again thought and matter in opposition to one another,
and these determinations we find joined to the former only in an
external manner.

The separation between God and the universe comes forward in Bruno in
this form: that, after having determined the universe as infinite, and
after having secured this infinity over against the finite that it
comprehends, by positing the finite as parts in the infinite, not as
parts of it, he separates that which is, as it were, infinite throughout
(absolutely infinite) from that which is infinite only as a totality.
The former is God, the latter the universe. God is absolutely infinite
because he is whole in the whole world and is in each of its parts in an
infinite and complete manner. The world is infinite as a totality
because it is without limit and spatial determinateness, but it is not
infinite throughout, because every part in it is finite and because in
its parts the opposites are not unconditionally united. God, who is
elsewhere designated by Bruno as natura naturans, as being in things, now
becomes a deus creans and is designated as the supernatural principle,
as the supersubstantial substance; in being determined as the highest
monad (individuality), as monad in the unconditional sense (monas
monadum), as the \emph{spiritual} abiding unity without becoming, he
receives the stamp of \emph{transcendent} personality. But this
distinction between God and the universe, which from the beginning rests
only \opage{16}on a postulate, becomes impossible to carry through
alongside the earlier determinations of the essence of the universe and
of its infinity, present at every point, and its indivisibility.

In that first separation, which, as it were, sets the infinite being at
variance with itself, all the other separations lie. The separation
between God as something supernatural and the existence of the world as
the natural has already been touched on. The notion of creation comes in
as nature is seen as the effect of the infinite transcendent cause; and
even if the production is posited as necessary, and thus indirectly
designated as an unfolding of essence that includes the unity of
essence, the expression creation nevertheless forces itself upon us
precisely because the cause is separated from the universe. And since
God, in the separation from the natural, the infinite body, essentially
becomes a spiritual unity, then, despite the endeavor to place matter in
him as a moment of his essence, there is nevertheless placed in God only
a spiritual source of it, a ``ratio'' of matter; and in the universe,
where the ether was to be at once the universal, all-pervading matter and
the living, acting force, the corporeal and the spiritual nevertheless
involuntarily part from one another. The determinations of the corporeal
remain unexplained on the basis of the monads, which, as images of the
essence of the divine, become spiritual forces, and between the soul and
the body a uniting intermediate link must therefore be inserted: the
ethereal vital spirit.

The influence from the past that we have shown in the conception of the
concept of the universe and of the relation of the spiritual and the
corporeal asserts itself likewise in the concrete conception of natural
things and in the theory of knowledge. It lay grounded in the
fundamental thoughts and fundamental direction of the period that,
consistently, it had to strive for a knowledge of concrete natural
existence on the basis of physical principles, and that it had to grant
the highest significance to experience and to free self-consciousness,
independent of foreign authority. This striving and this granting we do
find, especially in the \opage{17}second division of the period, where
one is not content, through the determination of the relation of the
material to God's essence, to affirm indirectly and in general the
self-groundedness of the natural, but where one seeks the principles of
concrete natural existence in the warm and the cold, fire and water and
the like, and where one, while laying claim to the independence of
knowledge from authority, to freedom of thought, points to a point of
departure for this free knowledge in nature, to the testimony of the
senses and of experience, and demands that the path of knowledge be laid
out in such a way that it leads from these to rational knowledge
(cf.\ in particular Telesio and Bruno). But the after-effect of the past
not only gives the systems a theological stamp at the beginning of the
transitional period, but never lets this stamp disappear completely,
even if it becomes successively weaker under the influence of the
natural sciences; it co-operates with the after-effect of the Greek
metaphysical conception of nature in depriving the physically named
principles of their physical meaning, and in giving the imagination the
upper hand over the understanding in the conception of nature, which for
the mind is still only anthropomorphized nature; it limits free
knowledge by faith in authority (the philosophers in the first division
of the period, Telesio, Campanella), and it lets the mystical
contemplation of the supernatural divine become the highest. This last
appears especially clearly in Bruno's worked-out theory of knowledge. He
shows the stages by which human knowledge raises itself from the
sensuous images of the ideas to the ideas, and from these to the
spiritual unity that is the ground of life of the All. Here the
necessity of successive progress is stressed; the finite, which is the
effect of the Godhead, is posited as that which is indispensable to the
knowledge of the cause; through nature thought is to find God. But all
at once this point of view is pushed aside, and even if a one-sided
theory of immanence and a one-sidedly determined concept of God assert
themselves in this, there is at the same time at work, as appears in
the execution, the after-effect of the medieval, supranatural
\opage{18}standpoint of transcendence. As God's indwelling in us is
accentuated in its infinity, it becomes the one decisive thing, over
against which finite conditionedness loses its significance: the degrees
of knowledge vanish in the infinite. By virtue of God's indwelling we can
see all in one intuition, with immediate evidence, in the manner of
intuition. In these determinations we still have expressions of
pantheism, even if its peculiar form is determined by influences from a
past whose relation to nature was other than that of the transitional
period. But the influence of the Middle Ages then betrays itself, as
% printed "Overgangperiodens" (defect logged in the transcription)
contemplation is expressly posited as a supernatural act, which is to be
``impossible and non-existent for whoever does not \emph{believe}.'' Just
as, according to the doctrine of Scholasticism, the supernatural
elevation of the human faculties is indeed posited as conditioned by,
but yet not as truly standing in teleological connection with, natural
development, so also this contemplation is to be prepared in us by the
development of the life of thought, but not to issue from it; it is to
come as a sudden effect of God, the ``universal intelligence'', which,
being determined as acting supernaturally, is involuntarily transformed
into the ``supranatural principle'', into the transcendent Godhead, in
the beholding of whom blessedness consists. In this beholding the soul
is in unity with God, and the soul's transformation into God is
completed in that \gk{μανία} of love, the ``heroic frenzy'', in which
knowledge and action meet, and the final goal of thought and of will, of
knowledge and of ethical life, is in unity.

\divrule

Since the systems in the transitional period, even where they lay claim
to independence and autonomy, essentially only reproduce the content of
thought and the manner of thinking of antiquity, and since they have
taken up into themselves opposed elements which they have not been able
to work together into organic unity, they cannot be regarded as
representing an independent fundamental form of \opage{19}philosophy.
But they acquire their great significance through the new spirit that
works its way forward through them to clarity, through what they have
prepared and striven toward in the opposition to the Middle Ages,
through the problems they have sharpened, through the material of
thought they have gathered. Through the heterogeneity of their elements,
through the unovercome dualism, and through the excessiveness of their
demand on human knowledge, there is called forth, as touched on above, a
skepticism, which first shows itself as a ferment in the most developed
systems, but at the close of the period also appears in isolation and
carried through, and which betrays its source by taking its motives
partly from the disagreement among human opinions, partly from the
discord that dualism brings into human nature, between senses and
understanding, body and spirit, and which is supposed to make science
impossible, partly from the impossibility of realizing the infinity of
knowledge, from which the impossibility of a knowledge of the
individual, which is a member of the totality, is derived. The
enthusiastic faith in the power of human knowledge with which the period
began is shaken by the very occupation with nature, which opens the eye
to man's natural conditions and natural limits, and gradually gives way
to a skeptical sober-mindedness, which leads to the consciousness of the
necessity of securing philosophy by a sure method and by taking the
point of departure from something firm and unshakable. With the demand
for this, modern philosophy proper begins.

\dblrule

\opage{20}
\begin{center}
  \textbf{\large Second Section.}
\end{center}

\markboth{Second Section}{Second Section}
\begin{center}
  \textbf{The Philosophy of Modern Times in Its\\
  Independent Formation.}
\end{center}

\divrule

\begin{center}
  \textbf{Chapter One.}
\end{center}

\begin{center}
  {\large The General Character of Modern Philosophy.\\
  Its Form of Development.}
\end{center}

\divrule

The relation of independently developed modern philosophy to the thought
of the transitional period has been indicated in the foregoing. Common
to both is the direction toward nature and toward self-consciousness.
Peculiar to the philosophy of modern times is not only the altered
relation to the thought of antiquity, of which it is, and knows itself
to be, independent, but also the position and treatment of its task.
Modern philosophy no longer has the enthusiastic but blind confidence in
the power of the spirit that was characteristic of the poetic thinkers
of the transitional period; it begins with a more sober consciousness of
the necessity of securing philosophical knowledge by a preceding doubt,
a knowledge whose task becomes the methodical unfolding of a coherent
and uniform content of thought, drawn from the reality of spirit and of
nature. For both the realistic and the idealistic form of modern
philosophy doubt becomes the point of departure, either, as in Bacon,
conceived as the liberation of the mind from prejudices in order to be
able to take up the object of experience purely and without distortion;
or, as in Descartes, conceived as thought's separating out of everything
that is not thought, in order to find itself through this separation.
And just as doubt becomes the common point of departure, the
consciousness of the freedom and independence of thought becomes the
\opage{21}common presupposition, and for both the idealistic and the
realistic main form of modern philosophy nature becomes, even if in
idealism thought begins by enclosing itself within itself, an essential
object of knowledge. Thought, which seeks itself in the \emph{reality}
of self-consciousness, is drawn irresistibly toward the objective real,
which asserts for itself an independent and essential significance, and
in the immediate relation to this real thought asserts its freedom from
authority, develops its forms and becomes, more clearly or more
obscurely, conscious of being unable to tear itself loose from its
natural ground.

The external circumstances under which modern philosophy developed were
favorable to its independence. Protestantism, which in the Peace of
Westphalia won the recognition of equal rights, broke the supremacy of
the papacy. The spiritual power of the hierarchy, which had long been in
decline, was decisively broken by the development of the political unity
of the state, by the development of the natural sciences and by the
coming-to-be of the vernacular literatures. The place of the Church as
the center of spiritual life was taken by a learned estate, which in the
17th and in part still in the 18th century had in Latin a universal
means of communication. The new spiritual movement was led not so much
by the universities, which had followed the current of the new age only
in part, as by the academies and in particular by independent men of
science. Philosophy was only rarely molested by the power of the state,
and then always found sympathy and always a refuge. It gradually grew
into a power that in the 18th century could attack the established order
in every domain and contribute in high degree to the French Revolution;
and even if since that time it has essentially stepped back into its old
position as calmly inquiring science, it nevertheless indemnifies itself
in this more modest but truer position by entering into a relation of
reciprocal action with all the sciences, and by having laid down in the
forms of religious, political and social life the germ \opage{22}of a
development in accord with its own essence, the essence of free thought.
---

Modern philosophy develops, in accordance with the opposition between
nature and spirit that was its presupposition (cf.\ Part 1, p.~8),
through a \emph{double} series of systems: a realistic and an
idealistic. From a double point of departure: from nature and from
spirit, it strives toward unity, and by the double point of departure
the character and essential object of knowledge and the end point and
result of the series are differently determined. The tendency that
proceeds from nature gets nature as its essential object, experience as
the form of knowledge, and leads at last to letting spirit vanish in
matter as a single form of existence of matter, and to positing
spiritless matter as the unity of existence. The tendency that proceeds
from spirit, from self-consciousness, posits, even if, in accordance
with the tendency of modern times, it is involuntarily drawn toward
nature, spirit as its highest and most essential object; it gets a
priori thinking as the form of knowledge, and its final goal becomes a
panlogism, which lets all actual being be transformed into thought-being,
which lets nature and the material become a vanishing form of being
within the ideal process of spirit, the eternally actual, sole truth,
and thus posits the natureless spiritual as the unity of existence.

More precisely, the two lines of development articulate themselves in
the following way. The \emph{realistic} series develops from its first
form as \emph{philosophy of experience} (Bacon) to an \emph{empiricism}
that denies a priori knowledge (Locke), and from this point it divides
into two branches, of which the one, through \emph{formal} (empirical)
\emph{idealism} (Berkeley), runs out into \emph{skepticism} (Hume), the
other, through \emph{sensualism} (Condillac), into \emph{materialism}
(Système de la nature). The \emph{idealistic} series divides into a
dogmatic and a critical division. It begins as \emph{dogmatic} idealism,
in that thought, trusting in the power of ``clear and distinct''
knowledge, \opage{23}without making the possibility of an objective
knowledge a problem for itself, turns toward the object in order to take
possession of it. Dogmatic idealism develops through the systems of
Descartes, Spinoza and Leibnitz; it is continued spiritlessly and
unproductively in Wolf and languishes in the 18th century. Kant, by
taking the problem of knowledge as his point of departure, founds
\emph{critical} idealism, within which there again separate out a form
critical in the narrower sense, with a direction toward the speculative
(Kant, Fichte), and a speculative form, which nevertheless bears within
itself, even if as hidden, the element that is conditioned by the
one-sidedness of the determination of the principle and is critical in
the one-sided sense (Schelling, Hegel). Against one-sidednesses in the
development of critical idealism there appears a reaction from the
philosophy of faith and from Herbartian ``realism.''

\divrule

\begin{center}
  \textbf{Chapter Two.}
\end{center}

\divrule

\begin{center}
  {\large The Realistic Line of Development.}
\end{center}

\divrule

\begin{center}
  \textbf{1. Bacon.}
\end{center}

\emph{Frants Bacon} was born in London in 1561. His father
was Lord Keeper of the Great Seal under Queen Elizabeth. In his twelfth year
he went to the University of Cambridge, where continued studies
soon brought him to see the deficiencies of the scholastic
philosophy. From the university he went with the English ambassador to Paris, and during a two-year stay there he wrote,
at nineteen years of age, his work de statu Europæ. His father's
death called him back to England, and since, as a younger
son, he was without fortune, he applied himself to jurisprudence,
soon acquired a respected name as a jurist, and was
\opage{24}consulted by the Crown on important questions. Alongside his
juridical studies, however, he continued to cultivate philosophy,
and already at that time he laid the plan for a reform of the sciences. The Earl of Essex favored him, but after the Earl's fall Bacon appeared as the Crown's prosecutor against him.
From 1595 Bacon had been a member of Parliament, but
only after Elizabeth, who harbored mistrust of his character, had died did he receive a state office. By attaching himself
to James the First's minister Buckingham he rose quickly
to the highest dignities. In 1617 he became Lord Keeper of the Great Seal,
in 1619 Lord Chancellor and Baron of Verulam, in 1620 Viscount of St.
Albans. But already in 1621 he was impeached before Parliament
for bribery, and sentenced to lose his offices, to a
fine of £4000 and to imprisonment; he was, however, pardoned
the two latter penalties, and after James's death and Buckingham's
fall his sentence was quashed and he was released from all consequences
of the judgment, though more, no doubt, because he was regarded as
sacrificed by the generally hated Buckingham than because
he was held to be innocent. After that time he lived
as a private man for the sciences alone, and died in 1626.

Bacon's two main works are the treatise de augmentis et dignitate scientiarum (published in a first draft in 1605, complete in
1623) and Novum Organon (1620). They form the first two members
of a comprehensive work which, as an Instauratio
magna of the sciences, was first to give a survey of the whole
articulation of science, next to develop the new scientific method,
then to collect the material of the sciences, especially of natural science: the material of induction, thereafter to work it up according to the method of induction, so that the laws of nature and the explanations of the phenomena
would be found, and finally to give directions for the practical application of the knowledge gained (scientia activa). Bacon's Silva silvarum sive historia naturalis, published after
his death, contains an extensive collection of material; to the last parts of
the work Bacon gave only few and imperfect contributions.

\divrule

\opage{25}Bacon founds the realistic tendency of modern philosophy, in that, recognizing the necessity of a complete reform of philosophy, he determines with clear consciousness
nature, which in the transitional period had drawn the spirit to itself as if by magic,
as the essential object of thought, and natural science grounded on
experience as the mother of all the sciences.
The philosophy of the past he judges on the whole as formalistic
and as unfruitful in a practical respect. A genuine natural philosophy has in his view not existed at all: it has
been corrupted either by logic (metaphysics), as in Aristotle, or by natural theology, as in Plato, or
by mathematics, as in the Neoplatonists. The whole condition of the
sciences is bad; their decline is essentially grounded in
their having torn themselves loose from experience and nature; more specific causes
% errata list B, p. 25: printed "medvirkende", corrected "speciellere"
Bacon seeks in the spiritualistic haughtiness of the past, in its superstition and blind religious zeal, in its one-sided occupation with theology, morals and politics, in its
belief in authority and its blind reverence toward antiquity,
and finally in its despondency before the great difficulties.
In order that science may be \emph{renewed} and become \emph{fruitful} for
man's dominion over things, there is required a return to nature and experience, and a new method by
which nature can be apprehended undistorted according to its real
essence, since only such an apprehension can in truth become
fruitful. ``The world is not, as has hitherto happened, to be forced
into the narrow sphere of the human understanding, but the latter
is to be stretched and enlarged so as to take in the image of the world
such as it is''. Therefore man must first free his
mind from all ``idols'', i.e.\ from all the deceptive images (notions) that do not express the nature of the objective but spring from something subjective: from human inclinations or
prejudices. Bacon distinguishes between ``idola tribus'', which are
grounded in human nature in general, e.g.\ in
its inclination to anthropomorphize, to want to find
order and likeness everywhere, to overlook the negative instances,
\opage{26}to one-sided abstraction, and the like; ``idola specus'', which have their
origin in individuality and upbringing; ``idola fori'',
which are called forth by language, and ``idola theatri'', which derive
from what has been impressed by tradition, from preconceived opinions
and accepted methods. When the mind is cleansed of these idols,
it is fit to take in its object undistorted. This
then happens through the application of the true method of knowledge:
\emph{the method of planned experience}. It takes its point of departure,
in order to gain a secure foundation, in the real particulars of nature,
and from these, after they have been methodically
ordered, one ascends step by step, with critical observation of the positive and negative instances, with emphasis on the decisive cases and with the use of experiment and analogy, by true induction up to the knowledge of the form, the uniting law that works in the material, and through
the graded series of forms one strives toward the unity of nature
(forma magna) as the theoretical goal. By this method,
which is the novum organon of philosophy, a true interpretation of nature is won (an interpretatio naturæ as opposed to an
antecipatio naturæ), and by it is effected the ``great
renewal'' of the sciences.

Bacon sets his method of induction against the syllogistic
method, which was that of Scholasticism. The syllogism rests on the
concepts contained in its premises; but since these
concepts can be determined reliably and exactly only through induction, the syllogism points back to the latter. The
earlier philosophy did indeed know a kind of induction, but
not the true and methodical one. One wanted, as it were in a rush,
to reach the most general propositions so as then to
arrive through syllogisms at the intermediate propositions and the particular propositions.
True induction, on the other hand, ascends quite successively from
the sensible and particular; it divides and separates nature
by the proper exclusions and rejections, and infers to
the affirmative determinations of a thing only when it has collected and
examined a sufficient number of negative instances and
\opage{27}excluded from the object all the determinations that do not
essentially belong to it. In order that the investigation shall not become
endless, a contractio inquisitionis is required, which takes place by
bringing forward the essential instances, the instances that have a
prerogative (thus Bacon cites the ratio of the weight of quicksilver to
that of gold as decisive for determining specific
gravity as dependent not on cohesion but on mass).
As one ascends successively to ever more general concepts, these are ``non notionalia, sed bene terminata, et
talia, quæ natura ut revera sibi notiora agnoscat, quæque
rebus hæreant in medullis'' [not notional, but well defined, and such as nature would acknowledge as truly better known to herself, and such as cleave to the very marrow of things]. The concept, the ``form'', expresses
the active nature of the thing, and includes, as what runs through the particular, the law of the particular. To the highest
form of nature (its forma magna) the way goes through the knowledge
of the ``simple'' (simplices) forms, which are the foundation for the apprehension of nature in its unity. They must not be apprehended as
separated from matter, but they must be found again in the heterogeneous parts of matter.

On the knowledge of the form rests the practical significance of science. He who has known the form has power
over nature and is in a manner omniscient. He can ``superinduce'' the form on all kinds of objects and perform
all the wonders that the alchemists claim to be able
to perform. ``Scientia et potentia in unum coincidunt'' [Knowledge and power coincide in one].
Science in its true form, as the image of being, founds man's dominion over the nature ``that is
conquered only by being obeyed'', and it serves man's benefit.
The theoretical and the practical are inseparably united
in it.

\divrule

In his \emph{encyclopedic ordering of the sciences}, which,
while it delimits them from one another, holds fast to their
organic connection, Bacon takes as the ground of division for ``doctrina humana'' the powers of the soul that ``depict the world'': \opage{28}memory,
imagination and reason. According to their difference
the ``globus intellectualis'' is then divided into history, poetry and philosophy. \emph{History}, which divides into historia naturalis and historia civilis (under which the history of art and of science is included), has for its
object individuals, not species. It describes facts, while philosophy comprehends causes. \emph{Poetry} likewise
has to do with individuals, but with feigned ones. It
divides into the narrative (epic), the didactic (parabolic)
and the dramatic. Lyric poetry is lacking in the division;
it is referred by Bacon to philosophy and rhetoric. \emph{Philosophy} occupies itself with general concepts,
their connection and separation, and in this is guided by the laws of nature and the evidence of the matter. It seeks, in contrast
to revealed theology, the natural causes, and divides
according to its object into natural theology, natural philosophy and
the doctrine of man. \emph{Revealed theology}, to
which belong the positive doctrine of God and the doctrine of a spiritual element
in man (as distinct from the natural soul), Bacon regards not as an object for the knowledge of reason
but for faith, and in general he keeps religion and philosophy strictly separated, chiefly in the interest of philosophy, so as
to protect it against encroachments on the part of the Church. His utterances
about the objects of faith often sound very ironic; thus when
he says: quanto igitur mysterium aliquod divinum fuerit
magis absonum et incredibile, tanto plus in credendo exhibetur
honoris deo, et fit victoria fidei nobilior [the more discordant and incredible, therefore, a divine mystery is, the more honor is shown to God in believing it, and the nobler the victory of faith becomes]. In his polemic against
fanaticism and superstition he sometimes strikes positive
religion as such. --- \emph{Natural theology} acquires in
Bacon a somewhat doubtful significance. It is not to go so far as to
state anything positive about the essence of God; it is only to be able to serve
to refute atheism --- inasmuch as the explanation by natural
causes needs to be supplemented by recourse to divine providence, --- not to ground faith. Nature reveals only God's power; it is not his image but his
work; there is no analogy between it and him. --- \emph{Natural \opage{29}\setcounter{footnote}{0}philosophy} becomes the real main science and the foundation of all
\emph{knowledge}. As \emph{theoretical} (speculative) natural philosophy it comprises \emph{metaphysics} and \emph{physics}
(whose auxiliary science mathematics\footnote{As the doctrine of abstract quantity, which is one of the forms of nature,
mathematics belongs under metaphysics.} is), as \emph{practical}
(operative) \emph{magic} and \emph{mechanics}, and what Bacon in his
encyclopedic survey places before the particular philosophical
sciences as a \emph{philosophia prima}, as ``the mother of all sciences'', in reality falls under it. In the Novum
Organon it is natural philosophy that receives the predicate of mother of all
sciences; a philosophia
prima is hardly mentioned any more, and where Bacon in this work, in speaking of natural philosophy, talks of analogies in the sciences, he uses
the same examples with which he had earlier illustrated the thought of
% printed "hvike" (defect logged in the transcription)
a ``first philosophy''. Its objects were to be partly
the axioms that are common to all sciences, partly such conceptual determinations as ens and non ens, possibile and
impossibile, simile and diversum, which Bacon designates as
relativæ et adventitiæ entium conditiones. He demands that
the treatment of these concepts take place secundum naturæ,
non sermonis leges, by pointing out the relations in nature. ---
\emph{Metaphysics} Bacon strives to treat physically, whereas before physics was treated metaphysically, and he in reality reduces it to physics. He does indeed distinguish it
from physics in this, that the latter is to consider the form as working in this determinate, particular matter, which is changeable,
metaphysics the general law that is independent of the
particular matter, but not of matter in general. Physics
is thus, e.g., to consider the whiteness of snow, metaphysics
the conditions of whiteness in bodies in general. But
this consideration of the general laws, which is to be
the essential part of metaphysics, is not easily distinguished from
the ``abstract doctrine of bodies'', which is a part of \opage{30}physics. Another ground of distinction between metaphysics and
physics Bacon places in this, that metaphysics investigates ends, while physics considers efficient causes,
whose territory it cannot leave without harm to itself. (Through
the doctrine of final causes metaphysics points over to
natural theology). Bacon's utterances about these causes are,
however, essentially an accommodation to the opinions of the time,
and in any case they leave nature to have its activity
for itself beside the guiding activity of the Godhead. His
real interest is to remove final causes, which he
apprehends only from the standpoint of an external teleology, from physics,
not to secure for them a positive significance through metaphysics. They are, he says, ``drawn more from the nature of man
than from that of the universe, and they have greatly corrupted natural philosophy''. The investigation of them is ``barren, without the fruit of life, like a nun consecrated to God''. Nature is to be apprehended
physically, by itself, in its concrete, material reality.
This becomes the task of \emph{physics}. It divides into the doctrine of
communia principia rerum, into the doctrine of the structure of the world and
into the doctrine of the various bodies, which again is either the
concrete doctrine of bodies (the doctrine of plants, animals and the like), or the
abstract doctrine of bodies (the doctrine of the general physical properties, e.g.\ heat, light etc.). --- Of the parts of \emph{practical} natural philosophy, \emph{mechanics} is the application of physics,
\emph{natural magic} the application of metaphysics. Magic is the science or art that from the knowledge of the
hidden forms derives astonishing effects, and, by bringing
the active and passive elements together, brings the power of nature
to light. Bacon is, with respect to it, still rather caught
in the opinions of his time, partly on account of his deficient acquaintance with the natural sciences. Thus he still believes in the possibility of making gold, and he does not unconditionally reject the older fantastic magic either. --- The philosophical \emph{doctrine
of man} has man for its object partly as a single
individual (philosophia humanitatis), partly as a member of society
\opage{31}(ph.\ civilis). The doctrine of the individual considers partly the soul, partly
the body. \emph{Psychology} has to do not with the supernaturally inspired soul, which is the object of theology, but
only with the sensible soul, which is a corporeal substance
that has its place in the brain, and is invisible only on account of
its fineness. The properties of this soul are voluntary motion,
and sensation. The sensations of the soul are conscious; unconscious sensations (perceptions) are found in all bodies,
and express themselves in attraction and repulsion and the like. --- The sciences of the use of the soul's faculties are logic and ethics.
\emph{Logic} is the doctrine of knowledge directed toward truth; \emph{ethics} is a direction for the will to act rightly.
The good is what benefits man, the single individuals
and mankind. Action for the common benefit is the highest
duty; the common benefit is in all cases of collision the final
deciding instance. While ethics concerns bonitas interna,
\emph{philosophia civilis} concerns bonitas externa in intercourse with
men, in business and in the commonwealth. What virtue is
for ethics, prudence is for politics. Statecraft proper, which teaches how to guide the state or the multitude toward ends of common benefit, Bacon will not carry out, but only indicates its
place. He wants it to be treated by statesmen, not by
theoretical philosophers or jurists.

In giving his encyclopedic survey of the sciences, Bacon indicates at many points sciences that
were still lacking. Among these are, besides a philosophia
prima, a metaphysics and a magia naturalis in his sense, a history of literature, a history of philosophy, a physical astronomy, a comparative anatomy, a philosophical
grammar and others.

\divrule

If in Bacon, in whom, given his whole historical
position, the problem of the method of knowledge comes forward
ahead of the problem of the principle of existence, we would seek, through the scattered
\opage{32}and aphoristic statements found in his writings,
his view of the ground of existence, we shall find a striving
for unity that nevertheless does not overcome dualism. If we would
find Bacon's true opinion concerning the ground of existence,
we must not take his utterances about what he makes the object of
revealed theology according to the letter. These
utterances are only accommodating, only calculated to protect his philosophizing against theological encroachments. As
revelation and the supernatural are separated off from philosophy,
philosophy avoids, within its own domain, the dualism between the
supernatural and the natural, and, in asserting the
latter for itself, it makes the former into something non-existent for philosophy, into something that it factually negates. But in nature matter becomes the fundamental determination. The ``first matter'' is an immediately positive fact that cannot be derived from anything else: the series of causes cannot go on to
infinity. But in order to have a real existence the
first matter must originally have form and motion. It produces everything and dissolves everything into itself; its sum total remains
constantly the same, but is distributed in different ways.
Matter's resistance to annihilation is not a slight force,
but the most active of all forces. Everything corporeal has a certain
faculty of observation or of notion and a certain faculty of choice.
Between bodies with senses and those without senses there
is an agreement in properties.

Through these determinations of matter the concept is striven for of a unity of the spiritual and the corporeal, which are never separated in nature, the sole
object of philosophy, and in this concept is
sought an original source of the activity by which
the multiplicity of natural things and natural forms unfolds.
But the opposition, which was apparently canceled in the apprehension of the concept of matter, comes forward again within that concept.
In matter there takes place as it were a division into a twofold moment, i.e.\ the forming, moving principle, and that which is the stuff (substrate) of this
active principle. The former is designated with an expression from
\opage{33}Paracelsus as spiritus innatus, so that the name points
to the old opposition between the spiritual as the
active and the material as the passive; the latter is designated as
tangibile: the palpable, which is the material of the activity.
Spiritus is indeed in turn determined as material, as a spiritus
aëreus aut igneus, so that the distinction is thus canceled in the name,
but only so as to let an unexplained dualistic distinction of essence
be posited within the material. From this material the essential determination of spirit, the end, is then excluded, and as
the concept thus becomes narrower than the compass of reality,
forms of being remain over, with respect to which Bacon at least does not show any possibility that they could be brought
under it, e.g.\ understanding and will. There is an
% errata list B, p. 33: printed "under hiin Bestemmelse", corrected "under det"
ambiguity in Bacon where he speaks of these. He distinguishes, as we saw, in man between the rational and the
sensible soul. The powers of the latter were voluntary (spontaneous)
motion, and sensation: powers that are sought everywhere in the
material. Will and understanding are not posited as its powers,
but neither is it said that they are the powers of the rational soul;
rather, with an indeterminate expression, that they are those
``of the human soul''. Bacon's real thought can, since the
supernatural has no reality for his philosophy, only be
this: to see understanding and will, which are given to consciousness
as facts, as determinations of the sensible soul;
but he does not express himself definitely and does not show the essential connection. Understanding and will do indeed have an analogy
with sensation and voluntary motion, and with what is
the form of these in apparently lifeless things: perception and
the force of attraction and repulsion; but just as little as Bacon
proved the possibility that the material could have sensation and voluntariness, just so little does he carry through the analogy of the
mental functions with the manifestations of force of matter,
let alone develop it into identity of essence; and, despite the predicates he ascribes to the material, there thus remains a dualistic
distinction.

\divrule

\opage{34}Bacon has not given, and has not wished to give, a finished philosophical system. He wishes to put into the hands of posterity the means by which to work, with united forces, toward such a system, and this means he seeks in his method, while by his encyclopedic outline of the sciences he indicates the domain that is to be drawn into research, and the order in which research is to proceed.

Bacon's historical significance lies in his method and in the consciousness with which it is applied. Rejecting in principle the Middle Ages' manner of thinking, he has indicated the basic features of the inductive method in a form freed from its being bound to the single domain of natural science, and by this method he has put a true physical conception of nature in the place of the metaphysical, theosophical or theological conception of the past; and he has done so with consciousness of the task of the new age and of the significance of knowledge of nature. He is not the first to have pointed to experience and induction as the way to a true knowledge of nature. Several of the philosophers of the transitional period, whom he knew and made use of, but whom as a rule he neither names nor acknowledges, had already done so, but without holding fast to what they had intimated; and Leonardo da Vinci, who was as admirable as a natural scientist as he was as an artist, clearly indicated the method of experience. But Bacon nevertheless becomes the real founder of the philosophy of experience; for in him the demand for the method is united with the consciousness of its connection with the whole conception of existence and of the aim of knowledge. The philosophy of experience is founded not merely by the method, but by this: that the nature to which the method is applied is posited as something essential, as something true that exists objectively, as something valid in itself in its concretion, and as a source of knowledge and truth, inasmuch as in its peculiarity it is apprehended only through itself, by way of the receptive relation of the mind, and in its being grounded in law it rejects the arbitrariness of the imagination and the encroachment of immediate metaphysical \opage{35}\setcounter{footnote}{0}knowledge, so that the mind first gains mastery over it through the recognition and understanding of its conformity to law, and only when it has thus been recognized and understood can win from it means for the development of human life. Bacon's empirical inquiry aims not at a collection of particulars, but at knowledge of the general laws of nature, and that as real, immanent laws of nature, as the expression of the natural things' own active essence. And in accordance with the thought of the inner connection in the laws and forms of nature, which are to be led back to a highest unity, he strives to see the sciences in their inner connection, as an articulated unity, grounded on natural science and emancipated from all theology.

The defects of his philosophy lie, apart from the remnant of dualism pointed out above, partly in his imperfect consciousness of the conditions and presuppositions of induction, of the synthetic and a priori element in it, and of its relation to deduction; partly in his insufficient carrying-through and imperfect application of the methodological principles; partly in a frequently striking want not merely of mathematical insight but of real knowledge of nature\footnote{Of many of the most important discoveries of his age, e.g.\ those of a Galilei, a Harvey, a Kepler, Bacon seems ignorant; others, such as Gilbert's, he does not understand. The theory of Copernicus he rejects as a fantastic whim.}, and in the frequent relapse, conditioned by this, into a fantastic conception of nature;
% printed "en Galilei's en Harvey's" (defect logged in the transcription)
partly, finally, in the one-sided conception and application of the principle of experience. Inasmuch as Bacon posits experience as the source of true knowledge, no restriction of it to a determinate domain of what is has thereby originally been posited. Bacon originally wishes, in agreement with his striving to lead what is back to a unity, to let experience have \emph{all} that is for its object: what is designated as the spiritual just as fully as what is called the corporeal; \opage{36}\setcounter{footnote}{0}he thus wished, like Kant later, to be able to make the human mind too, in its self-activity, together with what originally proceeds from this, an object of observation. But this striving is checked at once: experience immediately takes its direction toward the physical in the narrower sense, toward the external, the sensuous-natural, and turns toward the mind only insofar as it is considered from its natural side. The mind is considered only as ``sensuous soul'', or, insofar as understanding and will, with the ambiguity touched on above, are posited as faculties of the ``human soul'' and account is taken of them, this nevertheless happens only through a describing of their phenomenal forms of manifestation, and neither their fundamental forms, their source nor their laws are investigated. And as the direction of knowledge
% printed "Erkjendelseens" (defect logged in the transcription)
thus goes essentially outward toward the given natural phenomenon, and the task becomes the \emph{objective apprehension} of it, i.e.\ an apprehension \emph{by the object itself}\footnote{Ea demum est vera philosophia, quæ mundi ipsius voces quam fidelissime reddit et veluti dictante mundo conscripta est, nec quidquam de proprio addit, sed tantum iterat et resonat. [That alone is true philosophy which renders the voices of the world itself most faithfully and is written as it were at the world's dictation, and adds nothing of its own, but only repeats and echoes.]}, the determination of activity that Bacon had directly and indirectly (by ascribing to the mind the faculty to invent and to judge, to order the material of experience,
% printed "til ordne" (defect logged in the transcription)
to single out the decisive instances, to find analogies, etc.) ascribed to the mind recedes in him, and the mind's relation to the object of knowledge, the experienced, unconsciously assumes the character of a merely receptive one; which is also why Bacon can express the hope that his method will be able to level the difference between minds and make all equally fit to receive and know the truth. In this restriction of experience and in this conception of knowledge as merely receptive lies the germ of the development through which the philosophy of experience, for which experience was originally indeed an essential, but not the only, source of knowledge, transforms itself first into an empiricism that denies everything a prio\opage{37}ri, and then either into skepticism or into sensualism and materialism.

\divrule

\begin{center}
  \textbf{2. Hobbes.}
\end{center}

\emph{Thomas Hobbes} was born in 1588 at Malmesbury; he educated himself by philosophical studies at Oxford, and by travels in France and Italy, on which he came into closer contact with Gassendi, Mersenne and Galilei, and acquired a knowledge of mathematics and the natural sciences that enabled him, better than Bacon, with whom he stood in a personal relation, to assess the great discoveries of his time in the domain of the natural sciences, and the significance of mathematics for the philosophy of experience. After his return home he occupied himself much with politics, left his fatherland with the Stuarts and stayed again for a time in Paris, where he became acquainted with Descartes; but he returned to England when Cromwell had acquired a power that accorded with Hobbes's views on the significance of the head of state. His writings on the state (de cive, 1642, and Leviathan, sive de materia, forma et potestate civitatis ecclesiasticæ et civilis, 1651) procured him many enemies, especially among the clergy, whose persecutions compelled him once more to leave England for a shorter time. He died in 1679. His main work is his Elementa philosophiæ, 1655 ff., in which Part 1 treats: de corpore; Part 2: de homine; Part 3: de cive (the earlier writing is incorporated here). In addition his writing de libertate, necessitate et casu, 1656, may be mentioned.

\divrule

In Hobbes far-reaching consequences of the one-sidedly conceived principle of experience already appear, inasmuch as knowledge is determined by him as springing from the sensuous sensation called forth by the pressure of external objects on the organs, and the corporeal becomes the only object of philosophy, \opage{38}since its concept is extended to \emph{all} that exists. Philosophy is the comprehension, attained by correct thinking or correct inferences, of effects and phenomena through their causes, and of possible causes through the observed phenomena or effects. All its concepts are concepts of something finite; the infinite lies outside its domain. Its aim is to bring utility and advantage to human life.

As the doctrine of bodies, philosophy divides, according to the two kinds of bodies, the natural and those joined together by the will, into philosophy of nature and philosophy of the state. The philosophy of the state comprises ethics and politics; the philosophy of nature comprises logic (the doctrine of names and words), ontology or phil. prima (which treats of the fundamental concepts: space, time, body and its accidents: magnitude and motion, etc.), the doctrine of relations (applied mathematics and geometry), and physics proper (the doctrine of natural phenomena). Ethics and politics must be grounded on physics, and are even at times regarded as parts of it. In the demonstrations of philosophy Hobbes demands a strict methodical order, such that proofs of everything that follows are successively derived from the first definitions. The order of the sciences becomes: logic (as the doctrine of correct thinking, without which no science can be had), ontology, mathematics, physics, morals, politics.

Hobbes's \emph{philosophy of nature} has a purely mechanical character. The corporeal (i.e.\ that which subsists independently of our notion, by itself, and with extension) is composed of elementary parts, which are, however, not absolutely indivisible, and it is passive; every motion comes from an external cause that is itself moved, so that motion presupposes motion to infinity. Every change in bodies consists in motion; the different determinatenesses of bodies are only differences in their motion. The qualities apprehended in sense-sensation, such as colors and sounds, are not as such in the objects, but only in the \emph{individual}, sentient subject as phantasmata; in the objects they are only \opage{39}\setcounter{footnote}{0}as kinds of motion, even though everything corporeal has in itself a disposition to sensation.

In the doctrine of knowledge our sensations are defined as \emph{changes} in the one who senses, and as changes they are motions in the internal organs\footnote{``Mens nihil aliud erit præterquam motus in partibus quibusdam corporis organici.'' [The mind will be nothing other than motion in certain parts of an organic body.]}. The cause of sensation is the external object. From the sense organs the motion goes inward, to the innermost parts: to the brain and from there to the heart, and from the heart there then proceeds a counter-pressure, a motion outward, whereby the sensation becomes sensation of the object as \emph{external}, and the image or sensuous notion is formed. When the object is removed, the sensation becomes imagination, which is a fainter sensation. The persistence of impressions conditions memory, as well as the ability to judge about the images by distinguishing and comparison, and in general to apprehend them in their determinateness, which could not happen if only one single image were before consciousness. Recollection is supported, and the communication of what is apprehended becomes possible, by the signs that we \emph{arbitrarily} attach to notions in naming. The word becomes a sign for many objects that resemble one another, and it acquires the character of the universal, while the actual thing designated by the word is always a single thing. In definition we make clear which objects we mean to designate by the word. By means of word-signs we get our notions into our power, and this Hobbes, although he otherwise distinguishes the thing from our impression of the thing, by a subreption lets be one with our knowing what is represented. Our thinking takes place only in words, and its truth consists in the correct apprehension of the relations of words to one another (metaphysics is therefore reduced to the art of language). Thinking is a reckoning, an adding and subtracting (combining and separating) of notions by means of signs. Whereas sensation lacks universal validity, rational thinking, which \opage{40}draws inferences by means of signs, acquires this character, and inasmuch as Hobbes by such inferences arrives at the concept of a substance that is the unchangeable substratum of accidents, and at the setting up of general propositions about the communication of motion and the general connection of necessity, he ascribes objective validity to the results of his inferences.

When our definitions, which ought to state the cause or the mode of origin, are correct, our philosophical thinking is demonstrative in the same way as mathematical thinking. A priori we can comprehend through demonstration only things whose origin depends on our will, e.g.\ magnitudes that are determined only by lines that we can draw. With natural things, whose causes are not in our power, we must demonstrate a posteriori; yet, since physics rests on mathematics, there is also much in it that can be demonstrated a priori.
% printed "a priori, Politiken" (defect logged in the transcription)
Politics and ethics can be demonstrated a priori, since we ourselves are the authors of the compacts and laws that first create right and wrong.

In ethical respect Hobbes conceives man only as the single sensuous individual. The will is only a desire and is dependent on the objects, which call forth an inward-going motion in the organs that facilitates or hinders the vital motion in the blood and is perceived as a feeling of pleasure or displeasure, by which feeling desire and aversion are aroused. The will is free only in the sense that we can carry out what we will, but what we will is dependent on the external. What we desire is good, what we abhor is evil; both are always relative. The greatest good is self-preservation, the greatest evil death.

On this conception is grounded also Hobbes's peculiar view of the \emph{state}, by which, despite all one-sidedness, he has the essential merit of having once again led the bond of society back to factors that lie in our nature and that we recognize as such. Bacon's demand that politics be based on physics, Hobbes realizes. The state, as a collection of isolated individuals, he regards as having arisen \opage{41}not through an instinctive social drive, but through voluntary compact. Prior to the state comes the \emph{state of nature}, which is a condition of unrestricted arbitrariness, in which the individuals, who are all equal and each of whom has a right to everything --- by virtue of the right to self-preservation and happiness --- stand hostile toward one another in their egoism, so that there reigns a bellum omnium contra omnes. But precisely for that reason no one can enjoy his right, and it is as if there were no right at all. The recognition of what is contradictory and unbearable in this condition, and \emph{fear} of being overpowered, therefore lead the individuals, in order to be able to live a secure and satisfying human life, to give up the state of nature by a compact and to form a society of right, which is protected by all the individuals, whose egoistic nature cannot be changed, transferring all their rights to a government that by its unrestricted power can maintain \emph{peace} in society. All wills are subjected to one will, that of a sole ruler or of an assembly, though most consistently that of a sole ruler, and to this the individuals transfer their right and power and, as it were, their personality. The multitude thus becomes one person: the state, the ``great Leviathan'', the ``mortal God''. ``The state is contained in the person of the ruler'', the ruler is the people, the people therefore always rules. The head of state, with its absolute power, stands above the laws and determines what is good and evil, right and wrong. Since all right has been transferred to it, it has no obligation toward the multitude. Since, however, the state was originally founded for the welfare of the citizens, the ruler ought to provide for this. It comprises not only peace and security and the physical enjoyment of life, but also goodwill among the citizens and the development of science and of intellectual life. The ruler, however, bears no responsibility: he can conceive the welfare of the state as he wishes: what he commands is by that very fact right. In his will the unconditional freedom of the state of nature is preserved, the right of the stronger perpetuated.

To what the head of state determines belongs religious faith. Religion is an affair of state. Religion and superstition \opage{42}have this in common, that they are a fear of invisible powers, invented or accepted according to tradition. But religion is fear of the powers recognized by the state authority, superstition fear of those not recognized. A private religious view is unjustified over against the faith sanctioned by the state, and becomes revolutionary. Hobbes therefore himself submits to the ecclesiastical laws of his country, although he stands remote from the positively religious, and in his philosophy really has no place for the concept of God, which in any case is kept purely negative. ``The mysteries of faith one must swallow whole, like healing but bitter pills; if one chews them, they are as a rule spat out''. ---

In Hobbes those consequences of the principle of experience are anticipated that come to development in the philosophy of the 18th century. But his influence on the time immediately following has been very limited, essentially, no doubt, precisely on account of the far-reaching character of his views; yet it may also have contributed that precisely in his theory of knowledge there is a lack of complete consistency, and that by his contemporaries he was regarded chiefly as a politician. It is not Hobbes, but John Locke, from whom the decisive influence on the further development of realistic philosophy proceeds.

\divrule

\begin{center}
  \textbf{3. Locke.}
\end{center}

\emph{John Locke} was born in 1632 at Wrington near Bristol. At Oxford he originally studied medicine and the natural sciences. He was frightened away from philosophy by Scholasticism, which still reigned at the university; later he was drawn toward it by the writings of Descartes; yet for a long time he still occupied himself chiefly with studies in natural science. Through his connection with the Earl of Shaftesbury Locke was led into public life, but had to share the changing fortunes of his patron and go into exile in Holland for 5 years, from where he returned with William of Orange. Under him \opage{43}Locke again entered the service of the state, from which, however, he withdrew in the last years of his life. He died in 1704.

For his main work: Essay, concerning human understanding, Locke, at the request of some friends, made the first draft as early as 1670; but the work, after being reworked repeatedly, was not published until 1689. Of his other writings may be mentioned the three letters on toleration (1685), and: the reasonableness of Christianity as delivered in the scriptures (1696), which drew upon him strong attacks from the clergy, which he answered with success, as well as the book on education (1693).

\divrule

Although Locke does not connect his investigations with Bacon's, and
on the whole makes only little use of him, as likewise of Hobbes, his
philosophy nevertheless stands in close connection with Bacon's and
issues from the tendency that was founded by him. Whereas Bacon had
posited experience as the true instrument of knowledge, Locke, in
raising the question of the origin of our knowledge, will show that it
can have no other source than experience, and he thus furnishes a proof
of Bacon's postulate. His empiricism points back to Bacon and Hobbes;
but the fact that he leads the investigation back to the psychological
betrays that he is influenced by idealism.

The task that Locke, prompted by a fruitless philosophical discussion
among some of his friends, sets himself in his main work is this: to
investigate the extent of our cognitive faculty, because only a clear
determining of the limits of our understanding, of the relation between
knowledge and opinion, of the degrees of certainty of our knowledge and
of the ground of the difference in degree, can keep us from the
encroachments that lead to uncertainty and skepticism. But for the
attainment of this goal there is, in his view, only one way: namely, to
make knowledge one's object in such a way that one seeks, not to gain a
metaphysical or as it were physiological insight into the essence of
the mind, but to find the source of \opage{44}knowledge, and the
different ways in which the understanding comes to its concepts of
objects. In this investigation Locke also has the practical aim, while
showing the necessity of faith, of freeing faith from superstition and
intolerance, and of determining the natural limits of civil and
religious laws.

In his investigation of the problem Locke makes his way to his
positive results through a polemic against the doctrine of dogmatic
idealism concerning innate ideas (ideæ innatæ). Since idealism took its
point of departure from the mind in its self-activity, it sought a
content belonging originally to the active mind, and this was conceived
as a sum of universal concepts and principles that were to be found in
the mind as a ready-made store. Thus the concept of ``innate ideas''
was determined empirically, and a polemic against it from the
standpoint of experience became justified. In wanting to maintain his
proposition, that there are no innate ideas in the sense indicated,
Locke does so through a critique of the proofs by which they were
supported. One appeals to the fact that there are certain theoretical
and practical principles that have validity for all. Among the
theoretical principles, accordingly, the principles of identity and of
contradiction are cited. But experience shows that this is by no means
the case: that these principles should be universally acknowledged.
Children and the uneducated have no notion of such abstract, universal
propositions; they understand the concrete: that the sour is not sweet,
but the abstract proposition does not exist for them. One then says:
they have it in the understanding but do not know of it; but this is a
flat contradiction, since being in the understanding is synonymous
with being thought. If one would then give the matter the turn that
men, as soon as they used their reason, came to agree on these
propositions, this would, on the one hand, say precisely that they did
\emph{not} have them from the beginning; on the other, it would be
incorrect, for we know much else before them: what comes first is not
the universal but the individual impressions.
\opage{45}Even if, however, those principles were really acknowledged
by all at once as true, this would not prove that they were innate; for
there are propositions so simple that, when they are heard and taken in
by consciousness, they are immediately acknowledged as true, e.g.\ that
1 and 2 are 3.

What holds of the theoretical principles shows itself still more
easily with respect to the practical and moral principles, none of
which has such evidence as those. They all require a proof, are
therefore not clear in themselves, therefore not innate either, and, as
experience shows, they do not have universal validity either. Even the
highest principle: do as you would that others should do unto you, is
not merely not followed, but not even universally acknowledged. To say
that such principles are obscured is a poor evasion, which would lead
to their being most clearly present to consciousness in children and
the uneducated. In fact the practical principles are different, indeed
even contradictory, in different individuals and different nations;
insofar as agreement is found, it is brought about by tradition and
upbringing. Innate in all men is only the direction of desire toward
happiness, not anything known by the understanding.

With respect both to the theoretical and to the practical principles
it holds that they cannot be innate when the concepts that are their
elements are not. But in the most universal principles lie the most
abstract concepts, such as being, non-being, identity and difference
and the like, and these can be correctly formed only by strenuous
reflection. They are incomprehensible to children and lie as far as
possible from the sensuous feelings that are originally the sole content
of children's consciousness. Not even the notion of God is innate. Many
peoples do not possess it, and where it is found, it is infinitely
different among different persons, even within the same religion. But
the divine wisdom and power reveal themselves so clearly in creation
that a rational being \opage{46}must, by serious reflection on the
causes of things, be able to discover God, and once this concept had
been found by a few, it would have to be so generally evident that it
would have to be preserved.

When it has thus been proved that our ideas (notions, ``the immediate
objects of the understanding'') are not innate, the task becomes: to
state positively the source from which they derive. As this source
Locke states \emph{experience} in its twofold form as sensation and
reflection. Originally the soul is like an unwritten tablet, a ``tabula
rasa,'' that is to be filled in. And this filling-in is not conceived
in the way Aristotle, from whom the image of the tablet is taken,
thought it, namely such that the mind, awakened to activity by its
relation to the external, by its own energy transforms the original
bare possibility into actuality. On Locke's conception the mind during
the filling-in is passively receptive: it is the object of experience
that inscribes itself on the tablet and successively fills it in. Of
the two forms of experience, sensation relates to external objects,
reflection to the mind's own states and activities. Reflection is also
called inner sense, since as observation it is akin to sensation. It is
dependent on sensation, since we must first have received ideas through
sensation before we can direct reflection to them, as they are, as
being in the soul, and to what goes on in our soul during the reception
of such ideas. That the understanding behaves passively in the
reception of ideas holds just as fully with respect to the ideas that
derive from reflection as to those that derive from outer sense: it
does not produce them, but they are produced in it as in a mirror. The
ideas that thus come to the understanding without its active
co-operation are the \emph{simple ideas}, which form the whole material
of all our knowledge. Locke divides them into 4 classes: those we get
through a single sense: e.g.\ the notion of heat, light, sound and the
like; those we get through several senses: e.g.\ the notion of
extension, form, motion; those that arise through reflection: the idea
of \opage{47}thinking and of will; and those that derive both from
sensation and from reflection, such as the notion of pleasure and pain,
of power, of unity, of existence, etc.

Our various sensuous notions stand in an essentially different
relation to what exists objectively. Locke distinguishes between
\emph{primary} and \emph{secondary} qualities. Primary or original
qualities are those that are really in the object and cannot be
separated from it, such as extension, impenetrability, figure, motion
and rest, number, duration. The ideas these qualities call forth in us
are real images of the properties of things and have a likeness to
them. But things also have the power, without our noticing that it
happens by means of their primary qualities, to call forth in us such
simple ideas as designate not a property of the things but a feeling in
us. Thus, e.g.\ what we perceive as color is properly a configuration
of the parts of the body; what we perceive as heat, a kind of motion in
these; but we feel it not as form and motion but as color and heat.
That power in things is their secondary qualities (or their sensible
qualities). The ideas called forth by them are naturally not images of
the properties of things but only signs of them. Finally, things have
the power, through the particular determinateness of their primary
properties, to call forth such changes in other bodies as produce an
idea in us. This is as a rule called not a quality but a power or a
force (e.g.\ the sun's power to melt wax). Properly there is no
difference between this and the secondary qualities; for neither the
feeling of light and heat, nor the altered state of the wax, are
properties of the sun, but only effects of it. But in the case of the
secondary qualities we are, since we cannot observe the configurations
by which the body calls forth the notion of color, inclined to regard
both as one. Since, on the other hand, we can separate and compare the
sun and the wax, we can here distinguish between \opage{48}property and
effect. But properly the relation here is the same as in the secondary
qualities, and in any case one can, if one wants to distinguish, call
these immediately observable, the powers mediately observable
qualities.

How the motions called forth in the passive mind by the action of
external things on the senses can be converted into feelings and
notions, adequate and inadequate, Locke pushes aside as
incomprehensible, and points to God's omnipotence.

The material given to the understanding in the simple ideas is now
worked up by a combination through which the \emph{compound} (complex)
\emph{ideas} arise, and in the formation of these the understanding can
exercise a voluntary activity. The number of compound ideas is, by the
nature of the case, infinitely great; but they can be brought under
three classes: the ideas of modes, of substances and of relations. By
\emph{modes}, which can be either simple or mixed, according as they
are formed from one idea (e.g.\ the idea of modifications of color, of
number, i.e.\ the repetition of unity, of modifications of time and of
space and the like), or from several different ones (e.g.\ the idea of
beauty), Locke understands such compound ideas as always designate
something that one cannot think of as subsisting by itself, but only as
occurring in a substratum that is, as it were, the bearer of those
modi. The concept of \emph{substance} Locke treats in detail. He seeks
to show that it arises when the mind observes that certain simple
ideas, which may derive either from sensation or from reflection,
always go together, and cannot think of what these ideas represent as
subsisting by itself. It then accustoms itself to assume a substratum
in which what is designated by the ideas has its subsistence and from
which it proceeds, and this it designates by the word substance. In
this concept, therefore, there lies in the first instance only that it
is something, one knows not oneself what, that is thought as the bearer
of the qualities that call forth the simple ideas, i.e.\ of the
accidents. From this origin of the concept it does not follow, however,
that substance \opage{49}should not exist outside us; on the contrary,
precisely the idea of substance has its archetype outside us. What this
is, is absolutely unknown to us; we know only the attributes, which as
a rule are secondary qualities or powers. Substances Locke divides into
cogitative and incogitative (not into immaterial and material, because
it may well be possible that the thinking beings, the spirits, are
material and that the power to think can lie in matter). The attributes
of the former are thinking and motive power (will), those of the latter
solidity and mobility. These attributes we know, but what the two kinds
of substances are according to their essence is equally unknown to us.
--- The concepts of \emph{relation} arise through the comparison
between several notions, through the transition from the one to the
other. Such concepts are the concepts of cause and effect, of identity
and difference, etc. They are all merely a product of the
understanding.

The activity that the mind exercises in relation to ideas consists,
according to Locke, partly in a voluntary turning toward or turning away
of our attention where several notions present themselves; partly in a
retaining and a recalling of the individual notion; partly in a
comparing, separating and combining of notions; partly, finally, in
abstraction, by which universal concepts are formed, whose verbal signs
can be applied to all the notions comprised under them. The faculty of
abstraction Locke regards as a \emph{specifically human} faculty.

Through the linking of ideas \emph{knowledge} arises, which \emph{does
not reach beyond the domain of ideas}. It is a perceiving of the
compatibility or incompatibility, agreement or disagreement, of ideas
(according to the 4 main relations: their identity and difference,
their determinate relation, their coexistence, and their real
existence), and a \emph{judgment} about the truth or untruth of the
connection contained in the proposition. Language here acquires its
essential significance as the necessary expression of thought and as a
means of communication. Words become signs for things, appellatives
comprehensive signs for classes \opage{50}of individual things,
expressions for the common concepts that are formed through
abstraction. Such comprehensive signs are necessary on account of the
number of things. The common concepts are not images of anything real;
for only what is individual exists; but they are complex ideas, which
are only products of our thinking. The words species, genus, etc., thus
designate nothing actual, but only the relation of things to our
understanding.

Truth and untruth first come forth in the proposition, because only the
proposition expresses a combination or a separation of ideas. The
complete or incomplete certainty in the judgment of knowledge about the
compatibility of ideas determines the difference between
\emph{knowledge} (certainty) and \emph{opinion} (probability). Our
knowledge can be intuitive, demonstrative, or sensitive (i.e.\ a
knowledge of the existence and constitution of the external). Our
knowledge is real only when the connection of our ideas corresponds to
that of actual things. Our knowledge is strongly limited both with
respect to material and to spiritual beings. Uncertainty has its ground
partly in the lack of ideas, partly in the lack of a comprehensible
connection between them, partly in the lack of exact investigation and
correct designation of them. Where certainty is lacking, probability
takes its place as supplementing defective knowledge. The conviction it
gives is only an opinion; but this is not arbitrary, since the greater
probability has power over the understanding. \emph{Faith} in a
revelation likewise becomes a supplement to our knowledge. The content
of the faith of revelation can be in agreement with reason, e.g.\ the
existence of God, to which we can infer with evidence by our own reason
from the existence of the world and of ourselves, or it can lie
\emph{above} this knowledge by reason (as the doctrine of the
resurrection of the dead). On the other hand nothing can be accepted as
revelation that \emph{conflicts with} knowledge by reason by being
either self-contradictory or incompatible with clear and distinct
concepts, e.g.\ the existence of several gods. It belongs to reason to
judge of \opage{51}revelation; faith in it cannot have the certainty of
intuitive or demonstrative knowledge by reason. ---

The empiricism that Locke founds becomes in the 18th century of
preponderant influence on English and French philosophy, and also
gradually exercises an influence in Germany, where at first it had
found its counterweight in the Leibnitzian philosophy. Of the many
philosophical phenomena that point back to Locke we mention only those
in which essential consequences of empiricism, distinctive stages in
its development, have come forth.

\divrule

\begin{center}
  \textbf{4. Berkeley.}
\end{center}

\emph{George Berkeley} (born 1685 in Ireland, Bishop of Cloyne, died
1758. Main works: treatise of the principles of human knowledge, 1710;
theory of vision, 1709; three dialogues between Hylas and Philonous,
1713) draws from Locke's empiricism the consequence of formal
(empirical) idealism. He reaches it by the route of dissolving the
universal into the individual, the thing into its properties, the
properties into perceptions.

If ideas, as empiricism teaches, have originally arisen from sensation,
then, he concludes, since all sensation relates to the individual, they
must all be expressions of something individual. What are called
general concepts are only notions of something individual, individual
notions that for the ease of thinking are taken as representatives of
all similar individual notions. (Thus in the proofs of geometry the
single determinate drawn figure represents all other figures of the
same kind. The proof acquires general applicability to the individual
figures by disregarding one property or another of them; but this
disregarding is something other than forming an abstract idea that
excludes \emph{every} concrete determination: such an idea does not
exist.)
\opage{52}If, next, only what is actually sensed is something actually
given, then, just as little as the unsensed universal, can such a
concept as substance have reality, which, although it is not the object
of any sensation, was posited by empiricism as the substratum of the
properties and was ascribed objective reality. The concept of substance
is only a fiction; we sense only the properties; only they remain as
given, while the feigned substratum vanishes. ---

And when, finally, empiricism, with respect to the properties, had made
a distinction between primary and secondary properties and had seen in
the former determinations of the objective real itself, Berkeley denies
the validity of this distinction too. In what are called primary
qualities, too, it is impossible to disregard the perceiving subject;
e.g.\ in the notion of extension, which is always the notion of a
determinate extension, in the notion of magnitude, which is determined
only by inferences from our feelings of sight and touch, in the notion
of impenetrability and the like. Our sensation is here just as
essential a factor as in the perception of color, heat or cold and the
like. To conceive the primary qualities as belonging to objectively
subsisting material things outside us, which as it were imprint their
images on our mind, entangles us in the greatest contradictions, since
it is unthinkable that something absolutely heterogeneous such as
matter should be able to act on the mind, and that an activity should
be able to proceed from anything other than a willing, i.e.\ a
spiritual, being. All our notions of the properties of things, the
primary as well as the secondary, in reality express only our
feelings; but when in all qualities we have to do only with feelings,
then what we call properties are only our own affections, their being
is percipi, and what we call thing, object, exists, as the sum of the
properties, only in us. Nothing exists but spirits (representing and
willing beings) and their \opage{53}ideas and acts of will. The ideas
of something corporeal that arise in us without our voluntary activity,
whereas the images of the imagination are produced voluntarily by us,
are brought about in us by God, in whom they are as archetypes. What we
call nature designates only their connection, what we call law of
nature, the constant order of the combinations of ideas. The fixed
order of nature is necessary in order that we may understand nature and
through it God's will. Ideas are signs through which he instructs us
and determines our will. Through ideas we know the sensible, through
immediate certainty ourselves; by inference from their effects we form
a \emph{concept} of other spirits and of God. God's relation to the
finite active spirits Berkeley determines only uncertainly. His
conception at times acquires a pantheistic tinge. ---

In his theory Berkeley believes he has found a secure bulwark against
materialism, atheism and skepticism. His denial of matter develops with
consistency out of empiricism, but since he lets the whole content of
experience subsist in the mind in the form of notions as something
given from without, his idealism (immaterialism) becomes formal and
empirical. Experience now merely seeks its object in us instead of
outside us; but the content is the same. The question of the
possibility that the spiritual God, who is supposed to have the ideas
in an entirely different manner from us, namely without sensuous
feeling, can produce sensuous notions in the mind, is not raised, and
neither is the truth of the content of the notions inquired into. But
Berkeley has brought forward and sharpened a problem that the French
sensualism and materialism of the last century evaded, and to which the
renewed sensualistic tendency of our day has returned: the problem of
the existence of the things of sense outside our sensation, and of the
relation of their objective being to our perception of them.

\divrule

\opage{54}
\begin{center}
  \textbf{5. Hume.}
\end{center}

\emph{David Hume} was born in 1711 in Edinburgh. From the study of law
he turned early to the study of ancient and modern literature and of
philosophy. From 1734 to 1737 he stayed in Paris, where he wrote his
first work: A treatise on human nature, being an attempt to introduce
the experimental method of reasoning into moral subjects. It was
published after his return home to England, but attracted little
attention. On the other hand, Part 1 of his Essays, which appeared in
1742 and treated moral, political and literary subjects, was received
with approval. During a new stay abroad he reworked his first work, and
published, in 1748, a part of the reworking under the title: an enquiry
concerning human understanding (later included in the Essays as
Part 2). With this work he called forth strong attacks, but laid the
foundation of his great renown. As librarian to the Faculty of Law in
Edinburgh he composed in the years 1752---1762 his great History of
England. A treatise on the principles of morals (later printed as
Part 3 of the Essays) had already appeared in 1751, and in 1757 appeared
his natural history of religion (included in vol.~4 of the Essays),
which drew many attacks upon him. In 1763 he went to Paris as secretary
of the embassy, where he was received very flatteringly and came into
closer connection with Rousseau, who in 1766 accompanied him back to
England, where, however, their relationship was soon disturbed by
Rousseau's morbid irritability and mistrust. In 1767 Hume became
Under-Secretary of State; from 1769 he lived in scholarly otium in
Edinburgh, in which city he died in 1776. After his death appeared his
dialogues concerning natural religion, and his Essays on suicide and
the immortality of the soul.

\divrule

Hume relies on Berkeley, though without sharing his idealism, and
takes over from him the doctrine of ideas as expressions only of the
individual, of the subjective charac\opage{55}ter of our notions, and of
substance as a mere complex of notions. But he goes further than
Berkeley, in that he extends the latter's conception of corporeal
substances to the spiritual substance as well, the self, which for him
becomes merely a complex of rapidly successive feelings and notions, to
which we (we!) unjustifiably supply an underlying substrate; and in
that, by his critique of the concept of causality, he annihilates the
possibility of a necessary and universally valid synthetic knowledge,
% errata list B, p. 55: printed "en almeengyldig", corrected "almeengyldig"
and thereby draws the consequence of skepticism from empiricism.

In investigating the origin of our notions, Hume separates impressions
(impressions) from ideas (ideas). By the former he understands the
stronger and more lively feelings that we have immediately in our
sensation, our passions and our emotions; by the latter, the weaker
feelings that we have when the former are copied in memory or the
imagination. The ideas are thus conditioned by the impressions and are
their copies; the impressions are the first material of thought, and
thinking consists only in a transforming and combining of them, in
which the passive mind is governed by the natural law of the
association of ideas. Even the idea of God the understanding forms only
by transforming the human qualities, wisdom and goodness, through
extending them to infinity. The impressions become the criterion of
whether a philosophical concept has a real meaning. If we cannot state
from which impression an idea is derived, we must regard it with
mistrust.

As the objects of the human understanding Hume determines partly the
relations between ideas, partly matters of fact. To the first class
belongs the content of geometry, arithmetic and algebra, in general all
propositions that have intuitive or demonstrative certainty. Of these
propositions we have a purely a priori knowledge; they are known by the
activity of thought, independently of existence. The geometrical
propositions about the circle and the triangle would thus have validity
even if none of these figures existed in nature. The propositions known
in this way all refer to mag\opage{56}nitude and number and contain
\emph{identical} (analytic) judgments. Propositions that concern matters
of fact do not have the same certainty as the former; they do not give
the same feeling of a necessary truth. The knowledge that relates to
matters of fact and connects them has a \emph{synthetic} character. The
relations according to which it connects them are those that determine
the association of ideas: resemblance, connection in space and time,
and the relation of causality. Of these, connection according to the
relation of causality is of the greatest significance, since only this
connection lays claim to being necessary and universally valid. To the
origin and validity of the concept of causality, therefore, Hume
directs his investigation, since he would determine the certainty of
the knowledge that relates to the connection between matters of fact,
and would determine it by inferences. The causal relation is a relation
between things that differ, a synthetic relation. The effect is
entirely different from the cause, and so cannot be derived from the
latter's concept by a priori thinking, by inferences of reason, since
these move only within the identical. By inferences of reason alone we
cannot, e.g.\ decide whether a heavy body that is in free space will
fall down to the earth or move upward. This we learn only through
experience. But experience can give us only a rule, not a law that
determines with necessity, and it can show us only the succession of
matters of fact, a post hoc, not a propter hoc, not the inner uniting
bond (the force); for this is not observed in the external, and the
reality of the concept is not secured by what we observe in our
interior, where we have indeed a feeling of the will, but not of the
way in which it is supposed to set the organism in motion. But the
question then becomes: \emph{how the concept of causality arises in
us}, and what brings us, when the same phenomena repeat themselves, to
expect the same consequences that we have previously observed. The
answer is: the concept of necessary connection, or of force, which
could not be formed on the basis of the single case, is formed as
\opage{57}the cases frequently repeat themselves, and there forms in our
imagination a \emph{habit} of passing over from a phenomenon to its
usual consequence, and of expecting that, when the former shows itself,
the latter will also occur. We then \emph{feel}, in this transition
become habit, a new \emph{impression}, the impression of a connectedness
between the phenomenon and the consequence, and from this impression is
then formed the concept of \emph{force}, and of a \emph{necessary causal
connection}. The individual cases we observe are all alike; when,
therefore, a \emph{number} of cases can give what the single case could
not, this shows that it is precisely the repetition, and the habit
produced by the repetition, that is decisive. Habit, which conditions
the liveliness of the impression, here gives the \emph{belief} (belief,
as distinguished from religious belief: faith) in the objective
validity of the causal relation, just as the constancy of certain
sense-impressions gives the belief in the existence of external
objects; and this belief satisfies practically, although it cannot be
transformed into a knowledge.

In the investigation of causes the human mind is, apart from the
decisive point that it does not comprehend the objective connection
between cause and effect but only posits a connection in accordance
with a subjective need, limited in a twofold respect. The aim of
knowledge is to lead the causes of natural effects, found by way of
experience, back to the smallest possible number of general causes. But
this reduction has its limits: we reduce the causes of natural
phenomena to a small number, but the higher causes of these general
causes, the ultimate causes, lie outside the limits of our inquiry. And
since the concept of causality has its origin in habit, in the
repetition of experience, its application is justified only within the
domain of experience. An inference from what is given in experience to
what lies outside experience, such as God and immortality, is
unjustified. ``Whatever is not an investigation, stemming from pure
thinking, of magnitude or number, or an experience-based
inves\opage{58}tigation of matters of fact and existence, can only be
sophistry and illusion.'' ---

Since Hume, by his critique of the concept of causality, annihilates
the necessity and universal validity of synthetic knowledge, which
stand and fall with that concept, he thereby draws the negative
consequence of the standpoint that lets the passive mind receive its
content exclusively through experience, which in its isolatedness
relates only to the individual in its isolatedness, and thereby in
principle excludes the universal and the necessary connection. In place
of truth his skepticism leaves only probability. His conception of the
coming-to-be of the concept of causality, which in reality presupposes
what it would explain, leaves the problem of the limits of our
knowledge unsolved, and only Kant finds, from a new point of departure,
a way to its solution.

\divrule

\begin{center}
  \textbf{6. Condillac.}
\end{center}

\emph{Etienne Bonnot de Condillac} was born in Grenoble in 1715. Through
Voltaire's writings he came to know Locke's philosophy, and in his first
work, Essais sur l'origine des connaissances humaines, 1746, 2 vols.,
he still stands on Locke's standpoint. His own distinctive sensualist
standpoint, which he had reached under the influence of Berkeley, he
developed in his main work, traité des sensations, 1754, 2 vols., in
the treatise on animals, 1755, in his Cours d'études pour l'instruction
du prince de Parme (13 vols., 1755 ff.), and in his Logic and his
langue des calculs. He died a member of the French Academy in 1780. ---

In Lockean empiricism experience was determined as the only source from
which the mind's tabula rasa received its filling. But Locke had not,
after all, consistently annihilated all activity in the mind. He had let
reflection stand as a source of ideas beside sensation, and had granted
it, if somewhat waveringly, an activity; he \opage{59}had let complex
ideas be a product of spontaneity as well, and he had conceived the will
as an independent faculty beside the passive understanding. Moreover,
although he denied innate knowledge, he had spoken of innate drives, of
an innate striving after happiness, of an innate conscience, and had
thereby granted to the mind what his principle of knowledge must
consistently have led him to deny. Condillac, by contrast, makes
sensation the only source of all our ideas, and from it he lets, by an
involuntary development, not only all the functions of the thinking
soul but also desire and will proceed. This sensualism is a correct
consequence of empiricism. The mind, which, as originally empty, is
posited as merely receptive, relates through experience to an object
that has only the determination of being the other of consciousness,
since even the object of inner experience, the state of consciousness,
is conceived as something determined from without and effected from
without. The primitive form of knowledge then becomes the sensuous
feeling, whose involuntary transformation in the receptive mind
produces the other forms of consciousness, and the primitive and actual
content of knowledge becomes the external, i.e.\ the corporeal, which,
as corporeal, is the bounded, the individual, and more precisely: the
mechanically moved. For the spiritual element, which in Bacon remained
behind in the admittedly only postulated concept, animated matter, must
vanish from what becomes the material proper, as dualism asserts itself
within matter (cf.\ p. 32 f.); and as that element vanishes, this
material can be posited only as the mechanically moved, which does not
have the point of departure of activity (motion) within itself. This
conception of the genesis of knowledge and of its content is that of
sensualism; but in sensualism lies in turn the consequence of
materialism. For sensualism the mind is only receptive, more precisely:
receptive in relation to a mechanically determined external, i.e.\
receiving its content, its determinations, through impressions on the
senses, the organs of receptivity. For this dogmatic sensualism there
consistently re\opage{60}mains nothing but the sensed, i.e.\ material
objects, and the sensing subject, which, in its bare receptivity,
relates to the object as thing to thing, and becomes a dependent link
in the chain of things, as it passively receives the impressions of the
mechanically connected finite objects of sense, which impressions are
then deposited in consciousness, as in a container, as concepts, in
which there is nothing independent that points to an ideal region, but
in which only the sensuous is as the real. But as the subject becomes a
mechanically determined link in the series of things, it must come
under the common determination of the series; it becomes a sensuous,
material subject, and everything in existence, the subjective and the
objective, is now led back to the same unity: to nature, determined as
passive matter; and this concept of matter has become the
all-embracing one, which absorbs the mind into itself.

\divrule

In order to make vivid how all psychic processes proceed from
sense-feelings, Condillac uses the fiction that a statue, organized
internally like a man and animated by a mind that as yet has no ideas
at all, has one sense after another awakened, so that it successively
becomes accessible to all impressions from the external world, to which
it was originally closed. He begins with the sense of smell, because
one usually takes this for the sense that contributes least to
knowledge, although in reality there is no difference between it and
the other senses, the sense of touch excepted. The other senses, namely,
all have this in common, that through them man has only affections of
his soul, not a notion of anything external. The senses give, in the
first instance, feelings, not ideas. The statue originally feels
\emph{itself} modified. This first stage Condillac calls now
\emph{sensation}, now \emph{perception}, and this perception is
immediately \emph{consciousness} (conscience). ``Feeling and
consciousness are the same act under two names. Insofar as one
con\opage{61}siders it only as an impression in the soul, one may use the
name feeling; insofar as it informs the soul of its presence, one may
give it the name consciousness''. Consciousness coincides with
reflection in the widest sense. As long as the statue has only one
feeling, it is, as it were, absorbed in it; but when, although it still
has only one sense, it is affected more strongly by the one impression
than by the other, \emph{attention} enters, which is the turning toward
the more powerful impression. By attention the impression is imprinted
in the soul and leaves behind a trace in the brain, and the retention
of this trace is \emph{memory}. Only through this function does the
statue get \emph{ideas}, since these are the retained images of the
feelings. In feeling, the movement goes from the sense organ to the
brain; in memory, the whole movement takes place in the brain itself.
Now when the statue has the idea of a feeling and at the same time a
new feeling enters, it will be divided between the two, and this gives
immediately \emph{comparing} and \emph{judging}. The series of
judgments that rest on a series of comparisons is \emph{reflection} in
the narrower sense, whose content is thus only sensations (sensations).
But with memory yet another, higher function follows. The ideas
imprinted in memory form, namely, a chain that is determined by the
original connection and sequence of the perceptions, and the soul is
thereby compelled, when it observes an idea, to pass over to the one
connected with it. If now, in a series of ideas, some make an
agreeable, others a disagreeable impression on the soul, the soul
passes quickly over the latter, whereby the former move closer
together, and the soul grows accustomed to connecting them. Through
this association of ideas the soul now acquires facility in calling up
ideas by means of other ideas, and this faculty is the
\emph{imagination} (imagination, from image). Through it, it also
becomes possible to \emph{abstract}, i.e.\ to represent to oneself an
idea that appears as connected with another as independent. Abstraction
has thus proceeded from sensation as its \opage{62}source, and one sense
is sufficient to give abstract ideas. These are nothing other than
certain classes under which similar ideas and feelings are subsumed
according to their similarity. The abstract ideas that the statue will
have as long as only the sense of smell, or the sense of taste, of
hearing or of sight, is in activity, will be limited to a certain
circle, and will express only the kinds (classes) of the soul's
modifications and their unity. Thus there arise first the abstract
ideas of the agreeable and the disagreeable, and the idea of the self.

From the other senses the sense of touch now differs by the essential
peculiarity that through this sense, whose object is solidity
(impenetrability), we observe an exclusive relation between two things,
thus not merely a modification of our own state, but something
different from the one who feels, something really objective, an
external object. Only through the sense of touch do the other senses,
too, learn to relate to their perceptions as to something objective, to
localize them in space as determinations of external bodies. The
functions of the soul whose genesis we considered before now take on a
modified form. Attention and comparison are now directed toward
something \emph{objective} as such, and there is \emph{reflection} on
the various relations of \emph{the objects}. Ideas now form of a
multitude of \emph{qualities}, and as these are felt simultaneously,
there arises by a kind of fiction the idea of an underlying substrate,
and in this way the concept of \emph{substance} is formed. Properly
speaking, what we call substance is only an aggregate of many
qualities; by the name we indicate a union among them, a connecting
bond; the word has no other meaning: we do not know substance. Besides
the idea of substance, the statue, now that the sense of touch has
entered into activity, forms countless other abstract ideas, as it
dissolves the groups in which the qualities show themselves to it.
There also now arises, as the concept of idea is applied to the
perception that is perception of something objective, a distinction
between \emph{intellectual ideas} (traces of \opage{63}past impressions)
and \emph{sensible ideas}, which represent to us objects that at this
moment act upon our senses. The assumption of innate ideas stems,
according to Condillac, from a failure to recognize this difference.
--- Through the cooperation of the sense of touch, finally, the notion
of the \emph{self} becomes a different one, since the statue now
distinguishes the self from the feelings and from the bodily organs,
whereas before it coincided with its modifications, the feelings. The
knowledge of the self is limited, however, to this: that all feelings
belong to it. Here too we do not know the substance that is placed
beneath the feelings, any more than we know \emph{what} force and cause
are: our knowledge is limited to the qualities.

Besides the evidence that the sense-impression and the observation of
our own states give, there is another kind of evidence, which is
attained through \emph{reasoning}, the highest function of knowledge.
This evidence is grounded on the relation of identity, since a correct
reasoning consists only in leading a proposition back, by substitution
of terms of equivalent meaning (equal terms), to an identical one. A
proposition is evident by itself when it is only another expression
for: the same is the same.

Since, in the manner indicated, all theoretical functions are derived
% errata list B, p. 63: printed "idet saaledes", corrected "idet paa den angivne Maade"
from sensation, Condillac leaves the soul only a freedom consisting in
this: that it is supposed to be able at will to direct its attention
toward a particular point in the store of our ideas, or to divert it
from there. But this freedom becomes doubtful in turn, since the
difference between the soul's activity and passivity is determined by
whether the cause of what takes place in the soul is an external
% printed "er er ydre" (defect logged in the transcription)
or an internal one; for the internal one can be the mere after-effect of
the sense-impression.

The \emph{practical} functions are, like the theoretical, led back to
sensation. Every feeling is, as such, a feeling of pleasure or
displeasure. From this it follows that the soul, as it has the
recollection of a feeling and compares it with a present one, at the
same time compares vanished and present feelings of pleasure and
displeasure, and \opage{64}thereby arise \emph{desire} and
\emph{aversion}. Pleasure and displeasure become the mainsprings of
every action, the spur to our \emph{whole} development. Desire becomes
\emph{will} when the soul, in the recollection of a previously
satisfied desire, expects to be able to satisfy it in the future too.
Will is, namely, desire that is conscious of its capacity for
satisfaction. As the soul raises itself to will, it also has the ideas
of the good and the beautiful, both of which have their source in the
senses, since good originally designates what pleases smell and taste,
beautiful what satisfies the other senses. Later the meaning was
altered so that what appeals to desire is called good, what appeals to
the understanding, beautiful.

\emph{Understanding} (entendement) and \emph{will} can, as class
concepts (abstract ideas), be used as comprehensive expressions for all
the activities of the soul. They are themselves comprehended in a more
general concept: \emph{thinking}, which on the standpoint of sensualism
then becomes synonymous with sensing and feeling.

In sense-feeling man meets the animal, but there remains a great
difference of degree. This difference is conditioned partly by the fact
that man's sense of touch surpasses that of all animals; partly by the
number and nature of the needs and desires that occasion knowledge and
give the means of acquiring new ideas and combining them; partly,
finally, by the richer development of language in man. The animals too
have a self, but only what one may call a habit-self (le moi
d'habitude), which consists in the totality of habits, and which, acting
instinctively, suffices for the general needs of animal life; man has
in addition a self that reflects
% printed "reflécterer" (defect logged in the transcription)
(le moi de réflexion), and by which we are active even in unaccustomed
circumstances. The animals do indeed reflect too, but since the circle
of their activities is so restricted, these are repeated continually
and become habit, and thus the animals are restricted to the mere
habit-self, to instinct.

\opage{65}\emph{Language} is of the highest importance for the association
of ideas and for the formation of new ideas. If we did not form signs
for the ideas, these would be few. The signs can be accidental, natural
(such as cries) or arbitrary. When the use of the last has become
habitual to us, which it becomes to such a degree that we often forget
that in the word we have only a sign for the thing, not its essence,
the external object calls up its name or its sign in our recollection,
these signs call up other signs, etc. Just as every thing is an
individual, so the signs are originally individual designations, names.
But general concepts also come to be designated, and are then only
names for certain classes into which we arrange our ideas, according as
we have felt a need to distribute our knowledge in a definite order.
Indeed, only through designation and naming is it possible for us to
have abstract ideas, without which we can neither infer nor reason at
all. All directions for correct reasoning are properly contained in the
direction for speaking correctly, and speaking, abstracting and
reasoning are entirely the same. The animals have a language, to be
sure, but an imperfect one, and from this it follows necessarily that
they communicate to and learn from one another little, and can imitate
and learn from others only to a slight degree. ---

Despite his sensualism, Condillac, in this still influenced by
Cartesianism, sharply separates mind from matter. The fundamental
property of the soul is thinking, which cannot be the predicate of
matter, since matter, which has extension for its fundamental property,
is an aggregate, whereas the subject of thinking must be really one.
This then originally leads Condillac to posit feeling not as effected
but only as occasioned by the body, and to maintain that the soul,
according to its concept, could arrive at knowledge by itself; but since
the soul, as it is an object of experience, is inseparably connected
with a body (a connection that Condillac derives from the Fall), the
unity involuntarily asserts itself in the closer development of the
relation, and the soul's func\opage{66}tions flow together with the
movements of the brain. The difficulties that sensualism encounters in
the question of the existence of external things and of an adequate
knowledge of things, Condillac thinks to remove partly by appealing to
the immediate testimony of the sense of touch, partly by maintaining
that, even if we do not know the essence of things in and for itself,
or their ``absolute qualities'', we do know certain qualities of things
that underlie all the others we know, and which may be regarded as the
secondary essence of things (essence seconde). The knowledge of these
is sufficient in a practical respect.

\divrule

\begin{center}
  \textbf{7. Materialism} (Système de la nature).
\end{center}

The materialistic consequences of sensualism are developed, amid a
passionate polemic against the established order in state and church,
by Diderot, la Mettrie and the author of the Système de la nature. In
\emph{Diderot}, the famous editor of the Encyclopédie (1713---1784.
Writings: Pensées philosophiques, 1746;
% printed "plilosophiques" (defect logged in the transcription)
Proménade d'un sceptique, 1747; Introduction aux grands principes;
Lettre sur les aveugles, 1749; Pensées sur l'interprétation de la
nature, 1754; Entretien entre d'Alembert et Diderot, and le rêve
d'Alembert, 1769), the deism expressed in his earlier works passes
over, through a skepticism, into a pantheistic fusion of God and
nature, and from there into a disappearance of God in nature, while at
the same time the distinction he earlier maintained between the
material and the immaterial is lost, inasmuch as feeling, which becomes
conscious in the organic, and out of which thinking develops as out of
its germ, is ascribed to the molecules, which have an original life.
Sentient matter now becomes the sole existent; the I is conceived as a
product of the brain's feelings, freedom and immortality as an
illusion, virtue and vice as a result of fortunate or unfortunate
natural dispositions, the passions as the productive powers of action.
Yet Diderot still polemicizes \opage{67}against those who make
self-love the principle of morals, and he lets the good man be happy,
the wicked man unhappy, without regard to their outward fortunes.

La \emph{Mettrie} (1709---1751. Writings: Histoire naturelle de
l'âme, 1745; l'homme machine, 1748; l'art de jouir, 1750) proclaims
materialism and atheism in the crassest form. For him man is a
machine, in which thinking and will, both of which develop out of
feeling, are products of the activity of the brain-muscles (!), the
soul one with the brain. Man's freedom is an absurdity, death is the
end of ``the farce of his existence,'' sensuous enjoyment the goal of
his life, and the transience of life only a motive for forcing
enjoyment. Religion La Mettrie regards as a hindrance to the peace and
felicity of the human race, atheism as the condition of it.

The main work of materialism is the \emph{Système de la nature}
(1770), written, or perhaps more correctly edited, by the German Baron
\emph{Holbach} (1723---1789), a friend of Diderot, whose influence is
strongly traced in the work, and with whose contemporary works it
agrees strikingly on many points. In the Système de la nature the
materialistic world-view
% printed "Verdensankuelse" (defect logged in the transcription)
is carried through in all domains under constant polemic against
Christianity. Sensualism forms the point of departure, but a skepticism
is avoided, inasmuch as in sensation, which is determined as a motion
in us, one finds the guarantee not only of an external motion but of a
moving substance: matter, with its original properties. On this
foundation the theory develops dogmatically, by inferences from the
known to the unknown. The fundamental thoughts of the work are the
following. The whole of existence: the universe, shows us nothing other
than matter and motion. Eternal matter, with the original qualitative
differences of its molecules, has an energy by which it maintains
itself, and in it itself, not in a being outside it, lies the ground of
the motion inseparable from its essence. This motion has at every
point reality only as \opage{68}communicated motion. As a necessary
consequence of the essence of matter, motion, by which the
imperishable molecules are united and separated under ever-changing
modifications, and the infinite chain of causes and effects is
produced, is determined by the unalterable law of necessity. Nature,
the active whole that contains the sum of the unknown forces that
govern what exists, maintains itself eternally, without purpose, amid
the uninterrupted change of the parts it comprises and governs. Man,
as part of nature and as nature's work, falls under its
determinations. He is only matter, his consciousness is a result of
the feeling of matter, his intellectual activity the product of the
invisible motions of the brain-molecules, his faculties, inclinations
and passions the result of the physical mixture of the body: the
temperament. Freedom cannot belong to him: that would be a breach of
the world-order of necessity, of the infinite chain of causes and
effects, of communicated motions. Immortality he cannot have; for when
the body is dissolved, nothing remains. The notion of a spirit,
absolutely different from the corporeal, has arisen from ignorance of
the invisible molecular motions in the brain, of which we have a
feeling and whose visible effects we observe. Since the cause withdraws
from our observation, we feign a cause, which we equip with predicates
that are wholly different from those of the visible corporeal. But this
notion, by which we as it were double ourselves, is empty and
contradictory. The soul is determined as immaterial, indivisible,
without extension, invisible, but nothing determinate can be thought in
a being that is to be only a negation of that which gives a knowledge.
And it becomes inexplicable how the two absolutely opposed beings:
spirit and body, which have no point of contact, should be able to
stand in a relation of interaction. The separation of soul and body in
reality means only the separation of brain and body. Thinking and
willing are peculiar functions of the brain. Like the notion of
spi\opage{69}rit, the notion of God too has arisen from ignorance. The
evils under which men suffer, and whose true causes one was unable,
from ignorance of the laws of nature, to find, led to the fiction of a
God who was to be their author. By this notion one doubled nature, just
as by the notion of a spirit one had doubled man. But the notion of
God, which in reality does not explain a single natural phenomenon, is
in itself full of contradictions, confusing and frightening, and not an
object of comprehension. One could determine God's essence only by the
negation of the human attributes; but since so indeterminate a being
could not hold men's interest, one sought to bring it nearer to man by
anthropomorphic determinations. Thereby incompatible, contradictory
elements entered the notion. The true system: the system of nature, is
free of such fictions; it is atheistic, just as it is materialistic and
fatalistic. It presupposes culture and courage and is to that extent
not for all. The atheist is he who denies that there is a spiritual
being whose invented properties can only disquiet man, and of such
atheists there will always be more, the more a correct knowledge of
nature and a sound reason spread. To assume that religion is necessary
for the multitude in order to keep it in check is the same as to think
that one should give a man poison so that he shall not misuse his
powers.

In the practical domain the same fundamental view is carried through.
Morals becomes the application of physics to man; medicine, which heals
the body, becomes its most essential aid. What in the domain of physics
is called ``gravitation toward itself,'' \emph{vis} inertiæ, the drive
of self-preservation, appears in morals as self-love; attraction and
repulsion appear as love and hate; the differing physical disposition
becomes the necessary source of the good and the bad action. Remorse
does not prove freedom and accountability, but only expresses the pain
over an ac\opage{70}tion's bad consequences for us; were the
consequences of our actions always useful to us, we would never feel
remorse. The criminal we punish for the same reason that we dam the
river that devastates our fields. Interest is the sole motive of all
actions, the sole true counterweight to harmful inclinations. Morals
first becomes effective when, instead of condemning our self-love, it
enlightens us as to where our well-understood interest leads. Whatever
agrees with this, to that we have a moral obligation. When we thus see
that, in order to bring others to take an interest in our welfare, we
must also take an interest in theirs, we have seen the obligation to do
so. He who satisfies his own interest in such a way that the others,
led by their own interest, must contribute to furthering it, is a
virtuous man. Virtue is the art of making oneself happy through the
happiness of others. Interest unites men into a society that furthers
the welfare of all individuals. The goal is happiness, and happiness is
continuous pleasure in the widest extent of this concept. But the goal
can be reached only where religion no longer frightens and confuses men
and separates them by the hatred of fanaticism. The system of nature
gives man calm and resignation, while at the same time it gathers all
his powers upon this life, the only one there is, and thereby brings man
to do everything here to satisfy his interest, and, since this can
happen only when he brings his fellow men to take an interest in it,
makes him a good man. Despite this morality of interest, the Système de
la nature nevertheless grants virtue, even as unrewarded, an
independent worth, (which is, however, essentially grounded in the
self-esteem it affords), and in the hope of the human race's progress in
culture, which is to be effected by the new world-view, there is
involuntarily acknowledged a purpose for existence and a distinction of
man from nature, unalterable in its eternal cycle.

\divrule

\opage{71}In the Système de la nature the realistic line of development
of modern philosophy has found its positive conclusion, just as it
found its negative one in Hume's skepticism. The unity of existence
that was striven after is here brought about by denying a divine spirit
and by making the human spirit a form of existence within the bounds of
matter, a particular kind of motion of the molecules. Shortly after
these consequences have been drawn, the Revolution in France brings a
halt to the philosophical movement in that country, and during the
ensuing reaction thinking there is led onto other paths. But in the
natural sciences the materialistic view finds a refuge, and from their
domain it has in our day, reacting against the idealism of this
century, renewed the struggle against the dualistic view, without,
however, this renewed materialism, despite all the aids the natural
sciences offer it, having developed a richer philosophical content than
that which the French materialism of the last century already
possessed.

In materialism the sharpest opposition to the view of life of the
Middle Ages is expressed, and the opposition betrays itself also in
materialism's passionate hatred of the Christian religion. And yet in
its very fundamental concept, despite the striving to assert its
validity as all-embracing, there shows itself an after-effect of that
view of life through the dualism that was characteristic of it. This
dualism found its conceptual expression at the beginning of modern
idealistic philosophy, when spirit and matter were determined by
absolutely opposed, mutually exclusive, predicates, and the essence of
matter was exhausted by extension and passivity. This conception still
asserts itself in materialism after the opposite of matter has been
denied. One does indeed transfer the determinations that dualism
ascribed to spirit, and that facts forbid one to deny, from spirit over
into matter; but they remain standing as foreign and heterogeneous, and
in the carrying-through they recede before the old conception. That
spirit \opage{72}is absorbed in matter does not make matter spiritual
and active; it remains what it was: the matter whose determination was
given precisely by the opposition to spirit, as it were the refuse of
spirit. In the very extreme opposition to the dualism of the Middle
Ages, grounded in spiritualism, the influence of that dualism betrays
itself.

This is easily shown when we examine materialism's concept of
\emph{matter}, as it is determined in the Système de la nature.

At first glance it does indeed look as if the concept of matter has
become another than it was in the dualistic opposition of idealism:
matter is posited as in its essence inseparable from motion; a protest
is thus made hereby against the notion of a dead, merely passive
matter. The question: whence comes motion? is answered with: whence
should it come but from matter itself, the sole existent? But this is
for the present nothing other than an expression, positive in form, of
the negative: it does not come, as the Cartesians assumed, from God, the
active spirit; for God is, after all, only a fiction. To be more, the
answer would have to be grounded through an analysis of the concept of
matter that showed the possibility of that connection; but this has not
been done. In the original properties of matter
(``primary qualities''),
% printed "Qualiterer" (defect logged in the transcription)
as the Système de la nature enumerates them: extension, divisibility,
mobility, impenetrability and vis inertiæ, there lies no point of
attachment for motion, nothing that can make it explicable; and since
it is expressly said that motion at every point in nature is present
only, and can become actual only, \emph{as communicated motion}, it is
thereby said that motion is nowhere present in matter in an original
way, that it is thus everywhere foreign to matter, and cannot be
comprehended by its nature. The postulated inseparable connection is
thus dissolved again; matter retains only the determinations that the
mechanical conception ascribed to it: it becomes passive, dead matter,
mechanically moved from without.

\opage{73}Now if not even the physical concept of motion, which in
idealism was put in connection with spirit only by a postulate, could be
explained from this concept of matter, it will already be given
beforehand that the peculiar determinations of spirit will not be able
to be so, that the matter that was to comprise the whole of existence
must leave one main form of existence standing as incomprehensible: the
main form whose essential determinations were precisely originally
excluded from its concept.

This shows itself when materialism is to explain the fundamental
determination of the human spirit: consciousness. In order to do this,
materialism first reduces, in the same way as sensualism had already
done, all the forms of consciousness to feeling, and then, appealing to
what it calls a fact, that feeling is \emph{effected} by motions in the
nervous system that are propagated to the brain, and to the observation
that the brain is the ``seat'' of feeling, it thinks it has furnished
the proof that a material organ factually has the capacity to be able
to feel, in general to perform acts of consciousness, and then
concludes that it must have it because feeling is a property of
matter, either of matter in general, which is supposed to be able to
have it as ``dead'' or ``living'' feeling, or of particular
combinations of the material. But whichever of these alternatives one
chooses, one gets no explanation of feeling, but arrives at it only by
a postulate. If we take the first alternative: that matter as such
should be sentient, then one ascribes to the matter for which one could
not even vindicate motion a predicate that has no point of attachment
in the ``primary qualities'' of matter, and directly conflicts with the
essence of the material so determined. Feeling, in which the affection
of the organism is actively transformed, presupposes a self that feels
\emph{itself}, an active subjectivity that in its feeling as it were
doubles itself in an ideal manner. But the material is the passive,
the selfless, that has its being in the form of objectivity, the
atomistically single, \opage{74}in which such an ideal doubling, such a
unity of subjectivity and objectivity, cannot be thought. To obtain a
point of attachment for the notion of such a doubling and unity one
would in any case have to go beyond the single material thing and
posit a relation between different parts of matter. This would then
lead to the second alternative: to posit feeling as the result of the
combination of the parts of matter. But this way out too materialism
cuts off for itself. ``Matter,'' whose name gives the semblance of a
unity, is, strictly speaking, only a common designation, hence a
subjective abstraction, and the real that is comprehended under this
common name is the ungenerated and unalterable molecules, of whose
various combinations the forms of nature are supposed to be formed. But
these forms in their apparent unity are only aggregates, compounds of
smallest parts independent in and for themselves. These parts can, as
aggregated, act reciprocally on one another, i.e.\ communicate to one
another their communicated motion; but the stream of motion goes
constantly outward from each part; all activity becomes, as
communication of motion, an acting upon another, a relation to another;
a relation to itself, a reflection into itself in activity, which would
have to correspond to that self-doubling, is excluded by the
presuppositions. Materialism does indeed ascribe to the brain the
property of being able to ``se replier sur lui-même,'' of being able to
observe its own changes. But in this expression an ambiguity hides
itself. Strictly speaking, only this can be understood by it: that one
of the brain's molecules can act on another, and this other thus be
determined in its being by the first; and only as one substitutes, for
the aggregate of independent molecules, whose reciprocal influence can
never become a relation of selfhood, the concept of the whole and of
subjective unity, surreptitiously obtains the view of the organic
totality, and thereby unconsciously goes beyond the bounds of
materialism, is the semblance of that possibility produced.

\opage{75}Not even the elementary form of the self: feeling, is
materialism thus able to explain, because for its explanation it has
only that whose concept is formed precisely by excluding the
determinations of the self. And the explanation of this elementary form
of the self's activity, appearing already in animal life, would still
be the necessary, if in and for itself insufficient, condition for the
explanation of the specifically human functions of spirit: the
functions of consciousness. But when consciousness now stands as an
irrefutable fact that cannot be denied without immediate
self-contradiction, its possibility becomes an unavoidable problem, and
the impossibility of being able to solve it becomes decisive for the
judgment on the principle of materialism, which proves to be
insufficient and one-sided. And since the spiritual, which cannot be
set aside, stands unexplained, it unavoidably calls forth
inconsistencies in the system, by which a validity independent of the
material is involuntarily granted to it. This shows itself in the
Système de la nature both when, despite its sensualistic theory of
knowledge, it \emph{dogmatically} lays down what can never become an
object of sensation, and uses concepts that cannot be given by the
relation of experience to natural things and thus point to something a
priori; when, in its conception of history, it lets the spirit that
itself ranks under the laws of nature and the power of necessity have a
perfectibility and a progressive development that defies the
tautological cycle of nature, and reach its goal through an
enthusiastic struggle for freedom; and when finally, in the domain of
morals, it lets disinterested virtue and humanity gain a place within a
view of life whose fundamental stamp, according to the principle, could
only be that of self-love striving after sensuous enjoyment.

The fundamental concept of materialism has thus proved to be defective
and incapable of being carried through. It does not overcome dualism,
but only covers over dualism, inasmuch as, while it is to express the
whole, it in reality expresses only that \opage{76}which had been
posited as the one member of the dualistic opposition. The other member
of the opposition is not reconciled with this one, so that what is
determined as the whole would have become richer in content; but it is
pushed aside, cast away, and since it cannot be annihilated for
consciousness, it forces its way to a place anew and makes existence
divided. Materialism has its relative justification so long as it
points out the emptiness in those concepts of the spiritual that express
only negations of the natural,
% errata list B, p. 76: printed "Naturlige", corrected "det Naturlige"
and so long as, against dualism's splitting of existence and of the
human being, it asserts the impossibility of a connection and an
interaction between the \emph{absolutely} different. Its justification
ceases when it positively determines its principle, and now, with
denial of spirit, the only thing that is immediately certain to us,
posits matter, which we know only at second hand, by means of the
knowledge of the denied spirit, as the sole being.

\divrule

Both terminal points of the realistic line of development thus point
beyond themselves to a higher form of thought. The problem of knowledge
demands a deeper solution than it received from the Humean skepticism,
and the problem of the principle of existence, which materialism
answered one-sidedly, demands a solution that, in reconciling dualism,
recognizes the essential peculiarity of spirit and of nature. But
despite its defects realism preserves an abiding truth through its
principle of experience, whose comprehensive significance was developed
above (p. 34 f.), and through its energetic assertion of the
essentiality of nature.

\divrule

\opage{77}
\begin{center}
  \textbf{Chapter Three.}
\end{center}

\divrule

\begin{center}
  {\large The Idealistic Line of Development.}
\end{center}

\divrule

\begin{center}
  I. \emph{Dogmatic Idealism}.
\end{center}

\begin{center}
  \textbf{1. Descartes.}
\end{center}

\emph{René Descartes} was born in 1596 at La Haye in Touraine.
Part of his childhood he spent in fanatically Catholic Brittany; from 1604---1612 he attended the Jesuits' educational institution at La Flèche, and thereafter lived for a time in Paris, at first eagerly taking part in all the pursuits of the young nobility, later in complete retirement, occupied essentially with the study of mathematics, the only science that satisfied him by its evidence. From his 21st to his 25th year Descartes was in foreign military service, first under Maurice of Nassau, later under Maximilian of Bavaria, Tilly and Boucquoi. His aim in this, however, was not to train himself as a soldier, but to acquire knowledge of men, to acquaint himself with the peculiarity of foreign nations, to free himself from prejudices and to clear his sight. After leaving military service he continued his traveling life for several years and came to know almost the whole of civilized Europe. He was constantly occupied with this task: by fusing the method of logic with that of mathematics, to find a new method of ``deduction,'' by which the individual sciences and philosophy could be given the same certainty that mathematics had. In order to be able to work undisturbed on this task, whose solution he long attempted in vain, he chose to live as a private man and abroad. For 20 years, from 1629 to 1649, he stayed at various places in Holland, where his main works were published. They brought him a great name, but made him the object of fanatical attacks from Ca\opage{78}tholic, and still more from Protestant, theologians, who accused him of atheism and the like. At the invitation of Queen Christina of Sweden Descartes traveled to Stockholm in 1649, to found an academy of sciences and to give the Queen instruction in philosophy; but already at the beginning of the following year he succumbed to the harshness of the climate. ---
Of Descartes's writings we mention, in accordance with the aim of this work, only the philosophical ones, and merely touch on the fact that in the domain of mathematics he created a new science, analytic geometry, that he enriched physics with considerable discoveries, and that he was distinguished as an anatomist. Of his philosophical writings the first to appear was: Discours de la methode pour bien conduire sa raison et chercher la vérité dans les sciences, 1637. Four years later appeared his Meditationes de prima philosophia, the work he himself valued most highly, and in 1647 it was published in French with additions and corrections. In 1644 appeared his principia philosophiæ, divided into 4 sections, of which the first treats of the principles of knowledge, the second of the principles of material things, the third of the visible world, the fourth of the earth. A couple of years later Descartes wrote for Princess Elisabeth of the Palatinate the treatise des passions de l'âme, which was published in 1649. After his death there appeared in 1664 his traité de l'homme, and in 1701 two works, presumably dating from his earlier period: Regulæ ad directionem ingenii, and: Recherche de la vérité par la lumière naturelle (in Latin translation). Of Descartes's letters published after his death many have significance for his philosophy. --- Descartes's writings were placed on the Index librorum prohibitorum in Rome in 1663. Already earlier, in 1656, the Synod of Dordrecht had ordered that his views were to be discussed neither in writings nor in disputations, and the Delft resolutions of 1657 had denied his adherents access to ecclesiastical offices. The University of Paris forbade in 1671 that his doctrine be lectured on. These clerical precautions did not, however, prevent Descartes's philosophy from spreading rapidly far
\opage{79}beyond the borders of his fatherland, and from founding a new, independent development of thought. ---

\divrule

\emph{Descartes} agrees with Bacon in the complete breaking loose from tradition, in the striving to find a sure method, and in the direction toward reality, from which thought is to draw a content of truth. But Bacon's method becomes that of induction, which proceeds from the individual given, Descartes's that of deduction, which, from the first sure point of departure won in the self, would arrive with mathematical stringency at the remaining content of thought. And Bacon's thought seeks outward toward the reality that is given in nature with its unchangeable laws; he is mistrustful of the mind's reflection on itself, because here it must at once be reflecting and master its reflection; he fears that it will, ``like the spider, spin threads out of itself, admirable for their art, but of no use for anything.'' In Descartes, on the other hand, the mind seeks inward toward the reality it finds in itself as the thinking self-consciousness, the thought present to itself. In itself it seeks the source of truth; it therefore becomes present in its content of thought, and claims for itself the right of decision over what is to be truth.

To that point of departure Descartes lets thought return, as it, through doubt about everything (``de omnibus dubitandum est'' [everything is to be doubted]), separates from itself everything that is not its very essence, in which it does not \emph{immediately} recognize itself and have an unshakable certainty. Doubt extends not only to the testimony of our senses about the existence of the external, the material, the existence of our own body included, but also to the testimony of our reason, to the content of the sciences that occupy themselves with the general determinations of the material: extension, magnitude, number etc., indeed even to the simplest mathematical propositions, and, in a word, to all the concrete content of thought, since I can always imagine that an almighty being has
\opage{80}so arranged my reason that I am constantly deceived in its use. But as doubt thus separates off all this content of consciousness, there remains One thing that it cannot reach, and that is the very fact \emph{that I doubt}. But the activity of doubt is an activity of \emph{thought}, and as thought through doubt seems to deprive itself of everything, it finds itself, and the I, which is posited as synonymous with thought (consciousness) and excludes everything else from its essence, finds in itself, as thinking, being and certainty. The proposition: de omnibus dubitandum est, leads to the proposition: cogito, ergo sum, i.e.\ I think, and as thinking I am. Here is the irrefutable point of departure of knowledge; here thinking and being coincide: the thinking I's knowledge of itself is immediate certainty of being, but a certainty that does not yet reach beyond the I, determined as a res cogitans, as ``a being that doubts, thinks and wills, that has notions and feelings.''

At this point, where the first certainty has been won, there then arises the question how this first abstract certainty of the I about itself can be filled out and become a knowledge, and in particular a knowledge of the external, which the I in doubt excluded from itself. Descartes, having first laid down that every knowledge that is to have validity must have the same certainty, the same clarity and distinctness as the point of departure, takes the path of seeking what is found in our thought as an innate store of ideas, and, in order to get beyond the chasm the I has set between itself and the external, of procuring for himself through the concept of God a guarantee for a knowledge of the latter. We find, as we search our consciousness, beside the notions that present themselves to us as called forth from outside, and those that we ourselves form arbitrarily, other notions that bear the stamp of being \emph{innate}, e.g.\ the notion of the thing, of truth, of thinking, which we draw from our own being. But among these innate ideas is also the idea of a God, of an infinite and all-perfect being: an idea that readily shows itself to be the expression of something actually existing. For since
\opage{81}this idea, expressing an infinite reality, is present in the thought of a finite being, it follows with necessity that it cannot be a product of that thought, since the effect cannot have greater reality than the cause, but that it must have its origin from an infinite, actually existing being. It is wrong to take the concept of the infinite as merely expressing an abstraction from the finite, a negation of it; it is on the contrary a purely positive concept, which includes a higher reality than the notion of the finite, and is the presupposition of the latter, since I grasp something as finite and limited only through the thought of infinity. --- A proof of God's existence lies also in this, that I am conscious of myself as finite, as not existing through myself; for my existence then points back to a first author; the notion of this author, the divine, is just as original in me as my consciousness of myself. It is like a ``trademark'' that God has set on his work, and which perhaps is not something different from the work itself. For since God has created us, it is credible that we are created in likeness to our cause, and that the likeness to God, in which the idea of God is included, is known by me through the same thought by which I know myself. --- In the very idea of God as the infinite and absolutely perfect being we can, finally, \emph{immediately} find the certainty of God's existence, for it follows with the same necessity from his concept as it follows from the nature of the triangle that its angles together are = 2 R. God's essence and his existence are inseparable: ipse suum esse est [he is his own being]. It is properly not we who carry out a proof of God's existence, but God who proves his existence to us. --- Since God's existence as the infinite, almighty, all-perfect being has thus been proved, Descartes singles out among his perfections \emph{veracity}, which must belong to God, since the opposite would express a defect in understanding or will, and in God's veracity he now in turn seeks, by an unconscious circular argument, the guarantee for the truth of our clear and distinct knowledge.
\opage{82}Originally evidence was the criterion of truth; it was the self-confirmation of thought, its immediate certainty of itself in its content. But since the \emph{doubting} thought is taken as a \emph{finite}, imperfect thought, its very clear and distinct knowledge must be guaranteed, and it is then guaranteed by God's existence, which had obtained its certainty precisely through the clarity and distinctness with which it was thought. The source of our errors is now sought in this, that our will has a greater scope than our knowledge, and that it leads us to pass judgment even where insight is only imperfect.

By means of the guarantee that is found in God's veracity, our notions of the external, the material, which before had been made doubtful, now within certain limits regain certainty. We know clearly and distinctly that the affections of our faculty of sense presuppose an external and bodily object, and we therefore infer with certainty the actual existence of the material. The objective essence of this material is to be judged by thought, not by sensation. From the determinations we are accustomed to ascribe to the material we must separate off those that we ascribe to it only through a confusion of the sensuous impression with external existence, such as the determinations of color, tone, taste, smell, hardness, weight and the like, which are just as subjective as the feeling of pleasure and pain. What we clearly and distinctly know as an objective determination of the material is the determination that \emph{thought} states about it in apprehending it as its opposite: \emph{extension}, the expression for the abstract quantity posited \emph{outside} thought. The distinguishing mark of the extended becomes, in opposition to the activity of thinking, to be the purely passive, and its modifications are restricted to those that arise through divisibility, receptivity to form and mobility. The reality of motion cannot be derived from passive extension; Descartes must seek it in the author of matter, in God, who originally called forth motion, and in accordance with his immutability constantly lets the same quantum of motion be maintained in the universe. All motion
\opage{83}here becomes the motion communicated from body to body by pressure and impact. On the basis of these determinations of the material Descartes now builds, with consistent exclusion of final causes in favor of efficient causes, a \emph{mechanical physics}, in which everything is explained by the laws of communicated motion in their application to what is determined by magnitude and figure, and by means of these laws he finds, in his hypothesis of the double motion of the homogeneous parts of matter originally separated by God: each single part about itself, or several in a vortex about a center, and of the new division into heterogeneous elements that arises from the friction of the first parts against one another, an explanation of the phenomena of natural existence in the large and in the particular. The restrictions that are made at the end in favor of the authority of the Church are only a formal tribute that Descartes pays to the Church in whose bosom he was brought up, in whose outward customs he still takes part, even by a pilgrimage to Loreto, and with which he avoids conflict as far as possible; but from which his thought has freed itself, a thought that independently strives after a truth unfolded from the very essence of the human mind.

As matter is known clearly and distinctly as the really existing, separated from mind, and as at the same time mind is clearly and distinctly known as that which, in its determinateness as thinking, does not include the notion of the bodily but excludes it from its essence, the extended and the thinking stand over against each other with independent existence. They are therefore determined, in this mutual independence, as the extended and the thinking \emph{substance}, substance designating ``a thing that exists in such a way that it needs nothing else for its existence.'' In the strict sense, however, this definition of substance can apply only to one substance: the divine substance; the two mutually opposed substances are indeed independent of each other, but need God's concurrence in order to exist, and they then become substances in the sense that, in order to exist, they need nothing other than
\opage{84}God. \emph{God} is conceived as producing them both, but as akin only to the thinking substance, as \emph{pure mind}, without such forms of consciousness (modi cogitandi) as, like sensation, point to a connection with a body. The property of the bodily, divisibility, is an imperfection that makes it unfit to be a determination of the divine.

In God, the author of the substances, the explanation is now also sought of the \emph{union between soul and body}, which facts of experience compel us to acknowledge in man. Such facts as feelings and sensations, which seem explicable only by the soul's being affected by something material, or as the voluntary movements of our organism, which show it to us as determined by the soul, guarantee us a real connection. But according to the concept our body and our soul have nothing to do with each other. The body is a mechanically moved machine, not essentially different from those that men form, if more artful and more composite. It can be designated as an automaton formed by God, in that, once set in motion, it continues to move itself, but only according to mechanical laws, and without possessing peculiar vital forces or the capacity to exercise any psychical function. All forms of thinking and of consciousness, feeling included, belong to the soul, from whose essence extension and motion are excluded. Where we have only a body, as with the animals, in which, since they do not think, we cannot assume a soul, there we have only matter, and it is only apparently that feeling takes place in them. If one strikes an animal so that it cries out, the cry is not the expression of a feeling, but the animal cries out, just as the string sounds when it is struck, by a purely mechanical effect.

The connection of soul and body in man, which could not be explained by the nature of the merely heterogeneous substances, Descartes must then posit as brought about by the divine
\opage{85}omnipotence, which as it were by force brings the heterogeneous together. But the connection can, precisely because it is a connection of the contradictory, become only a \emph{composition}, and it becomes \emph{incomprehensible} to thought.

In the connection into which soul and body are brought with each other, they are supposed to be able to act on each other \emph{directly}, by an influxus physicus, so that the body's affections bring about corresponding states in the soul, and the soul's acts of will corresponding motions in the body. The body's affections Descartes conceives as propagated, by means of the nerves and of the currents of the animal spirits, a fine ethereal fluid secreted from the blood, to a small part of the brain, the pineal gland (glandula pinealis), which he makes the soul's special ``seat,'' and which is supposed to be suited to this as the only single part of the brain, hence as the only one that can gather the sense impressions into a single image; and from this seat, in which the soul apprehends the communicated motion, it is in turn, by determining the direction of motion of the animal spirits through an influence on the gland, to determine the motion in the bodily organs. Through the soul's indwelling in the body new motions thus come into the organism, but they are produced mechanically, like the motions that have their source in the body itself.

With this doctrine of the connection and interaction of the substances in man Descartes has, however, not been able to avoid coming into contradiction with his fundamental determinations. It lies in the concept of substance that it is self-subsistent and independent, hence complete through itself (``substantiæ nequeunt esse non completæ'' [substances cannot be other than complete]); but since the two substances in man are to be composed into a unity, they must give up their self-subsistence, and become in this relation ``substantiæ incompletæ.'' The essence of the psychical was determined by thought alone, with the exclusion of extension. But since the soul, as connected with the body, is to enter into contact with the material, it acquires a place in space, a ``seat'' determined by space, and the determination of extension is not
\opage{86}avoided by reducing this space to a minimum. Like extension, motion is the property of the material and excluded from the soul. But since the soul is united with the body, it is to be conceived as affected by the body's motions, and as itself determining a motion, which, according to the fundamental determinations of Descartes's physics, would presuppose that it was itself bodily and moved (``corpus non movetur nisi corpore contiguo moto'' [a body is not moved except by a contiguous body in motion]). Finally the body, whose essence it is to be divisible to infinity, is, in the union with the indivisible soul, to be ``in a certain sense indivisible'' as an organism, --- a thought, however, that Descartes does not carry through.

\divrule

With Descartes's philosophy, as with Bacon's, a new and independent main form of philosophy is founded. In him thought, with the consciousness that a wholly new beginning is to be made, that all given, untested presuppositions are to be cast away, seeks back to self-consciousness as the source of knowledge. The I, whose essence is thinking, would be present in its thinking; it would seek in the very evidence of thought, in the clarity and distinctness of knowledge, the criterion of truth. And from the first sure point that is won, it would derive the whole content by a sure method, which at every point preserves evidence through stringency in the connection of thoughts. It would arrive from the certainty of the self's being as thinking at a knowledge of reality, of what essentially is in its necessity; and in setting itself this task it lets thought, which has found itself in self-consciousness, turn immediately --- for the detour through God is only a detour through the infinite \emph{thought} and through what is guaranteed by our very thought --- toward what is, without critical testing of the relation of thought to objective being. The power of thought to know this being is presupposed \emph{dogmatically}.

In the carrying out of Descartes's principle a defect soon shows itself. His method only links together a given
\opage{87}content, but cannot unfold a content from the abstract point of departure: the I as res cogitans. By an analysis of the I as thinking one can at any rate arrive at concepts that express the subject of thinking and the act of thinking, insofar as it has the subject itself as its object, and at purely \emph{formal} logical propositions. To get beyond the I, Descartes must partly take up empirically notions which, like the notion of God, are posited as innate to the I only by a postulate, partly ascribe to the merely thinking I determinations such as sensation and feelings, which become incomprehensible on his dualistic presuppositions. Only when by means of these he gets beyond the chasm that separates the I from nature does his thought acquire a concrete content, and it now dwells with predilection on this natural, which, however, by being excluded from the I and from the divine, is posited as something lower. But the attraction that nature exercises over Descartes's thought points to something universal in the tendency of the modern age, which asserts itself in idealism as well.

The defect shown in Descartes stands in close connection with another main defect in him: the dualistic separation that is posited in his system from the beginning, and that is never sublated. For there is indeed talk of the union of the opposites in man, but the union becomes only the incomprehensible composition of the absolutely heterogeneous by an act of omnipotence that eludes thought; and God is indeed posited as the common author of the thinking and the extended, but the origin of matter from the God determined as pure mind becomes the incomprehensible, and the dualism therefore persists despite the notion of the common author. And as the divine substance, which excludes matter from its essence and as creating stands outside it, is also separated in being from the finite minds that are its creatures, it becomes in its infinity a third \emph{beside} the thinking and the extended, which also appropriates the name of substance and brings an ambiguity
\opage{88}into its concept, and there arises a new dualism between God and created existence.

In the dualistic conception of what is to be united into unity in the human being, and in the notion of finite substances subsisting beside the infinite and, according to the concept, only substance, lie the contradictions that become incentives for the further development of dogmatic idealism. The contradiction in the ideal and real connection and mutual influence of the absolutely opposed finite substances leads, as the dualistic conception is retained, in Descartes's school, in Geulincx and Malebranche, to seeking another explanation of the agreement between the psychical and the organic processes, an explanation that makes God the efficient cause, body and soul only occasions (causæ occasionales) for the divine activity. And the contradiction in finite substances being supposed to have a subsistence beside the infinite one brings Spinoza, by consistently carrying through the thought of the infinite substance, to let the finite substances vanish in it as its attributes, and their external opposition become an inner difference in the infinite One that comprehends all existence.

\divrule

\begin{center}
  \textbf{2. Occasionalism (de la Forge, Geulincx, Malebranche).}
\end{center}

The occasionalist consequence of Descartes's conception of the relation
between the two finite substances was first drawn by \emph{Louis de la
Forge} in his traité de l'esprit humain (1661). From the fact that the
essence of mind is thought, consciousness, he concludes that the
sensations and notions that arise involuntarily in mind in consequence
of its relation to the body cannot have their origin from mind itself,
since it would then have to be conscious of this. But as modifications in
a mind they cannot come from the corporeal either, which is absolutely
different from mind. They must then come from God, who is mind, and
the motions of the body become not the cause of the sensations and the
\opage{89}thoughts, but only the condition, conditio sine qua non, for
them, just as the will of mind becomes only the condition for the
motion in the organs. A similar consequence was drawn by the German
\emph{Clauberg}, who in particular accentuated it insofar as the body, as
the lower, is supposed to act upon the soul, as the higher.

\divrule

In a more fully worked-out form and with greater effect the occasionalist
doctrine was developed by \emph{Arnold Geulincx} (1625---1669. Writings:
Logica 1662, \gk{Γνῶθι σεαυτόν} [Know thyself] sive Ethica 1665,
Metaphysica 1691, Physica, 1698.). Descartes had, in order to support his
doctrine of a direct relation of interaction between soul and body,
appealed to experience, which was supposed to show that certain bodily
states were effects of certain mental causes and conversely. But
against this appeal to an experience Geulincx raised objection; he urged
that by it one confused what was fact with what was only an
interpretation of the facts. What was a fact of experience was only that
those states accompanied one another, or immediately followed upon one
another. Experience, on the other hand, does not say, and cannot say,
that they are causes of one another; this is precisely only an
interpretation of what is experienced, and an unjustified interpretation
besides. It shows itself as unjustified as soon as one holds fast to the
fundamental concepts, and the relation is then seen in another light.
This is easily seen on a closer consideration of the facts appealed to.
When, e.g.\ a motion of an arm or a foot takes place in temporal
connection with a resolution of the will, then a causal connection is
not only not conscious to me, but it is altogether inexplicable to me.
The will is an activity that arises and takes place solely in the soul;
the motion is a material change, which must have material causes; it is
therefore a process of an entirely different kind, and it is not
possible to say how the psychical process, which has no determination
of matter and motion, should be able to pass over into the bodily one.
It is just as impossible to under\opage{90}stand how I should be able by
my will to set my arm or my foot in motion as how I should be able by my
bare will, without mechanical causes, to move a foreign body. For even
if my soul stands in a different relation to my own body than to other
bodies, the relation still cannot be this, that it touches it, or in a
mechanical manner pushes and presses it, and thereby moves it. If I,
that is my soul, were really the origin of the bodily motion, then,
since the essence of my soul is thought, knowledge, I would have to be
knowing not only \emph{that} but \emph{how} the effect was produced; but
this, as is well known, is not the case. In an entirely corresponding
manner it stands with the opposite series of data, the bodily causes
that are supposed to be causes of psychical states. There is no
difficulty in understanding that external objects act through motion
upon our sense organs, and that the action upon these propagates itself
by means of the nerves through the body into the brain: here we are
constantly in the domain of the corporeal. But we can in no way
understand how the communicated motion, because it is propagated to the
brain, should therefore be able to be continued into the soul. Here all
comprehensible connection ceases; there opens, as it were, an abyss
between the heterogeneous. The notion, both as individual notion and as
general notion, shows itself as something entirely other than the
impressions on the sense organs and the motions, just as the sentient
being is something entirely other than the bodily organs. But when the
states of the soul and of the body thus cannot be causes of one another
or effects of one another, and there is nevertheless a factual
connection between them, then the ground of this connection must be
sought outside the opposites, in the divine cause of the soul and the
body, which has joined them together in our being and which constantly
sustains them. God has placed the two constituents of our being, in and
for themselves altogether different, in such a relation that my body
moves when I will, and, conversely, \opage{91}my soul has sensation when
my body is affected, without there taking place between them, however,
any direct causal connection whatsoever. What takes place in the body or
in the soul becomes only an \emph{occasion} for God to bring about
something corresponding on the opposite side; God alone is in this
relation the efficient cause, causa efficiens; it is he who in an
incomprehensible manner, by a miracle continued at every moment, keeps
body and soul, like two machines, in agreement with each other. Body and
soul are only occasional causes, causæ occasionales.

Not only by this consequence of Descartes's fundamental determinations
does Geulincx depart from him, but on other points too, and he here
approaches Spinoza so closely that only the presuppositions of the
religious notion hold him back from reaching the latter's standpoint. As
he strictly holds fast to the concept of the infinite, he arrives at the
proposition that the finite must be thought in and through the idea of
God, by whose limitation of his perfections it has come to be. God is
not limited by the finite, for all the perfection that the finite
contains in itself God possesses in the same or in a more excellent
manner (formaliter or eminenter). The human mind becomes a modus of
the divine Thought (mens), hence of God himself; as pure thought it
thinks not in the brain but in God. In him the minds have their union,
and seen from this point of view they are not creatures. They are
dependent on God in all their being, thinking and willing. Extension
Geulincx will not make a predicate of God, for the reason that it is a
res bruta, omni destituta cogitatione, but not because it should be
divisible, since, as infinite, it is indivisible in the strict sense.
But extension is nevertheless supposed to be eminenter, in a higher
form, in mind, from which, as purely passive, it must receive its
motion. Since extension is indivisible, what we call individual bodies
become only modifications of extension (matter). Extension becomes a res
singularis, and there are properly \opage{92}only two substances: a
creating one, God, and a created one, ``the corporeal \emph{world}'';
otherwise there are only modes of existence. The creation Geulincx
posits as something incomprehensible, like the agreement between the
states of the heterogeneous substances united in us; and to the domain
of the incomprehensible he must also reckon the relation between God's
goodness, which demands the moral order and abhors all sin, and God's
power, which, as effecting existence, effects everything, the sinful act
too, as well as the freedom and independence which, in direct
contradiction with his doctrine of man's absolute dependence on God, he
nevertheless seeks to vindicate for the individual, so as not to make
God the author of sin.

\divrule

In \emph{Nicolas Malebranche} (1638---1715. Writings: de la recherche de
la vérité, 1675; Conversations métaphysiques et chrétiennes 1677; Traité
de la nature et de la grâce, 1680; Traité de morale, 1684; Entretiens sur
la métaphysique et sur la religion, 1688; Entretiens d'un philosophe
chrétien et d'un philosophe chinois sur la nature de Dieu, 1708) the
occasionalist theory is developed still further. He goes further in
that he partly sharpens the opposition between mind and body so that
not even the thought of the body can proceed from the human mind,
since this mind, as the absolute opposite of matter, cannot have
matter in itself even in an ideal manner, and partly determines the
relation between the infinite and the finite so that the latter, as
something created, receives no principle of activity whatever in
itself, and so that the finite mind, as individual, can develop no
general knowledge out of itself. For from these determinations it
follows, first, that not even body in relation to body, or the created
mind in relation to mind, can be efficient cause, but can only
become occasional cause, while God is everywhere the one active agent in
mind and in nature. The body cannot even modify the direction of
motion, while mind is nevertheless granted the ability to influence
that of desire. --- Second, it follows from this that \opage{93}we,
since our mind is absolutely opposed to matter, and the material is
absolutely obscure to us, must intuit things in God, of whose infinite
being the ideas of all that is, i.e.\ its essence as God intuits it,
are only limitations. We behold the ideas of things in God, and to these
God attaches for us sensuous feelings, which are not found in him, in
that he lets the material become occasional cause. --- Finally it
follows from this that we must gain in God the general (universal and
universally valid) knowledge which is the presupposition of the
knowledge of the individual and finite, and which we, as individuals,
cannot seek in ourselves. God becomes the universal reason (la raison
universelle) and the place of minds. He is the principle in all our
knowledge and in all our willing. The final general end of this willing
is unalterably fixed. God himself is this end, and he himself constantly
awakens in the created the drive toward himself as a natural drive
toward the general good, i.e.\ blessedness, which is complete only in
God. ``All beings, even the Devil, love God''.

In God material nature is as if idealized: the idea of extension is in
God. In his last writings Malebranche comes close to placing real
extension too, according to its perfection, in God. Hereby he
approaches Spinoza's conception of the divine substance as comprising in
its infinity thought and extension. And Spinoza's conception of the
relation of the finite minds to God he touches upon when he states
that the ideas God has of the creatures, among which are also found
ideas of our souls, are God's very essence. From this it would follow
consistently that God, as the infinite reason, comprised the finite
thinking minds as his modi. But on this point, as in the determination
of the relation of extension to God, Malebranche's religious notions
prevent him from drawing the final consequences. The finite,
\emph{created} minds are to retain a subsistence as finite, and to
receive only that mystical contact with the divine in their knowledge,
and matter is to come to be through \opage{94}an incomprehensible act of
creation, through an arbitrary effect of the will of God, in which the
material nevertheless is only in an ideal, i.e.\ immaterial manner. The
one who, unhindered by the power of religious notions, draws the whole
consequence of Descartes's fundamental determinations is, for the first
time, Benedict Spinoza.

\divrule

\begin{center}
  \textbf{3. Spinoza.}
\end{center}

\emph{Baruch (Benedictus) de Spinoza} was born in Amsterdam in 1632. His
father was a merchant who, on account of religious persecutions, had
left his fatherland, Portugal. The son was early destined for the
learned estate, and had the famous Talmudist Morteira as his teacher.
Under this man, who in scholastic fashion sought a union of Judaism and
philosophy, he came to know several of the Jewish philosophers of the
Middle Ages, in particular Maimonides, Levi ben Gerson, Ibn Esra and
Chasdai Crescas, as well as kabbalistic writings, and there are points
of contact with these in his philosophy. Yet his concept of substance,
which carries his whole doctrine, has been formed, in its peculiarly
stamped form, through the critique of the determinations of the
Cartesian philosophy. Already early Spinoza felt unsatisfied by the
doctrine of the rabbis and estranged from the notions of the Jewish
religion. In vain his teachers sought by authoritative pronouncements,
and the Jewish congregation by offers of money, to lead him back to the
synagogue; in 1656 he was expelled from the Jewish congregation, and
through the influence of its elders banished from Amsterdam. Already
before this happened he had sought, outside the narrow circle of Jewish
science, to procure for himself the means to a richer scientific
development and to humanistic culture. From the freethinker, the
physician Frants van den Ende, he learned Latin, and through his friend,
the physician Ludvig Meyer, he was led into the study of the natural
sciences and of Descartes's philosophy; and philosophy now, once the
decisive break with his religious community had taken place, became the
occupation of his life. In the various places where he lived
\opage{95}after his banishment from his native city: first in the
vicinity of Amsterdam, with a friend who inclined to the doctrine of the
Arminians, then in the village of Rijnsburg near Leiden, after that in
Voorburg near The Hague, and finally in The Hague itself, he worked
uninterruptedly and energetically toward the goal he had set himself:
the knowledge of God and nature and of man's unity with the infinite,
and toward the calm elevation above the human passions which he saw as
conditioned by the clarity of knowledge. The means for his frugal
livelihood he earned by the making of optical lenses, in which he
possessed a great perfection. An offer of a philosophical professorship
in Heidelberg, which was made to him by the liberal-minded Elector Carl
Ludvig, he declined so as not to expose himself to unavoidable
conflicts. On the 21st of February 1677 he died in The Hague of
consumption.

Of his writings only two appeared in his lifetime. The first is: Renati
Des Cartes principiorum philosophiæ pars I et II, more geometrico
demonstratæ, to which the Cogitata metaphysica are joined as an
appendix. In this writing, published in 1663, which was in part worked
out for the use of a pupil whom Spinoza was instructing in the Cartesian
philosophy, he has, in the course of expounding that philosophy,
indicated the points where difficulties that are insoluble from the
standpoint of Cartesianism make necessary consequences that lead on to
Spinoza's own, and in the whole design of the exposition the way has
been prepared for his own conception. The second writing is the
\emph{Tractatus theologico-politicus}, published in 1670 and
epoch-making for the historical-critical treatment of the Bible, in
which he seeks to make room for a free philosophy by drawing, through a
demonstration of the different aims of positive religion and of
philosophy (piety and obedience, and truth and wisdom), a boundary
between the two that secures for philosophy a neutral domain, and by
demonstrating the compatibility of free thinking and free speech with
the aim of the state and with its peace. In the same year that Spinoza
died his Opera posthuma appeared, which, \opage{96}besides his main
work, the \emph{Ethica}, ordine geometrico demonstrata (in 5 sections:
de deo; de natura et origine mentis; de origine et natura affectuum; de
servitute humana seu de affectuum viribus; de potentia intellectus seu
de libertate humana), contain the two unfinished treatises: the
\emph{tractatus politicus}, in quo demonstratur, quomodo societas, ubi
imperium monarchicum locum habet, sicut et ea, ubi optimi imperant,
debet institui, ne in tyrannidem labatur, et ut pax libertasque civium
inviolata maneat; and the \emph{tractatus de intellectus emendatione},
et de via qua optime in veram rerum cognitionem dirigitur, as well as 74
\emph{letters} from and to Spinoza and, as an appendix, an unfinished
Compendium grammatices linguæ hebrææ. The Ethics, on which Spinoza
worked right up to his death, had already for a number of years been
essentially ready for publication and had been communicated in copy to
% printed "være færdig" (defect logged in the transcription)
some of his closest friends; but Spinoza held it back so as not to
become once again the object of his opponents' fanaticism. --- In the
most recent times a treatise originally written by Spinoza in Latin,
\emph{on God, on man and his blessedness}, has been found in a Dutch
translation and published by van Vloten (Amsterdam 1862). This treatise
marks a stage of development in Spinoza at which he has indeed arrived
at the concept of the unity of substance, and of thought and extension
as attributes of the one substance, but at which his peculiar doctrine
of the mere parallelism between the modifications of thought and of
extension is not yet clarified. In some dialogues that are incorporated
in the treatise there shows itself an influence of Giordano Bruno's
doctrine of infinite nature.

\divrule

In the point of departure of the Cartesian philosophy and in its
determination of the concept of substance, the presuppositions are given
from which the fundamental concept of the Spinozistic philosophy, the
concept of the one, infinite, divine substance
com\opage{97}prising all mental and corporeal existence, must
consistently be developed when thought is not hampered by the influence
of religious notions. Descartes had taken his point of departure from
thought in the form of the I; he had maintained the subsistence of the
created thinking substance over against the divine, and had rejected the
problems that threatened its independence. But since the I was
determined only by thought, to the exclusion of the natural conditions
of individuality, the consequence had to be drawn that makes the
independent subsistence of the individual form disappear, and lets the
singularity of the I be absorbed in the universality and infinity of
thought. When the essence of the I is exhausted by thought, then,
according to the nature of thought, which is to be the thinking of the
universal, the eternal, the infinite, and therefore to have the
determination of infinity in its essence, the thought of infinity is the
first positive thought of the I, through which the finite is first
thought, as the \emph{limitation} of the infinite. In the thought of God
as the infinite thinking substance, human thought \emph{originally}
thinks its own essence; only in that thought does it first think it, and
in this thought it is in unity with the thinking infinite. What
Descartes, applying the concept of substance not merely to the mental
and the corporeal as wholes but also to the single, individual mental
and corporeal, gave the name of finite thinking substances must, when
these consequences are drawn, change its name and become forms of
limitation, modifications, of the one infinite thinking substance, and
have their reality only in it.

The same now holds also of what in Descartes was designated as finite,
\emph{extended} substances. Descartes determines them as created by God,
and as existing, in the same way as the thinking ones, solely through
his continued preservation; but the essence of extension is not posited
in God, who is to have only thinking as his predicate. But since the
cause can be conceived only as cause of what is grounded in its nature,
of that whose determination it itself bears within itself, God, if he is
to be thought as cause of the ex\opage{98}tended substances, must
receive extension, as infinite, as his determination, and the finite
extended substances must become mere modifications of the infinite
extended substance.

The concept of the infinite substance comprising all being is reached
also when the contradiction in Descartes is to be removed, that the
thinking and the extended substance should each by itself be conceivable
through itself, without going beyond its concept, and thus be
independent according to their essence, while with respect to their
existence they should be dependent on God, the only one to whom
Descartes's definition of substance fully applies. The contradiction
disappears only when what Descartes designated as created substances,
the thinking and the extended substance as wholes, loses its subsistence
for itself and becomes determinations of the one divine substance, its
attributes, the infinite forms of being of the infinite. Only thus is
the ambiguity removed from the concept of substance, and only thus is
the concept of infinity realized, since the divine is not limited by
something independent existing outside it. The one, infinite substance
then includes thinking and extension as its attributes in such a way
that each of them in its infinity expresses its essence wholly, and,
under the mutual independence that is conditioned by the concept of each
being thought without that of the other, they are nonetheless in unity
insofar as they express one and the same substance.

This concept of the one, divine substance comprising all mental and
corporeal existence becomes, in its form peculiarly determined by the
historical presuppositions, the alpha and omega of the Spinozistic
philosophy, and on its basis there is developed, for the first time
since antiquity, a distinctively formed pantheism hampered by no notion
of transcendence.

The \emph{first} thought of the thinking mind, that in which it thinks
its infinite essence, which immediately affirms itself to thought, is
then \emph{substance}. Spinoza defines substance \opage{99}as that which
is in itself and is conceived through itself. It \emph{is}
\emph{in itself}, not in another, as the modifications are in substance.
It \emph{is conceived through itself}, inasmuch as its concept is not
derived from a higher concept, as that of motion is from that of
extension, but is itself the highest concept. But when substance is in
itself and is conceived through itself, it must also be through itself,
be the \emph{cause of itself} (causa sui); for the concept of an effect
includes the concept of the cause; that, therefore, which includes only
its own concept can only be the effect of itself. In substance's being
causa sui there lies, again, that its existence is not contingent but
\emph{necessary}, given in its essence. This unity of existence and
essence is \emph{eternity}, and since eternity is the unconditional
affirmation of a being's existence, \emph{infinity} follows immediately
from it, while everything finite is a partial negation. But the concept
of infinity in turn includes the \emph{unity} of substance, both in the
sense that a plurality of substances, and in the sense that a
divisibility of substance, is excluded.

Since substance in its unity is determined by infinity, which does not
express the negation of a positive finite but the positive essence of
the negative finite, and, as causa sui, by activity, we are hereby led
to the two concepts that together with that of substance make up the
fundamental concepts of the Spinozistic philosophy: the concept of the
\emph{attributes} and the concept of the \emph{modi} (the
modifications).

Infinity excludes limitation, because limitation, by separating
something off from what is limited, negates. But limitation takes place
only through what is of the same kind, not through the heterogeneous.
Extension thus cannot be limited by thought, or thought by extension.
Infinity can therefore, without being cancelled as infinity, include in
itself as it were a new infinity, a multiplicity which, on account of
the difference in kind, does not entail limitation, and such a
multiplicity becomes, as compatible with infinity, compatible with
unity. It is even demanded by in\opage{100}finity, since everything that
expresses essence and includes no negation belongs to infinity. This
multiplicity of essential determinations, which without cancelling
infinity is posited within the unity of substance, the attributes
express. ``The \emph{attribute} is that which reason cognizes
\emph{as constituting} \emph{the essence of substance}.'' The attribute
is of one essence with substance; it is conceived, like substance,
through itself, and is eternal and infinite. Yet the single attribute is
infinite only \emph{in its kind} (in suo genere), since of each
attribute we can deny all the others; it is not
\emph{unconditionally infinite} (absolute infinitum), since to the
essence of the unconditionally infinite belongs everything that
expresses essence and includes no negation. This unconditionally
infinite being is \emph{God}, i.e.\ ``a substance that consists of
infinitely many attributes, each of which expresses an eternal and
infinite essence.'' In the concept God the concept of substance has
first attained its \emph{complete} expression. Substance is the
infinite; only in the unconditionally infinite, therefore, does it
attain its concept. The infinite number of the attributes is an
expression of the infinite reality of the divine substance. Each of
these attributes is conceived through itself, since none of them is
produced by another, but all are equally eternal expressions of the
essence of substance. But because they are conceived independently of
one another, they do not on that account constitute several beings or
several substances.

In this doctrine of the multiplicity of the attributes in their relation
to the unity of substance, Spinoza is determined by his historical
presuppositions. The concept of substance, as Spinoza defines it, leads
only to the attribute in its infinity without any limitation, not to
infinitely many attributes. From other historical presuppositions one
could have arrived at positing thought or extension as the sole
attribute. But Spinoza's presupposition was the dualism between thought
and extension, whose unity was demanded, but was brought about by
Spinoza, who maintained their mutual independence, only through their
being posited as mutually independent expressions of
\emph{the \opage{101}same} substance. And since with thought and
extension, the attributes that we come to know immediately through our
own essence, a \emph{plurality} of attributes is posited, for Spinoza
the infinite reality of substance must express itself in the
\emph{infinite multiplicity} of attributes, and God becomes the
substance consisting of infinitely many attributes. But just as little
as Spinoza has in reality deduced the multiplicity of the attributes
from the concept of substance, just so little has he, as will be shown
later, been able to maintain the unity in this multiplicity.

Substance, which in its unity includes the innumerable attributes, now
receives its further determination through the \emph{concept of
causality}. Substance is causa sui, whereby there is expressed not
merely the negative, that it is not the effect of another, but whereby
Spinoza positively means to designate God's essence as
\emph{activity}. When the determination causa sui receives a positive
meaning, there is posited in the identity of what effects and what is
effected a constantly vanishing, or rather: an eternally vanished,
difference, and through this difference as such we arrive at the third
Spinozistic main concept: \emph{modus} (modificatio, affectio). ``By
modus I understand a form of limitation of substance, or that which is
in another, through which it is also conceived.'' It is an
``affection'' of the attributes, by which these are expressed in a
determinate, limited way. Since the modus is determined as being in
another and conceived through another, it is at the same time determined
as posited by another, namely by that in which it is and through which
it is conceived. Determined as the effect of another, the modus is
\emph{different} from substance, and since there is at once talk of an
infinite multiplicity of modi, and this concept is brought together with
the concept of single things (res particulares), the finitude that was
absorbed by substance seems to emerge from it anew and, even if it is
posited within substance, to cancel its infinity, which seems to divide
itself into the unbounded multiplic\opage{102}ity of single things. This
would, however, be a misunderstanding of Spinoza's thought. For him
there is no question of a transition from the infinite to a finite. What
shows itself to us as a finite has no being \emph{as} finite; only
substance in its infinity is, and the very knowledge of the nothingness
of the finite does not come to be --- is thus not subject to the form of
finitude --- but \emph{is} in the manner of eternity.

Spinoza's conception of the causality of substance points at once to
this conception. God is the cause of all things in \emph{the same sense}
(eo sensu) in which he is the cause of himself. He is causa immanens,
not causa transiens, i.e.\ he is a cause that remains in the effect, not
a cause that in the effect passes over into being another. Causality is
in unity with the essence of substance; through causality this essence
realizes \emph{itself}. But since the effect of the causal substance is
substance itself, the infinite causality is exhausted in the infinite
effect; no room is left for an effecting of a finite.

Spinoza's conception of the determinations of finitude: limitation,
plurality, divisibility and number, leads to the same result.
\emph{Limitation} is for Spinoza one with negation; it therefore becomes
unthinkable when one proceeds from the infinite attribute; for since
limitation can take place only through what is of the same essence,
% errata list B, p. 102: printed "Væsentlige", corrected "Væsenslige"
it would have to take place through the attribute itself, and the
attribute, as what limits, would have to be conceived as the negation of
what is limited, hence as its own negation. The concept of the attribute
therefore excludes the limitation through which a finite gets being,
and when the limitation is seen from the standpoint of a hypothetically
assumed finite, it becomes, as something merely negative, immediately
something vanishing. --- \emph{Plurality}, according to Spinoza, does not
lie in the concept of the thing, but presupposes an external cause. If
one pursues this proposition backward, then, as one goes back to the
first cause, this cause, which is substance in its unity, will exclude
plurality, hence finitude. --- The determinations \emph{divisibility}
and \emph{number} lead to contradictions: to an infinite series
\opage{103}that can form no unity, to an infinity of parts that cannot
succeed in making up a whole, to a greater and a smaller infinity; and
these contradictions are removed only when the notion of the finite
single things is removed; for only the reality ascribed to these as
finite calls forth those contradictions.

But since Spinoza has brought forward the concept of modi, he must seek
new determinations that secure against this concept, despite the
contradictions that have been shown in the being of the finite,
nevertheless bringing a finitude into substance.
% errata list B, p. 103: printed "dog ikke bringe", corrected "dog bringer"
Such determinations are partly the concept of modi as apprehended under
the form of eternity and necessity, partly the concept of modi as a
whole, as a universal individual, as natura naturata; partly the concept
of infinite modi. When reason, the true form of knowledge, thinks a
modus, it apprehends it not under the determinations of time and space,
which are only determinations of the imagination, but under the form of
eternity and necessity, in God, i.e.\ according to its truth. But when
the modus is apprehended in this way, it is seen as it were into the
infinite; it is apprehended not as single and finite but in its
immediate connection with its eternal, infinite cause; its essence
expresses the essence of substance, its existence becomes essential,
necessary existence. Necessary existence conflicts with the essence of
finitude; through the conception of modi as necessary and eternal one is
therefore led to the conception of the totality of the modifications as
an \emph{infinite unity}. Limitation is negation, not something
positive; what, when a finite is presupposed, separates modi from one
another in representation is, seen in thought, something vanishing; but
when this negative is removed, the closed unity of the homogeneous is
posited. This gives the determination of the \emph{universal
individual}, which is equal to natura naturata as a whole.
\emph{Natura naturata} is ``all that follows from the nature of God or
of his attributes, that is, all modifications of God's attributes
insofar as they are considered as things that are in God, and without
God \opage{104}can neither be nor be conceived.'' \emph{Natura naturans}
is ``that which is in itself and is conceived through itself, or the
attributes of substance that express an eternal and infinite essence;
that is --- God, insofar as he is considered as free cause.'' Of the
universal individual it cannot be said that it is composite, that it
consists of parts; for parts would presuppose a subsisting limitation,
composition the reality of limitation, but the boundary is posited
precisely as the non-real. Thus the concept of infinity is again joined
together with the concept modus. The \emph{infinite modus} under the
attribute of thinking is intellectus absolute infinitus, which in its
unity unites all the modi of thinking and is equal to idea dei; the
infinite modus under the attribute of extension is motus et quies. By it
is constituted ``facies totius universi, which, though varying in
innumerable ways, constantly remains the same.''

Through these determinations, which are to maintain infinity and to
remove finitude and its forms from substance, Spinoza's conception of
substance receives a fuller stamp. Substance, which, as causa sui, is
the unity of natura naturans and natura naturata, shows itself to
Spinoza's gaze in the fuller intuition as the infinite that unfolds
itself in the infinity of the modifications of extension and of thinking
as an \emph{organic} whole, which as a whole is unchanged under the
eternal movement of life, and is eternally in undisturbed unity, since
the modifications, with vanishing limitation, close together in its
infinity. This fuller image of the infinite will indeed vanish again
when the fundamental determinations of the system, which make a
multiplicity of single things inconceivable, are consistently carried
through; but it forms itself before thought with irresistible necessity,
because thought, which flees the emptiness of abstraction and seeks a
concrete content, only through such content gets movement, only through
it avoids annihilating itself. The finite is indispensable to thought,
if thought is not to sink together into abstract infinity and negate
itself; the finite must therefore unavoidably press forward and compel a
validity for itself, and we shall see that \opage{105}Spinoza, in his
definitive determination of the human mind as an eternal mode of
thinking in God, comes to ascribe to the individual being a significance
that conflicts with his presuppositions, and to connect immediately what
according to his fundamental determinations is contradictory.

\divrule

From the fusion of the concept of substance with the concept of God,
and from the strict adherence to the concept of infinity, there follow
for Spinoza essential consequences with respect to the conception of
the divine, and of the natura naturata that is posited by the activity
of the divine. From the thought of God there is removed, first, every
anthropomorphizing, finitizing notion, thus the notion of a personal
God and of God's becoming man. In God, next, essence and activity are
seen in inseparable unity: the necessarily existent expresses its
active essence in necessary acting; but since the power that
determines the working of the essence is the essence itself, the
necessity becomes inner, free necessity. Freedom is the necessary
self-determination of the necessary essence. On the more precise
determination of the form of God's working the concept of nature now
exerts its influence. The form for the activity of substance becomes
causality in the form of mechanism. Just as the necessity of essence
excluded a liberum arbitrium, so the form of causality excludes
purpose, which for Spinoza is already incompatible with the divine for
the reason that it presupposes a lack in the agent who \emph{strives}
after the goal, whereas in substance as causa sui the goal of activity
is eternally attained.

By the conception of the essence and mode of working of substance, the
conception of natura naturata is determined. Since, according to
Spinoza, ``Everything that has being is in God, and nothing can be or
be conceived without God,'' the law for the essence and working of the
infinite must be transferred immediately to what is in the infinite.
The law of natural necessity must also govern natura
\opage{106}naturata. Every determination to being and activity in
things must be seen as a determinateness by God; the order of things
becomes an eternal, inviolable order, determined by the necessity of
God's essence. In the existence bound by the bond of necessity no
independent point of departure for a working can form itself; there
can be no question of an individual freedom, a freedom of choice, in
the meaning in which these concepts are usually taken, and just as
little of a contingency, or of a mere possibility; for these notions
have their origin only in the limitedness of the human cognitive
faculty. Nor can such determinations have validity as do not express
nature in its necessary connection with its cause, but only in its
relation to the human imagination.
% errata list B, p. 106: printed "til Forhold", corrected "Forhold til"
Such determinations are not merely those that designate nature's
relation to the human senses, not nature's being in and for itself;
but also determinations such as order and confusion, beauty and
ugliness, perfection and imperfection, good and evil, which all have
their source in the imagination. For in nature everything proceeds
with eternal necessity and with the same perfection from God's
necessary essence; to no modus does anything belong other than what
lies in its concept, and this concept is determined not by its finite
relation to us, but by its necessary relation to God. The imperfect
and evil does not become anything real; it becomes only a negation,
which vanishes when that by which the negation is, is led back to God,
of whom all that is positive is the effect. Existence in its truth is
therefore to be seen under the point of view of eternity and
necessity, as expressing the eternal essence of substance. No
disturbance then becomes possible in this existence, no separation of
a temporal from an eternal life.

If natura naturata is separated by an abstraction from its cause, and
if it is conceived from the point of view of the hypothetically fixed
finite, the law of necessity becomes no less dominant. The detached
finite in its limited existence does not have in \opage{107}itself a
determining ground to be and to act: it must be determined thereto by
another. But this other becomes here, in the domain of the finite,
another finite, which again demands another finite cause, and so on
to infinity. By this abstract conception of the finite, then, the
determining of things by God is conceived as the determining of them
by one another, running on in an infinite series. Everything finite
thus becomes governed by the compelling law of an \emph{external}
necessity; but its determinateness by another finite, coming from
outside, is, when the individual is conceived in its truth, in its
unity with the infinite, transformed into a determinateness by this
infinite. And since the infinite is its own essence, the external
necessity becomes an inner one, becomes freedom. In God as causa
libera the individual modi, conceived under the form of eternity, in
unity with him, obtain their freedom, in the true meaning of the
concept.

\divrule

Since the attributes of substance, as conceived through themselves
without pointing to another, are mutually independent, while yet they
are all expressions of one and the same substance, and all express it
wholly, there must be between the modifications that proceed from them
a complete agreement, but in such a way that the modifications of each
attribute are independent of those of the others. Between the
modifications of the different attributes there obtains a
\emph{parallelism}, but no causal relation, no mutual influence. If
this is applied to the two attributes that are given in the human
being, and that are therefore the only ones we know: extension and
thinking, it leads to Spinoza's proposition: ordo et connexio rerum
idem est ac ordo et connexio idearum [the order and connection of
things is the same as the order and connection of ideas], and to the
determination that a body and its idea are una eademque res, sed
duobus modis expressa [one and the same thing, but expressed in two
ways]. For the modifications of extension are the bodily things
(corpora), determined only by motion, figure and position; the modi of
thinking are the ideas, which in their mutual relation are
\opage{108}conceived after the analogy of the connection and
composition of the bodily. Each modus within the one attribute has its
corresponding member in the other: to each body there corresponds a
thought (idea), and conversely. But what takes place in the one domain
is not brought about by the corresponding member in the other, but by
a preceding member in the same series. Ideas are determined only by
ideas, bodies are moved only by bodies. The idea, which, as a modus of
active thinking, is to be conceived as an active thought, is the mind
(mens) of the corresponding thing and expresses its real essence in
the form of thought. All things are, though in different degree,
animated. The degree of perfection of a mind is determined by the
quantity of ideas it comprehends, and corresponds to the degree of
composition of the body. The human mind is the idea of the human body
(idea corporis), and it is at the same time, as essentially active,
the knowing of itself in its activity: it is idea of this idea (idea
ideæ), certainty of itself in this determinate activity, knowing of
itself through the knowledge of the determinatenesses (affections) of
the body.

By this conception of the relation between the modifications of the
attributes, the dualism that Spinoza sought to overcome is not in
reality removed. Thought and extension, which, by being conceived each
by itself without including the opposition, showed themselves to be
fundamentally different, were to show themselves as expressions of the
unity of the same essence through the complete parallelism of their
modifications. But precisely the denial of every relation of
reciprocal action between their modifications betrays that no real
unity is present. The attributes are not derived from the essence of
substance, but posited \emph{immediately} as determinations of essence
by a \emph{demanded} unity. But since their opposition is preserved
within substance, since there too they remain the \emph{merely}
different, substance, which in its one essence is to unite this
\emph{merely} different, what is inconceivable through each other,
must be split apart. Were the unity of substance to be preserved, that
which made the opposites opposites would have to vanish; but that is
here precisely their \emph{whole} essence, and the unity would then
become \opage{109}the empty and contentless, which comprehended in
itself two empty, contentless abstractions, and such a unity would not
% errata list B, p. 109: printed "Attributer", corrected "Abstractioner"
be able to govern the opposition in concrete existence. But while with
Spinoza the attributes thus distinguish themselves dualistically from
one another, and their modifications run parallel without any contact,
he nevertheless does not avoid letting the concepts of thought and of
nature, which are gathered together in substance, mutually influence
each other. The concept of nature is determined through thought by
extension alone, with the exclusion of the sensuous determinations of
concretion, and extension receives the predicates of indivisibility
and activity, thus predicates that are taken from the mental and
have meaning only there, not in the natural reduced to extension.
Conversely, the concept of nature determines the concept of mind,
inasmuch as the activity of infinite thinking is conceived from the
point of view of necessary causality, not from that of purpose, and
inasmuch as the ideal natura naturata that is posited by this activity
is conceived in analogy with the mechanically concatenated, moved
corporeal. The relation of necessity of causality determines the
connection of the mechanically linked ideas, inasmuch as the
individual modifications of thought come under the concept of the
\emph{thing}, and in their sequence and composition are intuited in
analogy with the modifications of extension. And the influence of the
concept of nature receives its definitive expression in this: that the
concept of substance, which in its truth is the concept of God,
receives the name of nature, so that Deus and Natura become synonyms.
In the mutual determining of the attributes the influence of nature is
powerful enough so to determine the conception of the totality of
existence that this is brought under the law of natural necessity,
% errata list B, p. 109: printed "fører", corrected "føres"
even if natural necessity is sublimated into the metaphysical ground.
The infinite existent becomes the power of nature, determined by
necessity, immediately self-certain in its unity with thought.

\divrule

\opage{110}\emph{how the notion of a finite subsisting by itself
arises}, Spinoza answers by pointing to \emph{the imagination}, the form
of knowledge in the human mind that corresponds to the latter as an
individual, finite modus, and which therefore \emph{itself presupposes
finitude}. The imagination gives no \emph{adequate} ideas, i.e.\ such
as contain only what is clearly conceived from our essence alone. In
the imagination the mind shows itself as an individual member in the
chain of causes, affected from outside, always pointing beyond itself
to something determining. Its ideas therefore always include the idea
of the foreign thing by which it is affected, and which is itself
known by it only through the affections it produces. That is the
reason why the imagination can give only inadequate, confused and as
it were mutilated ideas, both of that which affects us, of our own
body and of our mind. It does not express what is pervasive and one in
things, but forms only confused general images. It does not conceive
things in their necessary connection with their eternal cause, but
apprehends them under the form of time and contingency, in untrue
separation from the cause. It leads to the contradiction that the
mind, which according to its concept is idea corporis and idea ideæ,
can apprehend its body and itself only confusedly. But from the untrue
kind of knowledge of the imagination, which contradicts the essence of
the mind as thinking, there is distinguished the true kind of
knowledge, \emph{the knowledge of reason} (intellectus). It does not
conceive things as isolated and contingent, but as effects of the
necessity of God's eternal nature, and therefore considers them not
under the form of time, but under the form of eternity. And in its
idea of things there is included the idea of God's eternal and
infinite essence, from which things follow. Its ideas are adequate,
since it conceives things through the common origin of all, whose idea
is included in the idea of each thing and is therefore adequate in the
human mind. And these adequate ideas are not merely true, but
accompanied by the certainty of their truth, for the idea as active is
not something mute, like a picture on a tablet, but a living act of
thinking, \opage{111}which affirms itself and immediately gives
certainty. For this kind of knowledge the illusion of the subsisting
finite has vanished; all being is for it encompassed by infinity, all
time taken up into eternity.

In this form of knowledge the human mind, whose essence is active
thinking, has its true being, freed from the bondage of the passions
(the passive states), and Spinoza's \emph{Ethics}, which cannot become
a doctrine of tasks to be realized by the independent will of a free
individual, but can only become the natural science of the human mind,
the consideration of the human being in the light of the eternal and
necessary laws of nature, and a development of what follows with
necessity from the concept of the human mind in its self-assertion,
--- shows this knowledge, which is the mind's highest activity, its
most perfect self-affirmation, as the eternal determination of the
mind. Since the mind, in the highest form of the knowledge of reason:
\emph{intuitive} knowledge (scientia intuitiva), sees things as being,
according to essence and existence, in and through God, as eternal
effects of his necessary essence, it sees itself also as a member in
the totality that proceeds with eternal necessity from the active
essence of substance, and as eternally being in unity with God. This
intuitive knowledge, intuiting the effect in its eternal cause, is
grounded in the very eternity of the mind; it follows as an eternal
consequence of this eternity, and thus does not come to be, but is
eternally. What, with respect to our knowledge, shows itself as a
development, a becoming, becomes a vanishing semblance, and since the
untrue knowledge has vanished as a non-being, the finite as finite has
vanished, infinity has closed together with itself, and the eternal
modi of thought form a unity, a res particularis æterna, in God's
infinite reason (infinitus intellectus Dei). In unity with the
infinite the mind knows itself as free and blessed: its being is the
very eternal being of substance; its intellectual love, freed from all
passivity, for God, the origin of its being, working and blessedness,
is one with \opage{112}God's very love of himself and of what is
included in his essence. The true knowing of the mind is, as the
expression of the truth of its active essence, the mind's highest
virtue, with which its blessedness is one. ``Summum mentis bonum est
dei cognitio, et summa mentis virtus deum cognoscere'' [The highest
good of the mind is the knowledge of God, and the highest virtue of the
mind is to know God].
% printed "cognitio, et summa mentis virtus deum cognoscere" (defect logged in the transcription; the truncated note does not name the word; Latin kept as printed)

\divrule

The significance of the Spinozistic philosophy is to be placed
essentially in this: that it has wanted to be only philosophy, and that
as philosophy it has been borne by one thought: the thought of the
infinite. With clear consciousness of its difference in principle from
the positively religious and from theology, it has wanted to assert the
right of thought alone to determine the truth, the validity of reason
as ``the word of God that is in our mind and can never be corrupted,''
and it has not only, as pure philosophy, removed the disruption in
existence that always becomes the consequence of the interference of a
theological element and of the transcendence inseparable from it, but
has also sought to remove the dualism between mind and matter, which
was the inheritance of the Middle Ages, by gathering both in the
infinite substance, which expresses its whole essence through each of
them, and to do away with the separation between the finite and the
infinite which, after the notion of transcendence had been removed,
would bring discord into the infinite and annihilate infinity itself.
This infinity Spinoza has energetically striven to hold fast, so that
in it he intuited an eternal fullness proceeding from its activity.

By this concept of the infinity of substance, Spinoza's philosophy
becomes the most correct expression of \emph{dogmatic idealism}. The
fundamental concept of this idealism, the expression for what was the
essential goal of knowledge, was \emph{substance}: the essentially
existent in its independence and necessity. But essential,
\emph{necessary} being showed itself, when its concept was
consistently thought through, as a being through itself and, as such,
as infinite and as unity. Spinoza's concept of substance, which
comprehends the whole of exis\opage{113}tence in its infinite unity, is
therefore the consistent expression of the fundamental concept of
dogmatic idealism.

But this concept of infinity in Spinoza suffers from an essential
defect, which prevents him from reaching the goal after which he so
energetically strives. Substance is determined by Spinoza as activity,
causality, and by this determination he wants to procure for substance
an infinite fullness, to see it as eternally unfolding itself in the
infinite totality of the modifications. But in Spinoza's concept of
substance the concept of being becomes from the beginning what
underlies and predominates; activity comes to relate itself to this
underlying only as a predicate subsequently added to a subject
presupposed in its finished being, and therefore the infinity that is
the essence of that being does not become determinate, but becomes
\emph{abstract}, resting infinity. It is not seen as infinity through
the activity by which it posits its determinations and sublates them in
its unity: it is seen as infinity prior to the activity, and therefore
the latter can posit only something vanishing. And the postulated
activity itself becomes inexplicable in the substance from which every
inner opposition is excluded, and which therefore can obtain no
principle of unfolding within itself, and the infinite multiplicity
that the activity is to posit becomes inconceivable, since substance
excludes limitation, because limitation is only negation and substance
cannot negate itself. In the exposition of the system there can
therefore be no question of a genuine deduction, of a dialectical
development of thought; rather, an apparatus is taken up by way of
postulate which, in the form of definitions, axioms and the like,
anticipates the essential determinations without deduction. And only
with the help of what has thus been taken up does one proceed from the
first immediately certain, through inferences in mathematical form.
But the stringent linking together links together only something
already given and presupposed.

Since substance in its abstract infinity has in itself no principle of
determination and development, thinking and extension, which are
posited as the forms of its being, cannot be \opage{114}derived from
its nature; but they are taken up as given and, by virtue of the demand
for the unity of existence, posited immediately as forms of being of
the same substance, without their concepts including each other or
pointing to each other, which is precisely expressly denied by the
definition of attribute. But in this way, as has already been shown
above, they come to unity only in that their difference vanishes in
the leading back to substance, not in that they determine themselves by
themselves to unity, so that the difference is preserved in this unity.
Their opposition is therefore not reconciled; the dualism persists in
spite of the powerless unity.

And just as little as the dualism between thought and extension is
overcome, is the dualism between the infinite and the finite in reality
overcome. The thought of infinity was Spinoza's first thought; his
energetic striving was: to hold fast this thought as the all-embracing
one. But in order that it could be held fast as such, and in order that
thought could at all gain a self-affirmation through it, it had to
contain the finite within it. Without this, thought must as it were
expire in the airless space of abstract infinity. Therefore the finite
in Spinoza forced from thought, as it were, a lingering, and the
intuition of the infinite active substance shaped itself into the image
of a totality moved within itself, encompassing an infinite
multiplicity, eternally unchanged as a whole, and through this image
the thought of the infinite organism was intimated. But since the law
that was to bring the multiplicity bound by unity into infinity stood
unconceived, and since the finite according to Spinoza was to be
conceivable only through another finite, the finite remained
unexplained, and since it came to be for the notion which itself could
be understood only through the being of the finite, a dualistic
separation showed itself between it and the infinite. The dualism was
not removed by the fact that the thought of finitude everywhere
dissolved itself into contradictions, and that the finite, since
limitation was conceived only as negation, was seen as something
constantly vanishing. For even if it \opage{115}is constantly brought
to vanish, it is nevertheless constantly posited anew, because only
through it does thought gain movement, because only through it does
thought avoid denying itself. And at the close of the system it asserts
itself in the most pregnant form, inasmuch as the human mind, the
individual limited modus of the infinite thought, even when it is
conceived according to its truth, under the form of eternity and
necessity, is held fast as an \emph{individual} modus. But the concept
of the eternal individual modus contains, seen from Spinoza's
standpoint, a contradiction; for from this standpoint the determination
of eternity is incompatible with limitation and individuality. In the
eternal unity with substance the individual modus, the individual mind,
would therefore have to vanish; the highest act of its knowledge, its
highest self-affirmation, would become an annihilation of the thinking
individual, a vanishing into the \emph{abstract}, contentless infinity
of substance. But the fact that this concept is held fast shows
thought's irresistible urge to fill the emptiness of infinity; and by
this concept, unjustified from the standpoint of Spinozism, there is
implicitly a pointing toward another determination of the fundamental
concept. Since the individual modus of thought in the intuitive
knowledge of God is posited as containing the infinite, the finite is
in reality posited as infinite in its finitude, and the concept of the
\emph{monad}: the \emph{individual} substance with inner infinity, is
given. In Spinoza the conception of limitation merely as cessation, as
negation, made the infinite abstract and brought the drive of thought
into contradiction with itself. The task is therefore set: to conceive
limitation as the self-limitation of essence, as its qualitative
self-determining in its self-active differentiation, and by this task
% printed "dennne" (defect logged in the transcription)
there is a pointing toward the concept by which \emph{Leibnitz}
integrates Spinoza's concept of substance: the concept of substance as
monad.

\divrule

\opage{116}
\begin{center}
  \textbf{4. Leibnitz.}
\end{center}

\emph{Gottfried Wilhelm Leibnitz} (the original Slavic form of the
name is Lubeniecz) was born in 1646 in Leipzig, where his father,
originally a jurist, was then professor of philosophy. His
extraordinary gifts developed early, and his boundless thirst for
knowledge procured him, already as a boy, an astonishing wealth of
learning. Already in 1661 he came to the university of his native town,
where as his professional study he pursued jurisprudence, but alongside
it busied himself eagerly with the study of philosophy, classical
literature, history, mathematics and the natural sciences. Early he
conceived the plan of uniting the philosophy of antiquity with the new
philosophy founded by Descartes. Barely 17 years old he acquired the
philosophical doctorate in Leipzig with a dissertation, de principio
individui, in which the fundamental thought of his philosophy is already
intimated. Afterwards he studied in Jena under the guidance of the
mathematician Weigel, and in Altdorf, where, at 20 years of age, he
became Doctor juris, then lived for some time in Nuremberg in
association with alchemists, and in 1670 entered the service of the
Elector of Mainz as councillor in the High Court of Revision (supreme
court). With the son of the chancellor Boineburg he traveled to Paris,
where he spent 4 years important for his scientific and personal
development, and came into contact with famous men of science, among
whom Huygens, who became his teacher in higher mathematics, and the
Cartesian Arnauld were the most significant. In Paris Leibnitz worked
out a memoir in which he develops for Louis XIV, without any result, a
plan drafted by him and Boineburg for an expedition against Egypt, by
which France's lust for conquest was to be diverted from Holland and
Germany. When Leibnitz's official position in Mainz ceased with the
Elector's death, he received in 1676 a call from the Duke of
Brunswick-Lüneburg and Hanover as librarian in Hanover. He traveled
there by way of London, where he made the acquaintance of Newton, the
chemist Boyle and the mathematician Collins, and by way of Holland,
where he came into personal contact with Spi\opage{117}noza. In Hanover
Leibnitz was also appointed historiographer and member of the Chancery
of Justice, and his activity soon extended to all the affairs of the
court and of the land. In order to make studies for the history of the
princely house, and on diplomatic errands, he undertook several longer
journeys in Germany and Italy. Through the Hanoverian princess Sophie
Charlotte, who was married to Elector Frederick III of Brandenburg (as
king: Frederick I), he brought about the founding of the Academy of
Sciences in Berlin, whose first president he became. At several
meetings with Peter the Great of Russia he made proposals to him for the
development of the sciences in his realm, and the plan for the Academy
of Sciences in Petersburg, which however was realized only after Peter's
death, derives from him. For a union between Catholics and Protestants
he worked eagerly, but without success; later he labored just as
fruitlessly for a union between Lutherans and Reformed. Some of his last
years, from 1712 to 1714, he spent in Vienna, where the Emperor raised
him to the rank of baron of the Empire. He died in Hanover in 1716.

Characteristic of Leibnitz is the astonishing many-sidedness of his
gifts and of his knowledge. He embraced with the same mastery almost the
whole circle of the sciences, and everywhere he is original and
productive. Through the rich experience of life that he acquired in his
manifold and various circumstances, the mediating and uniting tendency
in him, grounded in his harmonious many-sidedness, developed ever more
strongly. As a philosopher he seeks to see all the philosophical systems
as members of the universal philosophy that comprehends all times and
nations, and his striving aims at uniting the truth of all the different
systems in an ``eternal philosophy''.

Of Leibnitz's numerous \emph{writings} we mention here only the
philosophical ones, and merely touch on the fact that by his treatise
published in 1684, Nova methodus pro maximis et minimis, in which the
basic features of the differential and integral calculus are set forth,
he has given Newton's calculus of fluxions, whose fundamental
\opage{118}thoughts were known to him through Collins, a more perfect
and worked-out form and a greater applicability. Of the
\emph{philosophical} writings only two are of larger compass, namely the
Nouveaux essais sur l'entendement humain (1704), written against Locke's
theory of knowledge, and the Essai de Théodicée (1710), directed against
the skeptic Bayle, to which a Discours de la conformité de la foi avec
la raison forms the introduction. The other philosophical writings are
short treatises and letters. The most important of them are the
Monadologie, written at the request of Prince Eugène of Savoy (in Latin
translation: principia philosophiæ 1714), and Principes de la nature et
de la grace fondés en raison (1714?); besides these may be named:
Système nouveau de la nature et de la communication des substances,
aussi bien que de l'union, qu'il-y-a entre l'âme et le corps (1695),
with the accompanying Eclaircissements; --- De primæ philosophiæ
emendatione et de notione substantiæ (1694); --- Sur les principes de
vie et sur les natures plastiques, and some letters to Arnauld, Bayle,
des Bosses and Clarke. A detailed systematic exposition of his
philosophy Leibnitz has given nowhere.

\divrule

As in Spinoza, so in Leibnitz the concept of \emph{substance} becomes
the fundamental concept; but the concept of the infinite One becomes
individualized in him into the concept of the indeterminably many
substantial individual beings, and the Spinozistic equating of the
mental and the natural as attributes of substance is annulled in favor
of the mental, which now alone becomes the essential determination of
the substantial.

Everything that shows itself as something extended is, as divisible,
something composite, and the composite points back to the individual
elements of the composition. These are not atoms, for the atoms still
have an extension, and are therefore still an aggregate; but they are
in the strict sense simple, indivisible, really existing ``metaphysical
points'', ``points de substance''. These individual beings, which
Leibnitz, with an expression from Giordano Bruno
\opage{119}taken over by him, calls in his later writings
\emph{monads}, are the \emph{only} real existent. As the fundamental
determination of the monad Leibnitz posits (in connection with Glisson,
Cudworth and More) \emph{active force} (vis activa), and the activity is
understood immediately as the realization of an \emph{individual}
nature, so that the monads are originally intuited not as uniform atoms
but as infinitely individualized, and more precisely as forming an
infinite \emph{continuous} series of qualitatively different grades of
being. Since the fundamental determination, force, is a \emph{mental}
determination, the monads, as determined by it, become mental,
\emph{immaterial}, soul-like beings, ideal centers of force. As
immaterial they have in the strict sense the \emph{simplicity} that was
required for what constituted composite things; and as in the strict
sense simple, i.e.\ indivisible, they are in turn \emph{unreceptive to
an influence from without}, since the possibility of such an influence
would presuppose that they had parts, inasmuch as the external influence
would have to propagate itself successively through what is affected.
The individual monads thus have an \emph{independence}, and in this
they are self-active out of their \emph{inner} source of activity. But
this activity, by which they have an inner unfolding, remains within
themselves, because everything that is has the same independence; it
becomes an ideal acting, whose fundamental form is determined by a
\emph{representing} (perception), inasmuch as the individuality of the
monad sets it in an ideal relation to the totality of the other monads,
which mirrors itself in it according to the peculiarity of its
microcosmic nature. Each individual monad thus becomes a little world
for itself, a peculiar mirror of the whole of existence. And since the
monads are from the beginning determined as \emph{ideally} supplementing
one another in their individuality into an organic unity of existence,
which mirrors itself in each one of them, there is, despite their
independence of one another, an undisturbed \emph{harmony} preserved in
existence during their unfolding of life: a harmony that, conceptually
understood, becomes the unity, given by the inner relation of the
monads, in the infinite
\opage{120} \setcounter{footnote}{0}
multiplicity, and as such the expression of the perfection of the
universe, while in the form of the theological notion it is understood
as a harmony \emph{predetermined} (harmonia præstabilita), by the
calculating adjustment of the creating God, the primal monad, between
the activities of the monads independent of one another.

By this conception one and the same essential determination is given
for everything that truly is: the whole of existence is transformed
into \emph{psychical} being, dissolves into beings whose fundamental
determination is that of the mental, and which differ from one another
only by a difference of degree, determined by the degree of development
of representation (slumbering, dreaming, feeling, self-conscious and
thinking monads), and that in such a way that there is a continuous
transition in the differences of degree, and that the possibility of
infinite development stands open to all, even to the as it were
slumbering monads in the apparently lifeless. While existence acquires
fullness and richness through the infinite series of beings formed by
the monads and through their inner infinity, its multiplicity is seen
as permeated by the unity of harmony, which is still understood only as
the harmony of the mental individual substances, not as a harmony
between the mental and the bodily, since a bodily world and its
foundation, matter, have as it were not yet come into view.

In order that the phenomenon of \emph{matter} can be explained, there
must be found a point of attachment in the essence of the monad, the
only actual. Leibnitz finds it in the finitude and limitation that
follow from the nature of the monads as created and individual, which,
as it appears in the imperfection of their representing, in the mixture
of their clear and distinct representations with obscure and
confused\footnote{For L. the \emph{clear} representation denotes the
one that is sufficient, the \emph{obscure} the one that is insufficient,
to distinguish the object from others. Of the clear representations
those are called \emph{distinct} by the help of which I can distinguish
the individual determinations of the object, \emph{confused} those that
do not give me the possibility of doing so. The highest degree of
distinctness, where the analysis is carried to the last elements, is
\emph{adequate} knowledge or, when the elements are surveyed all at
once, \emph{intuitive} knowledge.} ones, which express only unclearly
the uni\opage{121}verse's infinity, points to the inhibitedness of their
fundamental activity, and so to a moment of \emph{passivity} in their
essence. The activity of the finite being becomes only a finite,
relative activity, hence an activity in which passivity is a moment.
This moment of passivity is, as it were, the \emph{possibility}
\emph{of matter} (materia prima) in the being itself. And this
possibility of matter is the condition for matter \emph{as object of
representation}. As such, matter (materia secunda) denotes \emph{the
confusedly intuited aggregate of monads}, hence only a semblance, a
\emph{phenomenon}, albeit a \emph{well-founded} one (bene fundatum),
namely founded in the essential and substantial. For a representation
distinct throughout, such as we think it in God, for a representing
activity that can comprehend the whole of existence in an adequate
knowledge, this phenomenon would vanish, since it would dissolve into
the essential that lies at its ground. The obscure unity of the
aggregate would show itself as clearly distinguished individual monads;
but for the imperfect, limited human representing the phenomenon
remains. Our sight is like the naked eye's view of the Milky Way, whose
countless stars for this eye flow together into continuous nebulae of
light; the divine sight is like the view of the armed eye, for which
the individual points of light show themselves as points in their clear
distinctness. Matter as a form of existence thus becomes a phenomenon of
the connection of the monads, existing only for finite sight, and the
fundamental determinations of the material: extension (space), motion
and time, are then understood as phenomena of the position of the
individual monads, of the real active force, and of the order of the
effects of force.

When these determinations are now applied to the \emph{relation
between soul and body}, they lead to the following conception. The body
is an aggregate of monads, as extended mass something
pheno\opage{122}menal. In this aggregate a central monad, the soul, is,
amid the uninterrupted change of the parts, the ruling principle, the
entelechy, and through it the aggregate becomes an organism with a unio
realis and with a \emph{substantial bond} between the elements. That
the central monad is ruling is to be understood not as if it intervened
in the being of the other monads through an influxus physicus, but
ideally, in such a way that it is the more perfect, more developed one,
the one that has the greater capacity for distinct representations,
hence greater activity, and the one to which the others are adjusted,
whose state furnishes a ground of explanation for the changes of the
body. It becomes the center for just these monads, i.e.\ it becomes the
soul in this body, because it represents to itself just these monads'
states and changes more distinctly than those of the other monads, and
so stands in the closest relation to them. Each of the monads to which
it stands in this closer relation nonetheless leads, according to its
nature as monad, an independent life in the organism too, and the
subordinate monads can in turn be understood as relative centers, as
souls for smaller wholes, for the individual parts of the organism. The
organism is a machine made by God, of which even the smallest parts are
again machines, a composition in whose smallest parts there is again a
world of living beings. We constantly reach, as it were, an infinite of
a higher order. Since in the body that belongs to a soul there is a
constant change of the parts, the soul constantly changes its body, but
gradually, in such a way that it is never deprived of all its organs at
once. Soul and body are inseparable in concept: there are no souls
absolutely separated from something bodily. Therefore there is no
generation or death in the proper sense; but what we call generation is
an evolution, an enlargement; what we call death is an involution, a
diminution. Not merely the organic body exists in the germ before
conception, but also the soul in this body; that is: the living being
itself, and at conception this living thing is merely disposed to a
greater development.

\opage{123}

According to this fundamental conception, which holds fast to the
psychical monads as the sole existent, and posits the soul and the
grouping of ideally ruled monads designated as body as inseparable in
concept, the harmony between soul and body becomes merely the simple
consequence of the harmony between the monads in general; for the body
is, after all, in its actuality only a sum of monads. But the conception
is essentially nuanced in that, as it were involuntarily, the bodily,
the material, determined as phenomenon, with its fundamental
determinations excluded from the essence of the monads, is granted a
place as a special real form of existence beside the mental. The bodily
world, ruled by the law of mechanism, then comes to form an
\emph{independent parallel} to the world of mind, ruled by the law of
teleology, of purpose. The bodily and the psychical become two
independent domains, and the doctrine of harmonia præstabilita is
modified so that it denotes the agreement, originally ordained by God,
between the two independent, heterogeneous worlds. The soul and the
body, now intuited as a unity, thus go their independent ways as two
coordinate reals; but the divine will has arranged both so that
everything that takes place in the body has its psychical reflection in
the soul, and conversely, everything that takes place in the soul its
mechanical counterpart in the body. Soul and body then become, as it
were, two clocks, which from the beginning are so exactly made by the
artist that, without mechanical connection and without later repeated
intervention, they continue to agree.

This transformation of the concept of harmony, by which the standpoint
of the idealistic doctrine of unity is again exchanged for that of
dualistic parallelism, is not accidental in Leibnitz, but is an
unavoidable consequence of the fundamental determinations of the system.
Leibnitz has from the beginning determined the fundamental concept, the
concept of the individual substance, in such a way that all the
fundamental forms of the material are excluded from it. But what is
excluded is something from which we cannot
\opage{124}get away, which we, in attempting to deny it, indirectly come
to affirm, because it is an essential determination in us. When it then
asserts itself, Leibnitz determines it as \emph{phenomenon} of the
ideal, which alone is actual. But for this determination to be
sufficient, the phenomenon would have to be explicable from the
essence, which, given the presuppositions, becomes impossible. When
Leibnitz attempts to deduce matter and the fundamental forms of material
existence from the essence of the monads, he reaches them only, and can
reach them only, by surreptitious moves. Where the emergence of matter
is to be shown in the essence of the substances, this happens only in
that the moment of limitation in the active monad, identified with that
of passivity, is converted by a postulate into the determination of
matter, as if through a recollection of the dualistic conception of
matter. And when space, motion and time are to be explained as
phenomena of the relation of the ideal, this happens only by a similar
surreptition, which simply inserts into the monads themselves what is
to be explained. Thus when space is to be the phenomenon of the
position of the monads, the concept of space is already anticipated in
the determination, position, and the monads, which as immaterial are
free of space, are brought into a relation to space at the same moment
as one has excluded it from their essence. In the same way the
explanation of motion and time presupposes what is explained. The
determination of force leads to the determination motion only when
force is from the beginning understood as moving force, but not when it
is understood in its original meaning as representation, as ideal
activity. And temporal succession is not explained by the ideal logical
sequence (consequence). But when the phenomenon with its peculiar
fundamental determinations thus in actuality stands underived and
unexplained, it unavoidably asserts, in its uncomprehended being, an
independence beside that which could not explain it, and in this
independence it no longer stands as a dependent phenomenon in relation
to the essence, but as a real in relation
\opage{125}to an essentially different, coordinate real. Body and soul
now become two independent worlds, each following its own laws, and the
agreement, the harmony, becomes the uncomprehended harmony of the
heterogeneous, whose author must be sought outside the beings. Leibnitz
does indeed intimate that the mechanical laws of the material are in the
last instance to have their ground in the teleological, and to be means
to its realization; but the intimation is not carried through, and
mechanism becomes an unexplained, independent power, ruling its separate
domain, beside purpose.

Thus in Leibnitz, in the determination of the relation between the
psychical and the bodily, there comes into view the dualism that was
latently present in the first, one-sided determining of the substantial
as something merely mental, to the exclusion of all the determinations
of the material, which abstracting thinking nevertheless has not the
power to transform into a non-being. They force their way involuntarily
into the intuition of the place, position and composition of the
monads, and from the sphere of the bodily, which as it were compelled
for itself an independent actuality, mechanism forces its way into the
conception of the psychical itself, inasmuch as the unconscious
processes of the obscure and confused representations are determined by
it, and it touches the very conception of harmony, inasmuch as this is
seen as an external adjustment of the monads to one another proceeding
from a divine will. ---

But if Leibnitz's original determination of the essence of the monad
thus has this defect, that it makes a dualistic split in existence
unavoidable, it leads on the other hand, through the independence and
activity that it asserts for all monads, to a conception of the
\emph{knowing of the mind} that energetically applies a necessary
corrective over against the one-sidedness of realism, albeit through an
opposite one-sidedness. Leibnitz must unavoidably, since his conception
of the essence of the monad must lead him to let \emph{all} our
representations have their source in the mind itself, enter into a
definite opposition to the theory of knowledge of Lockean empiricism,
which lets everything
\opage{126}come from without. When empiricism sets up the proposition:
``nihil est in intellectu quod non antea fuerit in sensu'' [there is
nothing in the intellect that was not previously in the sense], Leibnitz
adds the corrective: ``nisi intellectus ipse'' [except the intellect
itself], and since this intellectus is force, activity, and since its
\emph{essence} is to think, it has in itself, as a principium
cogitationis, the possibility, the disposition, to develop a content
out of itself, to give the ideas, which are originally in it virtually,
an actual being for consciousness. Sensation, the indistinct, confused
representation that mirrors the mind's relation to other things, is
indeed a necessary condition for clear and distinct thinking, a
necessary point of passage in our inner development; but just as little
as it itself has its source outside the monad, is it the \emph{source}
of thinking. It can occasion a conceptual knowledge, but the mind
produces this of itself; it can inform about the factual and individual,
but cannot give knowledge universal validity and necessity. The fact
that our mind is in possession of \emph{necessary truths}, such as those
of logic, metaphysics, mathematics and morals, points back to its
original activity, by which it produces the concepts (e.g.\ the
concepts of substance, of identity, of the good and of the true) on
which the necessity of those truths rests, and the principles that
govern all our thinking: principium contradictionis and principium
rationis sufficientis. Reflection becomes, on Leibnitz's conception, an
attention to what is originally in us, and the being of the necessary
truths in our mind becomes the condition for our being able to bring
them to consciousness.

As the mind's form of being Leibnitz determines the \emph{synthesis of
self-consciousness}. Self-consciousness is the thinker, the thought, and
the unity of both. The reflection into itself that takes place in
self-consciousness becomes possible because the mind has the capacity
to think the eternal and necessary truths, in which it does not relate
itself to anything foreign. Through consciousness and reason the drive
toward will and freedom, present in all monads, is raised in the human
mind, a drive whose final goal is happi\opage{127}ness, which is
realized in the universal harmony. The imperishability that is given in
the essence of the monad in general becomes, for the human mind, which
is an image not merely of the universe but of God, an \emph{immortality
of personality} as a member of a \emph{kingdom of grace}, with which the
\emph{kingdom of nature} stands in harmony, inasmuch as what happens in
it is always what is most serviceable for the mind. Despite the
imperfection inseparable from the concept of the creature (the
``metaphysical evil''), this kingdom of nature, the existing world, is
\emph{the best world}, and precisely for that reason God has, among the
countless possible worlds whose ideas are eternally the inner object of
his understanding, chosen it according to a \emph{moral} necessity, and
by his power given reality to the resolution of his will.

\divrule

Compared with Spinoza's philosophy, the Leibnitzian has essential advantages on several points.
By setting up the concept of the monad, which brings the determination of individuality into
substance, it has made possible the explanation of the multiplicity of existence. By the
determination of the infinitely many continuous stages of development in the independent monads,
existence is conceived as possessing an infinite richness and an unbroken continuity. By the
determination of self-consciousness as the form of being of the highest monads, mind is posited as
more than a vanishing modus; and by the concept of the inner infinity of the self-conscious, and in
general of the representing, monad, and of its ideal relation to the totality of monads, a point of
attachment has been won for the possibility of the being, for the subject, of the objective, which
by its nature is mental and \emph{of like essence} with the monad, as formed in representation by
the subject's activity. But the concept of the monad, as it is developed in Leibnitz, still leaves
essential difficulties behind. It has already been shown how it is not sufficient to remove the
dualism whose overcoming was a main task for Leibnitz. Another difficulty shows itself in that
Leibnitz, in \opage{128}his striving to preserve the religious notions, is not content with the
thought of the ideal unity of an infinite organism of existence, but transfers the concept of the
monad to God, as central monad, monas monadum, and lets him create, or as it were ``radiate,'' the
finite monads and pre-establish their harmony. For inasmuch as God, as a monad, thus comes to stand
beside the other monads, on the one hand God's omnipotence will annihilate the fundamental
determination of the individual monads: independent activity, while on the other hand the monad's
other fundamental determination: individuality, comes into conflict with the infinity of the concept
of God. If the relation between the divine monad and the created monads were thought through, it
would lead either back again to the fundamental view of Spinozism or forward to a standpoint higher
than Leibnitz's. --- The Leibnitzian concept of the monad, which becomes necessary through the
defects of Spinozism, has finally the defect that it is in reality only a postulated supplement to
Spinoza's concept of substance, not one won by a proper deduction, and that it occupies a doubtful
position between dogmatic and critical idealism. By the concept of the synthesis of
self-consciousness, which develops out of the concept of the monad, a pointer is given beyond the
sphere of the concept of substance, the fundamental concept of dogmatic idealism, to that of
subjectivity, the fundamental concept of criticism; and when the concept is thought through, it must
lead from dogmatism's immediate security in knowledge to criticism's reflection on the possibility
of knowledge. But Leibnitz's analysis of the form of being of the highest monads is not carried
through and pursued in its consequences, and it does not acquire retroactive force to transform the
concept of substance.

\divrule

\begin{center}
  \textbf{5. Idealism in the 18th Century.}
\end{center}

With Leibnitz the productive time of dogmatic idealism is concluded, and the span of time between
him and Kant marks its decline. While realism develops with strict \opage{129}consistency through a
series of advancing systems, reaches its final goal and extends its dominion, idealism appears in
the various countries without coherence, uncertain and strongly influenced by its opposite. We see
it represented in \emph{Germany} by \emph{Wolff} and \emph{the philosophy of the period of
Enlightenment}, in \emph{Scotland} by \emph{Reid} and his school, in \emph{France} by
\emph{Rousseau}.

\emph{Christian Wolff} (1679---1754; numerous writings in Latin and German on all parts of
philosophy) attaches himself closely, without originality, to Leibnitz, and carries out what the
latter had only sketched, in a system executed with pedantic carefulness by the mathematical method,
which is to comprise \emph{everything that has possibility} (i.e.\ what can be thought as free of
contradiction). But in the execution of the system Leibnitz's fundamental thought is trivialized:
representation is no longer the fundamental determination of the monad as such; the monad retains
only the determination of simplicity, as res (substantia) simplex; its concept is approximated to
that of the atom, and between the monads that become the molecules of the corporeal and those that
as souls are to form the center of an organism, and which now only by a postulate get back the
determination of representation and of consciousness, there forms an unbridgeable gulf, a dualistic
separation. --- Leibnitz's thought of the harmony between the monads becomes restricted to the
relation between soul and body, and is conceived as a hypothesis, so that the possibility of a
direct interaction is not excluded. Likewise teleology is conceived in an external manner. While
Wolff strives by his mathematical method to bind together the whole content of the system by the
necessity of proof, he is compelled, since his method is not an instrument for a real unfolding of
thought but can only bind together something given, to resort everywhere to experience and to take
up his determinations from it. --- In the division of philosophy, \emph{logic} becomes for Wolff, as
for Aristotle, the preparatory science for philosophy in general. \emph{Theoretical} philosophy:
\emph{meta\opage{130}physics}, he divides into \emph{ontology}, the doctrine of what is in general,
of the principles and the fundamental concepts; \emph{cosmology}, the doctrine of the universe;
\emph{rational psychology}, the doctrine of the soul as a single substance; and \emph{natural
theology}, the doctrine of God, his existence and attributes. \emph{Practical} philosophy Wolff
divides, after the example of the Aristotelian school, into \emph{ethics, economics} and
\emph{politics}, according as man is considered as man, as a member of the family, or as a citizen
of the state. --- The significance of Wolff's philosophy lies in its having ordered and systematized
the philosophical material at hand, and in its having, by its easily grasped, if not profound,
treatment of that material, procured it entry into general culture, while at the same time it
influenced the various sciences by its method and by its consistent carrying-through of a
\emph{rational} consideration of the world. Only with the spread of Wolffianism was the
scholastic-Aristotelian philosophy, with its faith in authority, definitively displaced from the
German universities.

\divrule

In the \emph{period of Enlightenment} in Germany, philosophy, which more and more assumes the form
of a popular philosophy, loses its systematic and stringent character, and with the Wolffian
philosophy and its reasoning of the understanding, which predominantly governs thinking, elements
from realism are eclectically mixed. In this span of time philosophy takes a one-sidedly practical
direction. As the infinite, which is incommensurable with the determinations of the understanding,
is posited as incomprehensible, interest gathers on the finite subject, on man, and the task of
philosophy becomes man's happiness; the value of knowledge is determined according to the
contribution it makes to this. Religion is conceived as a doctrine of happiness; God's existence
acquires essential interest because only a highest, all-good and almighty being can guarantee man's
happiness. One \opage{131}is content to demonstrate the existence of God, ``the highest being,''
% printed "hoieste" (defect logged in the transcription)
but willingly resigns a knowledge of God's essence. Only that God is, is known, not \emph{what} he
is; only his power and goodness are postulated as conditions of our happiness. The dogmas of
positive religion are pushed aside or interpreted rationalistically. Of essential interest, besides
God's existence, becomes immortality, the condition of an infinitely continued happiness; and in it
special weight is laid on memory being preserved, on the relations that made the individual happy
in this life being continued. ``Seeing our loved ones again'' becomes the catchword. Morality is
likewise treated eudaemonistically, and interest attaches not so much to the investigation of the
general nature of ethical activity, its laws and motives, as to the special relations of life,
which are treated by ``applied morality.'' Psychological self-observation acquires an extraordinary
significance; attention is directed especially to the life of feeling, and confidential
communications about one's own life, diaries, confessions, become a matter of fashion. Nature is
considered not for its own sake but from the standpoint of a finite, external teleology, which has
the whole existence of nature adapted by God in an external manner to man's utility.
Characteristic of this age with its abstract intellectuality is the lack of sense for what has
developed historically, a narrow-minded superiority toward the ``prejudices'' of the
``unenlightened'' centuries, a rejection of all authority. ---

The eclectic element in the philosophy of the period of Enlightenment comes out especially in the
manner in which Lockean empiricism is inconsistently joined with Leibnitzian propositions. An
unconscious and perhaps indirect (through Rousseau?) influence from realism comes out in the period
of Enlightenment's strong accentuation of the natural side in man, in its respect for natural
endowments, in its demand for naturalness in education and in life, in its stress on the
significance of the material conditions of life. The living sense for nature awakening in this age,
which often degenerates into sentimentality and \opage{132}affectation, presumably springs
originally from the same source. --- The most significant philosophical representatives of the
period of Enlightenment are Mendelssohn (1729---1786), Garve, Sulzer, Engel, Platner, Tetens and
Reimarus. \emph{Lessing} (1729---1781) shares its striving for freedom of mind, but by his
many-sided theoretical interest, by his sense for a deeper philosophy, such as it appears in
Aristotle and Spinoza, and by his eye for the development of the human race, especially its
religious development, through history, he is raised above its one-sidedness and superficiality.
For the Kantian philosophy the period of Enlightenment becomes a \emph{negative} preparation.

\divrule

In \emph{Scotland} \emph{Thomas Reid} (1710---1796) becomes the founder of a philosophical tendency
which seeks in the sound common understanding of men (common sense), independent of philosophy, a
foundation for philosophy and a defense against subjective idealism and against skepticism. Reid
polemicizes against the ``ideal system,'' i.e.\ the doctrine that we originally have only ideas,
and only later, by connecting them, come to a judgment about the reality of what is represented;
and he distinguishes in the function of sense between sensation, which expresses only a state of our
soul, and the perception immediately connected with it, which instinctively and originally includes
a judgment about the existence of the object; and as he seeks back, through experience and
analysis, to ``the original laws of our mind,'' he finds as the first in our thinking the primitive
judgments, independent of reasoning, conscious only through their application and unexplained as to
their origin, the principles of sound common understanding, which assure us immediately not only of
the existence of the object, but also of the identity and substantiality of the I, of the existence
of other intelligent beings, of the validity of the concept of causality, of freedom, of the
harmony between mind and body, and of an abiding order in phenomena, and, in \emph{practical}
respect, of imputability and obligation.
\opage{133}The principles set up by Reid, however, even though they are called laws, become only
the first in experience, that is, the ultimate abstraction of the empirical content of
consciousness; there is no going back, as in Kant, to that which prior to all experience conditions
the possibility of experience. The principles, which are taken up empirically as facts, are not
derived from knowledge of the essence of mind; they come to stand as postulates, which sound common
understanding, which in reality is only the expression of the generally prevailing opinion that by
virtue of habit apprehends itself as the normal, cannot maintain against the contrary pronouncements
of another ``sound reason,'' i.e.\ of another age's accustomed opinion. --- Reid's successors,
\emph{James Beattie, James Oswald, Dugald Stewart, Thomas Brown} and \emph{James Mackintosh}, have
modified and worked out his doctrine only in inessential details. The merit of the Scottish school
lies chiefly in special investigations in the domain of psychology and of morals.

\divrule

In \emph{France} eclectic idealism is represented by \emph{Jean Jacques Rousseau} (1712---1778).
Although influenced by sensualism, he seeks in immediate \emph{feeling}, in conscience, an innate
source of consciousness of good and evil, and in opposition to materialism he champions the
difference of mind from matter, the existence of a highest spiritual being, and the freedom of man.
He shares his age's striving after nature, and passionately condemns culture as the corruption of
nature through the misuse of superfluous powers, as the source of luxury, self-interest, injustice,
unfreedom and inequality. He demands (in his Emile) an education that lets the individual develop
by itself as nature's disciple, without regard to the existing customs and forms, which it only
learns to abhor. He sets up (in his Contrât social) the ideal of a state which, as having issued
from a social contract between individuals, is to unite the \opage{134}advantages of the state of
nature and of the social state, and has as its goal only the inalienable rights of man: freedom and
equality; inasmuch as the individual who enters it in order to win its protection through
subordination obeys only the indivisible sovereignty of the people, which embraces all power, and in
which he himself has an equal share with all others, and, even if he is compelled to obey the
infallible general will (la volonté générale), which aims at the common good, is only compelled to
freedom. He proclaims, finally, a ``religion of man'' (a ``religion of nature and of the heart''),
which has as its content only the existence of the highest unknown being, freedom, immortality,
tolerance, humanity, and the ``virtue'' independent of all ethical life, and which only in the
interest of the state is supplemented by a ``religion of the citizens'' fixed by the power of the
state, whose dogmas, however, are almost only moral and social, an expression of le sentiment de
sociabilité. --- Rousseau's essential significance is that of having conceived freedom as the
essence of man; and even if he conceives it only in its abstraction, and as little determines its
relation to God and to the power of nature as he can distinguish it from, or protect it against,
arbitrariness, he nevertheless becomes not only the active preparer of the French Revolution, but
also a forerunner of the deeper and more consistent ethical idealism that Kant founds.

\divrule

\emph{Dogmatic idealism} has its abiding significance in having led thought back to
self-consciousness as the source of truth; in having asserted something a priori in the mind as
active; in having arrived from the I at the infinite essence of the thinking I, and in having at
least striven to comprehend this in unity with the infinite essence of nature; and finally in
having, directly or indirectly, recognized the essentiality of the natural. Its essential defects
have already been touched on in what precedes. They lie in the unproven fundamental presupposition;
in the unovercome dualism, in the lack of dialectic in the concept of substance (since even in
Leibnitz the \opage{135}form of immanent development is not comprehended), and in the imperfection
of method conditioned by this lack. When in the span of time after Leibnitz we see dogmatic
idealism, which would continue and develop his thinking, trivialized and lost in the
unphilosophical, this conclusion is grounded precisely in the defect of method conditioned by the
peculiarity of the concept of substance. Since the concept of substance, thinking's immediately
certain point of departure, did not have in itself the possibility of a living unfolding, an
apparatus of definitions and axioms had to be taken up by way of postulate, with whose help thinking
could advance through inferences in mathematical form, the necessity of whose connection of thought
becomes an image of the necessity in the postulated activity of substance. But since the concepts
are not won by a deduction, but are taken up by way of postulate as something given, as ideæ
innatæ, they acquire validity for thought in a separation made by the understanding, in their fixed
apartness, and in this apartness they are employed in the reasoning of the understanding that is
clothed in the mathematical form. Even in Leibnitz, where the monad is seen as a principle of
development and where the essence of self-consciousness is made the object of a sharper analysis,
the relation of externality between the concepts is preserved, since the form of
\emph{development} is not comprehended and the immanent unfolding is not known, and the form of the
movement of thought does not become in him essentially different from Spinoza's. But in that
separation the thoughts become finite and abstract, and philosophy becomes a metaphysics of the
understanding, which must end in the unphilosophical. --- In the course of the development of
dogmatic idealism the infinity which from the beginning covered over finitude must be pushed back,
and the finite gain the upper hand, as we saw it in the period of Enlightenment. A current of
thought in which the finite, abstract thoughts went together into a rich, concrete unity of thought
is excluded at this standpoint by the immobility of the fundamental concept.

When this finitizing of thoughts points to the \opage{136}necessity of another form of idealism, a
pointer to such a form is also given by the uncertainty of the fundamental presupposition. In
dogmatic idealism, what gave it the character of dogmatism was the immediacy of knowledge, which
was grounded, not, as with the Greeks, in the immediate harmony of the oppositions of existence, but
partly in the immediate unity, excluding all skepticism, of individual thought with the absolute,
infinite substance, in which it meets nature; partly in the fact that, within the opposition of
mind and nature, the accent, consciously or unconsciously, is so laid on the ideal that the ideal,
as the primitive determinateness of essence of the absolute, in its overreaching of the opposition
implicitly has the determinateness of totality in itself, and in clear and distinct thinking is seen
as immediately expressing the essence of the opposition in its forms. But since the dualism is in
reality not canceled, even where the opposites form parallels within the unity of the infinite
substance, the opposition must assert itself \emph{as} opposition, and the immediate certainty must
be threatened by a doubt about the possibility of a knowing of the objective separated from
thought, --- a doubt which then, as the finitude of thought comes forward in the separation, will
also lead to doubt about the possibility of a knowledge of the principle in its infinity. From the
problem of the relation of thought to the objective, the real, there proceeds the form of idealism
that replaces the dogmatic: \emph{critical} idealism. Idealism arrives, at the point where the
skepticism proceeding from realism as it were wakes it from its dogmatic slumber, at this new form
by going back again to self-consciousness, to which Descartes had pointed and whose form of being
Leibnitz had indicated. It becomes critical idealism as thought turns from the object toward itself,
as it makes the presupposition of dogmatism into a problem, and would solve this problem through a
critical examination of the cognitive faculty, through an analysis of the content of consciousness,
which, as it vindicates for the mind what originally \opage{137}belongs to it, is to determine the
relation to the object through the right distinction between the factors of knowledge.

\divrule

\begin{center}
  II. \emph{Critical Idealism}.
\end{center}

\begin{center}
  \textbf{1. Kant.}
\end{center}

\emph{Immanuel Kant} (Cant) was born in 1724 in Königsberg,
to which the family had immigrated from Scotland. After having
gone through the Gymnasium in his native town, he came to the
University in 1740, where he studied philosophy, mathematics, physics
and theology. Through the physicist Knutzen he became closely
acquainted with Newton's writings; the philosophical lectures led him
into the Wolffian philosophy. When in 1746 he had finished his
university studies, he spent a number of years as a private tutor in
various places in the neighborhood of Königsberg; only in 1755 did he
habilitate as Privatdocent at the University there, and as such he
lectured on logic, metaphysics, morals, natural theology, philosophical
encyclopedia, mathematics, physics, physical geography and
anthropology. The circumstances during the Seven Years' War made Kant's
position at the University difficult, and only in 1770 did he obtain a
professorship in logic and metaphysics, after he had been appointed,
some years before, sub-librarian with an annual salary of 62 thalers.
As professor Kant displayed a great activity; his lectures, which were
clear and stimulating, gathered a numerous circle of hearers, and by
his Kritik der reinen Vernunft (1781) he secured for himself the place
of leader of a new philosophical direction, to which a numerous school
quickly attached itself. He constantly followed the various
philosophical movements in Germany, England, Scotland and France with
lively interest, and Hume's and Rousseau's writings in particular
exercised a strong influence on him. Under the Pietist Wöllner's
administration of the Ministry of Public Worship, Kant came into a
conflict with the government on the publication of his work:
``die Religion innerhalb der Grenzen der blossen Ver\opage{138}nunft''.
Kant was accused of ``distortion and debasement of many of the main and
fundamental doctrines of Holy Scripture and of Christianity,'' and it
was enjoined upon him to employ his reputation and his talents for the
benefit of ``the paternal intentions,'' just as all the theological and
philosophical teachers at the University were commanded not to lecture
on the said work. Kant declared, since he would not deny his conviction
in order to accommodate those ``intentions,'' that he, ``as a loyal
subject,'' would keep silent and in future refrain from giving lectures
or publishing writings on religion. Until 1797 Kant continued his
activity as a teacher; from that time on the increasing infirmity of
age compelled him to give it up. He died in his 80th year: 1804. The
fundamental traits of his personality were a genuine spirit of freedom,
permeated by moral earnestness, an unconditional honesty and love of
truth, and a strict and conscientious fulfillment of duty.

Kant reached his distinctive philosophical standpoint only after a long
development. It is first fully formed in his Kritik der reinen
Vernunft, which appeared in 1781, and the first decisive intimation of
it is found in a dissertation written in 1770 on the occasion of his
accession to the professorship of philosophy, while in Kant's
``Träume eines Geistersehers'' (1766) the doubt about the validity of
the \emph{whole} previous metaphysics already comes forward. In the
numerous writings that precede the last-named, Kant still stands on the
whole on the ground of the Leibnitz-Wolffian philosophy, even if in
particulars he departs from it under the influence partly of the
physics and astronomy of his time, partly of the empiricist and
skeptical philosophy. Of the writings from this period we mention only:
``Allgemeine Naturgeschichte und Theorie des Himmels (1755), in which
% unbalanced guillemet kept as in the Danish: the opener before
% "Allgemeine Naturgeschichte" is never closed.
Kant sets up a hypothesis about the origin of our solar system and
about the origin of the tangential motion of the planets, which in the
main agrees with the more fully worked-out and more rigorously
mathematically grounded one in Laplace; --- Versuch, den Begriff der
negativen Grössen in die \opage{139}Weltweisheit einzuführen (1763);
--- Der einzig mögliche Beweisgrund zu einer Demonstration des Daseyns
Gottes (1763); --- Untersuchung über die Deutlichkeit der Grundsätze
der natürlichen Theologie und Moral (1764); --- Beobachtungen über das
Gefühl des Schönen und Erhabenen (1764). In Kant's inaugural
dissertation touched on above, De mundi sensibilis atque intelligibilis
forma et principiis (1770), his distinctive standpoint is intimated in
that space and time, the forms of the world of sense, are conceived as
subjective, and the sensible as phenomenon; whereas the categories are
still conceived as objectively valid, and a knowledge of the
intelligible according to its essence is affirmed. After \emph{Kritik
der reinen Vernunft}, Kant's epoch-making work, had appeared in 1781
(the 2nd, revised edition appeared in 1787), the other writings in
which his distinctive doctrine was carried through followed one another
quickly. In 1783 appeared his Prolegomena zu einer jeden künftigen
Metaphysik, die als Wissenschaft wird auftreten können; in 1785 his
Grundlegung zur Metaphysik der Sitten; 1786: Metaphysische
Anfangsgründe der Naturwissenschaft; 1788: Kritik der praktischen
Vernunft; 1790: Kritik der Urtheilskraft; 1793: die Religion innerhalb
der Grenzen der blossen Vernunft; 1797: Metaphysische Anfangsgründe der
Rechtslehre, and: Metaphysische Anfangsgründe der Tugendlehre
(collected under the title: Metaphysik der Sitten, 1.---2. Theil). Of
more or less peripheral writings from this period we mention: Ideen zu
einer allgemeinen Geschichte der Menschheit in weltbürgerlicher Absicht
(1794), the treatise: zum ewigen Frieden (1795); der Streit des
Facultäten (1798) and the Anthropology (1798). By disciples of Kant
there were published, still in the last years of his life, his Logic
(1800), his Physical Geography (2 vols., 1802---1803) and his Pedagogy
(1803). Later his lectures on the philosophical doctrine of religion
(1817) and on metaphysics (1831) were made public. In the collected
editions of Kant's writings are found numerous letters from him, in
part of considerable interest.

\divrule

\opage{140}It is the concluding forms of dogmatic idealism and the
skeptical branch of the realistic line that, as the immediate
historical presuppositions, condition the form and character of Kant's
distinctive philosophy. Wolffianism was the philosophical school from
which Kant had originally proceeded; he stands close to the most
significant representatives of the period of the Enlightenment in his
cast of mind, in particular in the conception of the relation of the
religious to the moral; but at the same time he has been influenced by
the idealist thinkers in France and Scotland, and he has followed the
development of realism with interest. Its materialistic consequence,
however, only repels him, but its skeptical consequence gives him a
decisive incitement. It was, by his own statement, \emph{Hume} who by
his investigation of the concept of causality awoke him from his
``dogmatic slumber,'' and compelled him to make the problem of
knowledge, which dogmatism had skipped over, the object of a
penetrating critical investigation. By his analysis of the genesis of
the concept of causality Hume had been led to deny all necessary and
universally valid synthetic knowledge; our a priori knowledge had been
restricted to identical propositions. The mind, whose spontaneity had
been denied, had no content originally belonging to it and proceeding
from its active essence; it was merely to bring something given into
relation, not itself to be productive.

From this decisive point Kant now sets out in his \emph{Kritik der
reinen Vernunft}, whose task is, through an investigation of the
cognitive faculty, to determine the origin, extent and limits of our
knowledge, and to find the answer to the question by which reference is
made back to Hume's investigations: ``\emph{are synthetic judgments a
priori possible, and how are they possible?}''

What is compressed in this fundamental question of Kant's becomes clear
when we analyze the individual expressions. The question concerns
\emph{judgments}, not concepts, because only the judgment as assertion,
not the isolated concept, contains knowl\opage{141}\setcounter{footnote}{0}edge.
The question concerns \emph{synthetic} judgments, i.e.\ judgments whose
predicate, although it stands in a connection with the subject, does
not lie in the concept of the subject (e.g.\ ``all bodies are
heavy''), because only such judgments ``extend'' our knowledge, give it
a new content, whereas the analytic judgments, the ``judgments of
explication'' (e.g.\ ``all bodies are extended'', cf.\ Descartes's
concept of matter), which are only an application of the principle of
identity and of contradiction, state in the predicate only what is
identical with the subject or is found, through an analysis of its
content, to be already contained in it. Finally, the question concerns
synthetic judgments a \emph{priori}, because only the a priori
judgments, i.e.\ the judgments independent of experience, include
universal validity and necessity, whereas the a posteriori judgments,
i.e.\ the judgments of experience, lack this character. On the
possibility of synthetic judgments a priori, then, depends the
possibility of such an a priori science, independent of experience,
rich in content, determined by the stamp of necessity and universal
validity, a \emph{metaphysics} in the wider sense of the
word\footnote{By metaphysics in the wider sense Kant understands the
pure science of reason in general, the sum of all the doctrines that
derive from a priori concepts a real knowledge, i.e.\ a knowledge that
concerns a determinate object (not, like logic, merely the \emph{forms}
of our thinking). By metaphysics in the narrower sense he understands
the scientific knowledge of the supersensible.}, which Hume had made
doubtful. Several sciences lay claim to the character of such an a
priori science, and Kant's fundamental question includes the question
of the validity of their claim. One has to ask about the possibility of
a pure mathematics, of a pure natural science and of a metaphysics (in
the narrower sense). With respect to the first two sciences, whose
existence is guaranteed partly by their inner evidence, partly by
confirmation through experience, the question can at once become:
\emph{how} are they possible? A metaphysics does indeed also exist in
fact, and to that extent it might seem that here too one could at once
ask: how? but \opage{142}metaphysics does not have such an unchallenged
validity as the two other sciences, and the question with respect to it
must therefore run: \emph{whether} metaphysics is possible \emph{as a
science}, and if so: \emph{how?} and in what form? To these questions
correspond the main sections of Kant's Kritik der reinen Vernunft,
which subjects the fundamental problem to a ``\emph{transcendental}''
investigation, i.e.\ one that concerns not the objects but the
conditions of the possibility of an \emph{a priori} knowledge of them.

Kant now seeks the solution of his problem by the path of sharply
distinguishing between what in our knowledge, as something given from
without, derives from the object, and what is the original property of
the subject, grounded in the nature of its cognitive activity,
expressing its laws. This distinction is carried through by
distinguishing in \emph{all} knowledge between the \emph{matter} as the
empirically given and the \emph{forms} as the a priori, which belong to
the mind and gather the manifoldness of the given under their unity.
Kant concedes to the philosophy of experience that all our knowledge
historically begins with the empirical, and that the mind, something
Leibnitz had denied, has a side of receptivity; but in opposition to
the assertion that all our knowledge should derive from experience
alone, he maintains the self-activity of the mind as the source of a
priori forms for our sensation, our understanding and our reason: of a
priori intuitions, categories and ideas.

\divrule

In the \emph{first} main part of the critique of reason, the
\emph{transcendental aesthetic} (the doctrine of \emph{sense-}knowledge),
Kant investigates one of the ``two stems'' of our knowledge: the
receptive \emph{sensibility}, the faculty of intuitions (immediate
singular representations). That fundamental distinction between matter
and form comes forward here for the first time, in that in the sensible
intuition, in the singular representation, its matter, i.e.\ the
material given through sensation, through the action of the object on
our faculty of represen\opage{143}tation, which is something manifold,
empirical and a posteriori, is distinguished from the form by which
this manifold is coordinated into a unity, and in that it is stressed
that this form, as that which orders the elements of sensation, cannot
itself be sensation, but must belong to the mind as something a priori,
which is immanent in the mind and is ready to take up the matter into
itself and to supply the norm for its order. In every intuition, then,
there is a \emph{pure} element, i.e.\ one independent of experience,
and thereby the possibility is given of abstracting from everything
empirical. The remainder of the abstraction is then what is a priori in
the intuition, i.e.\ what sensibility contributes a priori to the
intuition, the subjective conditions of intuition: \emph{space} and
\emph{time}. Space is the form of outer sense, time of inner sense
(i.e.\ the intuition of ourselves and of our inner state) and mediately
also of outer sense, since all representations, as determinations of
the mind, belong to the inner state. Both are \emph{a priori} and
\emph{necessary intuitions}, not concepts. That they are intuitions
follows from the fact that the concept is a universal that always
comprehends the particular, whose mark it becomes, \emph{under itself},
whereas time and space are singular representations which include the
individual spaces and times as parts \emph{within themselves}. That
they are \emph{a priori} intuitions is shown by the fact that the
representation of them cannot be given to us through the observation of
coexistence and succession, since these relations presuppose space and
time, and that it cannot be borrowed from experience at all, because
the intuition of space and time is precisely the condition for any
experience at all to be made, for an object to come to be for us. That
they are \emph{necessary} intuitions is shown by the fact that it is
impossible to abstract from them, while one can abstract from
everything empirical. That they are \emph{pure} intuitions a
\emph{priori} is evident also from the fact that concerning them one
can determine much a priori in apodictic form, set up propositions with
the stamp of unconditional universality and necessity, e.g.\ that space
must have three dimensions, which would be
impos\opage{144}sible if they had an empirical origin. But from this it
also follows that they are \emph{only} \emph{subjective} forms, that
they cannot be immanent in things. As a priori and subjective they
ground the possibility of a \emph{pure mathematics}; for the subject,
which has these forms in itself a priori, can also a priori draw out of
itself such synthetic determinations as \emph{must} have universal
validity for everything that can be subsumed under those forms and
determined by them. Precisely in \emph{intuition} lies the possibility
that a priori judgments can become \emph{synthetic}, which they could
not by means of concepts without intuition alone. All the propositions
that merely concern the \emph{forms} of sensible intuition \emph{must}
be valid for the objects of the senses, but such propositions are
precisely the pure mathematical propositions. The geometrical
propositions rest, according to Kant, on the intuition of space, the
arithmetical on the intuition of time; all of them are, in his
conception, synthetic propositions.

But from the fact that those forms have \emph{only} \emph{subjective}
significance there follows the essential consequence that they have
validity only for the \emph{phenomenon}, i.e.\ for the thing as it
shows itself in our apprehension, \emph{not} for the \emph{thing in
itself} (das Ding an sich), i.e.\ for what affects our sensation, as it
is apart from our intuition. As such it is spaceless and timeless, not
determined at all by the forms of our knowledge. Only the phenomenon,
which \emph{must} always show itself to us in those forms, is an object
for sense, outer and inner; only to it do the laws of mathematics find
their application. This view Kant designates as \emph{critical} or
\emph{transcendental idealism}. It does not deny the existence of
things, which is precisely affirmed as necessary for sensation to come
about, but it maintains that things are for us only in the forms of our
representation, and that they are transformed by being taken up into
these merely subjective forms; and about the constitution of things it
therefore asserts only the negative: that they do not
\opage{145}exist in the determinations of the phenomenon: in the
determinations of space and time.

The result of the transcendental aesthetic can therefore be summarized
thus: with respect to \emph{phenomena} as determined in space and time,
\emph{synthetic} judgments a priori are possible \emph{because} space
and time are pure \emph{intuitions} a priori, i.e.\ with respect to
phenomena a \emph{pure mathematics} is possible.

\divrule

The \emph{second} part of the Critique of Reason, the \emph{transcendental
Logic}, subjects \emph{thinking} to an investigation similar to that to
which the transcendental Aesthetic had subjected the faculty of sense. The
transcendental Logic in turn divides into two divisions: the transcendental
\emph{Analytic} and the transcendental \emph{Dialectic}. The former treats
of the understanding, the latter of reason. The faculty of thinking is
determined in relation to sensation in such a way that the former
represents the mind's spontaneity, the latter its receptivity. Sensation
gave us, as we were affected by things, intuitions, singular
representations. Thinking first gives us, as it comprehends and combines
the intuitions, concepts, general representations. Just as sensation
formed the material given through \emph{affection} into a unity, so
thinking forms the intuitions given through sensation, as by its
\emph{functions} it orders different representations under a common one,
and in this way it forms \emph{judgments}. And the form of thinking that
expresses itself in judgment is the \emph{understanding}. The
investigation of the a priori in thinking thus becomes first an
investigation of the \emph{pure forms of the understanding}, the pure
concepts of the understanding, which are to make possible an a priori
knowledge of the understanding. These pure forms must, since the
understanding is the faculty of judging, be the \emph{forms} of the
combinations into which the material given in intuition, the matter of
judgment, is brought: the forms of synthesis, which the understanding
applies to the content. If one now looks away from the matter of
judgments (the content of the subject and the
pre\opage{146}dicate) and considers them only with respect to their form
for the understanding, one finds that they can be referred to 12 kinds,
which are gathered under 4 classes. According to \emph{quantity},
namely, judgments are singular, particular or universal; according to
\emph{quality}, affirmative, negative or infinite; according to
\emph{relation}, categorical, hypothetical or disjunctive; according to
\emph{modality}, problematic, assertoric or apodictic. These differences
do not concern the matter of judgment, which is taken from the empirical,
but precisely only the form of synthesis, and can therefore have their
ground only in certain norms, independent of the empirical and immanent in
the understanding, for the subsumption of the subject under the predicate
in judgment. In each of the 12 kinds of judgment (which Kant takes up
without any deduction from the current logic) the act of synthesis is
different; under each, therefore, there lies a different principle of
synthesis as ground. These different principles Kant calls pure concepts,
pure primitive concepts, or \emph{categories}. Of these there then are,
as corresponding to the 12 kinds of judgment: under the class of
\emph{quantity} the categories unity, plurality and totality; under that
of \emph{quality}: reality, negation and limitation; under that of
\emph{relation}: substantiality and inherence, causality and dependence,
and reciprocal action; under that of \emph{modality}: possibility and
impossibility, actuality and non-being, necessity and contingency. From
these fundamental concepts, of which the third in each class issues from
the combination of the two preceding, a complete system of transcendental
philosophy would have to derive the other a priori concepts of the
understanding given in them: e.g.\ from the concept of causality the
concepts force, action, passion, etc.

All these forms of the combining activity of thought Kant traces back to
one source, without, however, demonstrating their issuing from it. In
everything that originally belongs to the mind there is an expression of
spontaneity, a comprehending into unity. Already through the forms of
intuition, space and time, the empirical manifoldness of the sensuous
impression was gathered by an activity of the mind into the unity of
intuition; the
\opage{147}
forms of the understanding, the categories, link the intuitions together
into a unity of thought, and we shall see that the forms of reason, the
ideas, in turn link together the cognitions of the understanding into a
highest, all-embracing unity. But all these activities of unity and all
these expressions of spontaneity point back to a common source, to a
final point of unity, a final source of activity, and this Kant finds in
the \emph{original unity of self-consciousness}, in what he calls the
\emph{transcendental} --- non-empirical --- \emph{apperception}, the
\emph{original} synthesis of self-consciousness. This transcendental
apperception, which has no empirical content but is only that
fundamental form of self-consciousness, I think, I am, which is the
presupposition of all empirical consciousness of myself, just as the
derived forms are the presupposition of all knowledge of the object, is a
self-consciousness that precedes all experience. The \emph{empirical}
apperception, the consciousness of the self as filled with an empirical
content, is only a passive perceiving, by inner sense, of our inner
states, of ourselves \emph{as phenomenon}. The pure apperception, on the
other hand, is a pure act of spontaneity, and therefore has nothing at
all to do with sensibility, and so not with its manifold content either,
with the manifold changing states of our phenomenal I; but it contains
only this: \emph{that} I am, not a determination of the What of this
being; for the fact that: I \emph{am}, is equal to: I \emph{think}, is
not to give being a content of knowledge, since this: I think, designates
only the abstract form of thought, which bears in itself the possibility
of all knowledge, but no actuality of knowledge. This primitive unity of
self-consciousness is, as the source of the a priori forms, the source of
the unity by which the manifold of sensation first becomes the unity of
the object, and it is the source of the unity of the several objects that
the understanding combines. Since the categories are its forms of
combination, they are \emph{entitled} to be applied to what exists only
through it: the object as object, and the object as \emph{object of
experience}, i.e.\ as determined by the
\opage{148} \setcounter{footnote}{0}
\emph{objective} validity of the connection, i.e.\ by the universal
validity and necessity, by the \emph{conformity to law}, that exists only
through the \emph{thought}-forms of that original unity, and that is
distinct from the empirical connection of perception, supported by
\emph{sensation}, which gives only empirical, subjectively valid and
contingent judgments. The \emph{phenomenal} object of intuition and of
the understanding exists only through the \emph{forming} subject; it
therefore conforms to the latter's a priori forms (just as, according to
Copernicus's theory, the phenomenal motion of the vault of heaven
conforms to the motion of the spectator). As mere representations, the
empirical objects, the phenomena, stand under no other law of connection
than that which the connecting faculty prescribes to
them\footnote{``All the manifold in intuition must agree with the
conditions of the transcendental apperception, and thus stand under the
general functions of synthesis according to concepts.''}.

Through the categories, then, there will be able to be a synthetic
knowledge a priori with respect to all objects of experience, for which
the rules of connection \emph{must} hold, \emph{if} only it can be shown
that the empirical and sensuous intuitions can be subsumed under the pure
and intellectual categories. If this can happen, that knowledge becomes
possible, that is, there will then be able to be a \emph{pure natural
science}; for \emph{nature} designates precisely (as distinct from the
world of sense, the mere sum of phenomena) the complex of \emph{lawfully
ordered} phenomena, which exists only through the transcendental
apperception, and whose general laws are given a priori by the
understanding.

In order that this subsumption of the empirical intuitions under the
pure concepts of the understanding may be thinkable, there must be a
combining third, which comprehends in itself determinations of both. Such
a third is \emph{time}, which as \emph{a priori} form is homogeneous with
the categories, and as the form of everything \emph{sensuous} with what is
intuited, and which, in being determined a priori,
\opage{149} \setcounter{footnote}{0}
through the transcendental synthesis of the \emph{productive}
imagination, by the categories, yields the \emph{transcendental
schemata}\footnote{The \emph{schema} Kant distinguishes from the
\emph{image}. The schema is the expression of the general procedure of
the imagination by which it procures for a concept its image. Thus the
schema of the triangle is not the image of the determinate right-angled
or obtuse-angled triangle, but the general intuition of the rule that the
imagination follows in forming this figure.}, the sensuous analoga of the
categories, by whose help that subsumption becomes possible. Following
the order of the categories, these schemata express the time-series, the
content of time, the order of time and the scope of time. The schema of
quantity is the time-series, number, i.e.\ the representation that
comprises the successive addition of a (homogeneous) one to one. Of the
categories of quality, reality has its schema in time filled (by
sensation), negation in empty time, limitation in the degree of filling.
Of the categories of relation, substance has its schema in the permanence
of the real in time, causality in the rule-governed succession of the
manifold, reciprocal action in the simultaneity of the determinations of
the one substance with those of the other according to a general rule. Of
the categories of modality, finally, possibility has its schema in the
representation of a being at some time or other, actuality in the
representation of a being in a determinate time, necessity in the
representation of a being at every time.

That every object of perception must fall under the schema of time
follows from the fact that time is directly the form of the phenomena of
inner sense, indirectly of those of outer sense. As the \emph{given},
i.e.\ what is received through receptivity, is combined through the
schemata with the categories (i.e.\ with the forms of the activity of the
understanding), we arrive at an actually \emph{real} knowing, as distinct
from mere thinking, which moves in merely formal concepts. ``Concepts
without intuitions are empty, intuitions without concepts are blind.'' By
the necessity of the \opage{150}given for a knowing by means of the
categories, the limit of the application of the categories is determined.
They can be applied \emph{only} to \emph{objects of a possible
experience}, i.e.\ to sense-perceptions or phenomena. ``Sensibility
realizes the understanding while at the same time it restricts it.'' That
the categories do not acquire validity for the thing in itself, for what
is independent of our representation, the transcendental object, whose
reality Kant maintains against Berkeley, would also follow immediately
from their being, as a priori forms of the understanding, subjective
forms, which have no validity for what objectively is in itself.

But with respect to the phenomenon there is, then, inasmuch as the
transcendental schemata secure the possibility of combining the
heterogeneous, an a priori knowledge by means of the categories, which
constitutes a \emph{pure natural science}, whose content becomes
\emph{the general principles a priori}, the general laws, for the
connection of intuitions through concepts, to which \emph{every} nature,
i.e.\ every law-determined sum of phenomena, must be subject. The supreme
principle of all these a priori principles is: ``Every object stands
under the conditions under which a synthetic unity of the manifold in
intuition, and consequently experience, is possible,'' that is, it stands
under the concepts of the understanding. Since the principles are
properly ``judgments that state the schemata as predicates of the objects
of experience'' (Reinhold), there must be a parallelism between them and
the table of categories and the schemata, and we get 4 classes of
synthetic principles:

a) (corresponding to quantity) \emph{The Axiom of Intuition}: ``all
intuitions are \emph{extensive} magnitudes.''

b) (corresponding to quality) \emph{The Anticipation of Perception}: ``in
all phenomena the real, which is an object of sensation, has
\emph{intensive} magnitude, i.e.\ a \emph{degree}.''

c) (corresponding to relation) \emph{The Analogies of Experience} (i.e.\
rules that state how one must judge when per\opage{151}ceptions are
given). The general principle here is: ``Experience is possible only
through the representation of a \emph{necessary} connection of
perceptions.'' To the individual categories of relation correspond the
three particular principles:

\gk{α}) ``in all change of phenomena substance persists, and its quantum
in nature is neither increased nor diminished'';

\gk{β}) ``all alterations take place according to the law of the
connection of cause and effect'';

\gk{γ}) ``all substances, insofar as they can be perceived as
simultaneous in space, stand in thoroughgoing reciprocal action''.

d) (corresponding to modality) \emph{The Postulates of Empirical
Thinking}:

\gk{α}) ``what agrees with the formal conditions of experience (that
which concerns intuition and concepts) is \emph{possible}'';

\gk{β}) ``what hangs together with the material conditions of experience
(sensation) is \emph{actual}'';

\gk{γ}) ``that whose connection with the actual is determined according
to the universal conditions of experience, (its universal laws, such as
the law of causality), is (exists) \emph{necessarily}''.

These are the general laws of nature, which are given a priori by the
understanding, but from which the particular laws of nature, although
they stand under them, can \emph{not} be derived.

\divrule

After the question of the possibility of synthetic judgments a priori has
thus been answered as regards pure mathematics and pure natural science,
and after it has been explained how such judgments became possible in the
domain of these sciences, it remains for Kant to investigate the
possibility of a \emph{metaphysics} in the sense indicated above: of a
priori synthetic determinations concerning \emph{the supersensible}. This
investigation is undertaken in the \emph{second} division of the
transcendental Logic, the \emph{transcendental Dialectic}, which analyzes
\emph{reason},
\opage{152}
just as the Analytic analyzed the understanding and the Aesthetic
sensation. Reason, which, like the other faculties, is taken up as
something found in us, is determined as the faculty that by its
principles brings a thoroughgoing unity into the cognitions of the
understanding (the rules of the understanding), just as the understanding
brought unity into the phenomena. Its principles include, as principles
of an absolute totality, the \emph{unconditioned}, whereas the knowledge
of the understanding always related itself to something conditioned. This
unconditioned Kant designates as the \emph{idea}, which for him means a
necessary concept of reason that goes beyond every object of experience,
and therefore cannot express something that is in experience, but only a
demand, a maxim for the ordering of the cognitions of the understanding
into a unity of reason. The same concept is pointed to when reason ---
as distinct from the understanding, which was the faculty of judgments
--- is determined as the \emph{faculty of inferences}, inasmuch as the
inference, which presents something as conditioned by the premises,
points, through the necessary grounding of the premises, to the completed
totality of the conditions, and thereby to the unconditioned, which is
itself this totality, or includes it. From the forms of inferences Kant
deduces the ideas, just as he deduces the categories from those of
judgments. To the three kinds of inference, the categorical, the
hypothetical and the disjunctive, there correspond three ideas: the idea
of the unconditioned unity of the thinking subject (the soul), the idea
of the unconditioned unity in the series of the conditions of the
phenomenon (the world) and the idea of the absolute unity of all
realities (the Godhead as ens realissimum), and these ideas are the a
priori subjective forms of reason, just as the categories were those of
the understanding and the forms of intuition those of sensibility.

By the ideas, which, as maxims, as rules for the use of the
understanding, which point out beyond everything conditioned, beyond the
sphere of experience, and toward something \emph{unconditioned}, cannot
have objective reality, because no congruent object can be given to them
in the world of sense, there is unavoidably produced in us something like
an optical deception. Essentially they are only rules,
\opage{153}
tasks; but since, in order to solve a task, it is necessary to represent
the solution to oneself as realized, those maxims too can become
principles for pure reason only by one's assuming: when this conditioned
is given, then the whole series of conditions is given too, i.e.\ the
unconditioned is given. In order to think its task determinately, reason
thinks its idea as object, and there arises the unavoidable illusion that
the subjective necessity of a certain connection of our concepts should
be an objective necessity for the determination of things. This illusion
is unavoidable like the optical illusions, but can be made harmless if we
do not let ourselves be deceived by it. This we attain by demonstrating
that reason becomes \emph{dialectical} (sophistical) through that
transformation of the merely subjective, and this is the essential task
of the transcendental Dialectic, the critique of the ``\emph{transcendental
illusion}'' that arises when one would use the a priori forms of the
understanding and of reason for the knowledge of what lies beyond the
limits of experience. This critique becomes at the same time a critique
of the main parts of the earlier dogmatic metaphysics (the Wolffian),
inasmuch as this metaphysics' pretended knowledge of the supersensible was
grounded precisely on such an unjustified application of the categories
of the understanding to a non-empirical \emph{being}, to an unsensuous
and unconditioned object. Through the confusion of what is a norm for the
use of our understanding, a subjective condition of reason, with an
object for an extended knowledge of the understanding, there arose the
false dialectical inferences, which are ``sophistications of pure reason
itself'', and which the critique is to show to be false.

Of the 4 main parts of the earlier metaphysics, \emph{ontology} is
criticized indirectly in the Analytic's doctrine of the categories as
forms of our subjective thinking. Its 3 other main parts, rational
psychology, rational cosmology and rational theology, find their critique
in the Dialectic.
\opage{154}

\emph{Rational psychology}, which, starting from the thinking I's
consciousness of itself as thinking, would construct an a priori doctrine
of the \emph{soul}, made the soul into a res, a soul-thing, and determined
it as a substance that was simple and immaterial, that by virtue of its
simplicity was imperishable, that as intellectual preserved the identity
of the personality, and, while it stood in reciprocal action with the
body, was immortal. But in all these inferences Kant finds only
\emph{paralogisms}, inasmuch as in them one has confused the subject of
the act of thought, that determinationless ``I think'', which as
\emph{logically} identical \emph{subject} accompanies all my
representations and thoughts, with a \emph{psychical substance}, which has
simplicity and imperishability as its determinations, and inasmuch as one
has then posited this substance as a \emph{transcendental object},
although the concept of substance has validity only for the phenomenon.
That the thinking I is always the subject of thinking, that in
apperception it is a logically simple subject, that it shows itself to
itself as the same in the many acts of consciousness, and that in its
existence as thinking it distinguishes itself from other things outside
it, are simple analytic, or even identical, propositions; but they mean
something wholly other than that I \emph{as object} should be something
substantial subsisting for myself, that I should be a simple substance,
that this substance as thinking should preserve its identity through all
changing states, that it should be different in essence from that from
which it distinguishes itself, and finally, that it should be able to
have consciousness without external things and so continue its existence
without a body. Here there is everywhere a confusion of logical rules, of
statements about the logical form of the act of thought, with
metaphysical predicates, a confusion of the possibility of an abstraction
from my empirically determined existence, by which I arrive at the
indeterminate concept of a thinking being in general, with the
possibility that this thinking being can have a separate objective
existence. With respect to the knowledge of \emph{our soul}, i.e.\ of the
\emph{phenomenal} object of our inner sense, our inner experience, our
\opage{155}
empirical consciousness, we are referred to the perception of phenomena,
so that only an empirical psychology acquires validity, and rational
psychology, as an a priori doctrine, rests on a misunderstanding, whereas
it has a good significance as a discipline when it rejects all dogmatic
assertions about the essence of the soul, thus both the statements of
dualism and of spiritualism and the materialistic denials of freedom,
immortality, etc. As reason thus gives up all dogmatic pretensions, it
leads precisely to the domain where its power is not assailed at all: to
the domain of ends, which lies out beyond the limit of all experience. The
misfortune of speculation, to be so limited, might perhaps be precisely a
good fortune for man's practical destiny. In the practical domain the
% errata list B, p. 155: printed "Væsen", corrected "Væren"
being of the subject as intelligible can be maintained. ---

The fundamental concept of \emph{rational cosmology} is the concept of
the \emph{universe}. The idea of the world as the unconditioned unity of
the members in the series of causes and effects of phenomena originally
expressed only the \emph{demand} on the understanding: to seek such a
unity, a completed system. But when this demand is transformed into a
dogmatic assertion, there arise \emph{antinomies}, i.e.\ mutually
contradictory propositions, for which equally stringent proofs can be
given. Of such antinomies there will be as many as there are classes of
categories. Reason postulates unconditioned completeness, and since
experience gives only something conditioned, it directs the understanding
out beyond experience; but the understanding is bound in its thinking to
the categories, so that the demand on it becomes equal to: extending the
categories to the unconditioned. This can happen only with such
categories as contain a series of subordinated conditioned terms, for
which the unconditioned can be sought. These are the categories:
magnitude, reality, causality, contingency. The one world-idea therefore
appears in 4 transcendental ideas (``world-concepts''), which demand
absolute completeness with respect to composition, division, origination,
dependence. The opposed as\opage{156}sertions that come forth when these
demands are applied ``doctrinally'' or constitutively cannot be refuted
by any experience, and a proof can even be given for them, but the proofs
for the opposite stand against one another with equal strength. The
antinomies are the following: 1) \emph{Thesis}: the world has a beginning
in time and a limit in space; \emph{Antithesis}: the world is unlimited,
both with respect to time and space; 2) \emph{Thesis}: every composite
substance in the world consists of simple parts; \emph{Antithesis}: there
exists nothing simple; 3) \emph{Thesis}: besides causality according to
laws of nature there is in the world a causality of freedom, an absolute
spontaneity; \emph{Antithesis}: there is no freedom, but everything in
the world happens solely according to the laws of nature; 4)
\emph{Thesis}: to the world there belongs, either as part of it or as its
cause, an unconditionally necessary being; \emph{Antithesis}: there
exists no necessary being, either in the world or outside it as its
cause.

In these 4 antinomies the thesis is that of rationalism, the antithesis
that of empiricism. The proofs for them Kant constantly conducts
\emph{indirectly}, the thesis being proved by the fact that the endless
continuability asserted in the antithesis cannot be completed,
% printed "Antíthesen" (defect logged in the transcription)
the antithesis by the fact that the limit maintained in the thesis is
posited as arbitrary and as one that can constantly be overstepped. The
contradiction in the antinomies Kant would resolve by his critical
distinction between idea and concept and between thing in itself and
phenomenon. If ``the world'' is taken as transcendental object, as
noumenon, then in the first two antinomies both thesis and antithesis are
false, since the determinations of space and time, of unity and
compositeness, find no application at all to what is in itself, which
therefore excludes both of the opposed predicates. Only as a regulative
% printed "bégge" (defect logged in the transcription)
principle for our inquiry does the demand acquire validity: in the domain
of experience to regard no limit as the absolutely last. In the last two
antinomies both thesis and antithesis can be true, when the former is
referred to the thing in itself, the latter to the phenomenon. In the
world of phenomena everything is determined
\opage{157}
with necessity, and there is no unconditioned cause; in the domain of
what is in itself, on the other hand, freedom and something unconditioned
are thinkable as \emph{possible}, but their actuality must be proved by
another way than that of speculative reason. ---

Kant criticizes \emph{rational theology} by demonstrating the illusory
character of its \emph{proofs of the existence of God}. The idea of God,
as the idea of a unity (a sum total) of all realities, was a necessary
\emph{ideal} for reason, which seeks the unconditioned for conditioned
(limited) reality, but it was, like the other ideas, valid only as a
rule. But theology conceives it as a thing, and thereby there arises the
representation of an ens summum, or of God as an objective being, and the
ideal of the most real being is first realized, i.e.\ made into an
object, then hypostatized and finally even personified. Yet one has a
consciousness that this conversion is not justified without further ado,
and therefore proofs of God's existence are set up. But none of these
proofs holds good. The most important of them is the properly a priori
one, the \emph{ontological}, which infers from the concept of God to
God's existence; for on it the others lean. In it one first infers: the
\emph{necessary} being is that whose non-being is impossible, since this,
as a predicate of that being, comes into conflict with the subject. But
when one cancels (negates) that \emph{merely postulated} subject together
with the predicate, this conflict falls away. Therefore one then seeks to
show that there is at any rate one being (subject) that cannot be
negated, cannot be removed in thought, and that must therefore claim
necessary being for itself, namely the \emph{most real} being. The
concept of this being contains no contradiction, and is therefore
thinkable (possible). And since existence is a reality, it must, so it is
inferred, follow from the concept of the most real being that it must
have existence, and more precisely necessary existence. But against this
inference Kant protests. Existence, he says, is not a reality, if by
reality we understand that which gives a concept more content. ``A
hundred possible thalers are just as much as a hundred actual ones.''
\opage{158}
Existence, actuality, expresses only that the object relates itself as
\emph{given} to my experience: existence is an object's being \emph{posited}
with all its properties (predicates). Therefore, whatever the concept of
the object may contain, we must go out beyond it in order to ascribe
existence to the object; we arrive at existence only by presupposing
existence, and the actuality of existence is proved in the last instance
only by \emph{sensuous} perception. By this distinction of concept and
existence, which is justified as against rational theology, which by
God's existence understood not an existence in reason, but the existence
of a personal being outside thought, Kant thus denies the validity of the
main proof. But on the ontological proof, according to Kant, both the
cosmological and the physico-theological proof rest in the last instance.
The \emph{cosmological} rests on it inasmuch as, after having inferred,
by a leap out beyond experience, from the contingent existence of the
world to a necessary being as the author of the world, it now determines
this necessary being as ens realissimum, because only the necessary being
of the latter can be thought, so that it is thus in reality through the
postulate of the fundamental concept of the ontological proof, the
\emph{concept} ens realissimum, that one secures the necessary
\emph{existence}. And the \emph{physico-theological} rests on it
likewise, inasmuch as it infers from the contingency of the world-order,
known to us only incompletely, to an almighty and all-wise
\emph{author} of the world (not merely an \emph{architect} of the
world), converts this concept into the cosmological proof's concept of
the absolutely necessary being, and with this ends in the ontological
proof. The result is therefore that there is no speculative proof of
God's existence, and that, unless a practical need should lead to
postulating it, there can be no theology of reason.

This negative result of the critique, which transforms the ideas from
\emph{constitutive} principles, which were to lead to a knowledge of
non-empirical objects, into \emph{regulative} principles, into rules for
the guidance of the understanding, which demand absolute systematic unity
in empirical knowledge, absolute complete\opage{159}ness in the use of
the understanding in relation to the objects of experience, and which
therefore forbid empirical inquiry to stop at a determinate limit, is not
dangerous; for it is not negative in the sense that, e.g.\ God's
% printed "f. Fx." (defect logged in the transcription)
existence should thereby prove to be impossible. It implies, indeed,
precisely that this existence can as little be disproved as proved; for
the concept of God is not contradictory, and there is no question of an
object of experience. The ideal of reason, precisely by pointing, if only
as a demand, out beyond the world of experience, both in our
contemplation of nature and in our action, makes materialism impossible,
and the negative result of the critique makes room for the practical
doctrines that are to form the supplement to that negative. The final
result of the whole transcendental Dialectic can namely be formulated
thus: there can indeed \emph{not} be a metaphysics that would go out
beyond the domain of phenomena and know an intelligible \emph{being}, a
\emph{thing in itself}; the domain of \emph{theoretical} knowledge is the
domain of \emph{experience}, but this is the world of \emph{phenomena},
and only through the phenomenon do the categories, empty in and for
themselves, receive a content, does thought become knowledge; but it must
nevertheless be said that reason, with respect to a priori synthetic
determinations, is not restricted to the sphere of phenomena, the
\emph{domain of nature}; on the contrary, it has the power to go out
beyond it, and by its maxims a priori to determine what \emph{ought} to
be. This one might, as distinct from theoretical knowledge, call a
\emph{practical knowledge}. The domain of theoretical knowledge was that
of the \emph{concepts of nature}, but the tasks and demands go out beyond
the domain of nature and can be called \emph{concepts of freedom}, whose
sphere is the sphere of the \emph{noumenon}, of \emph{what is in itself}.
A metaphysics of tasks with respect to the intelligible can thus be
given, and through the negative result of the Dialectic there issues the
possibility of \emph{the metaphysics of morals}.

\divrule

\opage{160}Kant's Kritik der reinen Vernunft concludes with a \emph{Doctrine
of Method}, which, by distinguishing philosophical discursive rational
knowledge from concepts from mathematical intuitive knowledge through the
construction of the intuitions corresponding to the concepts, by
separating the domains of the understanding and of reason, and by giving
a division of the system of philosophy, would answer the question: in
what form a metaphysics, in the sense in which he acknowledges its
validity, is possible \emph{as science}, and would sketch the plan for a
complete systematic ordering of all the cognitions of pure speculative
reason. (Reason is here taken in the wider sense as the faculty that
contains the principles of a priori knowledge.)

\divrule

The results won in the earlier divisions of the critique of reason now
find their further development in the ``metaphysical first principles of
a natural science'', to which a ``metaphysics of nature'' was to have
been joined, and in the ``metaphysics of morals''. In the metaphysical
first principles of \emph{natural science} the transcendental principles
of the pure understanding, developed in the \emph{Analytic} and valid for
\emph{every} nature, find their application to a determinate
\emph{given} nature: \emph{corporeal} (outer) nature, through the
constructions of mathematics. ``Thinking nature'' withdraws from a
similar treatment, since it does not admit the application of the
mathematical, and natural science contains only so much science proper
as it contains mathematics. Of thinking nature there is only an
empirical consideration. The \emph{empirical} basic datum of corporeal
nature is \emph{matter} (i.e.\ that which is the object of the outer
sense), and it is considered according to the schema of the categories.
Since all its predicates are traced back to \emph{motion}, through whose
affecting of the sense the corporeal first becomes phenomenon for us,
pure natural science becomes the \emph{doctrine of motion}. By tracing
the phenomenon of matter as that which fills space back to forces
(attractive and \opage{161}repulsive force), Kant, while polemicizing
against hylozoism as ``the death of natural science'', asserts a
\emph{dynamic} conception of nature, which has become of essential
significance for the later philosophy of nature.

\divrule

In Kant's ``Grundlegung zur Metaphysik der Sitten'', in his ``Kritik der
praktischen Vernunft'' and in his ``Metaphysik der Sitten'' the results
of the transcendental \emph{Dialectic} find their execution, the first
two writings developing the law for moral action, the last presenting
the system of moral actions.

Since in the domain of the practical, too, a search is made for
something determining a priori that contains a universally valid and
necessary norm for the will, Kant finds it in the practical fundamental
law, which, in the form of a not hypothetical but absolutely
unconditioned \emph{categorical imperative}, announces itself to our
consciousness as a fact, through which pure reason asserts itself as
legislative. And just as in the theoretical domain the normative a
priori was the \emph{forms} of knowledge, so in the practical domain too
that by which the moral fundamental law of the will is determining
becomes the \emph{form} of \emph{validity as law}, i.e.\ of
universality; it expresses only the demand: ``Act according to such
maxims (i.e.\ subjective practical principles) as you can will should
hold as universally valid law'', (``act so that the maxim of your will
can always at the same time hold as principle of a universal
legislation''). All \emph{material} practical principles are
\emph{empirical}, and since they take their content from our desire and
our sensuous impulse, they are \emph{sensuous} and egoistic determining
grounds, aiming at our individual happiness, which make the will
\emph{heteronomous} (bring it under a foreign law), whereas that formal
unconditioned command springs from the \emph{autonomy} of the will. Only
determination by the consciousness of duty, not determination by
inclination, by impulse, \opage{162}is moral; morality is right action
out of respect for the moral law. This moral law of the will is the law
of \emph{our own} rational being, an ``inner, necessary legislation of
reason''. In its \emph{universal validity} it is determined by what is
\emph{universally human} in us, by ``homo noumenon'', ``humanity in the
individual man''. In this, that through it we are self-determining,
consists our moral dignity as personality and end in itself. It sounds
to us as \emph{law}, as constraining our natural inclination, because,
as sounding from the world of the intelligible, it \emph{obligates} us
as those who belong to the world of the phenomenon; because we, at home
at once in two worlds, receive our own command as a foreign command, and
thus become at once legislators and subjects of the law, who bow in
\emph{respect} before its majesty.

Through the categorical imperative there opens in the practical domain a
prospect of what lies beyond the limits that theoretical reason could
not overstep. The possibility of the categorical imperative is, since
the fact that we \emph{ought} presupposes that we \emph{can}, that we,
determined by the a priori, are independent of the natural law of the
sensuous phenomenon, conditioned by \emph{freedom}, which thus becomes a
\emph{postulate} of pure practical reason, by which the sphere of pure
reason and the application of the categories are extended beyond the
domain of phenomena, governed only by necessity, to that of the
noumenon, the intelligible, but in such a way that this extension
acquires only \emph{practical} significance and gives us no increase in
theoretical \emph{knowledge}. For the postulate is not a theoretical
dogma, but a presupposition that is necessary in consequence of an a
priori practical motive.
% "but": printed "meu" (defect logged in the transcription)
It does not extend speculative knowledge, procures no insight into a
``How'', but expresses only a subjective certainty about a ``That'',
inasmuch as through it the ideas of speculative reason, by their
relation to the practical, are posited as reality in a practical
respect, whereas before they stood at most as a possibility, a
hypothesis. As free I am an intelligible character, the transcendental
timeless ground of the whole \opage{163}series of my actions in time,
while these, as phenomena, as effects in the world of sense, are subject
to the law of necessity, and with regard to my ``empirical character'',
the sum of my phenomenal actions, there is no freedom. From freedom as
the first practical postulate one arrives in turn, inasmuch as the
highest moral task demanded by practical reason is determined as this:
by freedom to further the realization of the \emph{highest, complete}
good, i.e.\ the harmonious union of virtue and happiness, the
actualization of the law in nature, at postulating the conditions under
which alone this task can be accomplished: \emph{immortality} and
\emph{God}. The moral law demands the perfect accordance of the will
with the law, perfect virtue, \emph{holiness}, (the \emph{supreme}
good). But since man, as also belonging to the world of sense, only ever
\emph{approximates} to this accordance, an \emph{infinite} duration is
demanded, i.e.\ ``\emph{immortality}, in order that an infinite progress
in perfection may take place and the contradiction between demand and
reality be removed.
% unbalanced guillemet: the print opens a quotation before "Udødelighed" and never closes it
\emph{The existence of God} is postulated because it alone makes it
explicable that a harmony can be brought about between what happens in
consequence of a determination by the moral law, independent of all
natural motives, in the world of morality and reason, and what takes
place in nature, which is not produced by us and is therefore withdrawn
from our power, and by whose course in accordance with the end and
essential determining ground of the will happiness is conditioned. Only
an intelligence (a rational being) that is the author of nature can make
thinkable the harmony between the two separate spheres that follow
dissimilar laws. By these postulates practical reason asserts its
\emph{primacy} over the theoretical, which latter is as it were obligated
to seek to unite with its concepts the pronouncements of practical
reason that go beyond its domain without, however, contradicting it.

In Kant's ``Metaphysik der Sitten'', which comprises the \emph{Doctrine
of Right} and the \emph{Doctrine of Virtue}, the principle of the moral
is applied to \opage{164}determining the universal and a priori laws for
man's \emph{legal} and \emph{moral} action: for the action that in its
content accords with the demand of the moral law, and for the action in
which duty is the incentive, the disposition from which the action
proceeds thus according with the command of reason. The domains of the
legal and of the moral are sharply separated from each other by Kant,
and every relation that does not rest exclusively on a moral obligation
is treated as a purely juridical institution (thus marriage and the
relation of the state).
% "obligation is treated": printed ";" after "Forpligtelse" (defect logged in the transcription)

The \emph{Doctrine of Right} concerns actions, not according to their
end, but according to their form, insofar as human choice reveals itself
through them. The command of the universal law of practical reason, that
we are to act so that the maxim of our action is fit to be a principle
for a universal legislation, receives in the law of right, which
concerns the form of our outer actions, the expression: that we are to
act so that the freedom with which we assert our choice in our actions
can also be granted to all others in the same way. Right becomes ``the
sum of the conditions under which the choice of the one can coexist
with the choice of the other according to a universal law of freedom''.
All right is derived from the concept of lawful freedom. In this, that
right contains the universal conditions for all use of freedom, it lies
that with right is connected the authorization to compel individuals to
observe its commands, and rightful coercion becomes only a removal of
the hindrances that the choice of single persons would set against
universal freedom. Those conditions can be partly such as follow
immediately from the nature of a free being and therefore hold for all
such beings in their mutual intercourse, partly such as follow only from
the union of many men into a society. The former constitute
\emph{private right} (natural right), the latter \emph{public} (civil)
\emph{right}.

The detail of the Doctrine of Right is of lesser interest for our
purpose. In the development of public right, which Kant
di\opage{165}vides into the right of a state, the right of nations and
cosmopolitan right, the pure republican form of state (in which ``the
united people'', with whom sovereignty lies, asserts it through a
representative system that distinguishes the legislative power from the
executive and the judicial), is designated as the only rightful form of
state, i.e.\ the one corresponding to the idea of the state of right, as
the original rational form in which the law is self-ruling. The other
forms of state that have appeared historically are regarded as merely
provisional forms, which are gradually to dissolve into that one. In the
right of nations Kant posits peace among the peoples as the moral task,
because only in peace is the rightful condition present. And although he
concedes that the impossibility of bringing about the conditions for a
``perpetual peace'' makes it an unattainable ideal, he nevertheless does
not give up faith in the continued progress of the human race toward
this ideal. He has an unshakable trust that the a priori demands of
reason have the power successively to assert themselves, and that
nature, ``the great artist'', through the very inclinations in men that
immediately conflict with a universal rightful condition, will precisely
bring it about, inasmuch as the experience of the consequences of these
inclinations forces the egoistic wills toward that universal rightful
condition. The evil inclinations counteract themselves, and thereby room
is made for reason. A condition for the universal rightful condition
Kant sees in the development of republican forms of state, and in the
founding of a mighty federation for the protection of peace and of the
right of nations. The faith in the power of practical reason that comes
forward in Kant's conception of the development of the right of nations
comes forward also in his conception of the \emph{historical} in
general. He has a firm trust that humanity is in constant progress; the
opposite view seems to him to stand in direct conflict with the demands
of practical reason. Here too he stresses the cooperation of nature in
progress.

In the \emph{Doctrine of Virtue} the characteristic basic feature is the
ethical rigor with which the unconditionedness of the command of duty is
\opage{166}inexorably maintained and respect for the law alone is
recognized as the motive for an activity in conformity with duty, while
an action according to merely natural inclination is posited as
non-moral. By a free self-constraint impulse is to be mastered, and duty
becomes a ``necessitation to an end adopted unwillingly''.

Kant divides the \emph{duties of virtue} into the duties to ourselves
and to other men. \emph{Only} to men are there duties. What has been
called duties to what is lower than man, to the animals, which are not
persons, falls in reality under the duties to ourselves, namely under
the duty not to blunt sympathy in oneself. ``The duty to have religion''
is likewise not a duty to God, who is not given to us as an empirical
object, but a duty to ourselves, from whose reason the idea of God
proceeds. Faith in God can be demanded only as a condition or aid for
our morality, only on account of the practical effect that it \emph{can}
have to regard the voice of conscience as a divine voice, inasmuch as the
moral law can thereby gain greater strength.

\divrule

Kant's conception of \emph{religion} is in general conditioned by his
conception of the relation between the theoretical and the practical,
and by his upholding of the autonomy of the moral. Since the existence
of God, which for theoretical reason was only a hypothesis, first became
for practical reason a postulate, necessary in consequence of the
consciousness of the moral law, an object of a pure rational faith,
there could be no speculative theology, but only a moral theology.
Religion must, with respect to its matter, coincide with morals, and
differs from it only in this, that the idea of God, demanded by morals
itself, is added to the concept of duty, so that duties are conceived as
divine commands. For natural religion duty is the primitive; its faith
is a pure rational faith. In the historical forms of religion this faith
has, on account of the weakness of human nature, received an admixture
of a statutory element, \opage{167}whose determination is only to be a
vehicle for the former, and successively to disappear. The successive
transition of ecclesiastical faith to the sole dominion of pure
religious faith is the approach of the kingdom of God. Rational faith
must become the right interpreter of ecclesiastical faith, and therefore
Kant attempts in his ``Religion innerhalb der Grenzen der blossen
Vernunft'', in developing the moral rational faith, to bring by a
symbolic interpretation the principal dogmas of Christianity: the dogmas
of original sin, of the Son of God and the atonement, into accordance
with this faith as the true religion. The ecclesiastical dogmas become
moral doctrines. The dogma of original sin becomes the doctrine of a
radical evil, an original propensity, grounded in an intelligible
timeless act, in sensuous man to invert the moral order in the
incentives to action by letting egoistic desire take precedence over the
command of the moral law. The doctrine of the Son of God becomes a
doctrine of the man well-pleasing to God: humanity in its moral
perfection; the doctrine of the atonement through Christ becomes the
doctrine that the consciousness of an imitation of the exemplar in moral
disposition gives the practical faith of being well-pleasing to God, the
knower of hearts, who beholds the disposition, the intelligible
character. The archetype that is imitated is sought complete only in
pure reason; no example in outer experience can correspond to it, since
no experience reveals the inwardness of the disposition. The acts of
cult that the positive religions prescribe, but that are not expressions
of the moral commands, have only temporary significance; if the emphasis
is laid on them, it is a false service of God and priestcraft; the
\emph{Church} is according to its idea only an \emph{ethical} community
under the divine moral legislation.

\divrule

Through Kant's Kritik der reinen Vernunft and through his Kritik der
praktischen Vernunft a sharp distinction had been drawn between the
spheres of the theoretical and of the practical, between the domain
\opage{168}where the law of necessity and that where the law of freedom
held sway, between the domain of the ``concepts of nature'' and the
domain of the ``concept of freedom.'' An insurmountable barrier seemed
to separate them from one another. But this barrier Kant breaks through
in his third main work, in his \emph{Kritik der Urtheilskraft}, in which
he subjects the faculty of feeling (the feeling of pleasure and
displeasure), as a third fundamental faculty beside the cognitive
faculty and the faculty of desire (and as a ``middle term'' between
them), and the reflecting judgment related to it, as a distinctive
function beside understanding and reason (and as a ``middle term''
between them), to a critical examination, by which the question of an
a priori legislation for that faculty too is to be answered.

In the domain of the theoretical as well as in that of the practical
there is found a point of attachment for the concept mediating between
the domain of nature and the domain of freedom,
% printed "Fribedsgebetet" (defect logged in the transcription)
which, as the a priori determination of the reflecting judgment, gives
the feeling of pleasure and displeasure an a priori rule. In the domain
of the \emph{theoretical} the possibility had shown itself of an a
priori legislation of the understanding for nature. But this
legislation comprised only the universal laws of nature, and from these
the particular ones were not supposed to be derivable; they were
supposed to be findable only by way of experience, and therefore to be
contingent for our insight. If, however, they are to be called
\emph{laws} --- as the concept of a nature and of an order of nature
demands ---, then they must be regarded as \emph{necessary} according to
a principle of the unity of the manifold, even if this principle is
unknown to us. Such a unity in the empirical, a \emph{system} of
empirical laws, is demanded a priori by our understanding, because
otherwise there could be no thoroughgoing connection of empirical
cognitions into a totality of experience. Now the principle of the
\emph{reflecting} judgment (i.e.\ the judgment that seeks the universal
for a given particular, whereas the \emph{determining} judgment
subsumes the particular under a given universal according to the
universal laws of thought) is precisely this: that the particular
empirical laws, with respect to what in them is not determined by the
universal ones, must be \opage{169}regarded according to such a unity,
as if an understanding (though not our understanding) had given them
for the sake of our cognitive faculties, in order to make possible a
system of experience of particular laws of nature. This principle is,
in other words: the \emph{purposiveness} of nature; for purposiveness
expresses the determinateness of the manifold by the unity of the end,
the mastery of the real by something ideal, by thought. Purposiveness
is the a priori principle for reflection, and as such \emph{only} a
subjective principle, a maxim of judgment, not an objectively valid
concept of nature.

From the realm of the \emph{practical} there was a pointing toward the
same concept of purposiveness, inasmuch as the demand was made that
freedom should realize in the world of sense the end determined by its
laws (its ``Endzweck''). For this demand presupposes that nature must
also be capable of being thought in such a way that the lawfulness of
its form corresponds to the possibility of the realization of such
ends, that is, that nature is determined by purposiveness. Through this
concept there is a pointing, by way of the ideal unity it presupposes,
toward a connection between the domain of nature and the domain of
freedom.

Inasmuch as nature is apprehended under the point of view of
purposiveness, it is apprehended as corresponding, in a manner
contingent \emph{for our insight}, to the aim of our cognition. But the
attainment of any aim is bound up with the feeling of pleasure, and
when, as in this case, the aim is determined by an a priori formal
principle, then the feeling of pleasure too is determined by an a
priori ground valid for all, and that merely by the relation of the
object to the cognitive faculty, without regard being taken here to the
faculty of desire. ---

If in our representations we distinguish the \emph{aesthetic}, which
designates their relation to the subject, from the \emph{logical},
which determines the concept of the object, this leads to the
distinction between aesthetic pleasure and logical satisfaction, and
between aesthetic (formal, subjective) and logical (real, objective)
purposiveness. \emph{Aesthetic purposiveness}
ex\opage{170}presses only that the form of the object, immediately in
perception (apprehension), prior to any concept of the thing, stands in
agreement with our cognitive faculties, inasmuch as the representation
of it involuntarily sets the two faculties that the reflecting judgment
brings together: imagination (the faculty of intuitions a priori) and
understanding (the faculty of concepts), in harmony, quickens them to a
concordant activity, and thereby awakens a pleasure which, as
conditioned merely by the \emph{form} and its relation to the normal
proportion of our cognitive faculties, is \emph{necessary} and
\emph{universally valid}, and leads to the judgment of \emph{taste} on
the \emph{beauty} of the object. \emph{Logical purposiveness} expresses
the agreement of the form of the thing with the previously given
\emph{concept} of the thing, which contains the ground of this form; it
is judged, not like the former by taste, but by the understanding, and
according to concepts. Natural \emph{beauty} is the presentation in
intuition of the concept of formal purposiveness, the natural
\emph{end} the expression of the concept of real purposiveness.
\emph{Aesthetic judgment} is the faculty of judging formal purposiveness
through the feeling of pleasure and displeasure; \emph{(teleo)logical
judgment} is the faculty of judging the real purposiveness of nature by
understanding and reason. But it must not be forgotten, on the one
hand, that purposiveness \emph{everywhere} stands in relation to the
faculty of feeling, and, on the other hand, that it, whether it is
called formal or real, is not a constitutive but only a regulative
principle.

The critique of \emph{aesthetic} judgment analyzes the concepts of the
\emph{beautiful} and of the \emph{sublime}. The \emph{beautiful} is
determined, as it is considered according to the various classes of
categories, as that which awakens a free, \emph{disinterested} delight;
% printed "som det, vækker" (defect logged in the transcription)
as that which, without concept, pleases \emph{universally}, as if
according to a common sense (a uniform excitability); as the
purposiveness of form, which is perceived in the object \emph{without
the representation of an end}; and finally as that in which the delight
is \emph{necessary without, however, being grounded \opage{171}on a
concept}, and demands the assent of all to an aesthetic judgment that
is apprehended as an example of a universal rule which one cannot
state. By these determinations the beautiful is distinguished from the
\emph{agreeable}, i.e.\ that which in sensation appeals to the senses,
and from the \emph{good}. The agreeable is, like the beautiful, a
subjective determination; but the delight in it is bound up with an
interest in the existence of the object, and it has only individual
validity. The good has the character of universal validity in common
with the beautiful, but it pleases through reason by virtue of the
concept, and the delight in it is in the highest degree interested in
existence. The good enters into a peculiar relation to the beautiful
in that the latter can be apprehended as a \emph{symbol of the morally
good}, so that taste thus becomes a faculty of judging the sensible
rendering of moral ideas.

The \emph{sublime} Kant determines as the absolutely great, which
pleases immediately through a resistance to the interest of the senses.
It appears where imagination enters into relation, not with the
understanding, but with reason. Sublime, therefore, is everything whose
intuition calls forth the idea of infinity. But since this can happen
only inasmuch as the incongruity of every intuition of the imagination
with the idea is felt, on the one hand the sublime proper cannot be
sought in sensible form, but only in the ideas of reason; no natural
object can itself be sublime, but can only serve for the presentation
of the sublime; --- on the other hand the sublime cannot, like the
beautiful, produce a pure pleasure, but only a satisfaction accompanied
by displeasure. For since no presentation in intuition corresponding to
the ideas of reason is possible, there arises as it were a conflict
between imagination and reason; but since precisely through that
incongruity the ideas of reason are called forth more clearly before
the mind, that conflict calls forth a subjectively purposive state in
us, a feeling that we have a pure, independent reason, whose
preeminence comes forward through the inadequacy of the faculty that is
capable of surpassing all limits. The fe\opage{172}eling that thereby
arises is respect, namely for ourselves, but which by a subreption is
transferred to the object, inasmuch as we as it were lay sublimity into
the representation of it and therefore call it itself sublime.
What we, by transferring the determination of sublimity to the object,
designate by this predicate is either the \emph{mathematically}
sublime, whose greatness is that of extension, or the
\emph{dynamically} sublime, whose greatness is that of force. By both
kinds of it we are led back to the infinity in our reason. Everything
\emph{great} in nature, even that which surpasses every measure of the
senses, always remains only relative, and the imagination, which goes
beyond every given magnitude, still cannot give a counterpart to the
idea of reason in us, which precisely through that faculty's fruitless
striving after an infinite totality becomes alive in us. Everything
\emph{mighty} in nature is mighty only in relation to us as sensuous
beings, but it brings us to become conscious of the power in us that is
not nature, but by which we are superior to everything natural and in
relation to which all the might of nature vanishes. The sublime as the
absolutely great lies only in the subject's own determination. The
\emph{feeling} for it calls forth a disposition of mind that resembles
the moral one, and it presupposes a moral development. ---

All aesthetic judgments on the beautiful and the sublime, which are at
once subjective and yet come forward with a claim to universal
validity, rest unconsciously on a concept not realizable through any
intuition: the concept of a supersensible substrate of phenomena, by
which the subjective purposiveness of nature for the cognitive
faculties is grounded, and whose unity with the supersensible substrate
of the human faculties is, as it were, divined.

This meeting of mind and nature is intimated also in Kant's definition
of art and of genius. \emph{Art} is an intermediate sphere between the
spheres of nature and of freedom. In the products of fine art we have a
work of freedom, \opage{173}but one which, in its independence of the
constraint of consciously formed rules, shows itself as a pure product
of nature. \emph{Genius} is a talent given by nature, which
unconsciously produces the exemplary and gives art a norm. It is the
faculty of \emph{aesthetic ideas}, i.e.\ intuitions (representations of
the imagination) to which no determinate concept can become adequate,
but which, inasmuch as they add to a concept an unnameable wealth of
representations that cannot be contained in it, by the feeling of this
wealth quicken the cognitive faculties in their free play and join
spirit to the mere letter of language. But the production of these
ideas therefore cannot be a conscious deduction from concepts, but must
have the character of an unconscious gift of nature. The highest
aesthetic idea is the idea of purposiveness. Only the \emph{idealism}
of purposiveness, the conception of it as having a merely
\emph{subjective} validity, makes it possible that the question of
synthetic judgments a priori with respect to \emph{aesthetic} judging
can be answered affirmatively, and that their possibility can be
explained. But such synthetic judgments are the judgments of taste,
inasmuch as they attach the feeling of pleasure as predicate to the
intuition of the object, and pronounce this connection valid for all.
---

\emph{Teleological} judgment considers nature under the point of view
of \emph{objective} and \emph{real} purposiveness. To apply this point
of view to nature, and thereby to form the concept: natural end, our
judgment is led only by such objects of experience as we are unable to
comprehend
% printed "istand til begribe" (defect logged in the transcription)
as lawful from the point of view of the mechanism of nature, but only
when we regard the idea of the effect as determining the causality of
the cause, or, in other words, when we ascribe to a concept of the
object, which is as it were laid into nature, causality with respect to
the object, or, more precisely, represent to ourselves the possibility
of the object by analogy with such a causality according to ends as we
meet with in ourselves, and thus apprehend nature as acting \emph{in
the manner of art}. (Thus purposiveness, in \emph{both} its forms,
comes into relation with the \opage{174}\emph{artistic}.) Here again
there is talk only of a regulative principle for judging phenomena, not
of introducing, in an unjustified manner, by a concept of reason, into
natural science a causality that we take from our own being, in order
to ascribe it to other beings, which we would nevertheless not
acknowledge as of like kind with us.

Only in those cases am I compelled to regard a natural product as
produced \emph{in the manner of art} when it is at once its own cause
and effect, as e.g.\ a tree is when in growth it produces itself as
individual, in propagation itself as species, or, what coincides with
the first determination, when in it the parts, as to their existence
and form, are possible only through the whole, are determined by its
idea, and are reciprocally each other's causes and effects. But these
relations we meet precisely in the \emph{organic} beings of nature, and
these we must therefore regard as natural ends, under the point of view
of \emph{inner} purposiveness, which is indeed to be distinguished from
the concept that had played so large a role in the Enlightenment's
contemplation of nature: the concept of \emph{outer} (relative)
purposiveness, i.e.\ the usefulness of one thing as a means for
another, a relation on which no absolute teleological judgment can be
grounded, since the contemplation of nature can never tell us that a
natural thing is the last end of nature. In the organic everything is
to be apprehended at once as end and means. The form of organization
may be called contingent, in contrast to the form necessary according
to mechanical laws, and precisely the contingency of the agreement
compels us to apply the concept of natural end. But since it is first
in the individual case, by particular products of nature, that we are
led to assert a teleological conception and to go beyond mechanism,
this mode of view is unavoidably extended to the whole of nature, and
we form the idea of \emph{the whole} of nature as a system according to
the rules of ends, and under this idea all natural mechanism must now
be subordinated, so that nature is seen as determined by the
all-embracing inner (conscious or unconscious?) \opage{175}end, under
which everything mechanical enters as a serving means. Without denying
the principle of mechanism we assert another principle, by which we
widen our conception of nature. And we now see aesthetic purposiveness
too in a new light, inasmuch as it enters under the point of view of
the total purposiveness of nature.

The \emph{twofold} contemplation of nature, from the point of view of
mechanism and from that of teleology, is necessary for \emph{our}
reason, according to its peculiar constitution. We must \emph{judge}
all production of material things and their forms as possible according
to merely mechanical laws, but nevertheless we must, in judging certain
natural things, resort to the concept of end, the concept of causæ
finales. Were we to transform this subjective necessity into a dogmatic
assertion about the coming-to-be of things, it would lead to the
insoluble contradiction that the production of \emph{some} things
should not be possible according to those laws alone which were to
explain \emph{all} production. If, on the other hand, we regard the two
maxims for judging merely as regulative principles, they can subsist
side by side. We then continue the inquiry after the mechanical
connection as far as is possible for us, and do not deny that it
extends beyond the limit to which we can pursue it, but we acknowledge
the necessity of reflecting on certain natural products, and, on
occasion of them, on the whole of nature, according to another
principle, the principle of final causes. Our reason is not able, even
though it postulates the subordination of mechanism to teleology, to
know the unity of the two principles. Our discursive cognition, for
which the whole results in a mechanical manner from the cooperation of
the parts, and which proceeds from the particular to the
\emph{abstract} universal, cannot see how the whole, as a
synthetic-universal by which the particular, the parts, is necessarily
determined, can be the ground of the constitution of the parts, the
prius that explains them. We can only think that the
\emph{representation} of a whole contains the ground of its form and of
the connection of the parts; but that which has its \opage{176}ground
in the representation of it we call an \emph{end}, and this concept our
cognition cannot do without in the face of certain phenomena of nature;
but it will not be able to bring the consideration according to it into
unity with the consideration according to the necessity of mechanism.
Only for a differently constructed understanding: for an
\emph{intuitive understanding}, an \emph{intellectual intuition}, for
which intuition was not separated from concept, its actuality not from
the possibility of thought, and which could intuit in the whole the
ground of the parts and of their concord, would what we separate as
mechanism of nature and connection of ends be one, and the concept of
end in our conception would not become necessary for it. But such an
intuitive understanding, which passes over from the
synthetic-universal to the parts, and for which therefore their concord
is not contingent, we must think, because only by it can we explain to
ourselves the concord of natural objects with the a priori laws of our
judgment. We represent to ourselves the possibility that in the ground
of nature unknown to us the physical-mechanical connection and the
connection of ends in the same things could cohere in one principle.
And this ground of nature inaccessible to our cognition (the
``supersensible real ground of nature''), this \emph{intelligible
substrate} that hides behind the phenomenon of nature, is, as was
already intimated in the development of the concept of aesthetic
judgment, seen in unity with ``the supersensible substrate of all the
subject's faculties,'' and designated as the \emph{common} intelligible
substrate of nature and of the human mind, as the intelligible, the
supersensible, which is the common deepest ground of both.

In this concept, which becomes the highest concept in the Kritik der
Urtheilskraft, the central concepts in the Kritik der reinen and der
praktischen Vernunft have found their unfolding. The central concept in
the Kritik der reinen Vernunft was the concept of \emph{transcendental
apperception}, the original unity of spontaneous self-consciousness,
the source of all the a priori forms, and through them of all necessity
in our cognition. But this self-consciousness was raised above the
\opage{177}limitation of individuality that clings to the word,
inasmuch as it was apprehended as the \emph{pure form} of
self-consciousness and as such as what is identical in \emph{all}
self-consciousnesses, as ``consciousness in general,'' i.e.\ as the
substrate of individual consciousness, as the connecting unity in the
individual. And a unity was intimated between this subjective and the
objective, the thing in itself, inasmuch as the possibility was
intimated that the latter, the unknown X, might be the same in all the
phenomena of outer sense
% errata list B, p. 177: printed "i alle Phænomener", corrected "ved alle den ydre Sandses Ph[ænomener]"
and not different from the substrate that lay at the ground of the
phenomena of inner sense, and inasmuch as the significance of the
unknown thing for our cognition became the same as that of
self-consciousness: to bring the determination of necessity into it.
In the Kritik der praktischen Vernunft
% printed "Vernuft" (defect logged in the transcription)
the central concept was the concept of \emph{homo noumenon}, whose
essence is freedom, whose legislation, as a legislation of humanity in
the individual human being, is valid for all individuals. Here the
determination of the universal is brought out more distinctly, and the
problem of the relation of the mental to nature
% errata list B, p. 177: printed "Sandseliges", corrected "Aandeliges"
is touched more determinately by the demand that the moral should be
realized in nature, that freedom should acquire a causality in relation
to nature. But in the Kritik der pr. V., in accordance with the sharp
critical distinction between the domain of the concept of freedom and
that of the concept of nature, the realization of the moral in nature
was still posited as conditioned by a postulated third: an almighty and
omniscient God. This notion of God now gives way, in the Kr. der
Urtheilskraft, to the concept of that which does not, as standing
outside the oppositions, connect them only in an arbitrary manner, but
is their \emph{ground of unity}, even if it is determined in the first
instance negatively, in a form that anticipates the Schellingian
concept of indifference. Since the law of mind and the law of nature
have their origin in the same being, the realization of the moral in
nature becomes more than merely a demand, and when the Kritik der
Urtheilskraft in both its main divisions as it were issues in the
ethical, through the conception of the beautiful as the symbol of the
moral, and through the concept of an ``Endzweck'' of nature that is to
be sought \opage{178}beyond nature itself, in the moral, this does not
lead us back to the standpoint of the Kr. der pr. Vernunft; but that
fundamental relation between the two laws which is decisive for the
actuality of the moral is conceived speculatively, grounded in the
deeper underlying, the encompassing, whose unity of essence is seen not
merely as the source of the oppositions but also as their
reconciliation. With this concept Kant has touched the goal toward
which the development of critical idealism leads; but inasmuch as he,
true to his peculiarity, protests against every attempt of our
\emph{cognition} to draw it into its circle, and gives what is
intimated a merely subjective validity, he holds back the consequences
that are then drawn by Schelling and Hegel.

\divrule

The \emph{lasting significance} of the Kantian philosophy we must
seek in what gives it its basic stamp: in its character as criticism,
and as a criticism that points toward the speculative as its goal. As
Kant, with critical consciousness and with a deliberate, mediating
mind, at once poses in the domain of the theoretical the decisive
problem of knowledge which dogmatism had evaded, he definitively
asserts against the encroachment of empiricism, through the critical
distinction between the factors of knowledge, which lets empiricism
keep its \emph{relative} right, the self-activity of the knowing mind,
the validity of the a priori, and with it a universally valid and
necessary knowledge. Through the critical distinction between the
world of the intelligible and that of the phenomenon he asserts in the
domain of the practical, by positing freedom as the essence of mind
and the law of action, the purity and independence of the ethical, the
realization of the ethical demand as the final, unconditioned end. And
in his last great and brilliant work Kant again, by sharp
conceptual distinctions, makes possible an ideal conception of the
beautiful, and through the critical distinction the thought of the
deeper unity of what first had to be \opage{179}separated in order
that the essence of mind could be asserted here breaks forth more
clearly than in the intimations to which, in the earlier writings, the
consequences of the fundamental problem had led him. The thought which
he had come nearest to before when, in the Kr.\ d.\ pr.\ V., reason, as
the expression of our \emph{own} free essence, of the absolute in us,
was seen as legislating and determining, whereby there was pointed not
only to the \emph{concretion} of the idea but also, since the highest
end of freedom was to be realized in nature, to the deeper harmony of
nature and mind; --- that thought makes itself mightily felt in Kant's
conception of the artistic and of genius, and of the inner teleology of
the organism, and it as it were forces its way to a definitive
expression in the determination: ``the common intelligible substrate of
nature and mind.''

But the limitation and defect of the Kantian philosophy spring from
the same source as its merit. The \emph{critical} moment in Kant, which
was the necessary condition for the solution of his peculiar task,
nevertheless gains a \emph{one-sided} preponderance over the
speculative, which comes forward only through intimations that
knowledge is held back from pursuing. The one-sided sharpness of the
critical distinction first appears at the beginning of the Critique of
Reason in the fallacy that what shows itself as originally belonging to
the subject must therefore be posited as belonging \emph{only} to it
and as excluded from the essence of the objective, so that the
subjective and the objective are \emph{abstractly} distinguished. In
this first abstract distinction all the decisive consequences lie
hidden.

In the \emph{theoretical} domain this distinction, which includes the
sharp separation between the form and the content of thought, between
a priori and empirical, leads to this: that being escapes thought,
since what is in itself, as something without determination, withdraws
from the forms of intuition and of thought, which are nevertheless
involuntarily applied to it, since as something existing \emph{outside}
us it is to \emph{affect} our sensibility; and that this which is in
itself, and the unconditioned, which is its most pregnant expression,
is not posited as an object of theoretical knowledge even where it
gains the \opage{180}certainty of fact. Kant does not shrink even from
the harshest consequences of the distinction between what is in itself
and the phenomenon, as when he asks: do \emph{I} myself, as phenomenon
for inner sense, exist outside \emph{my} power of representation, in
time?, and when he draws the conclusion that, since the I as pure I
receives determinations only through the categories, it is therefore
not determined according to its being in itself, and this although the
categories, as a priori forms of thought that are conceived as
proceeding from pure self-consciousness, would by this relation have to
be the essential forms of expression of the pure I, determinations of
its essence, through which it was to come into a relation of
consciousness to itself. And these a priori forms, which are all
postulated to go back to the unity of self-consciousness, cannot in
reality be deduced from this unity in its want of determination, but
must be taken up empirically, and remain standing in separated
multiplicity, without being led back to a unity.

In the \emph{practical} domain the abstract distinction between the a
priori and the empirical, between the sphere of the noumenon and that
of the phenomenon, between the domain of freedom and the domain of
nature, leads to the abstract and individuality-less formalism of the
moral law (which only, as it were, afterwards receives an ---
illusory --- filling-out through regard for rational beings as ends in
themselves), to the contentlessness of duty, to the merely negative
determination of freedom, to the conception of the ethical task not as
an objectification of freedom in the concrete forms of ethical life,
but as an annihilation of impulse, which constantly remains only a
demand, an ``ought,'' that is never fully realized, since the cessation
of the struggle would be the cessation of morality itself. And while
in the theoretical domain the notion of an influence of the thing in
itself on our feeling brought contradictions into Kant's fundamental
determinations, so in the practical domain the demand for a causality
of freedom, of the intelligible, in relation to nature, while the
opposite causality is excluded, makes the relation between the spheres
unclear and leads to a \opage{181}mystical doctrine of the
supra-empirical self-determination of freedom as the starting point for
the determinateness of the phenomenal human. And the demand for the
harmony between virtue and happiness, the positing of which is an
inconsistency, since happiness has no relation to man as noumenon,
receives no solution for thought through the postulate of a God with an
incomprehensible causality.

And even where Kant, in his Kritik der Urtheilskraft, most definitely
intimates the speculative, the unity of opposites, the original
distinction asserts itself again in that the intimations receive a
merely subjective turn and their consequences are cut off. Kant's
developments concerning natural teleology remain without influence on
his metaphysics of nature; his intimation of the unity of nature and
mind in their ground is not brought into the forms of the concept. But
in Kant's intimations of the unity of opposites, which already
glimmered forth through the central concept of the Critique of Reason,
incitements are nevertheless given for the development of criticism
into speculative thinking, and the subsequent speculative systems have
their root in Kant's philosophy.

\divrule

Kant's philosophy was epoch-making upon its appearance. It gave the
incitement to a development of thought that, with unexampled rapidity,
led the whole idealistic line of modern philosophy to its conclusion,
and by the wealth of possibilities of thought that it contained, it
could also become the starting point for realistic consequences, for a
philosophy that forms the opposite of the idealism carried through in a
subjective direction. But it also had to call forth attacks, not only
from the earlier philosophical schools, which defended themselves
against what was new and transforming in it, but also from such
thinkers as did not feel satisfied by the limitation of our theoretical
knowledge that criticism taught, or by the practical substitute for it
that it offered. Those quickly \opage{182}silenced attacks had only a
subordinate interest; these, which proceeded from what has been called
the \emph{philosophy of faith}, had a more essential one, since they
touched weak points in Kant's doctrine and, albeit only in the form of
the postulate, applied correctives against one-sidednesses in
criticism. We therefore, before we pass on to the systems that further
develop what Kant founded, briefly discuss the most prominent
representative of the philosophy of faith: Friedrich Heinrich Jacobi.

\divrule

\begin{center}
  \textbf{2. The Reaction against Criticism from the Philosophy of Faith: \emph{Jacobi}.}
\end{center}

\emph{Friedrich Heinrich Jacobi} was born in Düsseldorf in 1743, and
died as President of the Academy of Sciences in Munich in 1819. Of his
works we mention: Etwas, was Lessing gesagt hat, 1782; Briefe über die
Lehre des Spinoza, 1785; David Hume über den Glauben, oder Idealismus
und Realismus, 1787; Jacobi an Fichte, 1799; Ueber das Unternehmen des
Kriticismus, die Vernunft zu Verstand zu bringen, 1801; Von den
göttlichen Dingen, 1811; the philosophical novels: Allwills
Briefsammlung, 1792, and Woldemar, 1799. In the introductions to
volumes 2 and 4 of the Sämmtliche Werke a comprehensive survey of
Jacobi's doctrine is given. ---

Jacobi's philosophy is, in contrast to the thinking that seeks
knowledge only for the sake of knowledge (``der reine Vorwitz''), by
his own declaration a philosophy with a definite \emph{tendency}. What
moves his thinking is the urge to discover God as the first ground of
all our knowledge and to find him again everywhere; it is a striving
``not after truth in general --- a non-thought --- but after a
\emph{definite} truth, satisfying head and heart.'' His philosophy is
characterized by the following utterances: ``One seeks and loves truth
for the sake of its content. A truth that kills, that annihilates man,
man can neither \opage{183}seek nor love.'' ``All men in their inner
being call something truth in advance, something which they are not
yet in possession of, after which they strive, and which they could
nevertheless not presuppose unless it were in one way or another
present to them. --- If this presupposing of truth is a merely
subjective illusion in rational beings, if they have no intuition of
it, not even the dimmest, then their inquiry is altogether
fruitless.''

It is this urge toward a definite truth, this ``\emph{interest} in the
supersensible,'' that prevents Jacobi from settling down in a merely
``demonstrative'' philosophy, whether this has the character of
empiricism or of dogmatic idealism. Jacobi denies that demonstration
can reach beyond the conditioned and finite to the unconditioned,
infinite and supersensible, and he declares himself in complete
agreement with Kant as regards the limitation of our knowledge (our
conceptual knowledge) to the finite. Demonstrating thought does not get
beyond the mechanism of nature, for it only links together the
infinite chain of conditioned conditions. At most it can, since the
principle of ground (principium rationis sufficientis), on which all
demonstration rests, really --- as a reflection on mathematical
demonstrations shows --- comes to the same thing as the proposition:
totum parte prius est [the whole is prior to the part], reach the
concept of a universe, but not the concept of a pretermundane,
``living'' God. Since all its proving and comprehending is an
explaining by a conditioning other, precisely the attempt to
demonstrate, i.e.\ to ground, the supernatural, the infinite and
unconditioned, would be one with making it something natural, finite
and conditioned, for what is grounded becomes dependent on the ground;
precisely as grounded, it cannot be something unconditioned. If
knowledge is only the conviction won by demonstration, then it is a
contradiction that there should be a knowledge of God's existence. A
God who was the object of such knowledge would be no God. ``It is in
the interest of science that there be no God.'' \opage{184}Thus the
demonstrative philosophy necessarily leads to atheism and materialism,
and it leads at the same time to fatalism, since in the series of
conditions there is no room for a freedom. Spinozism, with its
deification of the universe, is the consistent philosophy. It is in a
certain sense unassailable: what it denies cannot be strictly
philosophically proved, what it proves cannot be strictly
philosophically disproved. Not by the understanding can one get beyond
its circle, only by a leap, by a salto mortale from the understanding
to faith, from the proving knowing of the natural to the immediate
certainty of the supernatural, which reveals itself to us as a reality
in us and above us. The organ for the supernatural Jacobi designates
now as a sense, as the immediate sense for existence, now as a
feeling, now as intuition, now as faith, now as reason, inasmuch as
this concept, which in his earlier writings was not definitely
distinguished from that of the understanding, is now to designate the
faculty of immediately feeling the supersensible, the faculty of
presupposing the true, good and beautiful in itself with full
confidence in the objective validity of this presupposition. Only an
immediate sense reveals existence to us, and this holds not only of
spiritual, supersensible existence, but also of sensible existence,
whose certainty is given in the very act of sensation. The
understanding comprehends only relations, connections, only forms a
\emph{given} content; as a merely reflecting and analyzing faculty it
never reaches further than to a \emph{relative} being; it states only
that something is like something else in concept, but not its
substantial being, its real being, its being ``schlechthin.'' This
being makes itself known only to the immediate sense: to sensible
intuition and to reason. But the understanding is necessary for sense,
both as sensation and as reason. The understanding is the faculty that
\emph{forms} our cognitions, whose material is given to it: ``without
understanding we would possess nothing in our senses; there would be no
force that united them in itself; the sensible being \emph{itself}
would \opage{185}not exist. And in the same way, without understanding
we would possess nothing in our reason; the rational being
\emph{itself} would not exist.'' ``The \emph{consciousness} of a reason
and of its revelations is possible only in an understanding.'' But
still man is raised above the merely animal being solely and alone by
the property of reason. ``The religious feeling is the foundation of
humanity.'' ``Reason is not grounded on the power of thought, a light
that is kindled only later in the understanding, but the power of
thought is grounded on reason, which, where it is, shines for the
understanding and awakens it to the contemplation on which
investigation follows, distinct knowledge, \emph{science}.''

What is first and foremost shown to us through the sense for the
supersensible is \emph{our freedom}, our power over the nature that
surrounds and permeates us, our independence and \emph{personality},
and thereby our \emph{immortality}. Freedom and reason become for
Jacobi two expressions for the same thing: for the capacity to raise
oneself above the conditioned and above natural necessity. Without
freedom the beautiful, the true and the good would be a deception;
admiration, esteem and love could not be conceived. We have an
immediate certainty of the reality of our freedom, even though it
constantly remains incomprehensible to our understanding, and its
connection with natural necessity in our being an insoluble riddle. But
in the certainty of the ``creative freedom'' \textbf{in} us we have the
certainty of a personal being \emph{above} us, of an almighty God
exalted above nature, whose image man is. Man finds God because he can
find himself only in his finding of God, the primal source of our life
and of our reason. The freedom of which we become conscious in our
rational personality must be present superabundantly in the giver of
our mind, in God, if nature has issued from him and not he from nature.
God is God only as spirit, and spirit he would not be if he lacked
self-consciousness and personality. Our consciousness of freedom, our
ethical consciousness in general, and our religious faith mutually
support one another. ``When I am \opage{186}convinced of the
objectivity of my feelings of the beautiful, the sublime, the true and
the good, and of a freedom in me that masters nature, I am convinced
of God's existence, and when these feelings grow faint in me, faith in
God also grows faint.'' ``By a divine life man gains the certainty of
God. Without morality there is no religiosity.'' But conversely it must
also be said: ``Our ethical convictions all perish when the ethical
primal being vanishes for us as an ethical, that is personal, being,
who wills and works the good.''

The determinations of God as the free, living, self-active,
\emph{personal} God we thus derive from what proclaims itself in
ourselves as something supernatural in us. ``When God created man, he
theomorphized; it is therefore necessary that man anthropomorphize
when he would know God.'' But the conception of God by analogy with man
may nevertheless, according to Jacobi, lead neither to such a holding
fast to this analogy that a development, a becoming, came into the
divine personality, nor to an attempt to fix God's essence for
knowledge in definite concepts. \emph{That} God is, is certain;
\emph{what} he is remains hidden from us. What is the object of faith
is always a miracle and a mystery, incomprehensible to the
understanding, and this incomprehensibility is insurmountable. The
understanding has only a ``non-knowing knowledge'' of God. Faith is a
reflection of the divine knowing and willing in man's finite mind.
``Could we transform this faith into a knowledge, then what the serpent
in Paradise promised the lustful Eve would be fulfilled: we would be as
God.''

In the recognition of God as the free, personal and creative God,
Jacobi sees what is essential in Christianity. ``Christianity is
essentially anthropomorphistic, paganism cosmotheistic.'' Only
Christianity is therefore religion; outside it there is only atheism or
idolatry. Jacobi calls himself a Christian, but to believe in Christ is
for him to \opage{187}believe \emph{what} Christ believed, and \emph{as}
Christ believed. Revelation is for him the inner revelation in the
reason of men: the true God cannot appear outside the soul, and we know
no other God than the one who became man in us. Faith in a historical
God-man is for him
% errata list B, p. 187: printed "en historisk Christus", corrected "et historisk Gudmenneske"
religious materialism. Christ as God-man is the \emph{product} of
Christian faith; the pious disposition lays its ideal into the
historical phenomenon. Therefore Christianity, ``conceived in its
purity,'' can for Jacobi reach beyond the domain of the Christian
Church. Socrates can become a hero of faith, and ``the God of Socrates
and of Plato'' and ``the God of the Christians'' become identical.
Jacobi's religion becomes a dogma-less mysticism, which, both with
respect to the conception of the form of revelation and with respect to
the content of faith, parts from positive Christianity.

\divrule

By the fundamental thoughts of Jacobi that have been developed, his
\emph{relation to the Kantian criticism} is determined. He has
essential points of contact with it, inasmuch as both Kant and he limit
the domain of the knowledge of the understanding and exclude it from
the supersensible, and inasmuch as they both assert a validity for
faith, and through it secure for consciousness the higher, which reason
demands but the understanding does not comprehend. But Jacobi departs
from Kant in three main points. He disapproves, first, of the way in
which Kant grounds the bounds of our theoretical knowledge by his
distinction between the phenomenon and the thing in itself. Against the
Kantian idealism he sets a realism: the real being of objects is for
him immediately secured by sensation, whose truth must be held fast,
% errata list B, p. 187: printed "i Sandsningen", corrected "ved Sandsningen"
even if as an ``incomprehensible miracle.'' It is not enough, with Kant,
to posit something objectively existing that, as a ``thing in
itself,'' does not reveal itself to us through its phenomenon, and
whose affecting of us becomes unthinkable by Kant's own determinations,
by which the relation of causality that lies in the determination
affection, as a merely subjective form of thought, is excluded
\opage{188}from the sphere of what is in itself. ``One cannot,'' says
Jacobi, ``get into Kant's system without this concept of affection, and
one cannot remain in it with this concept.'' Consistently, Kant would
have had to come to the Fichtean idealism, which lets the I be the sole
source of the content of our consciousness. But this idealism, which
makes all objective reality into a subjective semblance, destroys all
truth and transforms all our knowing into the I's empty play with
itself. --- Jacobi deviates from Kant, secondly, in the conception of
faith. The Kantian rational faith is an assumption in accordance with a
\emph{practical} need: a ``faith of necessity,'' in Jacobi's
expression; its content can, says Jacobi, consistently be only
postulates, practical tasks, and Fichte is in this respect too more
consistent than Kant. But Jacobi's faith also has a theoretical
character; it is not merely a faith in something that \emph{shall} be,
but in something that \emph{is}. The idea is something that is
\emph{experienced}; the unconditioned is not a merely \emph{subjective}
demand, but a being, and only thereby does a real faith in it become
possible. The true philosophy must be a real theism, must be faith in a
being that does not become, but is, as a being complete from the very
beginning, pretermundane, subsisting for itself. --- Finally, in the
practical domain, while he recognizes the purity of the Kantian
morality, Jacobi opposes the contentless formalism of the Kantian moral
principle, ``the will that wills \emph{nothing}, this hollow nut of
independence,'' the conception of the ethical demand merely as a
commandment of law, and the abstract universality of the principle,
which violates the right of individuality and of personality. He
demands a recognition of the significance of the ethical fundamental
\emph{impulse}, of the significance of inclination, feeling and passion
for the individual peculiarity of ethical life, of the claim of the
heart and of the individual's royal right to a privilegium aggratiandi
over against the pure letter of the universal law of reason. Feeling,
subjective need, the immediately given impulse here too becomes for him
the decisive thing.

Toward the systems that issued from the Kantian one, Ja\opage{189}cobi
places himself in a sharper opposition. He grants \emph{Fichte} the
greater consistency in comparison with Kant, but just as we have seen
him protest against the one-sided idealism, so we see him lodge an
objection against the substitution of a moral world-order for a
personal and conscious God. Against \emph{Schelling} he objects that
his philosophy effaces the difference between necessity and freedom,
between nature and spirit, and is only a Spinozism veiled under
theistic forms, whose concept of a becoming God, of a God who reaches
from the less perfect to the more perfect, is an absurdity. Whoever
starts from nature finds no God: ``God is the first or he is not at
all.''

\divrule

The justification of this peculiar philosophy lies in this: that it
asserts the original presence of the absolute, which is of one essence
with us, in the human mind determined by reason, which has its being
from that absolute; that, against a one-sided knowledge of the
understanding, which loses itself in the finite and conditioned and
does not understand the essence of mind, it applies a corrective
proceeding from self-consciousness, ``brings facts to light''; and, as
regards its polemic against Kant, in its demonstration of the
contradiction in his starting point, and in its assertion of the right
and significance of concrete personality in the domain of the ethical.
But its defect lies in this: that it applies its correctives only in
the subjective form of the postulate, and that its one-sided
conception of the form of knowledge against which it polemicizes
prevents the content of immediate certainty from gaining concretion and
universal validity, and leaves an unreconciled dualism behind in
consciousness. It excludes the philosophy of the concept from the
knowledge of the unconditioned, because a comprehending of it is
supposed to be a grounding that reduces the unconditioned to something
conditioned, and it forgets, although it involuntarily makes use of
it, the way of knowledge that goes from the conditioned finite to that
which is the unconditioned origin of the finite, which is known through
what the given necessarily presupposes in order to be understood in
\opage{190} \setcounter{footnote}{0}
its dependent being. And since, after having made this exclusion, it
abstractly distinguishes between immediate certainty and mediate
knowledge, it misjudges the psychological genesis of the notions that
announce themselves with immediate certainty; it misjudges the
possibility of the conversion, through appropriation, of what is at
first mediate into something immediate; it forgets that the immediate
certainty of the supersensible is mediated by an elevation above
sensibility, and that the concrete unity which joins opposed
determinations together in itself contains them in such a way that
each of them is conditioned, mediated by the other, which is the same
as saying that the unity through this mediation closes together with
itself in a higher immediacy, thus at once a mediation and an immediate
relation to itself, and in mediation itself sublates mediation. Thus
the absolute is known as spirit only when it is comprehended as
mediating itself with itself within itself. The abstract distinction
between immediate certainty and mediate knowledge leads unavoidably to
the consequence that the content of faith either becomes only the
abstract, the indeterminate, since every \emph{concrete} determination
must lead into the domain of mediate knowledge, or that, where concrete
determinations are applied, this happens with the reservation that
they are to be the incomprehensible, and that by their immediacy they
must take on the character of the contingent, the arbitrary and the
subjective (just as in the practical domain, too, the immediate
feeling's filling-out of the abstract moral law becomes contingent and
individualistic). We have also seen how Jacobi, with respect to the
content of faith, moves in determinations that either immediately
express something abstract, or do so by not being able to be
explained, i.e.\ unfolded in definite concepts\footnote{Despite the
  urge toward a definite truth, Jacobi comes to dwell more on the
  problem of the forms of knowledge than on the unfolding of the
  content, more on the How of revelation than on its What.}. Thus for
the divine he has either only \opage{191}the abstract determination
\emph{that} God is and \emph{is} unconditioned, or, where he takes
determinations from the human, they may not, and are not supposed to,
be convertible into definite cognitions, and these uncomprehended
determinations then receive their interpretation only from the
subject's individual feeling and individual need. Every content can by
this way of immediacy assert itself with equal right, as the history of
religion shows. --- From the uncomprehended divine, moreover, no
comprehending of its relation to man and to nature can become
possible, and Jacobi himself admits that the assumption of a personal
God in no way makes the existence of the universe
\emph{comprehensible}. But in this way there comes to be in
consciousness a doubly separated content: a content for immediate
certainty, which grasps the reality of the supersensible but cannot
through it obtain an explanation of actuality, and a content for
mediate knowledge, which, as it remains within the connection of
nature as the all-embracing and excludes the supersensible, cannot be
refuted; and there is an unreconciled dualism present in the
consciousness that is to contain these opposites, to which Kant already
drew attention when he said: one can believe what the understanding
can neither prove nor disprove, but one cannot believe that whose
opposite one thinks the understanding can prove. In Jacobi's
personality this discord comes to light; he was ``pantheist with the
head, mystic with the heart.''
% printed "Jacobis's" (defect logged in the transcription)
In a letter to Hamann (from 1783) the dissonance within him finds its
characteristic expression: ``There is light in my heart, but when I
want to bring it into the understanding, it goes out. Which of the two
clarities is the true one: that of the understanding, which indeed
shows fixed shapes, but behind them only a bottomless abyss? or that of
the heart, which indeed shines upward full of promises, but does not
give a definite knowing? Can the human mind grasp truth unless those
two clarities unite in it into one light? And is this union
conceivable otherwise than by a miracle?''

\divrule

\opage{192}
\begin{center}
  \textbf{3. The Development of Criticism into Subjective Idealism: \emph{Fichte}.}
\end{center}

\emph{Johann Gottlieb Fichte} was born in 1762 at Rammenau in Upper
Lusatia, of poor parents. With the support of Baron v. Miltiz, a
Silesian nobleman, he came to the school at Schulpforta and in 1780 to
the University of Jena, where he studied theology. His first studies of
philosophy led him to Spinoza, who made an overwhelming impression on
him, without, however, satisfying him. On v. Miltiz's death he was
brought into a very depressed situation, lived for some years meagerly
in Leipzig by teaching, and then in 1788 accepted a post as private
tutor in Zürich, where he came into contact with Lavater and
Pestalozzi. In 1790 he came back to Leipzig and now became acquainted
with the Kantian philosophy, which brought him beyond the determinism
to which Spinoza had led him. Already the following year he wrote from
a Kantian standpoint his ``Versuch einer Kritik aller Offenbarung,''
which he laid before Kant in Königsberg, and by which he won the
latter's esteem and goodwill. When the work appeared anonymously in
1792, it aroused general attention and was generally taken to be by
Kant himself. As soon as the author's name became known, his literary
reputation was established, and in 1794 he received a call to Jena as
professor of philosophy. Here he now developed, in his lectures and in
a series of writings, his distinctive doctrine, which he originally
regarded as identical with the Kantian according to its spirit, if not
according to its letter, until Kant pronounced a sharply condemning
judgment on the subjectivism of the Fichtean Wissenschaftslehre.
Fichte's stay in Jena, at whose then flourishing university he held a
prominent position, was only brief. In 1799 an essay of his in the
journal edited by him and Niethammer gave occasion to an anonymous
denunciation for atheism. The government in Weimar intervened in the
matter and dismissed him when, in a letter to one of its members, he
had declared that he would leave the university if, as was intended, a
sharp reprimand were given him. Fichte \opage{193}now went to Berlin,
where he lived in association with the brothers Schlegel, Tieck and
Schleiermacher, and gave philosophical lectures to a numerous audience.
In 1805 he received from the Prussian government a professorship at
Erlangen, but lectured there for only one half-year, as the events of
the war brought him to leave the city. He went to Königsberg, where he
appeared as a lecturer and worked on his fiery ``Reden an die deutsche
Nation,'' which he delivered in Berlin in the winter of 1807---8, while
the French troops were still in the city. When the University of
Berlin, whose founding he had eagerly urged, because through a
reshaping of moral and scientific education he wished to provide the
condition for the regeneration of the German people, was established in
1809, he became professor at it, and unfolded a restless activity as a
teacher, working unceasingly and energetically to give his philosophy,
in which an essential transformation had already taken place at the
beginning of his stay in Berlin, an ever more determinate and more
satisfying form. The War of Liberation aroused Fichte's lively
enthusiasm. During it his wife, in nursing the wounded, had contracted
a nervous fever; Fichte was infected by her and succumbed to the
illness in January 1814.

Of Fichte's \emph{writings} we mention --- besides his Kritik aller
Offenbarung and the pamphlets written in Switzerland: Zurückforderung
der Denkfreiheit von den Fürsten Europa's, die sie bisher unterdrückten
(Heliopolis, im letzten Jahre der alten Finsterniss) and: Beitrag zur
Berichtigung der Urtheile des Publicums über die französische
Revolution, 1793; --- the essay: Ueber den Begriff der
Wissenschaftslehre oder der sogenannten Philosophie, 1794; Grundlage
der gesammten Wissenschaftslehre, 1794; Vorlesungen über die Bestimmung
des Gelehrten, 1794; Grundriss des Eigenthümlichen der
Wissenschaftslehre, 1795; Grundlage des Naturrechts nach Principien der
Wissenschaftslehre, 1796; Einleitung in die Wissenschaftslehre, and:
Versuch einer neuen Darstellung der Wissenschaftslehre, 1797; System
\opage{194}der Sittenlehre nach den Principien der Wissenschaftslehre,
% printed "Wissenschaftlehre" (defect logged in the transcription)
1798; Ueber den Grund unseres Glaubens an eine göttliche
Weltregierung, 1798; Appellation an das Publicum gegen die Anklage des
Atheismus, 1799; die Bestimmung des Menschen, 1800 (herein the first
traces of his later standpoint); Sonnenklarer Bericht an das Publicum
über das eigentliche Wesen der neuesten Philosophie, ein Versuch, den
Leser zum Verstehen zu zwingen, 1801; Grundzüge des gegenwärtigen
Zeitalters, 1806; Anweisung zum seligen Leben, 1806; Ueber des Wesen
des Gelehrten, 1806; Reden an die deutsche Nation, 1808. In the
Nachgelassene Werke (3 vols.) published by his son there are found,
besides several smaller essays, his introductory lectures to the
Wissenschaftslehre (from 1813), three different courses of lectures on
the Wissenschaftslehre (from 1804, in which the new standpoint is first
determinately developed, from 1812 and 1813), lectures on transcendental
logic (from 1812), on ``die Thatsachen der Bewusstseins'' (from 1813),
on the system of the doctrine of right (from 1812), on the system of
ethics (from the same year) and on the vocation of the scholar (from
1811). In the collected edition of Fichte's writings (8 vols.) there is
found an essay written against Schelling: Bericht über den Begriff der
Wissenschaftslehre und die bisherigen Schicksale derselben. The
lectures on the doctrine of the state, or: on the relation of the
primal state to the kingdom of reason (from 1813), were published some
years after Fichte's death, likewise the work: ``Ueber den Begriff der
wahren Krieges''.

\divrule

In Kant the critical distinction had been made between subject and
object, and the claims of both had been given their due, inasmuch as
knowledge was to draw elements from both of them. The I was to give
the form, the thing the matter, and knowledge was to be the
determination of this given matter by that form. But since the form was
the \emph{subject's} form, and \emph{only} the subject's, the object was
seen as transformed by its reception into it. The thing, as it showed
itself to our knowledge, clothed in those forms, was not the thing as
it is in and for itself. But this thing in and for \opage{195}itself,
which stood outside the forms of intuition and of thought, could only
become the unknown, determinationless x, the bare form of objectivity,
which really expressed only the moment of necessity in our thinking,
and in its abstraction from all determinateness stood close to the
abstract fundamental form of subjectivity. In the transcendental
apperception, which is precisely this abstract, there lay, after all,
on Kant's conception, the necessity of connection, inasmuch as the
synthesis of self-consciousness was the source of the syntheses of the
categories, the necessary connections of the objects of experience.
And by the intimations touched on earlier (p. 177) the thing in and for
itself was drawn still nearer to the unchangeable center of
subjectivity. But in Kant it nevertheless kept its subsistence, its
objective validity, and Kant rejected every attempt to carry through
what lay in those intimations.

But in that assumption of an objective x, independent of us and
affecting our sensibility, there lies, as Jacobi already showed, a
contradiction on Kant's own presuppositions, inasmuch as the intuition
of space, which is included in the determination of the thing's being
\emph{outside} the subject's consciousness, and the concept of
causality, which lies in the notion of the subject's \emph{affection}
by the thing, were to be only subjective forms, which have validity for
the \emph{phenomenon}, not for what lies behind it. And the very notion
of something existing outside consciousness includes in and of itself
a difficult problem. In it I attempt to go out beyond my consciousness
to ``something represented that is not represented'' (Fichte), to
abstract from my consciousness, but in the attempt I yet remain within
consciousness, inasmuch as the whole act takes place only through
consciousness, and the determination I give the thing is a
determination of consciousness. For I determine the thing, the
undetermined existent, as \emph{objective}, and mean thereby to express
a being independent of the subject and of thought, but this
determination is only a necessary \emph{thought} of the
\emph{subjective} consciousness, which seeks the \emph{cause} of the
phenomenon; it is in and \emph{for} the consciousness that comprehends
the opposition between the subjective and the objec\opage{196}tive,
which, by thinking itself as subjective, fixes the opposition between
the subjective and the objective in itself. The ``objective'' object is
thus an object of consciousness, is in my consciousness; but as a
determination in my consciousness it is a determination of my I, is my
determination. If this consequence is drawn and seen as the expression
of the \emph{whole} truth, then the object is thus absorbed into the I
that posits it. The object, which Kant had reduced to the abstract
determinateness of an independent being, loses this being too, and
comes to be seen as that which, as to form and matter, is posited by
the I, and thus not as an external, independently existing object, but
as the I's opposite within the I itself, as the \emph{not-I}, which is
only in the relation to the I and through the relation to the I.

This consequence, which \emph{Fichte} draws, had already been prepared
by \emph{Jacobi's} and G. E. \emph{Schulze's} critique of Kant's
concept of the thing in and for itself as affecting feeling, by
\emph{Reinhold's} attempt to find the point of unity for sensibility and
understanding, and by the sharp-thinking \emph{Salomon Maimon's}
determination of receptivity as unconscious spontaneity, and his
rejection, connected with this, of the thing in and for itself. Fichte
carries through what Kant intimated in his doctrine of the
``transcendental apperception'' and of the primacy of practical reason,
and, in removing the notion of the I's passivity through an external
object, he develops a \emph{practical} idealism which, through its
fundamental concept taken from Kant's ``transcendental apperception'':
the spontaneous I, acquires a \emph{subjective} character. Philosophy,
as Wissenschaftslehre, as the science of knowledge and of the sciences,
is to deduce the standpoint of life and of the individual sciences, the
standpoint of ``sound common sense,'' on which we relate ourselves to
objects, and to demonstrate the possibility of experience; it must
therefore begin by placing itself outside experience and by giving up
every thought of something existing outside the I, since it is
precisely to show how the I, on \opage{197}the standpoint of ordinary
consciousness, comes to assume such a thing. Its content is the
\emph{necessary} acts of the I, which become the foundation of the free
(arbitrary) ones, and it is by bringing the former to consciousness
that it becomes the foundation of all other sciences. It contains their
principles and grounds their scientific form. It itself acquires the
systematic form of science by proceeding from the absolutely first
principle, unconditioned with respect to form and content, certain
through itself, which conditions the form and content of all knowledge,
and by closing itself, starting from this, in a circle that returns to
the point of departure.

The absolutely first principle, which is to express the
\emph{fundamental act} (Thathandlung, not Thatsache), not occurring
empirically, by which consciousness first becomes possible, and which is
reached when one raises oneself by freedom above ordinary, empirical
consciousness to an \emph{intellectual intuition} of the ground of our
essence in its freedom and self-activity, becomes, since the I, which is
the ground of all phenomena, necessarily \emph{must posit itself} in
every positing, the expression of \emph{the I's, the subject-object's,
positing of itself} as \emph{pure} activity, \emph{pure}
self-consciousness, \emph{pure} knowledge. I-hood is this very
\emph{act} by which the self posits itself, posits its being in the
doubling as subject-object: the I is only in this positing of itself.
But the I that is here reached in intellectual intuition is not the
determinate, individual I, but the ground of individuality, the
\emph{absolute} (pure) I, from which the individual is first to be
derived; it is the pure form of I-hood, and it is absolute knowledge,
pure knowledge, not a subjective knowledge of something. The first
principle runs: \emph{The I originally and unconditionally posits its
own being.} To this first principle there is now joined, without being
deduced (which is first attempted at a later point), a second, which
states: \emph{The I posits a not-I as its opposite.} It forms the
antithesis to the first thesis. The positing of the I is here a
counter-positing, and that which by the I is \opage{198}opposed to the
I, the condition of its representing, no longer receives in this
relation of reflection the name of thing, but, as that which is only in
this relation, that of \emph{not-I}. Since the first two principles
stand as contradicting each other, because the not-I too is in the I
that posits it, a third principle is required, which can become the
synthesis of thesis and antithesis, and it becomes this by determining
the I and the not-I as mutually \emph{limiting} one another. This third
principle states: \emph{I posit in the I a divisible} (i.e.\
quantitatively determinable) \emph{not-I opposed to the divisible I}.
By it the not-I is determined as a \emph{barrier}, by which the I
\emph{limits} itself and thus first becomes a \emph{determinate} I,
while at the same time the not-I becomes a \emph{something}. This
principle, as the \emph{fundamental synthesis}, contains all other
syntheses in itself, and becomes the point of departure both for the
theoretical and for the practical Wissenschaftslehre. By the last two
principles the \emph{method} of the science is determined, which moves
forward through an analysis that constantly seeks out new antitheses,
and through a mediation of these by constantly new syntheses; by the
first is determined its concluding \emph{goal}: the I as \emph{idea},
as no longer individual. From each of them, when one disregards the
determination of content and reflects only on the form of the
proposition, a \emph{logical} principle is derived: from the first the
principle of identity: A = A, from the second the principle of
contradiction: not-A is not = A, from the third the principle of ``the
ground'': A is in part = not-A and conversely (inasmuch as every
opposite meets its opposite in a mark = X, every like is different from
what it is like in such a mark, and this X is in the first case the
\emph{ground} of relation, in the second the \emph{ground} of
distinction). Likewise, when in the principles one reflects only on the
character of the act in general, not on the determinate content, there
is derived from each of them a \emph{category}: from the first that of
reality, from the second that of negation, from the third that of
limitation. And in a similar way Fichte \opage{199}seeks, through a
reflection on the activity of the I, in the further methodical
progression through theses, antitheses and syntheses, to deduce the
remaining categories, taken up by Kant without deduction, as well as
the forms of intuition. In such a deduction he sees the condition for a
fulfillment of the demand that science makes with respect to its form.

In the third principle two propositions lie included: \emph{The I
posits itself as determined} (limited) \emph{by the not-I}, and:
\emph{The I posits itself as determining the not-I}. The first of these
forms the point of departure for the theoretical Wissenschaftslehre,
which is to answer the question: how the I comes to assume something
% printed "bevare" (defect logged in the transcription)
objective; the last for the practical, which is to explain how the I
(reason) comes to ascribe causality to itself. In the
\emph{theoretical} Wissenschaftslehre the I is seen in that
determinateness which seems to come from without, but which --- since
the not-I is recognized to have reality only insofar as the I is
passive, and this passivity of the I in turn, as a negative activity, is
seen as determined by its activity, by the limiting of its infinite
activity, which includes the mutual determining of the I and the not-I,
--- finds its true explanation through the concept of the
\emph{productive imagination}. For this expresses the ``independent
activity'' of the I, which, as setting limits and again cancelling
them, hovers between infinite and finite, and which, as the
presupposition of consciousness, \emph{unconsciously} produces (the
notion of) the objects to which knowledge (the theoretical mind, the
intelligence) relates itself through feeling, intuition and
understanding, and which it, in consciousness, which is conditioned by
the opposition to the object and by the faculty of abstraction that
appears in its highest form in reason, expressly posits as determining
the I. In consciousness as this notion of our representing, the point
of departure of the theoretical Wissenschaftslehre has been reached
again: the synthesis of the subjective and the objective, the I that
posits itself as determined by the not-I. But the theoretical
Wissenschaftslehre explains
\opage{200}\setcounter{footnote}{0}only that and how the I is
determined by the not-I, but does not answer the question: \emph{why}
the I lets itself be determined by it, what impulse drives it to this.
Only the \emph{practical} Wissenschaftslehre does this, which proceeds
from the proposition: The I posits itself as determining the not-I. It
shows that the absolute I, precisely as I, i.e.\ as that from whose
essence \emph{reflection} is inseparable, realizes its \emph{infinity}
only by positing, as the condition of reflection, a limitation which is
found by the I as intelligence, and which as practical it constantly
breaks through in infinite \emph{striving}. The infinite I thus posits
itself as limited, finite, \emph{in order} to be infinite \emph{as I}:
the theoretical relation is the condition, the \emph{means}, for the
practical; reason is theoretical \emph{in order} to be practical. The
opposition in the I that the second principle expressed is hereby
explained. What showed itself to consciousness as a world is thus now
seen only as the condition, posited by the I, of the I's activity, of
its being as I. And as Fichte passes from the concept of the pure,
absolute I to the individual I's, whose \emph{plurality} is first
deduced in the Naturrecht from the fact that an I can think itself as a
free subject only by being solicited to self-determination by
\emph{another} rational being, thus by positing itself as an I among
several rational beings, --- the world becomes the material for these
beings' \emph{fulfillment of duty}\footnote{``The only ``Ansich'' is our
duty, which according to the laws of sensuous representation transforms
itself into a world of sense.'' ``The changes in this are changes
in the I, together with which the intuition of our world changes.
% unbalanced guillemet as printed: the second quotation opens but is never closed
}, a barrier posited by the I, which is to be lifted by ethical acting,
thus in reality an \emph{ethical} postulate. The highest task of
morality becomes this: in infinite approximation to realize the demand,
set by the concept of the absolute I, for a self-activity (a
\emph{free} self-determination) for the sake of self-activity, for
which the power, produced through freedom, over the sensuous
na\opage{201}tural drive aiming at enjoyment (the activity of the
\emph{objective} in us), the barrier of our pure drive, its adversary
in the I, is the condition. By the concept of self-sufficiency the
moral will is to be unconditionally determined; every action is to be
so determined that it ``lies in a series by whose continuation the I
would have to become independent''; thus the pure I is presented in the
individual. The goal represented as \emph{ideal}, by which there is a
pointing back to the point of departure of the Wissenschaftslehre: the
subject-object, the I that was \emph{not yet} an \emph{individual} I,
is complete independence, complete liberation from nature, which would
be one with the return, \emph{abolishing all individuality}, to the
subject-object, the absolute I, which is no \emph{longer} an
individual, an I limited by an object, with the complete realization of
universal reason. But this goal, since the object and the natural drive
constantly persist, is precisely only an ideal, only an object of an
infinite approximation, an infinite striving. The conviction that moral
action, in virtue of the agreement, grounded in the absolute I, between
the laws of nature and those of the moral world, leads toward the goal,
is the faith in \emph{the moral world-order}. According to it each
individual I takes its determinate place and by its ethical acting
contributes its share to the realization of the goal, the universal end
of reason. This moral world-order, which as an active ordo ordinans
embraces the individuals and as it were gathers their working, is the
eternal law, the Absolute itself, the divine: the divine must not be
sought outside and above this law, which expresses the absolute power
of the good, the Absolute in its process of realization. Faith in God
is only the practically self-revealing faith in this order. The notion
of an \emph{existing} God, with the determinations of substantiality
and personality, is self-contradictory, since it brings sensibility,
limitation and finitude into the Absolute. The true religion is not
faith in such a God, but ``die Religion des freudigen Rechtthuns.''

\divrule

\opage{202}The philosophical system that Fichte developed in his
Wissenschaftslehre and in the writings that work out its fundamental
thoughts more fully is his real philosophical achievement; the later
form of his philosophy, which arose partly through a development of the
pantheistic moment in the Wissenschaftslehre, partly through an
influence from Schleiermacher and Schelling, was never developed in
finished shape and with scientific rigor. Its fundamental concept
becomes the \emph{Absolute}, which stands above absolute knowledge,
which Fichte now conceives only as the schema, image, phenomenon of
absolute \emph{being}, God, whose essential determination is eternal,
unbecome, unchangeable being-in-itself. The issuing forth of being into
the image existing outside God is determined as necessary, but as
incomprehensible, and as necessary is likewise posited the phenomenon's,
the image's, intuiting of itself \emph{as} image in the absolute I,
which is the only form in which knowledge \emph{exists}, and, since
existence is only in knowledge, in the notion of being, is thus the
only possible form of existence. In the image the being that in and for
itself is one shows itself as manifold, and as knowledge intuits itself
as I, it divides itself into a world of I's, which all have one and the
same common world of sense as object. From the I there is now derived,
as before, a priori the whole content of knowledge, but since the
domain of the Wissenschaftslehre is the domain of images, it becomes
not a science of being but phenomenology, and the determination of the
Absolute leads to this: that Fichte's new philosophy no longer, like
the older, concludes with morality, with its infinite ought, its
infinite approximating striving, but with religion's being in the
divine, which is seen as conditioned by finite I-hood's turning back
from multiplicity to unity, from the phenomenon to true being, by its
unconditional self-surrender to the divine, ``the eternally One,'' the
only reality, and is conceived mystically as a calm, intuiting
enjoyment, not as ethical struggle and action.

\divrule

\opage{203}Through Fichte's Wissenschaftslehre idealism had been
carried through more consistently, the contradictory doubleness in
Kant's point of departure had been removed, the method and form of
science had been indicated more sharply, and from the fundamental
concept of the I in its spontaneity an ethical view of life and of the
world, as lofty and noble as it was bold, had been unfolded. But since
Fichte, determined by his point of departure in Kant's Kritik der
reinen Vernunft, posited the \emph{I}, the form of self-consciousness,
% printed "Selvbevisthedens" (defect logged in the transcription)
as the Absolute, he did not avoid that this abstract and one-sided
determination of reflection of consciousness constituted an insoluble
opposition between the subjective and the objective reached only by a
postulate, an opposition that, despite all syntheses, constantly comes
back, and leads only to an infinite uniting and separating, because the
I as I is only a member of an opposition, and the not-I must therefore
constantly escape it, constantly crop up anew, uncomprehended, when it
seems absorbed by the I. It looks as if the I, which receives the rank
of the Absolute, gathers the opposition into its unity, contains the
objective that it posits; but the opposition comes back, inasmuch as
the I, \emph{in order to be I}, is to need the incomprehensible impulse
(``Anstoss'') from the objective, which always remains something
unexplained and brings a dualism into the interior of the I. --- By
that fundamental concept, likewise, a real solution of the task of
\emph{deducing} the whole content of the system had been made
impossible. Fichte develops a virtuosity of reflection which in
infinite mirroring lets the oppositions, which cannot be gathered into
unity, be refracted in one another, but which, when the task is set of
deducing the a priori forms as necessary determinations from the I,
produces only an apparent deduction from the standpoint of reflection
of consciousness, not an immanent development from the true unity of
the subjective and the objective. And where the concrete objective is
attempted to be deduced as conditions of the I, of self-consciousness,
where e.g.\ a world of sense, the I's corporeality, light and air and
the like are to be derived, the deduction becomes merely illusory, and
reaches its goal only through a tacit borrowing from
ex\opage{204}perience. --- Finally it becomes impossible for Fichte to
reach a true conception of natural reality. Nature becomes only stuff,
material, barrier (``die nicht zu construirende Schranke des Ich''),
means for the I's reflection, material for ethical acting, a receding
barrier whose determination is only to be broken through; it acquires
no independent significance as an objectivity that has an end in
itself.

Through the one-sidedness in the idealism to which Fichte's abstractly
subjective fundamental concept, which cannot explain the object in its
reality and concreteness, necessarily leads, a demand is made for a
realistic corrective; through the ambiguity in the fundamental concept,
which would express the Absolute by way of the abstract and one-sided
form of the I, the demand is made for another determination of the
fundamental concept, one that expresses the Absolute in its
absoluteness, as the all-embracing. That corrective is applied from a
critical standpoint by \emph{Herbart}; this transformation of the
fundamental concept takes place through \emph{Schelling}, by whom
critical idealism is formed into speculative idealism.

\divrule

\begin{center}
  \textbf{4. The Reaction of ``Realism'' against Subjective Idealism: \emph{Herbart}.}
\end{center}

\emph{Johann Friedrich Herbart} was born in Oldenburg in 1776. Already
at the Gymnasium he became acquainted with the Wolffian philosophy and
with Kant's main writings. When in 1794 he went to the University of
Jena, he became Fichte's auditor, and came into a personal relation with
him. By his teacher's personality he was awakened to independence, and
his distinctive tendency appeared early. It is already intimated in a
critique of Schelling's first writings, composed in 1796. Herbart now
studied Greek philosophy eagerly and subjected the prevailing
fundamental concepts of philosophy to a sharp, critical examination. The
years 1797---1800 he spent in Switzerland as a private tutor, occupied
with the study of pedagogy, mathematics, morals and psychology. The
first traces of his mathematical treatment of psychology go back to
this time. After spending \opage{205}another couple of years as a
private scholar, he became Privatdocent at Göttingen in 1802 and
extraordinary professor there in 1805. In 1809 he was called to
Königsberg as ordinary professor of philosophy and director of the
pedagogical seminar. There he unfolded a considerable activity as a
university teacher and author; his main works were written there. In
1833 he went once more to the University of Göttingen, at which he
taught until his death in 1841.

Herbart's most important \emph{writings} are: Allgemeine Pädagogik, aus
dem Zweck der Erziehung abgeleitet, 1806; Hauptpunkte der Metaphysik,
1808; Allgemeine praktische Philosophie, 1808; Psychologische
Untersuchung über die Stärke einer gegebenen
% printed "einen" (defect logged in the transcription)
Vorstellung als Function ihrer Dauer, 1812; Theoriæ de attractione
elementorum principia metaphysica, 1812; Lehrbuch zur Einleitung in die
Philosophie, 1813; Lehrbuch zur Psychologie, 1816; Gespräche über das
Böse, 1817; Ueber die Möglichkeit und Nothwendigkeit, Mathematik auf
Psychologie anzuwenden, 1822; De attentionis mensura causisque
% printed "auzuwenden" (defect logged in the transcription)
primariis, 1822; Psychologie als Wissenschaft, neu gegründet auf
Erfahrung, Metaphysik und Mathematik, 2 vols., 1824---25; Allgemeine
Metaphysik nebst den Anfängen der philosophischen Naturlehre, 2 vols.,
1828---29; Kurze Encyclopädie der Philosophie, aus praktischen
Gesichtspunkten entworfen, 1831; De principio
% printed "Getichtspunkten" (defect logged in the transcription)
logico exclusi medii inter contradictoria non negligendo, 1833; Umriss
pädagogischer Vorlesungen, 1835; Zur Lehre von der Freiheit des
menschlichen Willens, 1836; Analytische Beleuchtung des Naturrechts und
der Moral, 1836; Psychologische Untersuchungen, Hefte 1---2,
1839---40.

\divrule

Herbart's philosophy, like Fichte's, has issued from criticism and bears
its stamp, but it seizes and develops precisely that in Kant which the
Wissenschaftslehre had pushed back, sets aside everything that had led
Kant's successors to idealism and pantheism, and designates itself,
since it \opage{206}holds fast to the concept of the thing ``an sich,''
the \emph{real}, as \emph{realistic}.

Philosophy is defined by Herbart as the ``working-over of concepts,''
and from the main kinds of such ``working-over'' he derives, somewhat
forcedly, the \emph{division} of philosophy into the 3 independent main
parts: logic, metaphysics and aesthetics. The task of \emph{logic} is
to set up the norms by which concepts acquire clarity and distinctness,
and thereby to give insight into which concepts can be combined and
which not. The distinctness of a concept, that is the clear
distinguishing of its marks, is the condition of the affirmative
judgment; the clarity of a concept, i.e.\ its distinguishability from
other concepts, that of the negative; and both forms of judgment in
turn that of the syllogism, which is either a syllogism of subsumption
(1st and 2nd figure) or a syllogism of substitution (3rd figure). Logic
has to do only with \emph{finished} concepts, not with their
(psychological) genesis; nor does it concern itself with the actual
objects corresponding to the concepts, but only with the form of the
concepts and their mutual relation. Its highest laws are the principle
of contradiction and the principium exclusi medii: critical distinction
is the norm, the synthetic unity of opposites the contradictory; logical
thinking is purely analytic.

Of the two other main parts of philosophy, to which the essential
interest of Herbart's system attaches, \emph{metaphysics} has the task,
through a working-over of the concepts of experience (concepts that
express \emph{the given}), of \emph{correcting} them, i.e.\ of removing
the contradictions that they include in their original form, and which
appear precisely when the concepts are made distinct.
\emph{Aesthetics}, which also comprises ethics, has the task of seeking
the normative concepts to which an immediate determination of value,
independent of the reality of what is judged, attaches itself through
an approving or disapproving judgment, and which thus as it were
receive a \emph{supplement} through this.

\divrule

\opage{207}The point of departure for the investigations of \emph{metaphysics} is
what is found in common consciousness, the representations and concepts
that are originally formed in it of that which calls forth our
sensations. From what is thus found we first detach ourselves through a
\emph{skepsis}, which becomes the preparation for genuine philosophical
thinking, but also only a preparation. ``Every able beginner in
philosophy is a skeptic, but every skeptic as such is also only a
beginner.'' Herbart distinguishes a lower skepsis from a higher. The
lower doubts only that we, on account of our apprehension's being
conditioned by our subjectivity, obtain a faithful and unaltered image of
what things really are. The higher goes further and doubts whether all
the forms in which we apprehend the thing are anything at all other than
subjective forms; indeed, whether there \emph{is} anything at all other
than the subject's representations. Here, however, skepsis has reached a
point where a reversal must take place. If one denies the existence of
everything objectively real, then in any case the semblance, the
sensation, remains; but this semblance is inexplicable unless a being is
posited: ``so much semblance, so much indication of being.'' And as
regards the \emph{forms} of experience, that which is is indeed not known
by us in and for itself, but only in relations; yet those forms must
still be really given through the \emph{complexes} of sensation, and not
be merely subjective forms that we impress upon a formless material,
since it would otherwise be inexplicable why we, when we apprehend a
definite object, cannot arbitrarily apprehend it in any form whatever,
but feel ourselves bound to a definite form.

Thus there is in the apprehension of common consciousness at least a
partial truth; but the concepts it forms, and which stem from something
given, contain, on closer consideration, contradictions which thought
must remove through a working-over. The fundamental concepts to which
the collective concepts of common consciousness for the given lead back,
and to which the most essential problems attach, are the concepts of the
thing, of matter and of the I. Of these, the thing is regarded
\opage{208}partly as a unity of different properties, partly as
changeable. In the concept of the \emph{thing}, in which several
properties are supposed to ``inhere,'' lies the contradiction that what
is one is supposed to be a plurality. For when the thing, according to
its nature, its what, is supposed to \emph{possess} the properties, it
itself becomes something manifold, whereas as a thing it is supposed to
be one. In the concept of the thing's change there lies partly the same
contradictory connection of unity and plurality, partly new
contradictions, as the question is asked what produces the change.
Whether the cause is sought outside or in what is changed, or the change
is posited as causeless, there is a contradiction. The concept of the
external cause leads out into the infinite series, and contains the
contradiction that the cause is supposed to bear in itself a
determination that is foreign to its nature. The concept of the internal
cause, self-determination, is contradictory, because the one being must
be thought to split itself into something active and something passive.
The concept of a causeless change leads to the concept, conflicting with
experience and self-contradictory, of an absolute becoming that cancels
every quality. --- In the concept of \emph{matter} the contradictions
show themselves, as the material is determined as what is extended in
space, but through extension, through the many parts of space lying
outside one another, is split into a plurality which is nevertheless
supposed to be a unity. The division is moreover continuable to
infinity; we never reach from the unity to the last, single, parts; and
if we start from these, determined as something infinitely small, we
never arrive at the whole. The same contradiction as with what is
extended in space comes forward with what is extended in time (i.e.\
that which fills time). Here too the unity contains a manifold of parts,
and of such parts as are never reached in the infinite division. --- In
the concept of the \emph{I}, finally, there recurs partly, as the I in
its unity is seen as the source of the manifold representations, the
contradiction in the concept of the inherence of the manifold in the
unity; partly the determination of the I as that which represents
\emph{itself}, since this \emph{it\opage{209}self} is = that which
\emph{represents} itself, leads out into an infinite repetition of the
same determination, without the concept of the I really being reached;
partly a new contradiction arises, as the I, as self-conscious, is
supposed to be subject-object, unity of what represents and of something
represented, an image. --- These contradictions in the fundamental
concepts metaphysics is to remove through its rectifying, methodical
working-over of them. Of the main parts of general metaphysics,
\emph{ontology} occupies itself with the concepts of that which is and of
the thing; \emph{synechology} with the concepts of matter, of continuous
being in space and time, and of motion; \emph{eidolology} with the
concept of the I, the bearer of all the phenomena. To general metaphysics
is then joined special (applied) metaphysics, which comprises the
philosophy of nature, psychology and the philosophical doctrine of
religion (natural theology).

The conditions for the rectification of the concepts given through
experience are partly the determination of the concept of that which is,
to which the phenomenon points, and to which its contradictions compel us
to go back; partly such methods as can provide those concepts of
experience that cannot be discarded with the supplements by which the
contradiction in them is lifted and the given and its forms are
reconciled with the logical forms of thought.

\emph{Being} designates the absolute position, which expresses a
subsisting for itself, freed from all relativity, independent of our
thought and with absolute reality. Being is thus, strictly speaking,
always absolute being. What is thus posited is \emph{that which is}, the
real, a \emph{being} (ens). Being itself contains no quality; it
expresses only the manner in which a quale is posited. Not being, but
that which is, is the really actual. From the concept of being as the
absolute position, excluding all relativity, Herbart derives: 1) that the
quality of that which is must be absolutely positive or affirmative,
i.e.\ not contain a negation or a limitation, since every negation
includes a relativity by which \opage{210}the absoluteness would be
cancelled; 2) that the quality must be absolutely simple, i.e.\ in no way
contain a plurality, or a becoming (a change); 3) that the quality of that
which is must be absolutely indeterminable by concepts of quantity, so
that that which is, as such, can never be regarded as a quantum, as
divisible and extended in space and time. The concept of that which is as
the absolutely simple, spaceless and timeless, unchangeable real becomes
the fundamental concept of the Herbartian metaphysics, influenced at once
by the Eleatics and by Leibnitz; and the real is conceived, in order to
explain the multiplicity of semblance, as an \emph{indeterminable}, but
not infinite (for the infinite does not tolerate the absolute position)
\emph{plurality} of individual realia. Their peculiar quality remains
constantly unknown to us; we cannot know how that which is is in and for
itself; what we know is only the given, the phenomenon; we know the
former only through this phenomenon, which as phenomenon guarantees
something that is, and as manifold phenomenon the individual
multiplicity of that which is.

The \emph{methods} by which the contradiction in the concepts is removed
while the given is respected are the \emph{method of relations} (die
Methode der Beziehungen) and the \emph{method of contingent views}. The
method of relations is so called because it seeks that to which the
concept \emph{points} in order to be able to be thought without
contradiction, i.e.\ the supplementary concepts that are demanded by the
contradiction. The method of contingent views considers the same thing
thought in different relations to other concepts, demanded by the given
problem, (as e.g.\ the straight line can be tangent, radius, diameter,
etc.), and through these different relations into which the concept is
\emph{brought} by us, it provides the possibility that what is simple and
unchangeable can nevertheless show itself as different. By means of these
methods we go, as it were, behind experience and gain an insight into the
relations of the underlying realia, which explains the empirical
phenomenon with its contradictions. We procure for ourselves the
possibility of going from the given to the \opage{211}supersensible, and
from this, reconstructing by means of concepts free of contradiction,
back to the given.

The contradiction that was demonstrated above in the concept of a thing
with many properties is removed in \emph{ontology} by the method of
relations. It consisted in this, that the same thing was to be thought as
one and as manifold. This can be designated thus: the contradictory given
concept A contains the opposed elements m and n. Its contradiction is
lifted when m is multiplied and n is set = the combination of several m.
This is thus to say: the concept of the thing is supplemented by the
assumption that many single realia, each of which has its single quality,
are ``together,'' and that this their being-together, which must be
posited because the phenomenon of the constant grouping of the properties
points back to such a \emph{relation}, conditions the phenomenon of the
multiplicity of the properties in the unity of the thing. The phenomenon
of unity is produced by one of the many realia that are grouped together
into a thing taking the place of the center, becoming the point in which
the series converge. In this center the manifold semblance is united and
the unity shows itself. It takes the position of \emph{substance}; the
others, which point to it, that of marks, accidents.

The contradiction in the concept of the thing that \emph{changes}, and
that, as changing, is supposed to remain the same while its marks
alternate, leads partly back to the same point of view as the
contradiction in the concept of inherence, namely to seeing the thing as
a complex of realia which become the ground of the phenomenon of change
through an alternation in their being-together; partly it leads to the
concept, decisive for the Herbartian philosophy, of the
\emph{self-preservation} of the reals against disturbances. Our own I
seems to present inner changes, and the fact that a coming-together of
different qualia produces different phenomena, dependent on the
qualities, seems to indicate that each real acts upon the others, as they
mutually become the cause of a different semblance. The question then
becomes: what really \opage{212}takes place hereby in each real? The
answer to this becomes: everything that can take place in what absolutely
is, from whose nature a real change is excluded, is one kind of activity,
or rather an analogue of activity, namely the \emph{self-preservation} of
the real, which coincides with its being. As the real with its simple
quality relates itself differently to the different qualities, (cf.\ as
an analogy: the effect of a juxtaposition of colors), there arise
different \emph{contingent views} of its own quality, and the kind of the
self-preservation, and of the corresponding inhibition conditioned by the
difference in quality of what is together, seems to change. In no case
does a real disturbance and change of the real as such take place, but we
presuppose a possible disturbance, and through this arises the notion of
a peculiarly modified self-preservation, while the real merely continues
to exist neutrally. To this self-preservation, which is thus apprehended
as modifying itself only through the contingent views, and which in
reality becomes only the expression of the real's unchanged being,
Herbart restricts all real occurrence (``das wirkliche Geschehen'') in
the reals; and although the disturbance and the self-preservation become
only imaginations of the observer, although the mutual conditioning takes
place only in the subject's apprehension, this concept serves him not
merely to explain that the real, without its unity being cancelled, can
show itself in manifold changed shapes, but he also uses it, in the
domain of psychology and biology, for the explanation of phenomena that
presuppose its significance for the inner state of the real itself.

The contradictions in the concepts matter, space, time and motion
\emph{synechology} seeks to resolve by means of the same fundamental
concept of the real and the same methods. Herbart here develops his
theory of \emph{intelligible space}, in which the single realia must be
thought in such a way that their objective relations furnish the
explanation for the phenomenal, represented space in which our
sensations are objectified. The intel\opage{213}ligible space, to which an
intelligible time and motion correspond, Herbart wants to construct
without any intuition, and he wants, without mixing in the empirical
representation we already have of space, to let that concept arise in
abstract-mathematical construction, starting from the spaceless, the
mathematical point, and through an atomistic concatenation of such points
letting first the line and, from it, the other dimensions of space come
into being. He does not, however, by the nature of the case, avoid
anticipating the representation of space from the very beginning. When
the single realia are thought as being next to one another, i.e.\
outside one another but without any distance, this gives the straight
(``starre'') line; when the points are thought as passing over into one
another, the line acquires the semblance of continuity (becomes
``stetig''); when two directions are combined, the plane arises; when a
third direction is added, the space of the body. That passing over into
one another of the single points presupposes the divisibility of the
points, a fiction that Herbart seeks to justify by pointing to the fact
of incommensurability. When we think of two ``starre'' lines (two rows of
points) intersecting one another and forming two equal legs, and now draw
a hypotenuse between their end points, then this, being likewise joined
together from points, will, owing to incommensurability, not coincide
perfectly with its end points with the two points between which it is
drawn, but will be an infinitely small part too large or too small. We
are thereby led to let two of the hypotenuse's points partly cover one
another, and these points can now be sought anywhere on the line, which
thus becomes fluid, a \emph{continuum}, since we are compelled to set the
single points neither \emph{next to} one another nor \textbf{in} one
another, but to let them melt together in such a way that they are not
one, not two, but, while not strictly outside one another, form an
infinitely divisible whole. In continuity, then, lies infinite
divisibility, but what is divisible is precisely only continuous space,
which is nothing in and for itself, but only
sig\opage{214}nifies the relation in which the reals are apprehended as
standing to one another.

\emph{Space} thus becomes for Herbart not a merely subjective form given
in the human mind, but a form, necessary for every spectator, for the
apprehension of the mutual \emph{objective} relations of the single
reals; it is a semblance, but an objectively given one, which is not
% printed "objectiv" for "objectivt" (defect logged in the transcription)
carried over by us onto the relations of things, but from them onto us.
The relations and positions of the reals, their union and separation in
the intelligible space in which we must think that which is, correspond
entirely to the coming and going that takes place in sensible space, in
which we see things. And what holds of space holds also of \emph{time}.

\emph{Motion} and \emph{rest} likewise refer to the relation between the
reals. Every real in intelligible space is, when it is considered in and
for itself, or as being in space, at rest; but this rest can, in
relation to other realia, be motion. Since space is contingent for that
which is, it must also turn out that the beings can withdraw from the
relation it expresses. Rest is, as a velocity = 0, only one possibility
among infinitely many; motion will thus be the usual thing in the mutual
relations of the beings. Motion is not something that is a determination
in the being itself, not a real change in it, but it expresses a relation
between the reals which is mirrored in that form in a comprehending
consciousness. Motion is not a ``wirkliches Geschehen,'' but an
\emph{objective semblance}, that is, one that issues from the universal
conditions, not merely those lying in human nature, for comprehension in
space, from the meeting of the images in the spectator as such, in
\emph{every} intelligence.

Through the notion of the reals' self-preservations against opposed
qualia, and through the fiction of the divisibility of the point, which
makes possible an imperfect being-together (a perfect being-together is
``einfache Durchdringung'') of the reals and a partial interpenetration,
an explanation is won of the phenomenon of \emph{mat\opage{215}ter}. When
a single element a is partially penetrated by another b, the \emph{whole}
real a is thereby occasioned to a self-preservation (because the real, as
simple, does not have real parts, which are found only in space); there
thus arises in the real an inner state that does not correspond to its
outer relation, to which only a perfect being-together of the reals a and
b would correspond. This abnormal state is now to be cancelled, and this
happens through the \emph{attraction} of the elements. If the single real
enters at once into a similar relation to \emph{several} realia of
opposite quality, without being able to accomplish the requisite degree
of self-preservation over against them, then it must --- since as a real
it must in any case maintain itself --- itself seem to recede, or, what by
the relativity of motion is the same, to thrust back, \emph{repel}, the
others. When an equilibrium is brought about between attraction and
repulsion, the several realia enter into such a spatial relation that
they form a molecule, and the combination of several of these gives
visible \emph{matter}, whose different modifications are determined by
the strong or weak and by the equal or unequal opposition between the
qualities of the reals, and by the intensity of the qualities. On these
oppositions Herbart bases his \emph{philosophy of nature}. To explain the
physiological (biological) phenomena he must use the fiction that the
inner states brought about by the self-preservation of the reals, which
are supposed to form as it were an analogue to the representations in the
soul, can persist after what called them forth has been removed, that
they can come forward again after having been inhibited by others, and
that the real can call forth similar states in other beings with which it
comes into connection. Matter thus acquires an ``inner formability''
(Bildsamkeit). ---

The contradictions in the concept of the \emph{I} are treated by the last
part of metaphysics: \emph{eidolology}, which becomes the foundation of
\emph{psychology} and can be presented in connection with it. The
fundamental concept here too becomes the simple real, the \emph{soul},
whose \opage{216}self-preservations are the representations, which by
their mutual relations condition all psychological processes, including
the genesis of self-consciousness.

To the determination of the soul as a simple real points the fact that
our representations are gathered in the unity of one consciousness. They
arise as this real is set, through its relation to others, in the state
of self-preservation, and they are supposed to persist after what called
them forth has been removed. What Herbart, with respect to the other
realia, could conceive as a contingent view that did not touch the
essence of the reals but fell only in an observing and comprehending
consciousness, he must here, in the case of the psychical real, the
substrate of consciousness, whose self-preservations are
representations, give a significance for the real itself; ``the
contingent view'' must here become a ``wirkliches Geschehen'' in
consciousness. The connection of unity and multiplicity, which everywhere
else was removed as a contradiction, asserts itself there, just as with
the Eleatics the semblance of multiplicity and becoming asserted for
itself, in sensuous consciousness, the being that was denied it. Change
asserts itself through the change of the representations; an influence
of other realia acquires reality in the representations that they call
forth. In order that the semblance (the phenomenon) may be explained,
everything must be ascribed to the soul that was otherwise denied of the
real. It is only a superficial veiling of the contradiction when it is
said that, without the soul's unity being given up or a plurality being
posited in its essence, a plurality of self-preservations can be posited,
because the being-together with other, disturbing, realia that calls
them forth is external and contingent for the soul; for the substrate of
the consciousness that holds the many representations is precisely the
soul-real itself. Nor is the contradiction removed by the distinction
between the general simple self-preservation, which is constantly the
same, and its accidental modifications, the determinate
self-preservations; for the question is precisely how the general
self-preservation, which coincides with the unchanged being of the
soul-real, can in fact, for consciousness, ex\opage{217}press itself in
the different determinate self-preservations. By letting the
representations in their multiplicity hover above the real untouched by
them, one makes the whole relation a riddle, and one finds nowhere a
subject for whose consciousness this incomprehensible semblance could
show itself.

But once the plurality of representations given in \emph{experience} is
acknowledged, the consequence is that these representations, which are
gathered ``in one sphere, in the one self-preservation of the one
soul-real,'' must modify one another. Simultaneous representations that
are homogeneous form fusions; disparate ones, i.e.\ belonging to
different continua (e.g.\ green and bitter), complications; if the
representations are opposed to one another, partially or totally, they
inhibit one another in proportion to their opposition, since it conflicts
with the principle of contradiction that the opposed should be at the
same time in the same point. What is changed by this inhibition is not
the quality of the representation but its degree of strength, its
intensity; and what is inhibited is not annihilated --- as positive, the
soul's self-preservations cannot be annihilated by one another --- but
is only pushed out of consciousness, so that from an actual
representation it has become a \emph{striving to represent}. Although the
single intensities cannot be measured, the \emph{relation} between the
intensities of the representations can nevertheless be subjected to a
calculation, and through it an \emph{exact} expression can be found for
the laws of the processes of representation. \emph{Mathematics} thus
acquires an essential significance for psychology. The \emph{statics} of
representations treats the equilibrium of the representations, the
\emph{mechanics} their motion and their varying strength during this
motion.

In order to calculate the mutual inhibition of the representations, one
starts from the simplest case: two representations with equal intensity
% printed "Forestillnger" for "Forestillinger" (defect logged in the transcription)
and total opposition. In order that one of them should become
uninhibited, the other would have to be totally displaced. But since both
assert themselves with the same force, the \emph{sum of inhibition} is
distributed equally over both: each loses half its intensity. If the
opposition is set as only partial, the sum of inhibition
\opage{218}would have to be expressed by a fraction. If the
representations are thought as totally opposed but unequal in strength,
the sum of inhibition becomes equal to the intensity of the weaker, and
it is distributed over the two representations in inverse ratio to their
strength. Thus, however great the inequality in intensity may be, when
only two representations meet there always remains a residue of the
weakest uninhibited; it is not pushed entirely down below the threshold
of consciousness. For this to happen, more than two representations must
meet. The sum of inhibition then becomes equal to the total sum of the
intensities after deduction of the strongest, and as it is distributed
in inverse ratio to the intensity, the weaker representation can be
pushed wholly out of consciousness. It then sinks down below the
threshold of consciousness, but since it is transformed into a striving,
it can again, supported by other representations with which it is
connected, be raised up above it. When a representation is connected with
others that were more or less inhibited, it will, when it rises into
consciousness, seek to raise these up in \emph{a definite order}. On this
rests the mechanism of memory and our apprehension of the forms of space
and time, and of the categories. \emph{Abstract concepts} are formed
mechanically as, in the single distinct complications, what is alike
produces a total impression which pushes back what is unlike, without
annihilating it; they are formed, at a more developed stage of
consciousness, artificially, as the point of intersection of the
different series is singled out in such a way that the series of which it
reminds no longer unfold with it by way of memory. The representation of
the I, finally, which is indeed to be distinguished from the concept of
the soul-real, is conditioned by one mass of representations apperceiving
(appropriating) other, alternating masses. The I designates the shifting
point in which all our series of representations converge, and the
representation of the I, or self-consciousness, arises only through this
point's being distinguished from the series that meet in it. By the
alternating objective itself the I is lifted out of its representing of
it. I-hood rests on a manifold objective \opage{219}foundation, of which
every part is contingent to it in the sense that the remaining parts
would still serve the I as support if that part were taken away, but not
in the sense that the I could do without them all. The determinate
content of the I is constantly given by the alternating complexes of
representations; the empirical I has this constantly alternating content;
its relative identity lies in the consolidation of the apperceiving mass
of representations. The \emph{pure I} is a scientific abstraction, which
is prepared by the motion in our representations, by the spontaneous
classification of the inner states as sensations and actions, both of
which concepts point to one end- or starting-point that distinguishes
itself as subject from the objects, and which is completed when,
disregarding the content of our knowledge, we reflect only on knowledge
itself, so that it now becomes \emph{knowledge of knowledge}. It is this
abstraction of the pure I that, by a self-deception, has been taken as
independent and identified with that which is, whereas the soul-real, in
its unknown quality, is something that does not represent, neither
subject nor object.

With these determinations of the soul-real and of I-hood Herbart regards
the contradictions in the concept of the I as removed. The multiplicity
that contradicted the unity is reduced to the plurality of
self-preservations. The relation of the I-subject to the I-object has
found its explanation in the doctrine of the apperceiving and the
apperceived masses of representations, and by the method of relations
the I-object has been made congruent with the I-subject by being thought
as a plurality which, \emph{taken together}, becomes equal to the latter.
--- The notion of a plurality of \emph{faculties of the soul} Herbart
rejects, partly on account of the contradiction between the plurality and
the simplicity of the soul-real, partly on account of the contradiction
in the concept itself: faculty, power, which designates a cause that is
as yet only possible, and thus does not yet act as cause. What has been
called faculties of the soul are only concepts of kinds of psychical
phenomena which have been hypostatized. All psychical processes are
explained by the mecha\opage{220}nism of the representations.
Recollection is a reproduction of representations according to definite
laws. The inner sense is an apperceiving of new representations by older,
homogeneous masses of representations. The understanding is a complete
working of the series of representations that have formed themselves in
the soul in consequence of an influence of external things. Reason is a
concurrent activity of several complete series of representations, which
leads to a weighing of reasons and counter-reasons. Feeling arises when a
representation is in motion by means of relations of fusion and is
affected by different, supporting and inhibiting, factors, between which
it is, as it were, squeezed in. Desire is the rising of a representation
in struggle against obstacles. Will is a striving that is connected with
the representation of the attainability of what is striven for. Free
will, in the psychological sense, is the secured dominion of the
strongest masses of representations over single affections.

\emph{The relation of the soul to the body} Herbart must conceive in such
a way that the soul becomes the real that occupies the place of the
substance-real in the organism's whole system of realia. Its seat becomes
an unextended, but shifting, point in the brain (probably in the pons
Varoli). When an affection of the senses propagates itself through the
nerves to this point, the soul is penetrated by the realia that are
nearest to it; it then exercises a self-preservation against the
disturbance, and thus the determinate representations come about. It
produces movements in the organism by calling forth similar
self-preservations in the other realia.

\divrule

The Herbartian \emph{doctrine of religion} cannot be brought into
harmony with the fundamental determinations of his metaphysics. The
beauty and purposiveness in living beings lead him to assume a divine
intelligence which has brought the independent reals into this
relation to one another. But as soon as the essence of this God is to
be determined, one comes into \opage{221}contradictions with the
metaphysics. God must be determined as a single real, but then both his
influence on the other reals becomes inexplicable, and the purposive
ordering of these becomes the \emph{presupposition} of God's
intelligence, which can first arise only through a being-together with
reals so ordered. Herbart also concedes that no scientific doctrinal
edifice of natural theology can be set up. But the concept of God has
an essentially practical significance, and therefore we fill it out
with ethical determinations: wisdom, holiness, power, love and justice.
Thus it satisfies the religious need.

\divrule

Herbart's \emph{aesthetics} is treated by him as wholly independent of
metaphysics. Its object is the \emph{beautiful}, i.e.\ the
involuntarily pleasing, to which the judgments of taste refer. The
aesthetic judgment concerns not the matter of things, which can be
represented purely theoretically without pleasure or displeasure, but
their form. The task of aesthetics is to seek out the first, most
simple thing that pleases, and to let the more complex proceed from it.
That most simple thing can, since the isolated simple (Einfache) is
indifferent, only be the \emph{relation} of the simple elements (thus
something \emph{formal}), and general aesthetics must therefore state
the simplest relations which, clearly represented, awaken a desireless
satisfaction. Of the various doctrines of art (practical sciences)
into which aesthetics divides as it is applied to the various kinds of
the given, the doctrine of the morally beautiful, the \emph{doctrine of
virtue} (practical philosophy), is the most essential, because its
material is ourselves, our willing and acting. The moral distinguishes
itself from the rest of the beautiful as that which determines the
unconditional worth of the person himself. Ethics does not concern
itself with what the will is (with the nature of the will), but shows
\emph{when} the will is morally beautiful, or when it pleases. The
ethical judgments are involuntary judgments of an aesthetic kind upon
relations of \opage{222}will that please or displease. They are passed
independently of the reality of what is judged, and rest on the
\emph{practical} ideas, moral model concepts, expressions of the normal
relations of will, which as ideals with immediate evidence determine
the judgment and pronounce a ``Shall''. They are 5 in number: the idea
of inner or moral freedom (which expresses the harmony between the
individual's will and its moral judgment); the idea of perfection
(which expresses the right relation of magnitude in the manifestations
of will); the idea of benevolence or love (which expresses that the
will immediately makes the satisfaction of a foreign will its object);
the idea of right (which expresses displeasure at strife), and the idea
of equity or the idea of retribution (which demands the reward or
punishment of the intentional, beneficial or harmful action). On ethics
the \emph{doctrine of the state} is grounded. The state, which has its
origin in the need for social life and is according to its concept ``a
society protected by power'', has the task of bringing \emph{all} the
ethical ideas to harmonious manifestation, so that it becomes an
``animated society'', a society in which the ideas animate all, a
common will strives to give them an expression, and a common good
conscience makes known a condition that corresponds to what in the
individual was inner freedom: action in agreement with insight, the
moral judgment. This freedom, which is explained psychologically by the
fusion of a compact mass of representations with self-consciousness,
and which expresses the completed independence of the determinate
character, is the true freedom, whereas a transcendental freedom is an
absurdity, theoretically contradictory since it conflicts with the law
of causality, and practically harmful because it makes the development
of character, education to virtue, and imputation impossible, and
transforms man's acting into something absolutely contingent. ---

\divrule

The Herbartian philosophy forms a sharp opposition to the idealistic
systems which, together with it, issued from the critical philosophy.
Against the sub\opage{223}jectivism of the Fichtean philosophy it sets
the concept of an independent existent; against the pantheism of the
Schellingian, the concept of the individual reals; against speculative
idealism's assertion of the unity of opposites it critically maintains
the principle of contradiction as the sole ruling one; against all
these systems' proceeding from the concept of the absolute it sets the
given, with its manifold problems, as the basis and point of departure
of philosophical knowledge. Herbart's philosophy is everywhere
polemical; inasmuch as it sets itself the task, not simply of
continuing where the earlier systems have left off, but ``everywhere to
go down to the foundations, and through the sharpest criticism of them
to test whether they are able to bear an edifice of science'', it finds
contradictions in all the fundamental concepts, errors in the methods,
defects in the determination of the relation between the main parts of
philosophy. By his sharp criticism Herbart has in all domains given
incentives and placed decisive problems in a clearer light. In the
domain of psychology he has given an essential impulse toward a genetic
explanation of psychical phenomena and an exact determination of their
laws. But the defectiveness of his fundamental metaphysical concept
makes impossible for him consistency in thinking and the formation of a
system of lasting significance. The concept of the existent, which in
Herbart is determined not by way of experience but by an a priori way,
which gets its determinate form through a fallacy, --- inasmuch as the
concept of absolute position, which in reality expresses only the
independence of the existent in relation to our thought, is without
warrant converted into a determination of the nature and constitution
of the existent in and for itself, --- and which suffers from the
contradiction that the single quality which is to exhaust the concept
of the unconditionally existent, free of all relativity, can be
precisely only something relative and dependent, becomes a dead and
unfruitful concept. Since all development is excluded from the real,
and since only imaginary self-preservations against imaginary
disturbances are admitted as a ``wirkliches Geschehen'', existence
becomes lifeless and dead,
\opage{224}
its life becomes only a \emph{subjective} semblance, a shadow-image
that forms itself in a consciousness whose existence becomes something
incomprehensible, and in which the representations play an unclear
role, as forces in interaction, as themselves representing
(apperceiving), indeed as real beings which meet in a given space and
displace one another, are combined and separated. Only by fictions
which let what was an image in a spectator's eye become actual ``inner
states'' in the real can a semblance of an explanation of reality be
produced. Contradiction, which was everywhere banished, everywhere
crops up, and can be veiled only for a moment, because reality, with
its oppositions, which is to be explained, forces it into the abstract
fundamental concept. --- From the real, determined, through a
reminiscence of idealism, negatively-ideally, i.e.\ with the removal of
the fundamental forms of the material, the forms of actual reality are
sought in vain to be derived. From the spaceless, timeless and
motionless ``real'' space, time and motion can be deduced only by
subreptions; dualism is unovercome in Herbart as it was in Leibnitz.
--- Where the phenomenon of the organic leads Herbart to posit an
ordering God, this concept becomes incompatible with the determinations
of the metaphysics, and comes to presuppose what was to be explained by
it. --- Finally, as regards Herbart's practical philosophy, in its
isolation from the theoretical it cannot find in the immediate
aesthetic judgments the certainty that the \emph{knowledge} of essence
would be able to give it.

\divrule

\begin{center}
  \textbf{5. Speculative Idealism: a) \emph{Schelling}.}
\end{center}

\emph{Friedrich Wilhelm Joseph Schelling} was born in 1775 at Leonberg
in Württemberg, where his father was a pastor. His rich gifts revealed
themselves early, and already in his 16th year he became an alumnus of
the theological seminary in Tübingen, where alongside theology he
studied philology and philosophy. His first works, the master's
dissertation Antiquis\opage{225}simi de prima malorum origine
philosophematis explicandi tentamen criticum (1792), and the treatise
über Mythen, historische Sagen und Philosopheme der alten Welt (1793),
show him as influenced by Herder. Later he was led by the study of
Kant, Schulze, Reinhold, Maimon and Fichte to an adherence to the
Wissenschaftslehre, and from its standpoint he wrote: Ueber die
Möglichkeit einer Philosophie überhaupt (1795), Vom Ich als dem Princip
der Philosophie oder über das Unbedingte im menschlichen Wissen (1795),
Neue Deduction des Naturrechts (1795), Philosophische Briefe über
Dogmatismus und Kriticismus (in Niethammer's Journal, 1796) and
Abhandlungen zur Erläuterung des Idealismus der Wissenschaftslehre
% printed "des Wissenschaftslehre" (defect logged in the transcription)
(1796). The last work was composed in Leipzig, where Schelling stayed
from 1796 on, occupied with philosophy, natural sciences and
mathematics. In 1798 he became a lecturer at Jena, and there came into
close relations with Fichte, Schiller and the representatives of the
Romantic school, especially A. W. Schlegel. His distinctive standpoint
is now developed. Schelling remained in Jena some years after Fichte's
dismissal; from 1803---1806 he was then professor at Würzburg,
afterward a member of the Academy of Sciences in Munich and the
Academy's general secretary; from 1820---26 he lectured at Erlangen,
was, upon the founding of the University of Munich, called there again,
and worked there until 1841, when the Romantic, King Friedrich Wilhelm
IV of Prussia, drew him to Berlin to combat Hegelianism and to ground a
Christian philosophy. In Berlin he lectured for some years on his new
philosophy, which had developed successively out of his philosophy of
identity, under the influence of the mystical and theosophical systems
with which Schelling gradually became acquainted, and from which he
appropriated elements in his own way. His last years he lived in
lofty seclusion, occupied with the working-out of the long-promised
system of revelation. He died in 1854.
\opage{226}

Of his writings we name, besides those cited above, the following:
Ideen zu einer Philosophie der Natur, 1797; Von der Weltseele, 1798;
Erster Entwurf eines Systems der Naturphilosophie, 1799; System des
transscendentalen Idealismus, 1800; Darstellung meines Systems der
Philosophie (in Zeitschrift für speculative Physik), 1800; Fernere
Darstellungen aus dem Systeme der Philosophie (in Neue Zeitschrift für
specul. Physik), 1802; Bruno, oder über das göttliche und menschliche
Princip der Dinge, 1802; Ueber das absolute Identitätssystem und sein
Verhältniss zum neuesten Dualismus (in kritisches Journal der
Philosophie), 1803; Vorlesungen über die Methode des akademischen
Studium, 1803; Philosophie und Religion, 1804; Ueber das Verhältniss
des Realen und Idealen in der Natur, 1806; Darlegung des wahren
Verhältnisses der Naturphilosophie zur verbesserten Fichteschen Lehre,
1806; Aphorismen zur Einleitung in die Naturphilosophie, and
Aphorismen über Naturphilosophie (in Jahrbücher der Medicin als
Wissenschaft, 1806---1808); Ueber das Verhältniss der bildenden Kunst
zur Natur, 1807; Ueber das Wesen der menschlichen Freiheit (in philos.
Schriften, vol. 1), 1809; Denkmal der Schrift von den göttlichen Dingen
des Herrn Jacobi, 1812; Ueber die Gottheiten zu Samothrake, 1815;
Preface to Hubert Becker's translation of a portion of Cousin's work on
French and German philosophy, 1834, and the introductory lecture in
Berlin, 1841. After Schelling's death a collected edition of his
writings has appeared, of which the section Posthumous Writings
contains his new ``positive'' philosophy, whose most essential part is
the philosophy of mythology and of revelation. ---

\divrule

Given the defects that were demonstrated above in the standpoint of the
Wissenschaftslehre, a going-beyond it was made necessary. Fichte himself
made an attempt at this in the later form of his philosophy, but without
being able to give the attempt a strictly scientific execution. It is in
Schelling's \opage{227}system of identity that the development of
critical idealism beyond the Wissenschaftslehre is actually realized. The
infinite principle that in Fichte was posited as the common ground of the
finite subjects (I's) and of the objects, and that was nevertheless,
although standing above the opposites, designated as I, thus by a name
that expresses the one opposite, is freed in Schelling from this
ambiguous designation of the absolute I, and is now determined as the
\emph{absolute}, reason, which is the unity of the opposites. And at the
same time the designations of the opposites are changed. For the abstract
form of the I is substituted the concept of \emph{spirit}; for the
determination of the objective as material void of reason, the concept of
\emph{nature} working according to purposes; and nature and spirit are
seen as \emph{coordinate} forms of appearance of the absolute, of
\emph{reason}. They become parallel forms of being of this absolute, the
subject-object, constituted by the quantitative preponderance of
objectivity or of subjectivity, while the absolute itself is the
\emph{total indifference} of the subjective and the objective. The
philosophy that knows the sole being of the absolute, the identity of
the opposites, is the \emph{philosophy of identity}.

Schelling arrives at this new standpoint as he, who originally himself
stands on the standpoint of the Wissenschaftslehre, seeks to bring what
Kant had contributed to the conception of nature in his Metaph.
Anfangsgründe der Naturwissenschaft and in his Kritik der Urtheilskraft
into connection with that standpoint. He wishes originally to construct a
science of nature as applied theoretical philosophy, to deduce the
content of nature from the forms of development of self-consciousness,
and to bring the Kantian determinations of the opposed forces, whose
product matter is, and Kant's concept of the organic together with
Fichte's conception of the opposed directions in the activity of the I,
and of its intuiting of itself in its successive realization of its
infinity, in which it is its own cause and effect. But as Schelling makes
this attempt, nature soon comes, as if involuntarily, to be conceived as
an independent parallel to the I, --- even if it is \opage{228}only by a
postulate that this independence is reached ---, and the philosophy of
nature is then \emph{coordinated} with transcendental philosophy. And at
the same time as nature and I thus come to form an opposition of equal
right, expressed in parallel series, the concept of nature is
transformed, the concept of spirit replaces the abstract concept of the
I, and the unity of nature and spirit is sought in the absolute,
determined as reason.

Schelling develops the relation between the philosophy of nature and
transcendental philosophy by starting from what the task of philosophy
is. This task is, in a word: to explain knowledge. But all knowledge
rests on the agreement between something objective and something
subjective. One can then either make the objective the first, and ask
how a subjective comes to be added to it that agrees with it; or one can
make the subjective the first, and the question then becomes: how an
objective comes to be added that agrees with it. The first question is
answered by the philosophy of nature (speculative physics), the second by
transcendental philosophy. The philosophy of nature shows how the ideal
(subjective) proceeds from the real; transcendental philosophy, which
leads the real or unconscious activity of reason back to the ideal or
conscious, shows how nature is objectified as the visible image of our
intelligence.

In the domain of the \emph{philosophy of nature} Schelling had first (in
``Ideen zu einer Philosophie der Natur''), before he yet coordinated the
two philosophical sciences, given the deduction of matter, whose
fundamental forces, repulsion and attraction, were conceived as
objectifications of the opposed activities of the intuiting I: the
original, infinite activity, and the limiting one. And from the differing
relation of intensity between the fundamental forces he had then derived
a difference of quality between the materials in its generality, while
it was conceded that the determinate qualities by which objects are
distinguished from the general concept of the material bear a stamp of
contingency that prevents an a priori construction. In his
\opage{229}work ``von der Weltseele'' Schelling then sought the principle
common to inorganic and organic nature. The activities of inorganic
nature repeat themselves in a higher form in those of organic nature; for
both the same law of opposition (polarity) holds; the mechanical and the
organic form no insuperable opposition; all natural causes make up a
continuous series; the positive in the world is absolutely one. The
principle of life in nature, which through the continuity in all natural
forms knits the inorganic and the organic beings together into a total
organism, Schelling designates as the world-soul, without, however,
determining the concept of this principle of life more closely. In the
conception of the organism Schelling attaches himself to Kant, but seeks
to ground Kant's determinations. The inorganic he determines as merely a
negative condition of the organic; the positive cause of the organic, on
the other hand, he seeks in a higher principle, which governs the
processes of the inorganic, and whose individual phenomena are the
functions of life: reproduction, irritability and sensibility. This
principle becomes identical with the principle of life of the whole of
existence.

What is intimated in the work on the world-soul receives its more
definite form in Schelling's ``Erster Entwurf eines Systems der
Naturphilosophie''. The independence of nature now stands out more
clearly; nature becomes, as it were, its own subject, and the intuition
of ``the infinite productivity of nature'' takes the place of the
idealistic construction of nature from the I. The task of the science of
nature, now coordinated with transcendental philosophy, then becomes
this: to see in the various products of nature the monuments of a true
history of nature, and to see all natural productivity aim toward
intelligence. The individual natural things designate, as it were,
points of inhibition of the infinite productivity, determined by that
productivity itself. The dead and unconscious products in nature are only
attempts of nature to reflect itself; the so-called dead, inorganic
nature is in general an immature intelligence, and therefore the
intelligent character already shines, in an unconscious manner,
\opage{230}through in its phenomena. The highest goal, to become wholly
object to itself, nature reaches only through the highest and last
reflection, which is nothing other than man, or, more generally, what we
call reason, through which nature first returns completely into itself,
and whereby it becomes manifest that nature is originally identical with
what is known in us as intelligence and consciousness.

\emph{Transcendental philosophy} shows in its \emph{theoretical} part, in
which Fichte's theoretical Wissenschaftslehre is used throughout, how, in
the continued unconscious self-objectification of intelligence,
everything objective in nature comes to be for it, so that the stages of
nature correspond to the stages in the history of self-consciousness, as
projections of the developing knowledge. All the forces and forms of the
natural become reflections of the activities of the I. The highest form
of nature, the organic with its circulation that closes upon itself, is
the objectification of intelligence, which intuits itself in its
production, in the succession of its notions, and in this intuiting makes
the succession permanent, returning into itself. The series of organisms
is concluded with the one that intelligence must intuit as identical with
itself, in which it sees itself as really existing. Above this phenomenon
intelligence raises itself by abstraction, and with the demand for
absolute abstraction theoretical philosophy passes over into practical
philosophy. In the \emph{practical} part of transcendental philosophy,
which treats the \emph{conscious}, \emph{free} actions of the I that
bring forth the world of ethical life, Schelling again attaches himself
closely to the practical Wissenschaftslehre, in that he sees in the
interaction of the I with other I's the condition for its free
self-determination, and for its fully developed consciousness of
something objective \emph{outside itself}; in that he makes the demand on
the I that it determine itself according to the laws of freedom; and in
that he sees in the state the condition under which the demand of
reason, made with regard to securing the freedom of individuals in their
relations of interaction, can be fulfilled. The legal order in the state
\opage{231}is, as it were, a second and higher nature, issuing from free
and \emph{conscious} self-determination, which is to rule with the
inviolability of natural law for the sake of freedom. The state is an
objective organism of freedom. The successive development of the legal
order is the object of history. \emph{History} as a whole is a revelation
of the absolute that successively clarifies itself. Every individual
determined by reason is an integral part of God, and contributes to
giving the moral world-order existence. The condition for history's
approaching its goal is the harmony between the objective and the
subjective, the law-determined and the free, the unconscious and the
conscious, which is grounded in the higher, which is the \emph{eternal
absolute identity} of the opposites. The history of the revelation of the
absolute has 3 main stages. In the first, the period of \emph{fate}, what
rules shows itself as completely blind power, coldly and unconsciously
destroying even the greatest and most glorious, bringing about the
downfall of the noblest humanity that has ever flourished. In the second
stage, the period of \emph{nature}, it reveals itself as nature and
brings the unloosed forces at least under the mechanical lawfulness of a
law of nature. This period begins with the expansion of the Roman
Republic, which binds the nations to one another and brings the spiritual
possessions of each into contact with those of the others. In the third
stage, the period of \emph{providence}, what earlier showed itself as
fate and as nature will reveal itself as providence, and what earlier
seemed to be merely the work of those powers will be known as a
beginning, still imperfect revelation of providence. --- In the last part
of transcendental philosophy the union of the theoretical and the
practical, of unconscious and conscious activity, is sought. It is found
in the domain of the \emph{artistic}. In artistic production the I
intuits \emph{itself} in the unity of its unconscious and its conscious
activity; in the work of art the character of the product of nature and
that of the product of freedom are united, and the revelation of the
absolute, which in history was one perpetually becoming, is there
realized. In the \emph{beautiful} \opage{232}the infinite is presented in
the finite; the infinite opposition is resolved in the object. Artistic
production is possible only through \emph{genius}, which unites in itself
the infinite opposites and unconsciously puts the infinite into its work
and lets it shine back from it. Science in its highest form has its task
in common with art; but the manner in which the task is solved by them is
different. The task of art can be solved exclusively through the natural
gift of genius: art is therefore the highest union, in the form of
reality, of freedom and necessity. The philosophy of art is the necessary
goal of the philosopher; he beholds in the work of art the essence of his
science as in a magic and symbolic mirror.

Both of the main philosophical sciences thus lead to the point where the
unity of the opposites reveals itself. The two forms of existence form a
parallelism that points to one and the same grounding principle. This is
the absolute, reason, that which alone is, in which all is, \emph{the
absolute identity}. This \emph{absolute reason}, the total indifference
of the subjective and the objective, is the truth in and for itself; to
know things in and for themselves is to know them as they are in reason,
and this is the task of philosophy. The organ of philosophy is
\emph{intellectual intuition}; through it, in which we are in identity
with the absolute itself, the object of this intuition, we see things in
reason. What \emph{is}, is, insofar as it \emph{is}, the absolute;
insofar as it is finite, it is \emph{not}. All plurality and finitude is
only in the phenomenon; according to their essence things express only
the absolute; through intellectual intuition they are, as it were, seen
into the identity that stamps its wholeness in each form of difference.
Since the absolute identity is only \emph{as self-knowledge} of the
identity, because the form of self-knowledge is one with its essence, it
must, in order to know itself infinitely, infinitely posit itself as
subject and object. But since both forms express the identity, there can
be only a \emph{quantitative} difference between them; \emph{the
absolute} is posited \opage{233}in \emph{both}, only with a preponderance
of subjectivity (knowledge) or of objectivity (being). The whole of
existence can be represented under the schema of a magnet, which in its
two halves shows a preponderance of one of the opposed forces, but in
each smallest separated part unites them both, and as a whole is the
indifference of the forces. The quantitative difference comes to the fore
only where something is seen outside the absolute identity, according to
the phenomenon, as an individual thing; if things are seen according to
their ``Ansich'' and as totality, only the absolute identity is seen in
them. Each individual being is only a determinate form of \emph{its}
being; each presents \emph{it} at a determinate stage, in a determinate
\emph{potency}; each is therefore infinite in its kind, each is a
relative totality. The potencies of nature, gravity (matter), light and
the organism, each express the identity in its own way; likewise the
potencies in the ideal series, knowledge, action, art, which the science
of spirit presents. By the absolute method, philosophical construction,
things are thought in their ideas, in everything the unity of the ideal
and the real is intuited, everything is regarded as an expression of the
undivided perfection of the absolute and therefore as itself absolute and
eternal; it is known that all that is particular is in the absolute and
conversely, and that neither can the former be comprehended without the
latter nor the latter without the former, so that philosophy thus does
not in general go outside the absolute.

These are the fundamental features of Schelling's philosophy of identity,
in which he, just as in general at each individual stage of his
development he leans on other thinkers and takes up their thoughts,
renews Spinozism in content and form, albeit in such a way that the
influence of the Wissenschaftslehre is not entirely forgotten. As in
Spinoza, true knowledge in the philosophy of identity has only the
absolute for its object: it grasps it through an immediate intuition, in
which it becomes one with it, and comprehends all forms of existence from
it and in it; it constructs them as expressions of the \emph{undivided}
absolute. It is only reflection --- Spinoza's \opage{234}imagination, ---
that gives the finite and the manifold an untrue existence as separated
from the absolute.

But precisely the possibility of this finitizing reflection becomes for
the philosophy of identity, as once for Spinoza, an insoluble problem,
and it is by this problem that Schelling is forced to transform this
standpoint and to develop what he later called his positive philosophy.

After Schelling (in his Fernere Darstellungen aus dem Systeme der
Philosophie and in Bruno) had made a fruitless attempt to explain the
separation of reflection by the very concept of absolute knowledge as the
``real unity and ideal opposition of the real and the ideal'', in which
there was supposed to lie a \emph{possibility} of the finite's separating
itself from the absolute into a particular existence, he posits (in
Philosophie und Religion) the finite and the corporeal as having arisen
through a \emph{falling-away} (of the ideas or souls) from the absolute,
and in this falling-away, which is to be reconciled through history, he
sees the condition for God's complete revelation. The ``falling-away''
nevertheless becomes, from the standpoint of the philosophy of identity,
the incomprehensible, by which nothing is explained, since it already
presupposes finitude and individual existence. The finite must either
continue to stand as a semblance, whose origin and whose apprehending
subject become a riddle, or its principle must be sought in the absolute
itself, which then, as bearing this principle within itself, cannot in
the same way as before be conceived as the absolute identity, or be
placed in the earlier relation to the forms of existence. To an altered
conception of this relation the determination of the origin of the
sensible world by the falling-away, and of the teleological significance
of the latter, already leads unavoidably. This determination must
transform the philosophy of identity's coordination of nature and spirit
into a subordination of nature to spirit. The ideal, which originally had
the preponderance in Schelling, then gets it again, and nature becomes a
condition and a point of transition for the realization of spirit. The
absolute is then no longer exhausted in the universe, but acquires a
transcendence.

\opage{235}Schelling himself now states his task thus: to gain the
possibility of distinguishing God and nature without dualism, and of
finding a ground of explanation for evil without pantheism. The new
standpoint, which is just as strongly influenced by Jacob Böhm's
theosophy as the philosophy of identity was by Spinoza, first comes to
development in the treatise on human freedom (of 1809). In it the
explanation of cosmogony and of the historical process, with its discord
and reconciliation, is sought in a theogony. The eternal, impersonal
primal ground or unground (Jacob Böhm's ``Ungrund''), the indifference,
splits with necessity into two equally eternal principles: the ground and
the understanding. Just as everywhere the perfect proceeds from the
imperfect, so too in God. God must have in himself a ground of his
existence that is not merely logically but really distinct from his
existence: something dark and unconscious, \emph{nature in God}; that in
God which is not God. From the process of the two eternal principles is
to be explained not only God's becoming toward personality and toward
all-fulfilling love that unites the opposites, but also the possibility
of a world distinct from God and the origin and overcoming of evil. The
ground is the dark basis of reality in God, which conditions his
existence as personality. It is the dark craving of the eternal One, its
unconscious willing, its longing to give birth to itself; in its
darkness it is lawless and formless, a will that only divines the
understanding. But from its dark stirrings there proceeds in God an inner
reflexive representation, by which he sees himself in an image: the God
begotten in God, and by which God first, absolutely regarded, attains
actuality. This representation is the understanding, the \emph{word} of
the longing, and the eternal spirit, moved by the love that it itself is,
utters the word, so that now the understanding, together with the
longing, becomes freely creating will. The word is the light in the
darkness of the ground, the reflection that brings understanding into its
dark craving. Through the understanding, the ideal, the real in the
ground is formed, the forces are separated, a world comes to be, in whose
graded series of forms the successive \opage{236}overcoming of the ground
by the understanding reveals itself. Only at the highest stage, in man,
is the innermost hidden unity of the forces raised up to the light. In
natural beings the two principles, the ground and the understanding, the
principles of self-will and of universal will, of imperfection and of
perfection, do not yet attain perfect unity: this unity is first realized
in man, in whom therefore a higher principle, spirit, comes forth, and
who, while all natural beings are peripheral beings, is in God and
through this being in God has the possibility of freedom. But in man
self-will and universal will are not, as in God, inseparable; there is a
possibility of a separation: the possibility of evil stands beside that
of good. Good is the subordination of self-will to universal will; evil
is the former's rebellion against the latter, the striving of self-will
to be, as particular will, what it is only in identity with universal
will. Evil comes to be as self-will is solicited by that in God which is
not God, and man, by an intelligible act that lies outside time, freely
makes the choice that with necessity determines his essence and through
it the whole of his action in time. The ground solicits self-will in
order that there may be an independent foundation for the good, which can
be brought under the dominion of the good. In the history of mankind the
struggle of the particular will against the universal will, and through
it the struggle of the principles, presents itself; its conclusion is
the reconciliation, which is founded by Christ, in whom the principle of
light and of love appears personally, and is consummated by the complete
raising of the ground up to the light. As man is united with God, he
becomes the redeemer of nature, the mediator between it and God. Through
him God unites nature with himself, and is now, as spirit, all in all. In
spirit the existent is one with the ground of existence; spirit is the
absolute identity of both. God has developed from the impersonal unground
into existing personality, which can only truly be comprehended when
there is posited in God an inner opposition, a negative principle, a dark
ground of nature, which is to be transfigured by \opage{237}spirit,
something lower than the personal and destined for subordination, but
without which personality would be an empty and contentless word.

Schelling determines the relation between the earlier and the later form
of his philosophy thus: the former, whose pantheistic character gave the
development of thought the form of necessity, is designated as the
\emph{negative} philosophy, the pure science of reason, in which merely
logical thought has for its object the a priori, the whole world of
\emph{possibilities}, the prius of all that is, even of the Godhead, the
negative in all knowledge (``das nicht-nicht-zu-Denkende'', i.e.\ that
without which no knowledge is possible, not that which makes knowledge
actual). But this negative philosophy reaches only to the summum
cogitabile, to the \emph{concept} of God (God's ``what''), to which the
potencies of what is, das Seinkönnende (the potentia of being, its
presupposition, its subject), das nothwendig Seiende (being as actus, the
purely existent, the object) and their unity (the subject-object), lead
as to that which comprehends them; but it does not reach God's
\emph{actuality} (God's ``that''), the \emph{existing} God, who is the
idea, and who is first intuited as the necessarily \emph{existent}, the
potency-less being (actus purus, das unvordenkliche Sein). With the
existing God begins the \emph{positive} philosophy, which is given in the
transformed system. This positive philosophy, to which no immanent
transition leads from the negative philosophy that it supplements, rests
on an ``a priori \emph{empiricism}'', and as it follows, through a free
thinking determined by \emph{will}, the life-process of the absolute
\emph{existent}, conceived as personal and particular, it is to give a
scientific grounding of \emph{theism}. Within it the content of the
negative philosophy returns, but as a reality that is seen in the light
of freedom as the product of free creative activity. Against Hegel
Schelling now objects that he, in bringing the philosophy of identity
into logical form, forgot that the science of reason makes up only the
one (negative) part of philosophy, and that, by putting the logical
\emph{concept} of the living divine \emph{subject} in its place, he had,
by a strange \opage{238}fiction, to ascribe to the concept the
self-movement that belongs only to the subject, and to make the
impossible attempt to reach actuality by the logical road.

In Schelling's Philosophy of Mythology and of Revelation, published after
his death, the theogonic and cosmogonic process is developed essentially
in agreement with the treatise on freedom, only more fully worked out,
and with an altered and more definitely formed terminology, and with an
endeavor to let God, as the All-One, the absolutely free, bound to no
kind of being (the ``Ueberschwengliche''), the absolutely personal, stand
eternally above the process of the potencies. The 3 divine potencies,
which come forth on the ground of the ``unvordenkliche'' being, and which
in the absolute spirit are in undivided actuality, not as potencies but
as spirit itself, become, through a process whose beginning, by a
postulate of ``free thinking'', is posited as grounded in a divine
arbitrariness, as ``das Seinkönnende, das Seinmüssende and das
Seinsollende'', factors in the process of creation in \emph{God}, whose
conclusion is the original ideal man, in whose consciousness the unity of
the potencies is attained. By a new postulate of free thinking this
original man is posited as, through his freedom, disturbing the harmony
of the potencies, which in the eternal process of creation, through their
mutual separation and ``tension'' and its overcoming, have been raised to
personalities; and by this act the eternal world is then supposed to be
posited as finite, as determined in space and time \emph{outside} God,
for whose knowledge it becomes the means, and the second divine
personality, ``the Son'', who shared dominion over the world with the
Father, is supposed to be separated from the Father and again reduced to
the condition of potency, deprived of dominion over the blindly working
principle (the first potency), while this potency is actualized in
consciousness. World history now becomes the history of the second
divine potency. Having been reduced by the Fall to the condition of
potency, it shows itself in the period of paganism, of mythology, only as
a naturally working potency, which struggles with \opage{239}the dark
principle ruling human consciousness and successively gains power over
it, until at the close of this period it reveals itself personally as
independent lord over the being it has subjected to itself. But as it
subjects this being to the Father, and as through the incarnation in the
person of Christ and through death it gives up its independent being
outside God, it atones for the guilt, and together with the Spirit, whose
appearance it makes possible, it returns into God in self-acquired
personal divinity and restores the unity of the divine. The goal of the
process is then attained, as all being in the Son is reconciled with the
Father, and God, as lord over a being distinct from him, is revealed as
the absolutely free and blessed. --- The period of Christianity is the
period of revelation. Through Catholicism (the Church of Peter) and
Protestantism (the Church of Paul) it strives toward ``the Johannine
Church of the future'', the Church of free unity. As a John the Baptist
of Christianity, John, the apostle of the future, points toward the free,
\emph{philosophical} religion, which is prepared but not actualized by
Christianity, and which is reached through philosophy's deeper knowledge
of God, of the development of religion, and of the historical connection
of all things. The philosophical religion (the religion of the spirit)
becomes the consummation of religion, and it becomes so by uttering the
thought of esoteric Christianity, of the eternal gospel.

\divrule

In the philosophy of identity, compared with the Wissenschaftslehre,
the advance had been made that attention was again directed to the
real \emph{content} of thought, that nature was recognized as the
objectively existing rational, as an organic whole, and that the
absolutely true was determined as the unity of nature and spirit, of
the objective and the subjective. But since this point of unity, the
concept of the absolute, was not reached through a logical movement of
thought, but was postulated as the object of an immediate intuition,
the consequence was that the \opage{240}absolute, despite every
striving to give it fullness, became the determinationless, from which
the concrete forms of existence could not really be deduced, and in
place of an unfolding of thought there had to step a so-called
construction, a formal schematizing. The defect that was pointed out in
Spinoza's system is still present in the philosophy of identity: the
opposites are not united in a full unity, but vanish in an empty
indifferent; and although the relation between the forms of existence
points toward unity more definitely in Schelling than in Spinoza, in
that nature is seen as teleologically disposed toward spirit, spirit
as genetically conditioned by nature and, in its highest activity, as
it were returning to nature, and in that a parallelism between the
\emph{stages of development} of the two is posited, still no real
\emph{knowledge} of the \emph{unity of essence} is reached. As regards
the parallelism, the agreement becomes only the likeness of the
\emph{analogous}, and the construction of the potencies, which takes the
place of a comprehension of a real development, leads at most to
% printed "idethøiste" (defect logged in the transcription)
ingenious intimations, often only to an empty play. And the emergence
of the spiritual from nature becomes as little conceptually determined
as that unity of spirit and nature, the ``incomprehensible miracle''
of genius, with which the science of spirit concludes.

The knowledge of nature in the philosophy of identity, although the
system was essentially carried through on the side of the domain of
\emph{nature}, has in the main got no further than the general view
that Kant had indicated in great and brilliant strokes; a richer fusion
of empirical inquiry and speculation could not be expected from a
philosophy that possessed in intellectual intuition the magic means of
beholding all things in God according to their truth, and that saw the
highest expression of spirit in the work of art. The speculation on
nature that the philosophy of identity called forth soon tipped over
into pseudo-inspired fantasizing, to which Schelling himself already
gradually gives himself up. And the nature whose independent being, as
coordinate with spirit, is the fundamental thought of the philosophy of
identity becomes in reality, in \opage{241}Schelling, only a hypothesis
whose justification Fichte could deny. The independence of nature is
reached only by a postulate, and, despite the demand for this
independence, the peculiar forms of nature are for Schelling always
involuntarily converted into reflexes of the activities of the
\emph{self-conscious} spirit, of the I, and in this conversion their
content of reality vanishes.

In Schelling's later philosophy, which, though influenced by
theological motives and in its form strongly determined by its relation
to mysticism, has shown itself to have grown naturally out of the
earlier philosophy in the endeavor to solve the problem that the latter
did not master, namely the being of the finite as finite, a defect is
pointed out that lay in the one-sided sublimating of being and reality
by the preceding thought. But instead of seeking the corrective for the
one-sided idealism in real being and in the experience that relates
itself to the reality of that being, Schelling seeks the corrective in a
being of the imagination and in an ``a priori'' experience corresponding
to it, and the ``free thinking'' that is to become the organ of the
positive philosophy becomes, in its arbitrariness and lawlessness, only
a name for thought in the service of the imagination, and its movement
proceeds only through postulates. This positive philosophy, which was
proclaimed as transforming the whole of thinking, and which, despite
its gnostic poetic recasting of Christianity, was received with naive
enthusiasm by speculative theologians, has had only a passing
significance. Its construction of history, reduced to the history of
religion, is just as fantastic as the conception of the reality of
nature held by the pseudo-inspired philosophers of nature.

\divrule

\begin{center}
  \textbf{6. Speculative Idealism: b) \emph{Hegel}.}
\end{center}

\emph{Georg Wilhelm Friedrich Hegel} was born in Stuttgart in 1770.
At eighteen years of age he came to the University of Tübingen, where
for 5 years he studied theology and philosophy. \opage{242}Theology in
its official form satisfied him but little; but he followed with lively
interest the movement in philosophy that had proceeded from Kant, and
the world-transforming political events in France. He stood in close
connection with Schelling and Hölderlin, and the latter's enthusiasm for
Hellenism influenced him strongly and lastingly. After completing his
university studies, Hegel spent some years as a private tutor in Bern
and Frankfurt am Main, where he eagerly continued his theological and
philosophical studies and kept up a correspondence with Schelling. A
draft of a ``Life of Jesus'' from this time distinguishes from the
dogma of the Church, which is made the object of a sharp critique,
Jesus' own religion, which Hegel wants to comprehend from the preceding
religious development and the whole state of the world at that time,
and which he determines purely ethically and from a pantheistic
standpoint. At the beginning of 1801 Hegel came to Jena, where he
qualified as a university lecturer. He had by then made the draft of a
philosophical system in which the absolute was regarded as pure idea
(in logic and metaphysics), as nature (in the philosophy of nature), and
as spirit (in the philosophy of spirit, which in this draft comprised
only the ``System of Ethical Life'', ethical life being conceived in
unity with religion and speculation). In Jena Hegel published in 1801,
first, a work: Differenz des Fichte'schen und Schellingschen Systems
der Philosophie, in which, as distinct from the subjective idealism of
the Wissenschaftslehre, he designates Schelling's philosophy as a
subjective-objective idealism, which in the philosophy of nature and in
transcendental philosophy constructs the absolute in its two necessary
forms of existence. Schelling's standpoint Hegel acknowledges as his
own. Together with Schelling, Hegel published in 1802--3 the
``Kritisches Journal der Philosophie'', in which the introductory
essay, Ueber das Wesen der philosophischen Kritik, was written by
Schelling and Hegel jointly, and the essays ``Wie der gemeine
Menschenverstand die Philosophie nehme'', ``Verhältniss des
Skepticismus zur Philosophie'', ``Rückert und Weiss oder die
Philosophie, \opage{243}zu der es keines Denkens und Wissens bedarf'',
``Glauben und Wissen'', ``Ueber die wissenschaftlichen
Behandlungsarten des Naturrechts'', and perhaps ``Ueber Construction in
der Philosophie'', by Hegel. Gradually Hegel came to clarity about the
difference of his philosophy from Schelling's. In his lectures on the
history of philosophy he pronounces the judgment on Schelling's
philosophy that it is the completion of the development up to that
time, but that it lacks Platonic dialectic and logical grounding.
Against the philosophy of identity's coordination of spirit and nature,
Hegel already at this time maintains spirit as the higher, the
overarching, which as world-spirit gathers together the infinitely
differentiated unfolding of nature, its totality existing in the form
of externality, in the ideal point of unity of absolute knowledge. When
Schelling left Jena in 1803 and the personal contact between him and
Hegel ceased, the difference emerged ever more clearly, and in the
preface to Hegel's ``Phänomenologie des Geistes'', published in 1807,
it is expressed with a certain sharpness. Hegel's activity as a lecturer
in Jena, which at first won only slight recognition, and which comprised
the philosophical encyclopaedia, logic and metaphysics, natural law, the
history of philosophy, the phenomenology, and pure mathematics, was
interrupted in 1806 by the events of the war, and Hegel was compelled
for a couple of years to seek a livelihood as a newspaper editor in
Bamberg, until in 1808 he became director of a gymnasium in Nuremberg,
a position he held until 1816. In Nuremberg he wrote his little
philosophical propaedeutic, for use in the teaching of philosophy in the
gymnasium, and his epoch-making main work: Wissenschaft der Logik
(1812--16), in which what was formerly separated as logic and
metaphysics is united into one whole. In 1816 Hegel was called to
Heidelberg as professor of philosophy. There he published in the
following year his ``Encyclopädie der philosophischen Wissenschaften'',
an outline of the whole system. The range of his lectures was there
extended with psychology and aesthetics. In 1818 Hegel was called to the
University of Berlin, where he developed a great \opage{244}activity as
a lecturer, and soon gathered a numerous school about him, so that his
philosophy was at that time unconditionally the dominant one in
Germany. Here he also drew the philosophy of history and the philosophy
of religion into the range of his lectures. In Berlin he published his
``Grundlinien der Philosophie des Rechts, oder Naturrecht und
Staatswissenschaft im Grundrisse''
% printed "Staatswissenschft" (defect logged in the transcription)
(1821), and was an active collaborator on his school's organ:
Jahrbücher für wissenschaftliche Kritik. He died in 1831 of cholera.
After his death a collected edition of his works was arranged by his
disciples, in which his lectures on the philosophy of history,
aesthetics, the philosophy of religion and the history of philosophy
were also included, and the Encyclopaedia and the Philosophy of Right
were furnished with detailed additions from his lectures.

\divrule

Since Schelling had arrived at the determination of the absolute as the
unity of opposites, but had not arrived at the standpoint of ``absolute
knowledge'' by way of a conceptual development, nor had gone by such a
way from the absolute to the determination of reality, there remained
for speculative idealism the task: to develop the \emph{method}
corresponding to the content of the philosophy of identity, to arrive
by way of comprehending thought at a knowledge of the absolute idea,
and, once it was reached, to know the idea in its life-process
comprising all reality. This is the task that Hegel wants to solve: he
wants to find for the absolute content the absolute form of thought.
Whereas Schelling had demanded as the condition of philosophy an
intellectual intuition that was to be the property only of the
specially gifted, and as such took on the character of something
contingent, Hegel wants to ground the necessity of the absolute form of
knowledge, in that he methodically leads consciousness through the
lower forms of knowledge, which are dialectically sublated in virtue of
their finitude and with inner necessity lead on to the highest, the
absolute. And since for absolute knowledge the absolute receives the
determination not merely \opage{245}of \emph{substance}, i.e.\ of the
ground of the accidents, but of the \emph{subject} that posits the
opposition and through it closes together with itself as \emph{absolute
spirit}, he wants to comprehend the absolute in its unfolding that
grounds all existence, and to let the absolute content of knowledge
form itself with logical necessity into an \emph{all-embracing system},
which is at once the system of thought and of reality, and in which not
merely the principles of the immediately preceding systems but those of
all the preceding systems are to be bound together as moments in a
higher unity.

The first part of this task, the demonstration of the necessary genesis
of absolute knowledge, is to be accomplished in Hegel's Phänomenologie
des Geistes, which was originally designated as the first part of the
system; the second part of the task, the dialectical unfolding of the
total content of absolute knowledge, in the logic, the philosophy of
nature and the philosophy of spirit.

\divrule

The ``\emph{Phenomenology of Spirit}'' is the doctrine of the
development of spirit through the phenomenal forms of knowledge, i.e.\
through those forms of reflection of consciousness in which the knowing
spirit appears (erscheint) as relating itself to a \emph{foreign}
object, up to the absolute, true form of knowledge. This development is
at once that of the individual consciousness and that of the
world-spirit, the universal spirit, in which individuals have their
substantial ground. The phenomenology shows both through what shapes
humanity must pass before absolute knowledge becomes possible, and
through what states the individual must pass before absolute knowledge
attains reality in him. The dialectical development of consciousness and
the historical development of humanity are constantly set in parallel,
not always without harm to the clarity of the development and the truth
of the conception. In our exposition, in accordance with Hegel's
original thought, we place the emphasis on the development of the
individual consciousness, and disregard, where it can be done, the
historical phenomena that are drawn into the exposition. The
progressive, dialectical movement of that \opage{246}development toward
the concluding goal is brought about by this, that at every
phenomenological stage what is the object of consciousness is
transformed in the very relation to consciousness, and, as it is
transformed, alters the very form of consciousness. In each new
determination of the object, consciousness is ``für sich'' what it
previously was ``an sich'', i.e.\ its new object is its knowledge of
the earlier object. Thus there is progress until every contradiction
between the ``Ansich'' and the ``Fürsich'' of spirit, every limitation
and foreignness, is removed, until spirit, in absolute knowledge, in
which reality and concept, truth and certainty, have become one, attains
a knowledge of itself in the form of spirit, a comprehension of its
thoughts, the pure essentialities, as the objective content of reality.

The main steps of the development of consciousness are determined by
the main differences in the object of consciousness. This object is
either an object that stands over against the I, or the I itself, or
something objective that belongs just as much to the I, i.e.\ the
objective thought, the concept. The main stages in the development of
consciousness thus become: object-consciousness, self-consciousness,
reason. Rational spirit develops through the stages of objective spirit
(ethical life) and of religion to absolute knowledge.

The first form of the life of consciousness, the first stage of
\emph{object-consciousness}, is \emph{sense-certainty} in its immediacy.
Here the object is a single, external object, and consciousness
apprehends itself as a single, immediately perceiving thing. But
everything that is for consciousness and is determined by it becomes,
as it were, infected by the universality that is the fundamental
determination of the life of consciousness. In order to hold the object
of sense fast as single, as \emph{this} object, I say: it is \emph{here}
(in space), it is \emph{now} (in time); but ``here'' and ``now'' are
something constantly changing, and in this change something identical
and universal, and already as I say this Here and Now, express it in
the form of \emph{language}, the \opage{247}element of universality has
penetrated. Now and Here as \emph{represented} are something universal,
and the same holds of the I. The single cannot be \emph{said} at all.
Language, as the expression of consciousness, bears the stamp of
universality. And by another way universality penetrates into the
singularity of the object of sense, in that the object, which first
showed itself to me as the undifferentiated single, comes forward under
consideration with a plurality of determinations, and now becomes the
``einfache Allgemeinheit'' of this manifold. The object has thus altered
under consideration, and consciousness, which relates itself to it, has
been transformed with it: sensing consciousness has become
\emph{observing} consciousness, which has in its object a combination of
something universal and something sensuous-single, a mixture of
sensuous determinations and determinations of reflection. Observation
distinguishes between \emph{the thing and its properties}, and the
properties are what was previously the content of immediate sensation.
But the contradiction between the unity of the thing and the
multiplicity of the properties leads, whether one seeks to resolve it
by placing the unity in the objective and the plurality in the
subject's apprehension through the different senses, or by placing the
unity in the comprehending subject and the plurality in the object, to
consciousness's distinguishing between the object and its apprehension
of the object, between the object in and for itself and the object as
it is for another. As consciousness thus doubles the object, it is
reflecting consciousness, consciousness of the \emph{understanding}. It
distinguishes between the two determinations to which the dialectic in
the relation between the thing that is only in its properties and the
properties that are only by the thing (the substance) pointed: between
the \emph{inner} and the \emph{outer} of things (the phenomenon). The
inner is the abiding, the essence, the law, in relation to the
phenomenon, the changing, which as such is nevertheless determined by
law. But in this inner, this concrete universal, which is the thought of
things, their concept, the understanding becomes conscious that it has
to do only with its own determinations and with its own movement, and
thus consciousness has turned from a consciousness of an other, of an
objective \opage{248}object, into a consciousness of what is its own, of
what is like it, i.e.\ into \emph{self-consciousness}.

The drive of self-consciousness aims at realizing its concept and at
giving itself in everything the consciousness of itself. Its development
has 3 steps, according as the self-conscious subject relates itself to
external things, or to another isolated individual self, or to other
self-consciousnesses that it recognizes as equal and by which it is
itself recognized. In its relation to the selfless object of
object-consciousness, the self-conscious subject posits this object as
the negative, through whose sublation \emph{desiring} self-consciousness
asserts its validity as the true and finds \emph{momentary} satisfaction
in \emph{self-feeling} as a single self-consciousness. In its relation to
another self-consciousness, a relation grounded in the concept of
self-consciousness ``as a subject that is at the same time objective'',
self-consciousness seeks certainty of its validity in the other subject,
i.e.\ seeks to be recognized by it as free and as being for itself. This
recognition is won through a struggle in which the existence of the
self is put at stake, and which in the first instance leads to the
relation of \emph{lordship} and \emph{bondage}. The lord is recognized
as lord, as a free self; the slave, who prefers life to freedom, is
selfless, he has his self in his lord, his relation is that of fear and
obedience. But this relation of inequality is untrue and can only be a
passing one. The self, whose essence is freedom, must be recognized by
free self-consciousnesses, not merely by unfree, selfless ones; only
thereby can it really know itself as free. In the lord's relation to the
bondsman only his egoistic \emph{single} will is recognized by the
unfree will; the lord's will does not in this relation gain a relation
to itself as free, precisely because it is not mirrored as a universally
valid will in a free self-consciousness. Toward this goal there must be
striving. The lord must give up his egoism toward the bondsman, and the
bondsman must be educated to freedom through bondage itself. Through
obedience the sensuous desire that made him a bondsman is tamed; through
discipline \opage{249}the bondsman's self-consciousness is formed: he
gains inner freedom and, through his work, consciousness of his
justification over against the natural selfless. Thus the bondsman is
prepared to leave the condition of bondage, through a struggle for
freedom or through the lord's voluntary recognition of his freedom. The
relation between lords and bondsmen is then replaced by a relation
between self-consciousnesses that at once recognize one another and are
conscious of themselves in their justification as free and essential.
% printed "væsenslige" (defect logged in the transcription)
Self-consciousness has now become \emph{universal self-consciousness}.
The self now intuits itself not as an isolated and particular self,
different from the others, but as the universal self that is in and for
itself. Thus it knows itself in the other self-consciousnesses, in that
it recognizes their equality with itself, and is recognized as equal by
the others. It knows itself as universal, in that it knows its
reflection in the others, knows that they know it as themselves. The
self thus becomes the bearer of the consciousness of a larger
substantial whole, has universal will and universal consciousness. Here,
then, is a unity of the relation of object-consciousness and the
relation of self-consciousness: the self knows itself in its object, and
knows itself in this object, which has laid aside its isolated
singularity, according to its universal essence; it knows itself as an
essential self that is in and for itself. As such, self-consciousness
is \emph{reason}.

At this new stage conscious spirit has the certainty of its reality, in
that it has certainty that all reality is nothing other than itself,
that its thoughts are the essential determinations of the objective. The
task is now that the certainty of the objective validity of reason is
to be confirmed in everything concrete, in the whole world of natural
and spiritual reality, and thus filled out into \emph{truth}. This takes
place, in constantly advancing measure, through the stages of reason, of
spirit, of religion, and of absolute knowledge. \emph{Reason} as
\emph{theoretical} (``observing'') investigates inorganic and organic
nature, in order to find law and reason in it; it considers
self-consciousness in order to find its logical and psychological laws
and its relation to the immediate organic reality of the self.
\opage{250}It gains the knowledge that it alone is itself the object to
which it positively relates itself in its occupation with things, and as
rational self-consciousness now relates itself to the other as to its
own, it becomes \emph{practical} reason, and as such has its
self-realization as its aim. When the individual as practical, in order
to attain this realization, subordinates his individual reason to the
universal in the \emph{ethical disposition}, which is the self-conscious
reality of the universal substance, reason becomes \emph{spirit}. Spirit
designates reason as realized in an ethical world (in the family and the
state), in which ethical consciousness exists as the consciousness of
the community. Through the development of ethical consciousness from
the form of immediate, substantial ethical life, through the forms of
the reflected spirit ``become foreign to itself'', divided between the
world of culture and that of faith, and through the stages of the spirit
``certain of itself'', of the consciousness that finds the moral law in
itself, there is reached, as the ``moral'' subjectivity that fills out
the abstract command of duty from its ``conviction'', in order not to
set the law of egoism in place of the universally valid law, evil in
place of good, gives up its particularity before the unconditioned and
universal, the stage of \emph{religion}, the consciousness of absolute
spirit. The religious consciousness, which historically has run through
the stages of natural religion (the oriental religions) and of the
religion of art (the Greek religion), reaches its last stage of
development in Christianity; for in Christianity the shape of spirit
that is known by self-consciousness is self-consciousness itself; in it
the unity of the divine nature with the human is intuited. But even at
the highest step of religion there is still, since the form of
consciousness of religion is representation, with its externality and
contingency, which preserves the split between a beyond and a here, the
untrue separation between the moments of the absolute, an incongruence
between the true, absolute content and its imperfect form. The goal of
the development is reached only when content and form coincide in
\emph{absolute knowledge}, \opage{251}in which spirit \emph{comprehends}
its unity of essence with the absolute, whose process comprises all
reality, and thereby \emph{comprehends} itself as that whose
determinations are not merely subjective thoughts, but the very essence
of existence. At this stage knowing and being, certainty and truth, are
in unity; all incongruence is removed. Spirit has itself as object in
everything. --- Since \emph{science} is the result of a development in
which the earlier stages have by their finitude led beyond themselves,
these stages are preserved as moments in it; it is, as spirit's
knowledge of itself, the recollection of them as stages in the very
movement of spirit, a comprehension of the history of spirit, while they
are as its foundation, a realm of spirits, out of which its infinity
wells forth.

\divrule

Since the Phenomenology has led the knowing spirit to the standpoint of
absolute knowledge, where thought and being were grasped in unity, where
spirit recognized its realm in the whole world of reality, it has found,
in the recognized absolute unity of concept and reality, the
\emph{idea}, the absolute, which is the object of philosophy. The
absolute, which is reached as a result, not taken up through a demand;
which is apprehended through the concept, not through immediate
intuition, is precisely thereby apprehended in its self-movement, as
subject. The idea becomes the absolutely concrete, and as such it is the
\emph{process} of placing itself over against itself, and of being for
itself in this opposition. By this movement of the idea the division of
philosophy is determined. Philosophy articulates itself into logic, the
science of the idea in and for itself; into the philosophy of nature,
the science of the idea in its ``Anderssein''; and into the philosophy
of spirit, the science of the idea that returns into itself from this
otherness. The \emph{whole} of philosophy is thus \emph{science} of the
\emph{idea}, the absolute unity of concept and reality.

\emph{Logic}, the science of the idea in the abstract element of
thought, is the fundamental science, which includes within itself
metaphysics (ontology) and presents the system of the ``objective
\opage{252}thoughts'': the categories, which have validity for all
rational reality, the natural and the spiritual, whose shapes are only a
particular form of expression of those logical forms, which are
comprehended in concrete unity in the absolute idea. The task of logic
is: by the dialectical method, which mirrors the immanent movement of
the matter itself, to present the progressing development in which, from
the most abstract concept, through the unfolding of the oppositions and
their joining together into a higher unity, one arrives at the more
concrete, finally at the absolutely concrete, at the adequate expression
of the absolute, at the concept of the absolute idea, which joins the
whole circle of thoughts, all concepts, in their sublated opposition,
into its highest, living unity. This dialectic is to have for its medium
``pure thinking'', from which every element of intuition is to be
excluded, because such an element, as reflecting the externality of the
natural, bears the untrue separation within itself. In the element of
pure thought all externality is to have vanished, the stream of thought
is not to be checked by any foreign element, and the concrete idea is to
accomplish its process unhindered through the positing and sublating of
its oppositions. In the timeless process of the movement of the idea all
reality is enclosed: the dialectic of concepts is immediately the
dialectic of existence, the system of thoughts is the system of reality.
The logical thoughts are not ``only'' thoughts, but every other content
is an ``only'' over against this content of thought.

Hegel's ``dialectical method'', which is in essence the same in all the
domains of philosophy, but which in the domain of logic, owing to the
nature of the object, appears with the greatest transparency, is a
\emph{dialectic of identity}, which everywhere seeks the concrete unity
of the oppositions, the dialectical union of the opposed concepts, and
which would reproduce the necessary self-movement of the object, which
the subject merely follows as a spectator. Such a dialectic was
indirectly pointed to by the problems that the \emph{dialectic of
distinction} of the earlier philosophy could not solve. Among the
Greeks, where, in accordance with the peculiar basic stamp of their
\opage{253}thinking, the unity of oppositions everywhere, even in
Aristotle, had definitively to yield to the separation by which concepts
were to be kept free of contradiction, the concrete reality of
existence, its movement, becoming and life, remained unexplained: what
is, excluding negation from itself, became dead and abstract. In the
more recent pre-Kantian philosophy
% printed "Phisophie" (defect logged in the transcription)
the dialectic of distinction led to a metaphysics of the understanding,
which moved in finite thoughts and did not contain the true infinite. In
Kant, despite all his profound intimations, the a priori thoughts became
a circle of \emph{isolated} thoughts, which were not deduced from their
source: self-consciousness. But as they are determined as the
expression of the essence of the I, the unity of that essence points
with necessity to the unity of the thoughts; it demands the dialectic
that does not merely separate but unites, and lets thought
\emph{unfold} its determinations out of itself. Fichte would attempt
this in his Wissenschaftslehre; but his deduction of the categories
became only an apparent deduction from the reflective standpoint of
consciousness, from the ever uncomprehended opposition between the I and
the not-I, not a development out of the matter itself in its identity of
the subjective and the objective. Schelling, finally, intuited
immediately the unity of all oppositions, and demanded this intellectual
intuition of the one who philosophizes; but the intuition did not
receive the form of thought, and instead of a concrete unity we
therefore got the emptiness of indifference, instead of a systematic
unfolding of thought an external schematizing. It is thus the necessity
of comprehending the richness and life of reality, of thinking the
infinite in its true infinity and in its unfolding, and of asserting the
unity of the active thinking being, through which dialectic in the
Hegelian sense is pointed to; and by his dialectical method, which
states the ideal form of science, Hegel brings critical idealism to its
conclusion.

This dialectical method, which, in the transformation conditioned by
the difference of the standpoints, renews the Fichtean
\opage{254}constructing method with its theses, antitheses and
syntheses, comprises 3 moments: the determining of the thoughts, their
transition and dissolution, and their union in a concrete unity. The
concept (the idea), which is the essence of everything actual, is in
the first instance as the immediate, the self-identical. The
determination of it as such, and the distinguishing of it from other
concepts that is identical therewith, is the business of the
understanding; the understanding stops at fixed determinateness and at
difference; the bounded abstract counts for it as subsisting for itself
and as being. The activity of the understanding consists, namely, in
giving its content the form of abstract universality; its principle is
identity, relation to itself. Its significance is precisely the
apprehension of the object in its sharp determinateness, the clear
precision of thought, which counteracts the vague and indeterminate, and
which by its sharpness and firmness is fitted to become the foundation
for the full unity of the thoughts. The dialectical now comes forward in
that the determinations of the understanding, which in their
abstraction are finite, precisely as finite, by the nature of the
matter, dissolve of themselves, as it were come into flux and pass over
into their opposite. The finite contradicts itself in itself and
sublates itself by this contradiction: in its limit it has at once its
determinateness and its negation (determinatio = negatio). ``Everything
finite is this: to sublate itself.'' This is already clear to ordinary
consciousness, which knows that everything finite, in the domain of
nature and of spirit, instead of being something fixed and final, is on
the contrary changeable and transitory; and this is precisely nothing
other than the dialectic of the finite, by which, because in and for
itself it is the other of itself, its own negation, it is also driven
out beyond what it immediately is, and turns over into its opposite.
But when the dissolution is held fast merely \emph{as} dissolution, the
dialectic is only negative and leads to skepticism. The dialectic,
however, is by its dissolving to become the means to something higher.
Through it the one-sidedness and limitedness of the determinations of
the understanding are to show themselves as what they are: as their own
negation;
% printed "sig sig" (defect logged in the transcription)
\opage{255}but only as finite is the finite sublated, and the
dialectic, which in the actual is the principle of all movement, all
life and all development, likewise becomes the moving power in the
progressing knowing of science; it becomes the principle by which alone
immanent connection and necessity can come into the content of science,
just as in it there lies, in general, the true, not merely external,
elevation above finitude. --- Philosophy contains the skeptical as a
moment within itself, namely as that negatively dialectical element. But
philosophy does not stop at the merely negative result of dialectic, as
is the case with skepticism, which has its relative justification over
against dogmatism but cannot be anything definitive. Skepticism
misjudges its result, in that it holds it fast as mere negation, i.e.\
as abstract negation. But since the dialectic has the negative as its
\emph{result}, this negative, precisely as result, is at the same time
something positive; for it encloses within itself that from which it
results, as sublated (preserved), and is, as \emph{determinate}
negation, not without it. Seen from this point of view, it shows itself
not as abstract, one-sided, finite opposition either, but as a concept
which is higher and richer than the preceding one, because as the
negation or opposite of that one it has it for its determinateness, and
thus in its essence comprehends both opposites. ``The one thing needful
in order to gain scientific progress is the knowledge of the simple
logical proposition: that \emph{the negative is just as much
positive}.'' The concept of negation as \emph{determinate} negation is
the fundamental determination with respect to the third moment of the
logical method: the speculative or positively-rational. By this,
namely, the unity of the determinations in their opposition is
apprehended; the affirmative that is contained in their dissolution and
transition is known. The rational is thus, although as thought it has
the character of abstraction, also something concrete, a totality,
since it is a unity of different determinations, not a formal unity
devoid of difference. ``The speculative content cannot be expressed in
a \opage{256}one-sided proposition. If we say, e.g.\ that the absolute
is the unity of the subjective and the objective, this is indeed
correct, but one-sided insofar as only the unity is expressed here and
the accent is laid on it, whereas in reality the subjective and the
objective are not merely identical but also different.'' The unity is
the result of the \emph{process}, and preserves within itself the
movement of the process. The affirmative unity of the finite
oppositions, thus determined, is the true infinite. The infinite, which
would be made finite if it were to be thought \emph{outside} the
finite, or as comprehending a finite fixed as finite, or as being, in an
abstract manner, without anything finite, can be maintained in its
infinity only when it is thought as eternally determining itself to
finitude, as being infinite by and in its determining activity, while
the finite becomes a link in its process and is seen at every point to
bear infinity within itself as the moving power that leads it out
beyond itself into unity with the supplementing opposite.

The Hegelian dialectical method, which is to reproduce in the
consciousness of the thinking individual the self-movement of the
content thought, everywhere lets contradiction appear in order
everywhere to resolve it, and contradiction becomes precisely the
elastic force of thought and of life. Contradiction is the life of
finitude; its reconciliation, the mediation of the oppositions, is the
life of the idea. In the affirmation of reason negation is negated; it
is affirmation as the negation of negation, as infinite negativity.
Negativity
% printed "Negativiteteten" (defect logged in the transcription)
is the incitement to the process through which unity emerges as the
union of absolute negativity and immediacy. The true unity shows itself
to be not something merely immediate, nor yet something merely mediated,
but something mediated with and by itself, which precisely thereby is
the true immediate. ---

\emph{Logic} divides itself in Hegel, according to the moments of
immediacy, of reflection and mediation, and of the unity of
oppositions, into the doctrine of being, the doctrine of essence, and
the doctrine of the concept and the idea. The \emph{first} part of
logic, the doctrine of \emph{being}, \opage{257}thought in its
immediacy, takes its point of departure from the concept of \emph{pure
being} as the indeterminate, simple immediate, which precisely as such
must make the beginning, since the first beginning cannot be mediated
and concretely determined. It is the immediacy of indeterminateness that
we have in this thought; but in its pure abstraction it is the
absolutely negative and thus equal to \emph{nothing}, its opposite. But
the opposition is the wholly abstract and, as such, vanishing
difference. Nothing, as this immediate, equal to itself, is for its part
precisely the same as being was. The truth of being and of nothing is
therefore the union of both, their constant transition into each other:
\emph{becoming} (as coming-to-be and passing-away), the first concrete
thought, while being and nothing are empty abstractions. Now since in
becoming being, as one with nothing, and nothing, as one with being, are
only as vanishing, becoming, which was only through their difference,
must by this contradiction within itself collapse into the unity in
which both are sublated; its result is thus \emph{determinate being},
the tranquil unity in which being has negation within itself as
\emph{determinateness}, as \emph{quality}. Determinate being, as
reflected into itself in this determinateness, is \emph{that which
determinately is}: \emph{something}. It is bounded in its reality; as
bounded it has negation within itself as otherness (Anderssein).
Something becomes \emph{an other}, i.e.\ changes, because this otherness
is its own moment, and this other is itself a something that is changed
in turn, and so on to infinity. But in this infinite process, the ``bad
infinity'', only the contradiction is continued, that the finite is at
once sublated and posited again, that it becomes, in an infinity,
something and its other; but the contradiction is resolved when it is
seen that something, in its passing over into other, only goes together
with itself, since the opposites have the same determination. This
relation to itself through its other is the true, inner
\emph{infinity}, the restitution of being through the negation of
negation: \emph{being-for-itself}. In this determination lies the
ideal\opage{258}ity of the finite, the sublatedness of finitude in true
infinity as its essence. --- That which is for itself is \emph{one},
monadic: it excludes all other ones from itself and is excluded by them;
but since what it excludes is again ones, repulsion is at the same time
attraction; being is, just as much as it is plurality, the sublation of
plurality through the coalescing of the ones with one another; it is at
once discrete and continuous. This being is that of \emph{quantity}, of
``sublated quality''. --- From these first groups of categories, whose
deduction can give a notion of Hegel's application of the method, the
dialectical development now advances further. We can here only state
the main stages by way of report. The main categories of
\emph{quantity} are pure quantity, determinate quantity (the quantum)
and the quantitative ratio. Through the last, which again gives the
quantitative a qualitative determinateness, the transition is formed to
the categories of \emph{measure}, as the expression of the unity of the
qualitative and the quantitative, of the dependence of the former on the
latter; and since this unity, which under the change of the
quantitative is constantly restored by the sublation of the quality
leading to a new quality conditioned by the changed quantity, in this
way becomes a \emph{mediated} unity, which comprehends within itself the
main forms of being, the immediacy of being is sublated, and the
transition is made from the sphere of being to that of \emph{essence}.
The \emph{second} main part of logic is then the doctrine of
\emph{essence}, thought in its determinateness of reflection. As the
sphere of transition between the spheres of being and of the concept,
that of essence shows the imperfect joining of immediacy and mediation;
it is the sphere of reflection, in which contradiction is posited, since
essence, as the being that is mediated with itself through the negation
of itself, and in relation to which immediate being is only semblance,
is relation to itself only in being relation to an other. In this
sphere we have to do with such determinations of reflection as each
point to the other and cannot be thought without it (e.g.\ essence and
semblance, positive and negative, ground and consequence, cause and
effect, etc.).
\opage{259}The categories of essence fall under 3 main groups: those of
\emph{essence as grounding existence} (the pure determinations of
reflection: identity, difference and ground; existence and the thing);
those of the \emph{phenomenon}, i.e.\ of ``semblance filled with
essence'' (the world of the phenomenon; content and form; the relation
between the whole and its parts, between force and its expression, and
between the inner and the outer); and those of \emph{actuality}, the
unity of essence and phenomenon (possibility, actuality and necessity,
the relation of substantiality, the relation of causality, and
reciprocal action). Reciprocal action is the highest form of necessity;
through it, as the infinite negative relation to itself, it comes to
light that everything individual is a phenomenon of the substantial,
which determines itself to individuality and in its distinguishing is
the self-identical essence, the totality transparent in itself. Hereby
the transition is formed to the concept, which in its freedom is the
truth of necessity and of substance. The \emph{third} part of logic,
the doctrine of the \emph{concept}: thought in its developed being
for itself (Beisichsein), the substantial power being in and for itself
(anundfürsichseiende), free actuality, whose moments, which it joins
together within itself, are identical with the whole, divides into
the doctrine of the subjective or formal concept, the doctrine of
objectivity and the doctrine of the idea. The first division develops
the moments of the \emph{concept}: those of universality, particularity
and individuality; its self-separation in the judgment, its mediated
unity in the syllogism. Concept, judgment and syllogism are here
regarded not as forms of our thinking, but as forms that have validity
for all reality, as the expression of the unity of the actual with
itself, of its unfolding into its moments and of its joining of these
together into a totality. The syllogism is to be the form of everything
rational, the circuit in which the conceptual moments of reality are
mediated. The immediacy posited through the mediation of the syllogism
is \emph{objectivity}, the reality of the concept, whose development
through the categories of mechanism, chemism and teleology, taken in a
general metaphysical sense, is given by the second division. The
category of teleo\opage{260}logy, since in it the concept, in the
realization of the immanent purpose, posits itself as the essence of the
object, being in and for itself, which determines itself in the object
and relates itself to itself, forms the transition to the third
division: the doctrine of the \emph{idea} as the unity, being for
itself, of the concept and objectivity. Through the ideas of life, of
cognition and of will (of the true and of the good), the idea
determines itself as the \emph{absolute}, all-truth-comprehending,
self-thinking idea, the truth that knows itself, as the pure transparent
form of the concept, which intuits its content, the system of the
logical, as itself. The absolute idea thus encloses the systematic
totality of the logical, comprehended into unity; and the unity is not a
resting unity, but as absolute negativity the idea is process and
overreaching subjectivity. What during the dialectical development
showed itself as presuppositions of the idea is now seen as moments
posited by the self-determining of the idea, of the absolutely first;
and these moments are known, in the being-for-itself of the idea, in
their significance and validity within its totality. The dialectical
method of the development is here, at the standpoint of the idea,
converted into the speculative method: the form for the self-determining
of the idea and for its joining together with itself through the
determinations it posits.

Since the idea is the absolute unity of the pure concept and its
reality, it is, in the immediacy of \emph{being} and as totality in this
form, \emph{nature}. But since the idea is the absolutely free, its
determining of itself to nature is free, a \emph{free positing} of the
immediate idea, the idea in the moment of otherness; and in its
determining the idea is with itself, so that nature becomes only the
mediation for the free, \emph{eternal} joining together of the idea
with itself as spirit.

As immediate being, \emph{nature} is the idea in the form of otherness,
as having stepped out of itself, and externality, being-outside-one-another,
become its characteristic determinations. It is the unresolved
contradiction: to be idea, but not as idea; the concept is in it only as
something inner, its determinations (moments) have the semblance of an
indifferent subsistence side by side, and nature therefore shows in its
existence only con\opage{261}tingent determinateness and external
necessity, and it is powerless to hold fast the concept in its
irrational manifoldness. But since, on the other hand, nature is
nevertheless the form of existence of the \emph{idea}, it is, according
to its essence, a living whole, whose \emph{concept} unfolds itself
through a series of stages and, through this unfolding and
transformation, as externality successively vanishes before
being-in-itself and subjectivity, arrives at the point where, through
the sublation of natural life, spirit emerges as the truth and purpose
of nature and as the true actuality of the idea. But this unfolding is
an unfolding in the sense of the concept, not in that of time; for
nature has no history, and its stages do not emerge through a natural
production of one another, but from the inner idea that constitutes the
ground of nature. The philosophy of nature comprehends nature as reason,
in that it knows the development of the concept of nature, and thus
frees nature from its untrue existence. According to the main stages of
the development, the philosophy of nature divides into mechanics,
physics and organics. \emph{Mechanics} is the doctrine of nature in its
abstract being-outside-one-another, in its infinite dispersion, which
has the unity of form outside itself as ideal, as being only
``ansich'', as only sought --- the doctrine of matter and its ideal
system. In it the determinations of quantity prevail. It develops the
concepts: space, time, place, motion, matter, impact, fall, weight,
universal gravitation, and the realization of the latter in a system of
heavenly bodies. \emph{Physics} considers nature in its particularity,
as a manifoldness of qualitatively determined materials, whose
determinateness of form and difference are immanent to them. Under it
belong the doctrine of the free physical bodies (light), the elements,
the elemental process, specific gravity, cohesion, sound, heat, form,
(magnetism and crystallization), the particular properties of bodies
(light and colors, smell and taste, electricity) and the chemical
process. The latter points, through the relativity of the chemical
materials, to the unity that comprehends them as moments in its
totality, i.e.\ to the unity of the organic. \emph{Orga\opage{262}nics}
is the doctrine of nature in its subjectivity (organic nature), in
which the real differences of form are joined together into an ideal
unity that is for itself. The organism is conceived, as in Kant, as an
end in itself, as the totality producing itself through the taking up
and the excretion of inorganic materials and through the mutual
conditioning of its organs. Organics develops the concepts: the
geological organism, the plant world and the animal world. In the plant
world the sexual process is the highest, since through it the
individual becomes the representative of the species and, being negated
in its immediate subsisting for itself, serves the preservation of the
species. In the animal the point is reached at which the individual
organic being, in its immediate external reality, is at the same time a
self reflected into itself and negating externality: at which organic
individuality becomes sentient subjectivity, becomes soul; but animal
subjectivity is not yet for itself in the form of universality, but only
in individual, contingent determinateness: it is not yet self-conscious
and thinking subjectivity; it has only the \emph{feeling} (sensation) of
its self and of the genus, which in the death of natural individuals
shows itself as the universal power that sublates them. Liberation from
boundness to the individual is first reached within the sphere of
spirit, in human consciousness, which \emph{knows} the universal that
has existence in it.

The emergence of \emph{spirit} from nature as the latter's truth
determines it as the absolute prius of nature, since nature and matter
sublate themselves as the untrue. Spirit is, according to its universal
concept, the idea that is for itself and has, as it were, returned to
itself; it is identity with itself as absolute negativity. Its essence
is therefore freedom, its determinateness self-revelation; its
possibility is immediately absolute actuality. As spirit, in the
realization of its concept, raises itself from the form of immediacy,
of naturalness, to the absolute unity of concept and reality, it
determines itself from \emph{subjective} spirit, i.e.\ spirit that
unfolds itself from its natural form to knowledge of itself as free,
into \emph{objective} spirit, i.e.\ spirit that as free produces a
\opage{263}world as the realization of freedom, and, in its highest
form, into \emph{absolute} spirit, i.e.\ spirit in the unity, being in
and for itself and eternally producing itself, of its concept and its
objectivity. Subjective and objective spirit are \emph{finite} spirit,
and their finitude is the barrier that absolute spirit posits in order,
by its sublation, to be absolutely revealed in its freedom. Through this
realization of the concept of spirit the natural vanishes as something
independent; it becomes a means, existing only for finite spirit, of
raising itself up to its concept. From the vanishing relation to nature
spirit as subjective raises itself to the relation to itself; and as it
realizes its concept through freedom, so that substantial freedom
attains existence as subjective will and as objective actuality and
necessity, it arrives, through the world-historical process, in which
the individual national spirits are moments, by casting off their
limitedness, at the actualization of the world-spirit as knowledge of
absolute spirit as the eternally actual truth.

More precisely, the development of the concept of spirit proceeds
through the following stages. The \emph{first} main part of the
philosophy of spirit, the doctrine of \emph{subjective spirit},
considers in anthropology, phenomenology and psychology the main forms
of this spirit as soul, consciousness and spirit. The object of
\emph{anthropology} is immediate spirit, spirit as natural spirit, the
soul; the ideal unity of the organic body. The reality of spirit is at
this stage immediately existent, given by nature. Through the overcoming
of this natural determinateness and through the ideal tearing loose from
it, the soul raises itself to ideal identity with itself, to a
being-for-itself in which it distinguishes itself from the external,
i.e.\ to consciousness, the form of reflection of spirit.
\emph{Phenomenology}, which here comes again as a member of the system
and as such only demonstrates the development of individual
consciousness in respect of its \emph{form}, considers the determining of
consciousness to self-consciousness and to reason, and, since reason is
the substantial nature of spirit, leads over to \emph{psychology}, which
considers spirit \opage{264}as such, spirit that has the certainty of
relating, in the object, only to itself, to its own determinations, and
that confirms and fills out this certainty by making itself, through a
development, into what it is according to its essence. Through the
stages of development
% printed "Udviklingtrin" (defect logged in the transcription)
of intelligence (theoretical spirit): intuition, representation and
thinking, and through the stages of will (practical spirit), of thinking
that determines itself to action: practical feeling, impulse and the
will that seeks something universal, spirit determines itself as free
\emph{spirit}, i.e.\ spirit that knows itself as free and in its
self-determination wills its essence, freedom. The actualization of the
content of the freedom determined by reason determines spirit as
\emph{objective} spirit.

In the doctrine of \emph{objective} spirit the legal and the moral,
which in Kant and Fichte formed an abstract opposition, are conceived as
moments in the higher unity of ethical life; and the philosophy of
right, the science of right as the realization of freedom, of its forms
of existence, divides into the doctrine of formal right, i.e.\ the right
of abstract personality, the doctrine of morality, the right of the
subjective will, and the doctrine of ethical life, substantial freedom
existing as subjective will and as objective actuality and necessity.
Under \emph{formal right} the concepts are developed: person, property
(as the external sphere for the person's relations), legal agreement
(Vertrag), violation of right, and the restitution of right in
punishment, which is conceived from the standpoint of retribution as the
application to the criminal himself of the law set up by him through his
action. \emph{Morality} Hegel conceives one-sidedly as the standpoint
where the will bears within itself the unresolved opposition between a
willing of the command of duty and a determinateness by sensuous
impulse, which makes the realization of the former an infinite task, a
perpetual ought, and where the abstract command of duty is to receive
determinateness in the concrete case through conscience, through the
contingently subjective, which can lead just as well in the direction of
evil as of good. The contradiction of the moral consciousness is
demonstrated through the development \opage{265}of the concepts: purpose
and action, intention and the welfare of the individual, conscience,
good and evil. Through the demand that the abstract good receive
concrete determinateness and abstract subjectivity be filled out with
this determinate content, the transition is formed to the sphere of
\emph{ethical life}, which expresses this concrete unity of the good and
the subjective will: ``Ethical life is the concept of freedom become an
objective world and the nature of self-consciousness.'' At its
standpoint the individual is seen in harmony with the ethical powers
that assert themselves in laws and institutions and that, as the
individual's ethical substance, govern its life. The doctrine of ethical
life articulates itself into the doctrine of the family, of civil
society (in which individuals pursue their particular ends, but the
realization of these, the securing of their subsistence, their welfare
and their right, is conditioned by the universal, by the social
relation), and of the state. Hegel treats the last part with particular
fondness; it comprises the doctrine of the internal right of the state
(the constitution), of the external right of the state (the mutual
relation of states as individuals) and of world history. Hegel conceives
the significance and task of the state in the spirit of antiquity. The
state is for him ``the actuality of the ethical idea'', ``the
self-conscious ethical substance'', ``the ethical spirit unfolded into
an organic actuality''; it is ``absolute end in itself''. It is the
highest duty of individuals to be members of the state, and the state
has the highest right over against individuals, and makes the strictest
demands for a merging of their ends and interests in it, for the
patriotic disposition (the willing of the end of the state become
habit), which draws its content from the substantial realized in the
state, and in turn becomes a living bearer of the institutions of the
state. Against an individualistic conception of the state Hegel, with
vigorous one-sidedness, asserts the right and significance of the
substantial over against subjective arbitrariness. --- Since the state
is still bound to natural conditions, and since the individual state and
the individual national spirit in its particularity has a
limita\opage{266}tion and a finitude, states and national spirits must
enter as moments into the process of \emph{world history}, which shows
the dialectic of the finite national spirits, and in which each single
stage is represented by the people that is the bearer of the principle
of the epoch and therefore the dominant one, until its time has run out
and a new people, as representative of the higher principle of a new
time, takes the hegemony and executes history's judgment of the world.
The \emph{philosophy of history} comprehends that process as the one in
which the essence of the human spirit, freedom, wins its true form and
universal actuality, and the \emph{universal} thinking spirit of world
history, the \emph{world-spirit}, actualizes its knowledge of absolute
spirit as the eternally actual truth, in which knowing reason is free,
for itself, and necessity, nature and history are only as serving its
revelation and as instruments of its glory.

The subjective consciousness of \emph{absolute spirit}, which is itself
a moment of absolute spirit, raises itself from the form of intuition to
that of feeling and representation, and from this to that of thought;
and to these forms correspond \emph{art}, \emph{religion} and
\emph{philosophy}. \emph{Aesthetics}, as the doctrine of the beautiful,
i.e.\ the absolute in sensuous existence, the idea as actual in the form
of the bounded phenomenon, develops in its first part the metaphysics of
the beautiful, and determines the concept of the beautiful, in
particular of the beautiful in art, of the ideal, and of artistic
productivity. In its second part it treats the development, conditioned
by the different relation between the moments of the ideal, idea and
material, of the particular \emph{forms of art}: the \emph{symbolic}, in
which form is unable to permeate the material completely; the
\emph{classical}, where the ideal content is fully and harmoniously
expressed in the sensuous form; and the \emph{romantic}, in which the
depth of the heart, the infinity of subjectivity, spirit's breaking
through the material, leads to a new incongruence between form and
material, a going of art beyond itself, but in the form of art. In the
third part of aesthetics \emph{the system of the individual arts} is
treated: architecture,
sculp\opage{267}\setcounter{footnote}{0}ture, painting, music and
poetry. The principle of division is sought in the concept of the work
of art itself, and the arts are paralleled with the forms of art in such
a way that architecture becomes preeminently the expression of the
symbolic, sculpture of the classical, the three other arts of the
romantic form of art. The transition from the sphere of art to that of
religion is conditioned by the incongruity between the absolute content
and the finite, bounded, sensuous form of intuition in which art brings
it to consciousness.

In \emph{religion} the consciousness of the absolute has the form of
\emph{feeling} and of \emph{representation}: the sensuous element of art
has receded before the inwardness of the heart, in which the absolute
has presence; but since the absolute here does not yet receive the
adequate form of the concept, but is had only in the undeveloped state of
feeling or, as unfolded, in representation, this form, in which a
remnant of the sensuousness of intuition is preserved, has the effect
that the moments of the absolute are separated as independent shapes and
events, and that its eternal process receives the stamp of externality
and of time, even in the highest form of religion: the absolute
religion, Christianity\footnote{Christianity is the ``revealed
religion'', in that what is the essence of religion, namely that God
knows himself through man's consciousness, that he is as spirit, in his
congregation, is known in it, i.e.\ has become revealed. Its main notion
becomes the unity of the divine and the human essence, God's complete
self-communication to finite spirit, which through the overcoming of its
finitude and selfishness actualizes the unity of essence.}, to which
religion develops through the natural religions and through those
religions (the ``religions of spiritual individuality'') which, like the
Jewish, the Greek and the Roman, conceive God as subject, though not yet
as infinite subjectivity. The movement of the absolute idea within
itself, by which, through its self-separation, it emerges as absolutely
for itself existing \emph{spirit}, which knows itself in man, whose
knowledge \opage{268}of God becomes his knowledge of himself in God, is
separated for representation into processes falling outside one
another, each of which by itself presents the absolute content. The
eternal process of the divine Trinity is separated from the process
whose moments are the creation of nature and of finite spirit, the
breach of the created with the divine, and the reconciliation through
the historical God-man, which is appropriated through faith and the
negation of finitude in the true cultus. Since the death of Christ, as
the cessation of his sensuous existence, makes a spiritual conception of
him possible, and what was intuited in him receives a spiritual
existence in the consciousness of the congregation, the external
reconciliation is transformed, for the subject who gives up his
finitude, into an inner one, and thought knows, in what representation
separated, the inseparable, eternal, single self-revelation of absolute
spirit. In this knowing the form of consciousness of religion, whose
significance and development the \emph{philosophy of religion}
determines, is replaced by that of philosophy.

Since in \emph{philosophy} the absolute is the object of comprehending
thought, it is freed from the finite forms of art and religion, and
raised to the absolute form adequate to the content. Through the
sublation of those finite forms, phenomenologically necessary, the
thinking human spirit arrives in philosophy at comprehending, by
speculative thought, the eternal idea, being in and for itself, in the
process comprehending nature and finite spirit, in which it
\emph{eternally} produces and enjoys itself as absolute spirit. In this
comprehending, the semblance of nature and of finite spirit as
subsisting independently, and as conditioning absolute spirit, vanishes;
all reality is seen as the idea's own self-realization, spirit's
knowledge of it as its knowledge of itself in man, whose knowledge of the
absolute becomes his knowledge of himself in the absolute. The eternal
self-separation of the idea is eternal union: the temporal process is,
seen from the standpoint of the idea, eternally sublated in the eternal
process. --- While the \emph{system of philosophy} gives the moments of
the concept of philosophy, which at once concludes \opage{269}and
comprehends the system, the \emph{history of philosophy} presents the
becoming of this concept through the historically given series of
systems, which Hegel would conceive as developing out of one another
with logical necessity, and whose principles, which in their order of
succession are to correspond to the order of succession of the main
concepts of his system, he sees as moments that are all comprehended
into unity in the ``absolute philosophy''.

\divrule

The significance of the Hegelian system, and the source of the
extensive influence it has exercised, is essentially to be sought in
its method, by which Hegel, with an immense energy of thought, let an
edifice of concepts grow up more comprehensive, more fully worked out,
more firmly joined together and resting on a broader metaphysical
foundation than any earlier one, --- and in the rich working-out of
the sciences of objective and absolute spirit, in whose domain Hegel
everywhere, with a comprehensive and surveying gaze and with a
sympathetic understanding of the great and the beautiful, grasps the
essential, and deepens knowledge by letting it penetrate to what moves
most deeply, to what governs the whole development. In this system the
idealistic line of development of modern philosophy has found its
conclusion. The unity of existence, to which idealism sought to lead
the oppositions back, is here reached. All existence is here seen as
grounded in the process of the idea by which it eternally produces
itself as spirit: thought and being, reason and actuality are joined
together into a unity that lays claim to having removed all dualism.
But yet it turns out, as was already intimated earlier (Part 1, p. 9),
that the unity at which idealism arrives is still only one-sided,
brought about at the expense of being, of the real. In Hegel, thought
as ``pure thought'' is to comprehend all existence, to be in itself the
unity of thought and being. But the unity in which the reconciliation
of being and thought is already \opage{270}\setcounter{footnote}{0}brought
about in the domain of the logical, before being comes forth as the
being of nature, in the determinations of reality, is precisely only
the unity of thought and the abstract thought of being, the most
abstract determination of thought, a unity of thought with itself, in
which being as actually real being finds no place\footnote{Only by a
play with the concept: the immediate, is an image of real being brought
into the movement of thought. In this way nature is slipped into a
moment of the process of the logical idea: into the process's joining
together of the oppositions in immediate unity.}, and Hegel's system
has therefore rightly been designated as \emph{panlogism}. Although the
plan of the system, inasmuch as the idea through nature determines
itself to spirit, to the concept that is first the absolute expression
of the absolute, that toward which all religion and philosophy, indeed
all world-historical development, has striven (Hegel's Encycl. § 384),
--- had to demand an independent positive significance for nature, in
order that a fuller determinateness of spirit could unfold from it, yet
nature \emph{as} nature becomes only an untrue form of existence, one
that exists only for man, the finite spirit, and that \emph{only}
obtains truth as spirit, as sublated nature; its being is only to
vanish; it is posited as a nothing; in its richness and concrete
multiplicity is seen only the sign of its impotence to hold fast the
concept, and its life is designated, in accordance with this impotence,
as anxious and sickly.

But this conception of nature, and the reduction of the unity of
thought and being to the unity of thought with itself in its
abstraction, involves decisive consequences for the system. It brings
about, in the first place, that the concepts of the absolute idea and
of absolute spirit must in Hegel flow together, since the idea already
as the prius of nature becomes the absolute unity of concept and
reality (objectivity), and that absolute spirit must therefore become
the \emph{natureless} spirit, whose concept does not overcome dualism,
because natural being in its determinateness of reality becomes what is
incommensurable with the logical-metaphysical \opage{271}concept, what
comprises an ``existing irrational,'' thus something that is not
comprehended from the sole actuality of the idea. It prevents, next,
the carrying-through of the method by which the system was to be built
up with logical necessity, and whose moments and fundamental form Hegel
has stated. And it leads finally, despite all the significance that the
temporal development of the human spirit seems to acquire in Hegel, to
a volatilization of the actual life of the human spirit, of its ethical
life and its historical development.

Since in Hegel the logical pure thought in its abstraction, which does
not contain the real in its actuality --- a relation Hegel himself
indirectly acknowledges by letting nature be impotent to hold fast the
concept --- separates off from itself every element of intuition,
because only thought freed from all elements of sensibility is to be
able to have the infinite mobility that is the condition for the
dialectical joining of concepts into unity, this determination of
thought, grounded in dualism, affects the dialectic through which the
mutual relation of the logical concepts is to be determined, and the
character and connection of the concepts that are the content of the
empirical sciences. It turns out that beneath Hegel's dialectic of
identity, in which the determinateness of subjectivity of the ideal
principle is mirrored, the critical distinction still hides itself, and
that the distinction of the understanding is unovercome. Since thought
as pure thought pushes intuition aside in order to close itself up
within itself, the dialectic is thereby deprived of the incitement from
the real, concrete actuality that is the condition for this: that human thought,
which by its nature stands in relation to the totality of existence, and
which relates itself to itself only through its other, can,
through a true unfolding of thought, actualize the inner unity among
its forms, which in their mutual relation must reflect the relations in
the other of thought. When the movement of pure thought is to get
going, this happens only in that, in a hidden and inconsistent manner,
the element of intuition, which mirrors the relation to the other of
thought, is \opage{272}called in to help. But the unity of the concepts
that is brought about by a dialectic kept in motion by an element that
is excluded in principle, and that is taken up only in an external
manner, can only be illusory. And when the dialectic then advances from
the pure logical concepts to the concepts of concrete existence, it
becomes unavoidable that these concepts, while abstract thought carries
out with them the process of volatilization that would make them mere
reflexes of the logical categories, nevertheless fetch from intuition
and experience an element of reality, in which the distinction of the
real, incongruent with logical thought and unresolved in thought, is
mirrored; and through this determinateness by the foreign real the
concepts show themselves as brought only to an illusory unity: the
distinction asserts itself again.

As regards the conception of the actual life of the human spirit, it is
determined by the conception of the absolute idea that lets the idea,
\emph{as} the highest expression of ``pure thought,'' be the unity of
thought and being. From this conception of the idea it follows, namely,
as a direct consequence, that the \emph{eternal} process of the idea, in
which nature and the existence of finite spirit, whose independence is
only a semblance, are eternally taken up as sublated moments, must be
posited as including all actuality and all truth. Consistently, no
other true existence can be posited than that which is in the element
of metaphysical eternity; all that has subsistence in the forms of time
and space must be seen as something \emph{only} untrue, which with
eternal necessity is sublated in the absolute. The natural forms and
natural content of human life can thus become only something invalid
and untrue, and the only true being of the human spirit must become the
contemplative life of thought, freed from the limitation of
individuality, in the element of eternity: that form of thought in
which the absolute is known as that which alone is, in its eternal
process, and the thinking subject knows itself as a moment in the
absolute, in whose knowledge of itself man's knowledge of the absolute
and of himself as a moment in the absolute is taken up.
\opage{273}Over against the view of the eternal metaphysical process of
the idea, which is not a striving toward a goal but the
self-affirmation of what eternally is, the historical development of
the human spirit, whose essentiality Hegel sometimes stresses so
strongly that one has been able to ascribe to him the doctrine that
absolute spirit comes to a knowledge of itself only through the
development of the human spirit, loses its independent significance and
shrinks to being the way to the actualization, in the human spirit, of
that form of knowledge by which the human spirit knows itself as
enclosed by the idea. And over against the view of the necessity of the
process of the idea, in which every opposition is eternally canceled,
the problems of actuality of the ethical-religious life lose their
sharpness and their weight. No true emphasis can fall on the free
self-determination of the will, on its free relation in time to the
ethical tasks, on its concrete actualization of the human being as
natural-spiritual in a world of ethical life. Over against the
necessity of the eternal and eternally finished process, the movement
of finitude, even where it concludes in evil, becomes only a jest,
eternally forgotten in the process of the absolute, in which all the
seriousness falls. Through the thought of the absolute the individual
can take itself out of the contradictions of finitude, which can obtain
no significance for its life of eternity, and in the individualityless,
contemplative unity with the absolute it has absolution for all
imperfection and all guilt, and reconciliation with its essence.
Thought is action; to comprehend the absolute, the absolute action.
Life is volatilized, in that man in his nature-determined existence is
to live in the abstraction of pure spirituality, in time to take time
back into eternity; the ethical and religious relation to actuality is
sublimated into a timeless relation of thought, in which the thinking
subject nevertheless cannot find rest, because in its yearning for
eternity it must be thrown back upon the actuality that is not taken up,
as reconciled, into that eternity. ---

\opage{274}As the conclusion of modern idealism, Hegel's system forms
the parallel to the dogmatic conclusion of realism: its natureless
spirit is the counterpart of materialism's spiritless nature, and it
fully unfolds the thoughts that we have found in germ in the founder of
critical idealism. The concept of the I, which we already saw in Kant
growing out beyond itself, is in Hegel unfolded into the concept of
absolute spirit, which in its eternal process comprises nature and the
finite consciousnesses. And Kant's conception of that spontaneity of
the mind which idealism, as it becomes clearly conscious of the essence
of knowledge, as it gains consciousness of its knowing, must definitively
assert, and of the relation of the a priori to the unity of active
self-consciousness, receives in Hegel a corresponding development. The
result of the spontaneity of the mind, which Kant determined as the sum
of the separate forms of the merely subjective a priori, whose
proceeding from the unity of the self he asserted but did not
demonstrate, becomes for Hegel the \emph{system} of thoughts, enclosed
by the concept of the idea, in unity with actuality, which the
self-conscious human spirit unfolds through the dialectical method; and
as the standpoint is shifted from self-consciousness to the absolute,
in which self-consciousness comprehends itself as a moment, and the
system of thought is seen in its relation to this absolute, its
unfolding becomes one with the eternal unfolding of all actuality by
absolute spirit, determined as absolute subjectivity, and spirit
becomes, as the living source of the system of thought, the origin of
all actuality.

But when modern idealism has thus in the Hegelian philosophy reached
its goal and unfolded its plan, it has already been shown that the
concluding concept by which this system, which proclaimed that the
rational was the actual, the actual the rational, and which would see in
all natural and spiritual existence the unfolding of the grounding
idea, sought to express the unity of the oppositions of existence, the
unity of thought and being, spirit and nature, expressed only a
one-sidedly conceived unity: one in which the moment of reality, the
moment of nature, was volatilized, and over against which
\opage{275}it therefore had to assert itself in its form of actuality.
The contradiction that, owing to the hidden working of dualism, pervades
the whole of modern philosophy in its development is thus still present
at this concluding point, just as it was present at the conclusion of
realistic philosophy, and hereby a demand is then once more placed upon
philosophy: the demand that it shall attempt, by another way than those
hitherto trodden, the solution of the task that designates the
fundamental task of philosophy, determined by the essence of the human
mind: the knowledge of the unity of existence through the knowledge of
the principle as that which in its concrete unity joins together the
fundamental oppositions of existence. And the whole preceding
historical development points toward the way by which the solution is
to be attempted, and which the philosophy of the future must follow if
it would overcome the dualism that the preceding main philosophical
forms were not able to master.

\divrule

\begin{center}
  \textbf{Conclusion.}
\end{center}

\divrule

The fundamental form of philosophy that shows itself in its necessity
through the deficiencies of the preceding stages of development is,
according to what was already shown in the Survey (Part 1, p. 9 f.), to
be conceived as the last main form of philosophical development, and as
the one that bears within itself the germ of an unfolding through an
indeterminable span of time. The former followed from the juxtaposition
of the various possibilities for determining the relation of the
oppositions of existence; the latter from the significance that is
attributed to the knowledge of real actuality in its concretion. When
it has been shown that neither the determination of the relation
between the fundamental oppositions of existence as immediate unity,
nor as abstract opposition, nor \opage{276}the one-sided, exclusive
holding fast to the one member of the opposition, to which both the
dissolution of immediate unity and the consequences of abstract
opposition could lead, and which in reality is only a concealed form of
dualism, can be carried through, then only one determination of the
relation still remains possible: its determination as the unity of
opposites. And since this determination expresses that the real and the
ideal, nature and mind, while they assert themselves in their
peculiarity \emph{as} opposites, are joined together into concrete,
living unity in the principle, which can be thought only when the
essence of the principle is the source of the opposition, and the
principle is therefore, in its determinateness as absolute subjectivity,
an ideal-real principle, then by that determination of the relation of
the fundamental oppositions \emph{the essentiality of the real} is
asserted, and the knowledge of it in its infinite richness is made the
task of philosophy. ---

As the last possible fundamental form of philosophy, the demanded
\emph{philosophy of the future} must realize in completeness the task,
with respect to the content and form of philosophical knowledge, that
has been realized partially, and with the imperfection that partial
realization brings with it, in the main forms historically at hand. It
must then comprehend and preserve within itself all that is essential
and abiding in the preceding forms; for only under this condition will
what has not been realized in them be able to be actualized in it. And
since it comprehends within itself that essential and abiding result of
the earlier stages of development, it will not be able to relate itself
to what precedes and prepares it abstractly, as opposition or as the
unfolding of a single direction, and so will not be able to repeat the
relation of the earlier stages to one another; but it must relate
itself to what historically precedes it as the unity of all the
oppositions and as the completion of all that precedes.

This relation to the earlier stages must show itself most directly in
the determination of the \emph{principle}. Since it was the dualism
that at the close of the second main stage again \opage{277}came
forward as unovercome by which the demand for a new main form of
philosophy was made irrefutable, the most immediate task for this form,
which is to assert the unity of opposites, must be to transform the
concept of nature and the concept of mind in which dualism had
originally found its expression, and through which it continued to
assert itself. In place of the \emph{concept of nature} that had been
formed from the presupposition of nature's abstract opposition to mind,
and that determined nature only by what designates the real as it were
at its greatest distance from the unity and energy of the ideal, namely
by passive extension, there must come a concept that does not conceive
nature according to its poorest and most abstract determination, but
from the point where, in its highest phenomena, it expresses its
\emph{complete} essence most richly and most clearly: from the point of
view of the \emph{organism}, of the living. And since nature here shows
itself as determined by something ideal, for which the material
elements become serving moments, the one-sidedness of speculative
idealism must in turn be avoided, which lets the real evaporate and its
forms of being vanish into those of mind. The natural must be conceived
in its \emph{essentially valid} materiality; the fundamental
determinations of the real, the determinations of space, of time and of
matter, must be seen in their essentiality, in their decisive
significance for the thinkability of the fundamental determination of
the ideal: the determination of energy.

To the transformation of the concept of nature there will then
correspond an analogous transformation of the \emph{concept of mind}.
In place of the concept of mind that posited it as torn loose from all
the determinations of the real and as the abstract opposite of the
natural, as consisting wholly in thought and consciousness, there must
be substituted a concept formed in agreement with the experience that
shows the individual existence of mind as borne by, and in its
development in time conditioned by, a real organic foundation, and its
conscious mental life as going in unity with an unconscious psychic
life, by which it is connected with the objective reason of nature; so
that mind, in its freedom and ideality, is thus seen in its
con\opage{278}\setcounter{footnote}{0}nection with nature, as the
latter's completion and as standing in an essential relation to the
real. From an unprejudiced experience, whether nature or mind be taken
as the point of departure, one will arrive at the same essential unity
of the determinations of ideality and of reality under their opposition
as by the logical analysis of the conceptual relations.

Such a transformation of the concepts of the fundamental oppositions,
by which each of them is seen as pointing to that of its opposite and
as bearing it within itself, includes the determination of the
\emph{principle} as the \emph{concrete unity} that embraces both
opposites, as the \emph{ideal-real} principle of unity of existence. As
the \emph{active unity} of the opposites the principle can be thought
only when it is determined as \emph{absolute} \emph{subjectivity}, and
subjectivity is then seen as that which has actuality only in its
\emph{necessary} and \emph{original} relation to otherness, to the
opposite of subjectivity: the \emph{objective}, the real. Both
determinations receive, as necessary, equal originality, but conceptual
priority belongs to the ideal, since subjectivity becomes the ground of
explanation for itself and for objectivity, which is through it, but is
\emph{as the condition of its own being}. From this point of view,
then, the real is seen not as a vanishing moment, but as that which, in
its peculiar determinateness of being, is the \emph{abiding} condition
for the actuality of the ideal: it is seen as that which, in the
ideal's eternal positing of itself, is posited in order to be
idealized, but at the same time as that which is idealized in its real
actuality\footnote{
Cf., with respect to the closer development of this, my work:
Om det Religiøse i dets Enhed med det Humane, p. 65 ff.; 86 ff.}.
The principle, in being seen in its \emph{ideality}, is by the same
token seen as that which grounds \emph{real} existence and determines
it teleologically; it is seen as the \emph{reason of existence}, by
which existence \emph{is} as \emph{real} and subsists in \emph{real}
actuality.

\opage{279}
From this determination of the principle there then follows the
determination of the \emph{essential} \emph{object} of philosophical
knowledge: that this object can no longer, as at earlier stages of
development, where the principle was conceived incompletely and
one-sidedly, be a single, one-sidedly separated-out fundamental form of
existence, or existence in a determinateness of abstraction by which it
shrinks into the abstract being of the principle; but that it must be
the concrete, ideal-real whole of existence, unfolded in relative
independence, in its living connection with the grounding and
organizing principle. In this determination of the essential object it
is implied, on the one hand, that an essential emphasis must fall on
the single, the concrete, the individual, which in its \emph{real}
being steps out into relative independence and, as an \emph{integral}
member of the totality, acquires in its specific peculiarity essential
significance for the whole, and therefore is not to be seen as a merely
vanishing moment, but asserts itself in the peculiarity of its essence.
But on the other hand it is implied that the real, individual
existences, precisely as members of the indivisible totality, must be
seen as such as cannot be separated out in isolation, but must at every
point, even where a selfhood, through self-consciousness and will,
consciously opposes itself to the principle, reveal their living
connection with the whole and with its ideal principle of unity. In its
relation to its object, thought then acquires at every point a concrete
content of actuality, and everywhere asserts actuality as actuality;
but the content of actuality must everywhere show itself to thought in
the light of the concrete principle of unity, which clarifies it
through and through and idealizes it, but does not consume its
actuality.

The \emph{philosophical method} that is to correspond to the stated
determination of the principle in its relation to organized existence,
in which the separation of realities and the ideal unity both come into
their own, must be a \emph{dialectical} method in the \emph{complete}
meaning of this concept. It cannot be a dialectic of distinction that
overlooks the \opage{280}unity of essence, the inner infinity, nor yet
an abstract dialectic of identity that contents itself with unfolding a
system of abstract a priori concepts which do not contain the real; but
it must be a \emph{dialectic of identity} that unfolds a \emph{system of
concepts of actuality} having \emph{the necessity and transparency of
the a priori}. The task thus becomes an assimilation of the real in its
essentially valid peculiarity, an assimilation that coincides with the
unfolding of the a priori in concretion. The way to this becomes a
combination of the true method of experience, i.e.\ of the method of
experience completed and corrected by the recognition of the a priori,
with the immanent conceptual development of dialectic. Since that
method of experience leads to the point where the real, through the
concrete knowledge of its essential determinations and essential
relations, is comprehended as an ideal moment in an organic totality,
at this point its result will be able to be converted into dialectical
development of thought, assimilated by the a priori unfolding of
thought, and, through a constant relation of reciprocal action between
thought and the real in its living actuality, transformed into
something ideally produced by the subject, as the givenness of the
empirical is converted into total production of the formed content out
of the inner infinity of subjective thought. In this way the real gets
the right to assert itself according to its peculiarity, in order to be
taken up into the concept in concrete shape. The unity of the concepts
then becomes one in which actuality is contained: it acquires not
merely a concretion within the abstract, but a concretion in which the
essence of the real is preserved. The dialectical unity of the concepts
rests on the knowledge of actuality, and since \emph{a priori}
knowledge enters into relation with \emph{empirical} knowledge, the
boundary between them becomes, in accordance with the relation of the
human mind, universal in its essence, to the total organization of
existence, a \emph{constantly movable} boundary, as the a priori
successively assimilates from itself the empirical known according to
its concrete actuality, and at each stage of \opage{281}this process,
in its relative abstraction, gives an expression, true in principle and
as it were striving of itself toward its integration, of the actual,
which at every point reveals the principle and the totality. In
dialectic's development of its concrete content, the advancing but not
consummated appropriation of the real must come forward in such a form
that what is not yet appropriated is so linked in the development to
what has been appropriated that those points are brought out in which
ideality as it were shines through, and which thereby point toward the
more comprehensive appropriation. The philosophical exposition will
then surround its formed content with a material not yet converted into
the concept, but with such a material as betrays the possibility of
this conversion, since it is already, as it were, shone through by
single rays of the explaining thought.

What determines the method of philosophy at this standpoint determines
also its \emph{relation to the particular, empirical sciences}. Since
the determination of the principle includes the recognition of the
essentiality of the real, philosophy must stand in a constant relation
to the totality of this real. But thereby it will be in a constant
relation to the particular sciences, which relate themselves each to
its own domain of what is, and through their peculiar methods prepare
the real material of knowledge for a philosophical knowledge. The
particular sciences, which have their peculiar form and peculiar
method, conditioned by the peculiarity of their material, acquire
through their peculiarity an independence that philosophy must
recognize; but this independence can only be a relative independence.
The relativity is grounded in the merely partial and relative
realization of the concept of science in the particular sciences, which
relate themselves to something partial and rest on presuppositions that
they do not themselves ground. By virtue of its principle, which is not
a merely relative or presupposed one, but a universal one, grounded
through itself; by virtue of its object, which is not a limited one,
but a universal one; by virtue of the kind of its knowledge, which does
not merely \opage{282}relate itself to the object, but to knowledge
itself and its fundamental relations; and by virtue of its method,
which comprehends as moments within itself the method of experience,
the common determination of the peculiar methods of the individual
empirical sciences, and the exact deductive method of mathematics, and
unfolds the content of knowledge into a unity determined by principle,
--- philosophy becomes the adequate realization of the concept of
science, and must therefore have the primacy in its relation to the
particular sciences, and become their regulator. It becomes this
through its consciousness of the validity of the relative principles
and of the fundamental relations of knowledge (cf.\ Part 1, p. 1 f.);
but since it is to control the conceptual determinations of the
particular sciences and their applications of their method, it cannot
seek to force upon them a kind of knowledge foreign to them, but must
precisely let their kind of knowledge, preserved in its peculiarity,
lead, while it is secured against overstepping its limits, to results
that philosophical knowledge appropriates and assimilates through a
method that is not foreign to the method of those sciences, but has
precisely appropriated it as a moment. In this relation philosophy
shows itself in its inner unity with the particular sciences, and
precisely as these develop in their whole peculiarity, they become
conscious of their common task and their common means, and work, in a
living reciprocal action, consciously toward a \emph{unity of science}
that is realized through philosophical knowledge. The organization of
the empirical sciences, animated by philosophy, as such a living unity,
in which the individual members have their relative independence, must
become the goal of the historical development of philosophy and of the
particular sciences. ---

If, finally, we seek the \emph{view of life} that would correspond to
the conception of the principle within this historical main form of
philosophy, the determinations of the principle lead directly to the
determination of its peculiarity.

\opage{283}
If the principle is determined as the \emph{ideal-real} principle, thus
as the principle for which the real has \emph{essential} significance,
then the view of life too must assert for the real an essential
significance for human life, and set up the demand that man shall
develop according to \emph{the completeness of his essence as
natural-mental}, and, through such a development, in which the real
natural foundation is consciously formed through and through, give the
will a content of actuality.

If the principle, as ideal-real, is determined as that which grounds an
\emph{organic} ideal-real totality of existence, in which the human
being is ordered as an \emph{integral} member in the \emph{relative
independence} that the essentiality of the determination of reality
makes possible, and of which man, by virtue of the universality of his
essence, can become conscious, then the view of life must express this
by setting up the demand that man, in asserting himself in the
peculiarity of his complete, natural-mental being, shall in his
self-assertion become conscious of himself and will himself as a
\emph{member} of the totality in which he is harmoniously ordered, and
as determined by the principle that is the totality's, and thus, out of
his independence, freely determine himself to self-surrender. This
means, in other words: he shall, in asserting himself in his peculiar
being, assert himself \emph{ethically} in it. The ethical thus becomes
the demand on human life, and this ethical is determined in such a way
that it can embrace human life in its concrete actuality.

If, finally, the principle is determined as that which, in its joining
together of the opposites into concrete unity, \emph{excludes} all
\emph{dualism} from itself, then the view of life must express this by
conceiving the human being, which is in unity of essence with the
principle, as freed from all dualism, by seeing it as undivided, in
inner unity with itself in the harmony of its fundamental functions.
The unity of the principle can be reached only when what is the
phenomenological point of departure for the knowledge of it, finite
existence and the finite human mind, harbors within
\opage{284} \setcounter{footnote}{0}itself no opposition that cannot be
sublated. And conversely, from the concept of the principle as true
unity, a concept that, while it is reached from the finite, asserts
itself for thought through itself and bears certainty within itself as
the concept of what is grounded in itself, the unconditioned, no
abiding dualism in what is through the principle can be derived when
thought descends from that concept to the definitive knowledge of it.
The unity of the principle secures the unity of the human being and of
the whole of existence, shows every dualistic division in the one as in
the other to be untrue.

That unity in the human being, which has already been seen from the
side that the natural in man is not dualistically separated from the
mental, is here to be seen from the side that it expresses the inner
unity of the \emph{human mind} with itself in the harmony of its
fundamental functions. That view of life must then assert this harmony,
in asserting the relation of unity between cognition and action in
their true form, between knowledge and faith, between the ethical and
the religious. It must assert the relation of unity between
\emph{cognition} and \emph{action} because the willing that, as
ethical, has the realization of the essence as its goal, and is thus
the willing of the essence, must everywhere be guided by the cognition
in which the essence becomes manifest to itself, and because, therefore,
true cognition must become the condition for the clearly conscious
willing of the ethical, and for its concrete unfolding into a whole
that embraces the life of the actual personality\footnote{
Cf.: Om det Religiøse i dets Enhed med det Humane, p. 5---29.}.
It must assert the relation of unity between \emph{knowledge} and
\emph{faith} in its truth, because what is posited, in accordance with
a need of the essence, as the object of faith will not be able to
contradict the knowledge that is the true expression of the essence;
and because knowledge bears within itself the determination of faith,
just as it bears within itself that of devotion, and thereby shows
itself to be in essential unity with the religious according to the
latter's true meaning\footnote{
Ibid., p. 43---51.}. Finally, it must assert the unity between the
\emph{ethical} and the \emph{religious}, because the will too,
\opage{285} \setcounter{footnote}{0}in its truth, as ethical willing,
bears within itself the essential determinations of the religious,
faith and devotion; because the ethical choice, by which the
personality constitutes itself as ethical, comes into being only as the
relation to the absolute is posited; and because the ethical relations
everywhere include this relation, so that the religious absolute
relation becomes that which embraces all the relative relations, is
immanent in them, and, in giving them \emph{absolute} significance,
unfolds through it a \emph{concrete} content, and gains the shape of
actuality\footnote{
Ibid., p. 51---56.}.

In asserting the inner unity between the fundamental activities and
fundamental relations of the human mind, the view of life asserts what
is the condition for that \emph{reconciliation} of the mind with itself
in its actuality which cannot be thought so long as the mind is divided
by conflicting powers; for the consciousness divided within itself must
then at every moment be torn apart by doubt about the truth of
cognition, about the validity of action. It asserts what is the
condition for a life that is \emph{human} in the full sense of the
word, for the true \emph{humane} life, in which \emph{all} the powers
of the life of the mind work undisturbed, in conscious, willed harmony.
And this humane life is, as \emph{ethical} life, everywhere in unity
with the religious: the reconciliation of the mind in it is
\emph{religious} reconciliation.

\divrule

In the exposition of the history of philosophy we have stopped at the
systems that conclude the realistic and the idealistic line of
development of modern philosophy, and we have designated the main form
of philosophy that must succeed the philosophy of modern times as the
philosophy \emph{of the future}, that is, as \emph{not} \emph{yet}
historically actualized. A glance at the philosophical phenomena that
have appeared in rich multiplicity after \opage{286}those concluding
systems will show the justification of this designation.

The philosophical attempts that have appeared after the Système de la
nature with \emph{realism} as their foundation and point of departure
are either, like the \emph{materialism} of our time, mere repetitions
of the thoughts of the previous century, without any really new content
of thought; or they are consequences that differ from their
presuppositions only in extending the polemic against the a priori in
thinking to Kant's conception of it, and partly in the narrower
limitation of the domain of knowledge, and still fall under the same
principle. The forms of empiricism and sensualism that have appeared in
our century in \emph{England} and \emph{France} fall under the same
fundamental determination as realism in general, founder on the same
difficulties, and are guilty of the same inconsistencies. They have a
significance in having made contributions to the theory of the method
of experience, in having applied this method to the psychological and
the historical, and in having collected a rich material for
demonstrating the real conditions of the genetic development of the
psychic functions. But everywhere, in their consideration of these
functions, they have made what is only one side of the matter into the
whole matter, and have thereby foundered on the same difficulties
against which the older forms of realism struck when they would explain
the mental. The consequence that must necessarily strike every system
resting in fact on the foundation of atomism, the annihilation of
knowledge through the annihilation of the universal, they escape only
by arbitrarily keeping thought away from the concluding concepts and by
arbitrarily and inconsistently inserting a relative universal. And
despite all the apparent consistency with which they seek to remove the
a priori even from the domain of the mathematical (\emph{John Stuart}
\emph{Mill}), it nonetheless asserts itself involuntarily, and in the
most significant representative of this direction, in \emph{Comte}, the
founder \opage{287}of the ``positive philosophy'', which, by keeping our
thought away from first causes and from final causes, confines us to
the observation of the connections of phenomena and to the calling
forth of these by experiment, it is indirectly acknowledged through the
change of method in his second main work, ``Système de politique
positive'', which has human life as its object.

The philosophical attempts that have \emph{idealism} as their
foundation are, where they are not simple continuations, as in the
schools that have attached themselves to the more significant systems
and that, the Herbartian excepted, are now for the most part dissolved,
either merely eclectic attempts, like the philosophy that in France
itself calls itself \emph{eclecticism}, like the \emph{Italian}
idealist philosophy of our century, and like the \emph{Boströmian}
philosophy in our neighboring country, which ties together thoughts
drawn from Neoplatonism and Gnosis with thoughts from more recent
systems, notably Leibnitz's, and by its one-sided spiritualistic
conception of absolute being is compelled to bring in multiplicity and
the fundamental determinations of reality through borrowings from
experience, and, in explaining the phenomenon of the material, which
cannot be done away with, to violate the principle of contradiction
that governs the system; --- or they mark a reaction against that
monism of thought which has found its most consistent expression in the
Hegelian panlogism, but a reaction that has nonetheless not been able
to lead to the creation of systems that could gain abiding significance
as the expression of a higher fundamental form of thinking.

Among these philosophical attempts, those have been of lesser
significance in which the reaction against the Hegelian system took
the form of an assertion of the religious element of representation
against the supremacy of comprehending thought. In these
half-theological theories there has been too much unclarity and
prejudice for anything whole and abiding to be able to come of them.
Even where a truly philosophical endeavor to win \opage{288}a fuller and
more concrete concept of God than the Hegelian concept of absolute
spirit was joined with the religious interest of feeling, as in so
serious a thinker as C. H. \emph{Weisse}, there still remains an
after-effect of religious representation that prevents the attainment
of full clarity and consistency of thought.

Of greater significance have been the philosophical attempts in which
the reaction against the monism of logical thought came forward through
an assertion of the significance of \emph{being} and of
\emph{experience}. In them there has been a consciousness of the defect
in the development of idealist philosophy from the point at which, in
Kant, the moment of reality, which originally received its significance
as the objective presupposition of knowledge, stepped into the shadow of
the moment of ideality. And it is characteristic of the strength with
which this reaction has asserted itself in the post-Hegelian period
that the endeavor to assert being and experience appears also, if in
unclear and uncertain forms, within the half-theological philosophizing.

In the endeavor to assert the significance of being and of experience,
men meet who, with the Kantian philosophy as their common point of
departure, have arrived at the most opposed consequences. We have in
the foregoing discussed \emph{Schelling's} \emph{positive philosophy}
with its ``a priori experience'', and \emph{Herbart's}
``\emph{realism}'', originally directed against Fichte's subjectivism,
but then polemicizing against idealist philosophy in general. We touch
here on only two more men, \emph{Feuerbach} and \emph{Lotze}, in order
to point out with respect to them what justifies our judgment on the
latest forms of philosophy.

In \emph{Feuerbach} the critique of the one-sidedness of panlogism is
sharp and fine, and the corrective is pointed to with clear
consciousness. But as Feuerbach pursues with unbending energy the way he
indicates, it leads him to the opposite one-sidedness, and, despite all
the difference that remains between his philosophy, which cannot wholly
forget what Kant has secured, and vulgar materialism, he does not
escape \opage{289}meeting with the latter in the impossibility of
explaining the facts of mental life, and indeed already of organic
life, and of solving the problems these bring forth; and he ends in an
individualism that becomes the ruin of all thought and knowledge.

\emph{Lotze}, who was originally influenced by Hegel's doctrine, has
formed his peculiar philosophy, a ``spiritualistic realism'', through a
combination of Herbartian and Leibnitzian fundamental thoughts with the
concept of a substantial unity whose moments the monadic beings are,
and which conditions their law-determined reciprocal action and their
harmonious ordering. He takes his point of departure in the phenomenon
in order, by the way of experience, through the exact determining of
the mechanical reciprocal action of the elementary forces, to arrive at
a knowledge of the conditions without which the being or becoming of no
form of existence can really be causally explained. For all natural
existence, in his conception, mechanism becomes the form; all phenomena
without exception are to be explained according to its laws; all
actualization of something ideal becomes possible only according to
these. But the forms of space, of time and of motion are only the forms
in which, for finite consciousness, the order and relations of the
spaceless and timeless real beings (i.e.\ beings forming an independent
center of acting and suffering) in the intellectual connection of the
absolute, and their inner states, express themselves; matter is
materiality, i.e.\ the form for the revelation of a supersensible real,
and everything mechanical, the whole system of the mechanical
reciprocal actions between the elements, is only the expression of
something higher that reveals itself to us through it. Through the
world of laws and forms the world of purposes reveals itself. From what
is given as an actuality we are everywhere referred back to the ideal
that is its ground; what \emph{is} is to find its explanation in what
\emph{ought} to be: in purpose. But although all true actuality is thus
led back to the ideal, and in the last instance to the absolute that
embraces the feeling monads existing for themselves, Lotze nonetheless
does not bring his doctrine to a really concluded unity. His strength
lies \opage{290} \setcounter{footnote}{0}essentially in the middle
regions of thought, in the treatment of the individual concrete
problems; when his thought is to reach what is first and last for
knowledge, it becomes uncertain and is unable to gather the scattered
threads. Even if the unity of the absolute, of the highest actuality,
is postulated, this absolute is nonetheless supposed for our thought
constantly to split into a threefold of forms (forces): law, the real
with its active force, and purpose, which are said not to be knowable by
us in their unity, even though we must necessarily postulate that the
ground of law and of force must lie in purpose. The relation of the
individual monadic beings to the absolute that encloses them also
becomes uncertain in Lotze: now they appear as independent, now they
are wholly dominated by the unity of the substance. And finally, since
the monadic beings are conceived spiritualistically and determined as
the purely punctual, the unextended, the immaterial, there remains for
Lotze the same insoluble difficulty that was demonstrated above with
respect to Leibnitz: to comprehend the fundamental forms of the
material from that which excludes them from its essence. Only by
surreptitious moves like those of Leibnitz does Lotze arrive at the
determinations indispensable to him. His doctrine, which lacks the
ultimate, truly uniting point of unity, can only veil, not remove,
dualism\footnote{
With respect to the philosophical phenomena in Germany in the last
decades I refer to Dr.\ H. \emph{Høffding's} clear and thorough work:
Philosophien i Tydskland efter Hegel. Kbhvn., 1872.}.

Thus this glance at the most significant philosophical phenomena that
have appeared after the conclusion of modern philosophy's two main
forms\footnote{
If we have passed over \emph{Schopenhauer}, it is because in his
philosophy we can see only an incoherently and dilettantishly
constructed theory, which at every decisive point helps itself out with
mere postulates, is contradictory in its whole design, and has only
through accidental circumstances been able, for a time, to become a
matter of fashion.} confirms what was stated above: that there are as
yet only approaches to, and intimations of, a new
\opage{291}fundamental form of philosophy, but no real realization of
such a form. The tendencies that govern the attempts to form a new
philosophy are in most cases unclear, or they suffer from a
contradiction that makes them impossible to carry through. But through
the whole historical development of philosophy through the two
concluded main forms, the necessary task for thought is clearly
designated, and through the earlier development itself the way is at
the same time intimated that alone can lead to its solution.

\dblrule


\chapter*{Misprints and Corrections.}
\markboth{Misprints and Corrections}{Misprints and Corrections}

[The leaf of \emph{Trykfeil og Rettelser} bound at the end of this part
carries three lists: one for the first part (\emph{1ste Deel}), one for
this part (\emph{2den Deel}, 17 entries), and one for Brøchner's
\emph{Bidrag til Opfattelsen af Philosophiens historiske Udvikling}
(1869).  Since they correct the Danish wording, they are not translated
here; they are transcribed in full in the Danish edition.  The English
text above follows the corrections of the list for this part at each of
its sites, where the Danish edition keeps the printed reading.]

\end{document}
